% ترجمه تولیدشده با هوش مصنوعی است. اصلاحات و حروف‌چینی را OpenAI Codex — GPT-6 Astra با سطح استدلال Ultra تولید کرده است. نام امی نوتر و هم‌نویسندگان هر اثر حفظ می‌شود. ترجمهٔ انسانی یا داوری علمی مستقل انسانی به ناشر نسبت داده نمی‌شود.
% حقوق متعلق به صاحبان آن‌هاست. توزیع رایگان با هدف غیرتجاری و صرفاً آموزشی بر فرض ناشر دربارهٔ نبود مخالفت صاحبان حقوق استوار است، نه بر اجازهٔ احراز‌شده. دسترسی رایگان به‌خودی‌خود مجوز انتشار اثر مشتق نیست. صاحبان حقوق می‌توانند از پیوند زیر اعتراض کنند؛ ناشر بنا به درخواست آنان مطالب مربوط را برمی‌دارد. این بیان مجوز تازه‌ای اعطا نمی‌کند و آثار دیگران را به مالکیت عمومی نمی‌سپارد.
% https://github.com/KokunoYumeto/modern-latex-manuscripts/issues/new
\documentclass[11pt]{article}
\pdfvariable suppressoptionalinfo 767
\errorcontextlines=200
\usepackage{fontspec}
\usepackage{amsmath,amssymb,mathtools,array,booktabs,graphicx,adjustbox,pdflscape}
\usepackage[a4paper,margin=16mm,includefoot]{geometry}
\usepackage[english,bidi=basic,layout=sectioning lists contents footnotes captions tabular graphics]{babel}
\makeatletter
\@ifundefined{babelsubrl}{\DeclareRobustCommand\babelsubrl[1]{\leavevmode{\bbl@textdir\@ne#1}}}{}
\makeatother
\babelprovide[import=ar]{arabic}
\babelprovide[import=fa,main,onchar=ids fonts]{persian}
\babelfont{rm}{TeX Gyre Pagella}
\babelfont{sf}{TeX Gyre Heros}
\babelfont{tt}{Latin Modern Mono}
\babelfont[arabic]{rm}[Path=fonts/,UprightFont=Scheherazade-Regular.ttf,BoldFont=Scheherazade-Bold.ttf,ItalicFont=Scheherazade-Regular.ttf,ItalicFeatures={FakeSlant=0.15},BoldItalicFont=Scheherazade-Bold.ttf,BoldItalicFeatures={FakeSlant=0.15},Renderer=HarfBuzz,Scale=1.08]{Scheherazade-Regular.ttf}
\babelfont[persian]{rm}[Path=fonts/,UprightFont=Scheherazade-Regular.ttf,BoldFont=Scheherazade-Bold.ttf,ItalicFont=Scheherazade-Regular.ttf,ItalicFeatures={FakeSlant=0.15},BoldItalicFont=Scheherazade-Bold.ttf,BoldItalicFeatures={FakeSlant=0.15},Renderer=HarfBuzz,Scale=1.08]{Scheherazade-Regular.ttf}
\usepackage[hyperfootnotes=false]{hyperref}
\hypersetup{unicode=true,hidelinks,hypertexnames=false,pdflang=fa-IR,pdfauthor={},pdftitle={آثار امی نوتر — متن فارسی موجود},pdfcreator={},pdfproducer={}}
\makeatletter
\pdfstringdefDisableCommands{\let\babelsubrl\@firstofone}
\makeatother
\setlength{\parindent}{1.25em}
\setlength{\parskip}{0pt}
\setlength{\emergencystretch}{10em}
\sloppy\raggedbottom\allowdisplaybreaks
\providecommand{\A}{\Delta}
\providecommand{\D}{\Delta}
\providecommand{\R}{\varrho}
\providecommand{\N}{\operatorname{N}}
\providecommand{\modu}[1]{\;\operatorname{mod} #1}
\providecommand{\mcell}[1]{\ensuremath{\begin{gathered}#1\end{gathered}}}
\providecommand{\glqq}{\textquotedblbase}
\providecommand{\grqq}{\textquotedblleft}
\providecommand{\srcspaced}[1]{{\bfseries #1}}
\providecommand{\srcnumdisplay}[3][3.4em]{%
  \par\smallskip\noindent
  \babelsublr{\begin{minipage}[t]{#1}\vspace*{0pt}\srcLTR{#2}\end{minipage}%
  \begin{minipage}[t]{\dimexpr\linewidth-#1\relax}\vspace*{0pt}\centering
  \(\displaystyle #3\)\end{minipage}}\par\smallskip}
\newcommand{\srcRTL}[1]{\relax\ifmmode\text{\mbox{\babelsubrl{\foreignlanguage{persian}{#1}}}}\else\babelsubrl{\foreignlanguage{persian}{#1}}\fi}
\newcommand{\srcLTR}[1]{\relax\ifmmode\text{\mbox{\babelsublr{\foreignlanguage{english}{#1}}}}\else\babelsublr{\foreignlanguage{english}{#1}}\fi}
\newcommand{\srcfn}[2]{\begingroup\renewcommand{\thefootnote}{\srcLTR{#1}}\footnote{#2}\addtocounter{footnote}{-1}\endgroup}
\newcommand{\srcfnmark}[1]{\begingroup\renewcommand{\thefootnote}{\srcLTR{#1}}\footnotemark\addtocounter{footnote}{-1}\endgroup}
\newcommand{\srcfntext}[2]{\begingroup\renewcommand{\thefootnote}{\srcLTR{#1}}\footnotetext{#2}\endgroup}
\providecommand{\Lsys}{\mathfrak{L}}
\providecommand{\Jdom}{\mathfrak{J}}
\providecommand{\Gdom}{\mathfrak{G}}
\providecommand{\rcdots}{\,\cdots\,}
\providecommand{\pair}[2]{(#1\,|\,#2)}
\providecommand{\eps}{\varepsilon}
\RequirePackage{tikz}
\newcommand{\NoetherRosetteGlyph}{%
  \begin{tikzpicture}[x=0.1363584905660377pt,y=-0.1363584905660377pt,baseline=(current bounding box.south)]
    \path[use as bounding box] (0,0) rectangle (108,108);
    \fill[even odd rule]
    (33.5,84) -- (32.5,83) -- (30.5,83) -- (29,81.5) -- (29,76.5) -- (30.5,75) -- (31.5,75) -- (32.5,74) -- (34.5,74) -- (36,72.5) -- (36,71.5) -- (35.5,71) -- (33.5,71) -- (33,70.5) -- (33,69.5) -- (30.5,67) -- (28.5,69) -- (26.5,69) -- (24.5,71) -- (23.5,71) -- (22.5,72) -- (21.5,72) -- (20.5,73) -- (18.5,73) -- (17.5,72) -- (16.5,73) -- (16,72.5) -- (16,68.5) -- (18.5,66) -- (19.5,66) -- (20,65.5) -- (20,64.5) -- (26,58.5) -- (26,56.5) -- (24,54.5) -- (24,53.5) -- (17,46.5) -- (17,45.5) -- (16.5,45) -- (15.5,45) -- (15,44.5) -- (15,42.5) -- (15.5,42) -- (17.5,42) -- (19.5,44) -- (20.5,44) -- (21,44.5) -- (20,45.5) -- (21,46.5) -- (21,47.5) -- (22,48.5) -- (22,49.5) -- (29,56.5) -- (29,57.5) -- (27,59.5) -- (27,60.5) -- (26,61.5) -- (26,62.5) -- (24,64.5) -- (24,65.5) -- (22,67.5) -- (22.5,68) -- (24.5,68) -- (25.5,67) -- (28.5,67) -- (29.5,66) -- (31.5,66) -- (32.5,65) -- (33.5,65) -- (34.5,64) -- (36.5,64) -- (37.5,63) -- (38.5,63) -- (39.5,62) -- (41.5,62) -- (42.5,61) -- (44.5,61) -- (45.5,60) -- (46.5,60) -- (47.5,59) -- (48.5,59) -- (49,58.5) -- (49,57.5) -- (49.5,57) -- (51.5,57) -- (52,56.5) -- (50.5,55) -- (49.5,55) -- (48.5,54) -- (47.5,54) -- (46.5,53) -- (45.5,53) -- (43,50.5) -- (43,49.5) -- (42,48.5) -- (42,47.5) -- (41.5,47) -- (40.5,47) -- (39.5,48) -- (33.5,48) -- (32,46.5) -- (32,44.5) -- (34,42.5) -- (34,40.5) -- (33.5,40) -- (32.5,40) -- (31.5,39) -- (26.5,39) -- (26,38.5) -- (26,35.5) -- (27.5,34) -- (31.5,34) -- (33,35.5) -- (33,36.5) -- (34.5,38) -- (37.5,38) -- (38,37.5) -- (38,36.5) -- (41,33.5) -- (41,31.5) -- (40,30.5) -- (40,27.5) -- (39,26.5) -- (39,25.5) -- (38,24.5) -- (38,23.5) -- (37,22.5) -- (37,19.5) -- (37.5,19) -- (38.5,20) -- (39.5,20) -- (41,21.5) -- (41,24.5) -- (42,25.5) -- (42,27.5) -- (43,28.5) -- (43,30.5) -- (44,31.5) -- (44,35.5) -- (45,36.5) -- (45,37.5) -- (45.5,38) -- (46.5,38) -- (47,38.5) -- (47,40.5) -- (48,41.5) -- (48,42.5) -- (49,43.5) -- (49,44.5) -- (50,45.5) -- (50,48.5) -- (53,51.5) -- (52,52.5) -- (52,55.5) -- (53.5,57) -- (56.5,57) -- (58.5,59) -- (59.5,59) -- (60.5,60) -- (61.5,60) -- (62,60.5) -- (62,61.5) -- (63,62.5) -- (63,64.5) -- (63.5,65) -- (64.5,65) -- (65.5,64) -- (68.5,64) -- (69.5,65) -- (73.5,65) -- (74,65.5) -- (74,66.5) -- (75,67.5) -- (75,68.5) -- (75.5,69) -- (76,68.5) -- (76,66.5) -- (76.5,66) -- (78.5,68) -- (80.5,68) -- (81.5,69) -- (84.5,69) -- (85,68.5) -- (84,67.5) -- (84,65.5) -- (83.5,65) -- (82.5,65) -- (82,64.5) -- (82.5,64) -- (84.5,64) -- (86,65.5) -- (86,66.5) -- (88.5,69) -- (89.5,69) -- (90,69.5) -- (90,70.5) -- (91,71.5) -- (91,72.5) -- (93,74.5) -- (91.5,76) -- (90.5,76) -- (88.5,74) -- (87.5,74) -- (86.5,73) -- (85.5,73) -- (83.5,71) -- (82.5,71) -- (81.5,70) -- (80.5,70) -- (79.5,69) -- (77.5,69) -- (74.5,72) -- (73.5,72) -- (73,72.5) -- (73,74.5) -- (73.5,75) -- (78.5,75) -- (79.5,74) -- (80.5,74) -- (81,74.5) -- (81,78.5) -- (80,79.5) -- (80,80.5) -- (78.5,82) -- (73.5,82) -- (73,81.5) -- (73,79.5) -- (72,78.5) -- (72,77.5) -- (71.5,77) -- (70.5,77) -- (69.5,76) -- (68.5,76) -- (68,76.5) -- (68,77.5) -- (67,78.5) -- (67,79.5) -- (66.5,80) -- (65.5,80) -- (65,80.5) -- (66,81.5) -- (64.5,83) -- (63,81.5) -- (63,79.5) -- (62,78.5) -- (62,77.5) -- (61,76.5) -- (61,74.5) -- (60,73.5) -- (60,72.5) -- (59,71.5) -- (59,68.5) -- (58,67.5) -- (58,64.5) -- (56,62.5) -- (56,61.5) -- (55.5,61) -- (54.5,62) -- (53.5,62) -- (52,63.5) -- (52,64.5) -- (51,65.5) -- (51,66.5) -- (50,67.5) -- (50,68.5) -- (49,69.5) -- (49,70.5) -- (48,71.5) -- (48,75.5) -- (46,77.5) -- (46,79.5) -- (44.5,81) -- (42.5,81) -- (42,80.5) -- (44,78.5) -- (44,77.5) -- (43.5,77) -- (41.5,77) -- (39.5,75) -- (38.5,76) -- (37.5,76) -- (37,76.5) -- (37,78.5) -- (36,79.5) -- (38,81.5) -- (38,82.5) -- (37.5,83) -- (34.5,83) -- cycle
    (65.5,22) -- (63.5,22) -- (63,21.5) -- (63.5,21) -- (65.5,21) -- (66,21.5) -- cycle
    (55.5,28) -- (53.5,28) -- (53,27.5) -- (53.5,27) -- (54.5,27) -- (57.5,24) -- (59.5,24) -- (60,24.5) -- (57.5,27) -- (56.5,27) -- cycle
    (64.5,26) -- (64,25.5) -- (64.5,25) -- (65,25.5) -- cycle
    (63.5,29) -- (63,28.5) -- (63,27.5) -- (63.5,27) -- (64.5,27) -- (65,27.5) -- cycle
    (58.5,55) -- (57,53.5) -- (57,52.5) -- (56,51.5) -- (56,49.5) -- (57,48.5) -- (57,46.5) -- (60,43.5) -- (60,42.5) -- (59.5,42) -- (58.5,42) -- (58,41.5) -- (59,40.5) -- (59,39.5) -- (61.5,37) -- (62.5,37) -- (63.5,36) -- (65.5,36) -- (66.5,37) -- (68.5,37) -- (70,35.5) -- (69,34.5) -- (69,32.5) -- (68,31.5) -- (68.5,31) -- (74.5,31) -- (75,31.5) -- (75,32.5) -- (76,33.5) -- (76,35.5) -- (75,36.5) -- (75,37.5) -- (72,40.5) -- (72,41.5) -- (71,42.5) -- (72,43.5) -- (72,44.5) -- (72.5,45) -- (73.5,45) -- (74.5,46) -- (75.5,46) -- (76.5,45) -- (78.5,45) -- (79.5,44) -- (80,44.5) -- (80,47.5) -- (79.5,48) -- (77.5,48) -- (76.5,49) -- (74.5,49) -- (73.5,50) -- (71.5,50) -- (70.5,51) -- (68.5,51) -- (67.5,52) -- (65.5,52) -- (64.5,53) -- (63.5,53) -- (62.5,54) -- (59.5,54) -- cycle
    (86.5,50) -- (85.5,50) -- (85,49.5) -- (86,48.5) -- (86,46.5) -- (85.5,46) -- (82.5,46) -- (81.5,47) -- (81,46.5) -- (82,45.5) -- (82,44.5) -- (82.5,44) -- (84.5,44) -- (86.5,42) -- (88.5,42) -- (89.5,41) -- (92.5,41) -- (93,41.5) -- (93,42.5) -- (94,43.5) -- (94,44.5) -- (93.5,45) -- (92.5,45) -- (88.5,49) -- (87.5,49) -- cycle
    (44,42.5) -- (43.5,42) -- (43,42.5) -- (43.5,43) -- cycle
    (45,43.5) -- (44.5,43) -- (44,43.5) -- (44.5,44) -- cycle
    (80.5,60) -- (79.5,60) -- (79,59.5) -- (79,56.5) -- (81,54.5) -- (81,53.5) -- (82.5,52) -- (83,52.5) -- (83,53.5) -- (82,54.5) -- (82,55.5) -- (81.5,56) -- (80.5,56) -- (80,56.5) -- (80,57.5) -- (81,58.5) -- (81,59.5) -- cycle
    (55.5,54) -- (55,53.5) -- (55.5,53) -- (56,53.5) -- cycle
    (43,65.5) -- (42.5,65) -- (42,65.5) -- (42.5,66) -- cycle
    (61,67.5) -- (60.5,67) -- (60,67.5) -- (60.5,68) -- cycle
    (47,69.5) -- (47,68.5) -- (46.5,68) -- (46,68.5) -- (46,69.5) -- (46.5,70) -- cycle
    (62,68.5) -- (61.5,68) -- (61,68.5) -- (61.5,69) -- cycle
    (67.5,97) -- (65.5,97) -- (64,95.5) -- (64,94.5) -- (61.5,92) -- (60.5,92) -- (54.5,86) -- (53.5,86) -- (49.5,90) -- (48.5,90) -- (47.5,91) -- (46.5,91) -- (45.5,92) -- (45,91.5) -- (45,90.5) -- (45.5,90) -- (46.5,90) -- (48.5,88) -- (49.5,88) -- (54.5,83) -- (55.5,83) -- (60.5,88) -- (61.5,88) -- (62.5,89) -- (63.5,89) -- (64,89.5) -- (64,90.5) -- (64.5,91) -- (66.5,91) -- (67,90.5) -- (66,89.5) -- (67,88.5) -- (65,86.5) -- (65,85.5) -- (65.5,85) -- (66.5,85) -- (67,85.5) -- (67,86.5) -- (68,87.5) -- (68,90.5) -- (69,91.5) -- (69,94.5) -- (70,95.5) -- (69.5,96) -- (68.5,96) -- cycle
    (40.5,96) -- (36.5,96) -- (36,95.5) -- (37,94.5) -- (37,91.5) -- (40.5,88) -- (41.5,88) -- (42,88.5) -- (41,89.5) -- (41,92.5) -- (42,93.5) -- (42,94.5) -- cycle;
  \end{tikzpicture}%
}
\providecommand{\noethpIIrosette}{\makebox[9.44444pt][c]{\raisebox{-3.85pt}[0pt][0pt]{\NoetherRosetteGlyph}}}
\providecommand{\NoetherRosetteInline}{\makebox[17.72671698pt][c]{\raisebox{-3.85pt}[10.87671698pt][3.85pt]{\NoetherRosetteGlyph}}}


\newenvironment{NoetherWideTable}[2]{%
  \newgeometry{layoutwidth=#1,layoutheight=#2,margin=16mm,includefoot}%
  \paperwidth=#1\paperheight=#2\pagewidth=#1\pageheight=#2
  \begin{otherlanguage}{english}\begin{landscape}%
}{%
  \end{landscape}\end{otherlanguage}
}
\AddToHook{env/NoetherWideTable/after}{%
  \restoregeometry
  \paperwidth=210mm\paperheight=297mm\pagewidth=210mm\pageheight=297mm
}
\newtagform{NoetherLTR}[\srcLTR]{(}{)}
\usetagform{NoetherLTR}
\begin{document}
\selectlanguage{persian}
\begin{center}
{\Large\bfseries 1. دربارهٔ تشکیل دستگاه فرم‌های یک فرم سه‌تاییِ درجهٔ چهار.}\par
\vspace{1.2em}
نوشتهٔ \srcspaced{امی نوتر}.\par
\vspace{0.8em}
(برگرفته از رسالهٔ دکتری نویسنده.)\par
\vspace{1.2em}
\srcLTR{Sitz. Ber. d. Physikal.-mediz. Sozietät in Erlangen 39 (1907), S. 176--179}
\end{center}

آثار \srcspaced{گوردان}، \srcspaced{مایزانو} و \srcspaced{پاسکال} به دستگاه فرم‌های یک فرم سه‌تاییِ درجهٔ چهار می‌پردازند\srcfn{1)}{\srcLTR{\srcspaced{P. Gordan}, Über das volle Formensystem der ternären biquadratischen Form $f=x_1^3x_2+x_2^3x_3+x_3^3x_1$ (Math. Annalen Bd. XVII, S. 217 bis 233. 1880); \srcspaced{G. Maisano}, 1. Sistemi completi dei primi cinque gradi della forma ternaria biquadratica e degl' invarianti, covarianti e contravarianti di sesto grado (Giorn. di Battaglini XIX). 2. Sui covarianti indipendenti di $6^\circ$ grado nei coefficienti della forma biquadratica ternaria (Rend. Circ. Mat. di Palermo I. 1887); \srcspaced{E. Pascal}, Contributo alla teoria della forma ternaria biquadratica e delle sue varie decomposizioni in fattori (Memoria premiata dalla R. Accademia delle scienze fisiche e matematiche di Napoli. 1905).}}. آقای \srcspaced{گوردان} دستگاه کامل فرم ویژهٔ خودریخت (\srcspaced{automorphic form}) $f=x_1^3x_2+x_2^3x_3+x_3^3x_1$ را، که از 54 ساخته تشکیل می‌شود، بر پایهٔ اصولی همانند اصولی بنا می‌کند که برای دستگاه‌های فرم در حوزهٔ دوتایی به کار برده بود.

آقای \srcspaced{مایزانو} برای فرم عمومیِ درجهٔ چهار، فرم‌ها را تا مرتبهٔ پنجم\srcfn{2)}{در اینجا از \glqq{}مرتبه\grqq{}، مرتبه نسبت به ضرایب، و از \glqq{}درجه\grqq{}، درجه نسبت به متغیرها منظور است.}، با احتساب خودِ مرتبهٔ پنجم، به دست داده است؛ همچنین، به روش آقای \srcspaced{گوردان} برای فرم مکعبیِ سه‌تایی در جلد نخست Math. Annalen، چند ناوردا، هموردا (\srcspaced{covariant}) و پادهموردا (\srcspaced{contravariant}) با مرتبه‌های بالاتر را به دست داده است. آقای \srcspaced{پاسکال} با استفاده از نتایج مایزانو، عمدتا مسئلهٔ تجزیهٔ فرم درجهٔ چهار به عامل‌ها را بررسی می‌کند\srcfn{1)}{آقای \srcspaced{مایزانو} در یادداشت کوتاه دوم می‌کوشد وابستگی خطیِ سه هموردای دارای مرتبهٔ 6 و درجهٔ 6 را ثابت کند. در مقابل این برهانِ نه‌چندان کامل، آقای \srcspaced{پاسکال} گمان می‌کند استقلال خطیِ سه هموردای دارای مرتبهٔ 6 و نیز استقلال سه ناوردای دارای مرتبهٔ 9 را ثابت کرده است. بااین‌حال، محاسبات من نشان داده است که در واقع میان سه هموردا یک رابطهٔ خطی و میان سه ناوردا رابطهٔ خطیِ دیگری وجود دارد؛ فرمول‌های صریح، با نمادگذاری \srcspaced{Pascal}، چنین‌اند:
\[
0=20\Omega_1+6\Omega_2-3\Omega_3+4f\cdot C_1-12f\cdot C_2-4A\cdot\Delta.
\]
\[
0=4A^3-15A\cdot B+30C-90D+9E.
\]}. 

\srcspaced{هدف پژوهش‌های من تشکیل دستگاه فرم‌های فرم عمومیِ سه‌تاییِ درجهٔ چهار است؛ در آغاز فقط شالوده‌های اصلی را ارائه می‌کنم، یعنی یک \glqq{}دستگاه کامل نسبی\grqq{} به دست می‌دهم}\srcfn{2)}{برای این نام‌گذاری بنگرید به \srcspaced{Gordan-Kerschensteiner}, Vorlesungen über Invariantentheorie, Bd. II, S. 227.}\srcspaced{.}

ایدهٔ بنیادیِ تشکیل دستگاه فرم‌های سه‌تایی همان ایدهٔ حوزهٔ دوتایی است. از یک دستگاه کامل نسبیِ نخستین ــ یعنی دستگاه فرم‌هایی که از حوزهٔ دوتایی گرفته شده‌اند ــ بر پایهٔ قانونی معین به دستگاه‌هایی با پیمانه‌های هرچه بالاتر می‌رسیم؛ این روند ادامه دارد تا دستگاه یک پیمانه ــ با درنظرگرفتن پیمانه به‌عنوان فرم پایه ــ متناهی و شناخته‌شده شود، یا پیمانه‌ای به فرم‌هایی کاهش‌پذیر باشد که ناورداها در آن‌ها عامل‌اند. در حالت نخست، با ترانس‌وکشنِ دستگاه کامل نسبی بر دستگاه پیمانه (\srcspaced{transvection; Überschiebung})، دستگاه کامل مطلق پدید می‌آید؛ در حالت دوم، دستگاه کامل نسبی و دستگاه کامل مطلق یکی می‌شوند. چون \srcspaced{هیلبرت} متناهی‌بودن دستگاه‌های فرم را به‌طور کلی ثابت کرده است، این روش ناگزیر باید پایان پذیرد.

در حالت ما، پیمانهٔ $(abc)$ در دستگاه کامل نسبیِ نخست را می‌توان به پیمانه‌های زیر کاهش داد
\[
\Delta=(abc)^2a_x^2b_x^2c_x^2\quad\text{و}\quad \nu=(abu)^4;
\]
از این دو، دستگاه کامل نسبی به پیمانهٔ $\nu$ به‌آسانی به دست می‌آید. اکنون فرم‌های زیر را به‌عنوان رشتهٔ پیمانه‌ها برمی‌گزینیم:
\[
\nu;\qquad \nu^{(\nu)}=(\nu\nu_1x)^4=s_x^4,
\]
\[
\nu(s)=(ss'u)^4\ldots
\]
و دو فرم درجهٔ دو را که هنگام تشکیل دستگاه کامل نسبی به پیمانهٔ $s$ پدید می‌آیند، یعنی $u_\varrho^2$ و $t_x^2$، به‌عنوان پیمانه‌های دیگر می‌افزاییم. سپس می‌توان نشان داد که پیمانهٔ پس از $s$، یعنی $(ss'u)^4$، به پیمانهٔ $(\varrho,t)$ کاهش‌پذیر است. چون دستگاه مشترک دو فرم درجهٔ دو متناهی و شناخته‌شده است، پایان یادشده در بالا به دست آمده است. دستگاه کامل نسبی به پیمانهٔ $(\varrho,t)$ شامل 331 ساخته است. ــ

اکنون توضیحاتی دیگر دربارهٔ روش به‌کاررفته می‌دهم.

فرایند بنیادیِ ساخت فرم‌ها، یعنی \srcspaced{فرایند تاخوردگی}، در حوزهٔ سه‌تایی چنین تعریف می‌شود:

اگر در یک حاصل‌ضرب نمادی
\[
s_x^m t_x^n u_\sigma^\mu u_\tau^\nu
\]
جفت‌عامل‌های
\[
\begin{array}{@{}c@{\qquad}c@{\qquad}c@{\qquad}c@{\qquad}c@{}}
 & s_xt_x & u_\sigma u_\tau & s_xu_\sigma\ \text{یا}\ s_xu_\tau & t_xu_\tau\ \text{یا}\ t_xu_\sigma\\[0.45em]
\text{به‌ترتیب با} & (stu) & (\sigma\tau x) & s_\sigma\ \text{یا}\ s_\tau & t_\tau\ \text{یا}\ t_\sigma\\[0.45em]
\text{تاخوردگی} & \mathrm{I} & \mathrm{II} & \mathrm{III} & \mathrm{IV}
\end{array}
\]
جایگزین کنیم، آنگاه فرم‌های حاصل با تاخوردگی از فرم اصلی پدید آمده‌اند\srcfn{1)}{\srcspaced{Gordan}. Math. Annalen, Bd. XVII, S. 219.}.

همواره می‌توان $s_\sigma=0;\ t_\tau=0$ گذاشت. دربارهٔ رابطهٔ میان تاخوردگی‌های منفرد، قضیهٔ زیر برقرار است:

\srcspaced{تاخوردگی‌های I و II تاخوردگی‌های بنیادی‌اند و تاخوردگی‌های III و IV را می‌توان، مستقل از ترتیب ترکیب، از آن‌ها ساخت. به بیان دیگر: برای ساختن همهٔ فرم‌هایی که با تاخوردگی از یک عبارت معین پدید می‌آیند، فقط باید تاخوردگی‌های I و II را به کار برد}\srcfn{2)}{قضیه در حالتی استثنا دارد که $n$ و $\mu$، یا به‌ترتیب $m$ و $\nu$، همزمان صفر باشند؛ در این صورت \glqq{}رشتهٔ فرم‌ها\grqq{} به عضو قطری کاهش می‌یابد.}.

\glqq{}رشتهٔ فرم‌ها\grqq{} را یک فرم آغازین همراه با همهٔ فرم‌هایی تعریف می‌کنیم که با تاخوردگیِ آن در خودش از آن پدید می‌آیند. بنا بر قضیه، می‌توان رشتهٔ فرم‌های $s_x^m t_x^n u_\sigma^\mu u_\tau^\nu$ را در یک آرایش مستطیلی قرار داد. با پیشروی در آن به موارد زیر می‌رسیم:
\begin{enumerate}
\item با حرکت یک ستون به راست، همهٔ فرم‌هایی که با تاخوردگی I از فرم‌های مجاور پدید می‌آیند،
\item با حرکت یک سطر به پایین، همهٔ فرم‌هایی که با تاخوردگی II از فرم‌های بالایی پدید می‌آیند، و در نتیجه همهٔ فرم‌های حاصل از تاخوردگی.
\end{enumerate}

با این فرض‌ها، می‌توان قضیه‌های مربوط به کاهشگرها را در کلی‌ترین صورت و با تعریف زیر بیان کرد
\begin{center}
\glqq{}کاهشگر، رشته‌ای از فرم‌هاست که خودِ رشته کاهش‌پذیر است\grqq{}
\end{center}
می‌توان آن‌ها را به صورت زیر بیان کرد:

\begin{center}
\srcspaced{اگر فرم آغازینِ یک رشتهٔ فرم‌ها از این جهت کاهش‌پذیر باشد که یکی از اعضای آن با تاخوردگی به‌وسیلهٔ یک کاهشگر پدید آمده است، و اگر فرم پایانیِ رشته با تاخوردگی از همان عضو پدید آمده باشد، آنگاه کل رشتهٔ فرم‌ها کاهش‌پذیر است.}
\end{center}

دو روش کاهشِ دیگر نیز باید نام برده شوند:
\begin{enumerate}
\item آنچه \glqq{}کاهش دوگانه\grqq{} نامیده می‌شود؛ یعنی کاهش یک فرم به دو شیوه برای به‌دست‌آوردن \srcspaced{رابطه‌ای میان} فرم‌های مرتبهٔ بالاتر،
\item \glqq{}تاخوردگی با فرم‌هایی که به عامل‌ها تجزیه می‌شوند\grqq{}، که تا حدی \srcspaced{ویژگیِ} کاهش به‌وسیلهٔ کاهشگرها و تا حدی \srcspaced{ویژگیِ \glqq{}کاهش دوگانه\grqq{}} را نشان می‌دهد.
\end{enumerate}
\clearpage
\setcounter{footnote}{0}

\providecommand{\Afield}{\mathfrak A}
\providecommand{\Bfield}{\mathfrak B}
\providecommand{\Cfield}{\mathfrak C}
\providecommand{\Sdomain}{\mathfrak S}
\providecommand{\Kfield}{\mathfrak K}
\providecommand{\Rfield}{\mathfrak R}
\providecommand{\Xreal}{\mathfrak X}
\providecommand{\Yreal}{\mathfrak Y}
\providecommand{\qtext}[1]{\glqq #1\grqq{}}
\begin{center}
{\Large\bfseries 2. دربارهٔ تشکیل دستگاه فرم‌ها برای فرم سه‌تاییِ درجهٔ چهار\srcfn{*)}{گزیده‌ای کوتاه از مقدمه و فصل اول در \glqq{}Sitzungsber. der phys.-med. Sozietät Erlangen 1907\grqq{}، صص. 176 تا 179، منتشر شده است.}.}\par
\vspace{0.8em}
نوشتهٔ دوشیزه \emph{\srcLTR{Emmy Noether}} در ارلانگن.\par
\vspace{1.2em}
\srcLTR{Journal f. d. reine u. angew. Math. 134 (1908)}، صص. 23--90 و دو جدول
\end{center}

\vspace{3em}

\begin{center}
\textbf{مقدمه.}
\end{center}

آثار \emph{گوردان}، \emph{مایزانو} و \emph{پاسکال}\srcfn{**)}{\srcLTR{\emph{P. Gordan}, Über das volle Formensystem der ternären biquadratischen Form
\[
f=x_1^3x_2+x_2^3x_3+x_3^3x_1.
\]
(Math. Annalen, Bd. XVII (1880), S. 217--233.)
\emph{G. Maisano}, 1) \emph{Sistemi completi dei primi cinque gradi della forma ternaria biquadratica e degl' invarianti, covarianti e contravarianti di sesto grado}. (Giorn. di Battaglini XIX.)
2) \emph{Sui covarianti indipendenti di $6^\circ$ grado nei coefficienti della forma biquadratica ternaria}. (Rend. Circ. Mat. di Palermo I, 1887.)
\emph{E. Pascal}, \emph{Contributo alla teoria della forma ternaria biquadratica e delle sue varie decomposizioni in fattori}. (Memoria premiata dalla R. Accademia delle scienze fisiche e matematiche di Napoli. 1905.)}} به دستگاه فرم‌های وابسته به فرم سه‌تاییِ درجهٔ چهار می‌پردازند. آقای \emph{گوردان} دستگاه کامل فرم ویژهٔ خودریخت \(f=x_1^3x_2+x_2^3x_3+x_3^3x_1\) را، که از 54 ساخته تشکیل می‌شود، بر پایهٔ اصولی همانند اصولی بنا می‌کند که برای دستگاه‌های فرم در حوزهٔ دوتایی عرضه کرده بود.
آقای \emph{مایزانو} برای فرم عمومیِ سه‌تاییِ درجهٔ چهار، فرم‌ها را تا مرتبهٔ پنجم، با احتساب خودِ مرتبهٔ پنجم، می‌سازد\srcfn{***)}{منظور از \glqq{}مرتبه\grqq{} درجه نسبت به ضرایب و منظور از \glqq{}درجه\grqq{} درجه نسبت به متغیرهاست.}؛ همچنین چند ناوردا، هموردا و پادهموردای مرتبه‌های بالاتر را به روش آقای \emph{گوردان} برای فرم سه‌تاییِ درجهٔ سه در جلد اول Math. Annalen به دست می‌دهد. آقای \emph{پاسکال} با استفاده از نتایج \emph{مایزانو} عمدتا به مسئلهٔ تجزیهٔ فرم درجهٔ چهار به عامل‌ها می‌پردازد\srcfn{*)}{آقای \emph{مایزانو} در یادداشت کوتاه دوم می‌کوشد وابستگی خطیِ سه هموردای مرتبهٔ شش و درجهٔ شش را اثبات کند. در برابر این برهان نه‌چندان کامل، آقای \emph{پاسکال} گمان می‌کند استقلال خطیِ سه هموردای مرتبهٔ شش و نیز استقلال سه ناوردای مرتبهٔ نه را اثبات کرده است. بااین‌حال، فرمول‌های (13)، § 11 و (22)، § 17، یادداشت، در این مقاله نشان می‌دهند که واقعا یک رابطهٔ خطی میان آن سه هموردا و رابطهٔ خطیِ دیگری میان آن سه ناوردا وجود دارد؛ این روابط در آنجا صریحا داده شده‌اند. بنا بر اطلاع آقای \emph{پاسکال}، این تناقض از خطایی عددی در عبارت پادهموردای $p$ او (که اینجا $\varrho$ نامیده می‌شود) در ص. 46 مقاله‌اش ناشی است.}.
\emph{هدف بررسی‌های زیر، ساخت دستگاه فرم‌ها برای فرم عمومیِ سه‌تاییِ درجهٔ چهار است؛ در این مقاله فقط مبانی اصلی عرضه و یک \glqq{}دستگاه کامل نسبی\grqq{} ساخته می‌شود}\srcfn{**)}{برای نام‌گذاری، نک. § 3a.}\emph{.}
این مقاله پیوندی نزدیک با مقالهٔ \emph{گوردان} دارد؛ بااین‌حال، اصولی که آنجا فقط به‌صورتی کاملا کلی آمده بود، باید جزءبه‌جزء بسط داده می‌شد. با کمک قضیه‌ای دربارهٔ ارتباط تاخوردگی‌ها و با معرفی \glqq{}رشتهٔ فرم‌ها\grqq{} (§ 1)، قضایای کاهش به‌دقت صورت‌بندی و تکمیل می‌شوند (§ 3)؛ درحالی‌که ساخت بازگشتیِ بسط‌های ویژه (§ 2) ابزار محاسباتیِ اجرای واقعیِ کاهش‌ها را فراهم می‌کند.
اندیشهٔ بنیادیِ ساخت دستگاه‌های فرم سه‌تایی همان اندیشهٔ حوزهٔ دوتایی است. از نخستین دستگاه کامل نسبی ــ یعنی دستگاه فرم‌های برگرفته از حوزهٔ دوتایی ــ طبق قاعده‌ای معین به دستگاه‌هایی با پیمانه‌های پیوسته بالاتر می‌رسیم، تا سرانجام دستگاه یک پیمانه ــ با درنظرگرفتن پیمانه به‌عنوان فرم بنیادی ــ متناهی و شناخته‌شده شود، یا پیمانه‌ای به فرم‌هایی کاهش یابد که ناورداها را به‌صورت عامل دارند. در حالت نخست، با ترانس‌وکشنِ دستگاه کامل نسبی بر دستگاه پیمانه، دستگاه کامل مطلق پدید می‌آید؛ در حالت دوم، دستگاه کامل نسبی و دستگاه کامل مطلق یکسان می‌شوند. به‌سبب متناهی‌بودن دستگاه‌های فرم که آقای \emph{هیلبرت} آن را به‌طور کلی اثبات کرده است، این فرایند ناگزیر پایان می‌یابد.
در حالت ما، پیمانهٔ $(abc)$ نخستین دستگاه کامل نسبی (§ 4) به پیمانه‌های زیر کاهش می‌یابد:
\[
\Delta=(abc)^2a_x^2b_x^2c_x^2\quad\text{و}\quad \nu=(abu)^4
\]
(§ 5)، و سپس دستگاه کامل نسبی به پیمانهٔ $\nu$ به دست می‌آید (§ 6). این نتایج پیش‌تر در مقالهٔ \emph{گوردان} برای فرم عمومیِ درجهٔ چهار آمده‌اند؛ بااین‌حال، شیوهٔ استنتاج ما آسان‌تر به فرم‌های درجهٔ بالاتر تعمیم می‌یابد.
اکنون فرم‌های زیر (§ 7) را به‌عنوان رشتهٔ پیمانه‌ها برمی‌گزینیم:
\[
\nu,\qquad \nu(\nu)=(\nu\nu_1x)^4=s_x^4,\qquad \nu(s)=(ss'u)^4,\ldots .
\]

هنگام ساخت دستگاه کامل نسبی به پیمانهٔ $s$، دو فرم درجهٔ دوِ ظاهرشده در دستگاه، یعنی $u_\varrho^2$ و $t_x^2$، به‌عنوان پیمانه‌ها ملحق می‌شوند (§ 9 و § 10) و در نتیجه دستگاه کامل نسبی به پیمانهٔ $(s,\varrho,t)$ ساخته می‌شود؛ سپس نشان داده می‌شود (§ 17) که پیمانهٔ $(ss'u)^4$ به پیمانه‌های $\varrho$ و $t$ کاهش‌پذیر است. بدین‌ترتیب دستگاه بالاترِ بعدی به دستگاهی کامل نسبی به پیمانهٔ $(\varrho,t)$ تبدیل می‌شود. چون دستگاه مشترک دو فرم درجهٔ دو متناهی و شناخته‌شده است و در حالت کلی از 20 ساخته تشکیل می‌شود\srcfn{*)}{نک. \emph{Clebsch-Lindemann}, Vorlesungen über Geometrie. (Bd. I, S. 288 ff.)، دستگاه دو فرم هم‌درجه‌گرد. \emph{Gordan}, Über Büschel von Kegelschnitten. (Math. Ann. Bd. XIX. S. 530)، دستگاه دو فرم پاددرجه‌گرد.}، با ساخت دستگاه کامل نسبی به پیمانهٔ $(\varrho,t)$ پایان پیش‌گفته حاصل می‌شود. ترانس‌وکشنِ این دستگاه بر دستگاه $(\varrho,t)$ برای ساخت دستگاه کامل مطلق به بعد واگذار می‌شود.
به‌عنوان دستگاه کامل نسبی به پیمانهٔ $(\varrho,t)$، 331 ساخته به دست می‌آید که در جدول پیوست برحسب درجه‌شان نسبت به متغیرهای $x$ و $u$ مرتب شده‌اند.
در پایان، نمایی کلی از نمادهای معرفی‌شده می‌آوریم:
\begin{align*}
f&=a_x^4,\\
\theta&=\theta_x^4u_\vartheta^2=(abu)^2a_x^2b_x^2,\\
K&=K_x^6u_k^3=(a\theta u)u_\vartheta^2a_x^3\theta_x^3,\\
j&=u_j^6=(a\theta u)^4u_\vartheta^2,\\
\Delta&=\Delta_x^6=a_\vartheta^2a_x^2\theta_x^4=(abc)^2a_x^2b_x^2c_x^2,\\
N&=N_x^8u_\eta=(a\Delta u)a_x^3\Delta_x^5,\\
\nu&=u_\nu^4=(abu)^4,\\
H&=H_x^2u_\eta^4=(\nu\nu_1x)^2u_\nu^2u_{\nu_1}^2,\\
L&=L_x^3u_l^6=(\nu\eta x)H_x^2u_\nu^3u_\eta^3,\\
g&=g_x^6=(\nu\eta x)^4H_x^2,\\
\sigma&=u_\sigma^6=H_\nu^2u_\nu^2u_\eta^4,\\
s&=s_x^4=(\nu\nu_1x)^4,\\
Z&=Z_x^4u_\zeta^2=(ss'u)^2s_x^2s_x'^2,\\
\varrho&=u_\varrho^2=\theta_\nu^4u_\vartheta^2,\\
t&=t_x^2=a_\eta^4H_x^2,\\
i&=a_\nu^4,\\
J&=s_\nu^4.
\end{align*}
\begin{center}
{\Large\bfseries فصل اول. \quad قضایای کلی دربارهٔ فرم‌های سه‌تایی.}
\end{center}

\begin{center}
\textbf{§ 1.}\quad \emph{فرایند تاخوردگی. رشته‌های فرم‌ها.}
\end{center}

فرایند بنیادی برای ساخت فرم‌ها، \emph{فرایند تاخوردگی} است که در حوزهٔ سه‌تایی به صورت زیر تعریف می‌شود. یک حاصل‌ضرب نمادی داده شده باشد:
\[
s_x^m t_x^n u_\sigma^\mu u_\tau^\nu .
\]
اگر جفت‌عامل‌های
\[
s_xt_x\qquad u_\sigma u_\tau\qquad s_xu_\sigma\ \text{یا}\ s_xu_\tau
\qquad t_xu_\tau\ \text{یا}\ t_xu_\sigma
\]
به‌ترتیب با
\[
(stu)\qquad (\sigma\tau x)\qquad s_\sigma\ \text{یا}\ s_\tau
\qquad t_\tau\ \text{یا}\ t_\sigma,
\]
تاخوردگی
\[
\text{I}\qquad\text{II}\qquad\text{III}\qquad\text{IV},
\]
آنگاه فرم‌های حاصل با تاخوردگی از فرم اصلی پدید آمده‌اند\srcfn{*)}{\emph{Gordan}, Math. Annalen Bd. XVII S. 219.}.
عبارت $s_x^m t_x^n u_\sigma^\mu u_\tau^\nu$ را می‌توان یا حاصل‌ضرب واقعیِ دو فرم $S\cdot T$ دانست، یا یک فرم واحد با چند ردیف نماد:
\[
s_x^m t_x^n u_\sigma^\mu u_\tau^\nu=A_x^{m+n}u_\lambda^{\mu+\nu}.
\]
در حالت نخست از تاخوردگیِ فرم $S$ با $T$ سخن می‌گوییم و در حالت دوم از \glqq{}تاخوردگیِ فرم $A$ در خودش\grqq{}. برای اختصار، در ادامه همواره شرط زیر را فرض می‌کنیم ــ شرطی که در واقع برای همهٔ فرم‌هایی که بعدا بررسی خواهند شد برقرار است:
\srcnumdisplay{(a)}{\begin{gathered}
s_\sigma^\lambda=0,\\
t_\tau^\varkappa=0
\end{gathered}}
برای هر $\lambda$ و $\varkappa$ ناصفر و در ترکیب با هر تاخوردگیِ دلخواه دیگر. این شرط بیان می‌کند که فرم $s_x^m t_y^n u_\sigma^\mu u_\tau^\nu$ یک \glqq{}فرم نرمال نسبت به فرایند امگا\grqq{} است؛ یعنی با اعمال فرایند $\Omega$ هم نسبت به متغیرهای $x$ و $u$ و هم نسبت به متغیرهای $y$ و $v$ صفر می‌شود.
بنابراین شرط (a) هیچ محدودیتی ایجاد نمی‌کند، بلکه همواره می‌توان آن را برقرار کرد\srcfn{*)}{نک. \emph{Gordan}, Math. Annalen Bd. V S. 104.}.
دربارهٔ ارتباط تاخوردگی‌های منفرد، قضیهٔ زیر برقرار است:
\emph{قضیهٔ I. تاخوردگی‌های I و II تاخوردگی‌های بنیادی‌اند و تاخوردگی‌های III و IV را می‌توان، مستقل از ترتیب ترکیب، از آن‌ها ساخت. به بیان دیگر: برای ساختن همهٔ فرم‌هایی که با تاخوردگی از یک عبارت معین پدید می‌آیند، فقط باید تاخوردگی‌های I و II را به کار برد.}
\textbf{برهان:} با توجه به (a)، از قضیهٔ همانی‌ها، یا قضیهٔ حاصل‌ضرب ماتریس‌ها، نتیجه می‌شود:
\[
(\widehat{st}\sigma x)=-t_\sigma,\qquad
(\widehat{\sigma\tau}su)=-s_\tau,
\]
\[
(\widehat{st}\tau x)=s_\tau,\qquad
(\widehat{\sigma\tau}tu)=t_\sigma,
\]
\[
(stu)(\sigma\tau x)=s_\tau+t_\sigma-u_x\cdot s_\tau t_\sigma .
\]
بنابراین با ترکیب تاخوردگی‌های I و II، تاخوردگی‌های III و IV پدید می‌آیند؛ هم وقتی تاخوردگی II بر تاخوردگی I، یعنی بر عامل‌های $(stu)u_\sigma u_\tau$، اعمال شود و هم وقتی تاخوردگی I بر تاخوردگی II، یعنی بر عامل‌های $(\sigma\tau x)s_xt_x$، اعمال شود.
\emph{رشتهٔ فرم‌ها}\srcfn{**)}{نک. \emph{Clebsch}, Über eine Fundamentalaufgabe der Invariantentheorie (Abh. der Gött. Ges. d. Wiss. Bd. XVII): § 17 و § 18، پایان. \glqq{}رشتهٔ فرم‌ها\grqq{} با \glqq{}دستگاه هم‌ارزِ به‌معنای خاص کاهش‌یافته\grqq{} که آنجا معرفی شده است، در اصلِ ترتیب‌دادن برحسب فرم‌های بالاتر تفاوت دارد.}\setcounter{footnote}{9} را یک فرم آغازین همراه با همهٔ فرم‌هایی تعریف می‌کنیم که از تاخوردگیِ آن فرم در خودش پدید آمده‌اند، و رشتهٔ فرم‌ها را با فرم آغازین نشان می‌دهیم. در اینجا \glqq{}فرم بالاتر\grqq{} هر فرمی است که تاخوردگی‌های درونیِ بیشتری داشته باشد؛ در میان تاخوردگی‌های هم‌ارز $s_\tau$ و $t_\sigma$ نیز باید یکی را به‌طور قراردادی بالاتر برگزید. از قانون ساخت رشتهٔ فرم‌ها نتیجه می‌شود:
1) هر فرمِ رشتهٔ فرم‌ها نسبت به نمادهای فرم آغازین خطی است.
2) برای فرم‌های آغازینی که هرکدام دو ردیف نمادهای پادهموردا دارند، رشتهٔ فرم‌ها به‌طور یکتا تعیین می‌شود؛ برای فرم‌های آغازین با ردیف‌های نماد بیشتر، باید با ممتازکردن تاخوردگی‌های ویژه از میان تاخوردگی‌های مرتبط به‌وسیلهٔ قضیهٔ همانی‌ها، آن را به‌طور یکتا نرمال‌سازی کرد.
طبق قضیهٔ I، رشتهٔ فرم‌های $s_x^m t_x^n u_\sigma^\mu u_\tau^\nu$ ($m>n;\ \mu>\nu$) را می‌توان در طرح مستطیلیِ زیر مرتب کرد:
\begingroup\small
\[
\begin{array}{c|c|c|c|c}
s_x^m t_x^n u_\sigma^\mu u_\tau^\nu
& (stu)
& (stu)^2
& \cdots
& \noethpIIrosette\\[0.3em]
(\sigma\tau x)
& t_\sigma,\ s_\tau
& t_\sigma(stu),\ s_\tau(stu)
& \cdots
& \vdots\\[0.3em]
(\sigma\tau x)^2
& t_\sigma(\sigma\tau x),\ s_\tau(\sigma\tau x)
& t_\sigma^2,\ t_\sigma s_\tau,\ s_\tau^2
& \cdots
& \vdots\\[0.3em]
\vdots&\vdots&\vdots&\ddots&\vdots\\[0.3em]
(\sigma\tau x)^\nu
& t_\sigma(\sigma\tau x)^{\nu-1},\ s_\tau(\sigma\tau x)^{\nu-1}
& t_\sigma^2(\sigma\tau x)^{\nu-2},\ t_\sigma s_\tau(\sigma\tau x)^{\nu-2},\ s_\tau^2(\sigma\tau x)^{\nu-2}
& \cdots
& \vdots\\[0.3em]
\vdots
& t_\sigma(\sigma\tau x)^\nu
& t_\sigma^2(\sigma\tau x)^{\nu-1},\ t_\sigma s_\tau(\sigma\tau x)^{\nu-1}
& \cdots
& \vdots\\[0.3em]
\vdots
& \vdots
& t_\sigma^2(\sigma\tau x)^\nu
& \cdots
& \vdots
\end{array}
\]
\endgroup
\begingroup\small
\[
\begin{array}{r@{\quad}c}
\noethpIIrosette\srcfnmark{*)}
&
\begin{array}{c|c|c}
(stu)^n & \cdots & \cdots\\[0.3em]
t_\sigma(stu)^{n-1},\ s_\tau(stu)^{n-1}
& s_\tau(stu)^n & \cdots\\[0.3em]
t_\sigma^2(stu)^{n-2},\ t_\sigma s_\tau(stu)^{n-2},\ s_\tau^2(stu)^{n-2}
& t_\sigma s_\tau(stu)^{n-1},\ s_\tau^2(stu)^{n-1}
& s_\tau^2(stu)^n .
\end{array}
\end{array}
\]
\endgroup
\srcfntext{*)}{در اینجا، و به همین قیاس در ادامه، بخش پایینیِ نشان‌گذاری‌شده با \NoetherRosetteInline{} باید به جایگاه بالای سمت راستِ طرح که با همان نشان مشخص شده است متصل شود.}
در اینجا، اگر
1) یک ستون به سمت راست پیش برویم، همهٔ فرم‌هایی را به دست می‌آوریم که با تاخوردگی I از فرم‌های مجاور پدید می‌آیند،
2) یک ردیف به سمت پایین پیش برویم، همهٔ فرم‌هایی را به دست می‌آوریم که با تاخوردگی II از فرم‌های بالایی پدید می‌آیند،
و در نتیجه همهٔ فرم‌هایی را که با تاخوردگی پدید می‌آیند به دست می‌آوریم.
هر فرم دلخواه از کل رشتهٔ فرم‌ها، یعنی $t_\sigma^\kappa s_\tau^\lambda(stu)^\mu$ یا $t_\sigma^\kappa s_\tau^\lambda(\sigma\tau x)^\nu$، آغاز یک رشتهٔ فرم‌های تازه است. این رشته از همهٔ فرم‌های طرح مستطیلی تشکیل می‌شود که فرم $t_\sigma^\kappa s_\tau^\lambda(stu)^\mu$ یا $t_\sigma^\kappa s_\tau^\lambda(\sigma\tau x)^\nu$ رأس گوشهٔ بالای چپ آن‌هاست و عامل $t_\sigma^\kappa s_\tau^\lambda$ را دارند.
یادداشت پایانی. قضیهٔ I در حالتی استثنا دارد که $n$ و $\mu$ (یا به‌ترتیب $m$ و $\nu$) صفر شوند؛ در حالت نخست تاخوردگی‌های I، II و IV دیگر وجود ندارند و III باقی می‌ماند، و در حالت متناظر تاخوردگی‌های I، II و III دیگر وجود ندارند و IV باقی می‌ماند. در این حالت، عضو قطریِ طرح زیر را رشتهٔ فرم‌ها می‌نامیم:
\[
\begin{array}{c@{\qquad}c@{\qquad}c}
\begin{array}{ccccc}
s_x^m u_\tau^\nu &&&&\\
&s_\tau&&&\\
&&s_\tau^2&&\\
&&&\ddots&\\
&&&&s_\tau^\nu
\end{array}
&\text{یا، در حالت متناظر،}&
\begin{array}{ccccc}
t_x^n u_\sigma^\mu &&&&\\
&t_\sigma&&&\\
&&t_\sigma^2&&\\
&&&\ddots&\\
&&&&t_\sigma^\mu
\end{array}
\end{array}
\]
هماهنگ با نبود تاخوردگی‌های I و II، بسط رشتهٔ فرم‌ها برحسب قطبی‌های فرم در این حالت همواره به جملهٔ آغازین کاهش می‌یابد؛ بااین‌حال، قضایای مربوط به کاهشگرها همچنان برقرارند.
\begin{center}
\textbf{§ 2.}\quad \emph{بسط‌ها برحسب قطبی‌های رشتهٔ فرم‌ها}\srcfn{*)}{نگاه کنید به Study, Methoden zur Theorie der ternären Formen (Teubner 1889) II. § 7. برخلاف استنتاج آنجا، استنتاج ما روشی برای تعیین محاسباتیِ سریع ضرایب عددیِ $C_{pr\varrho}$ به دست می‌دهد. اگر بسط کلِبش--اشتودی را در حالت‌های a)، a') ویژه‌سازی کنیم، به صورت زیر درمی‌آید:
\[
(stu)^\mu u_\tau^{\nu-\varkappa}v_\tau^\varkappa s_x^m t_x^{\,n-\lambda}(tuv)^\lambda
=\sum_{p,r,\varrho} C'_{pr\varrho}
\bigl[s_\tau^\varrho(stu)^{\mu+r-\varrho}\bigr]_{\widehat{uv}^{\lambda-r+\varrho}v^{\varkappa-\varrho},\,w^p,\,(w=\widehat{xw})},
\]
و بنابراین ساده‌سازی‌ای را مجاز نمی‌کند که در بسط ما، وقتی می‌دانیم همهٔ جمله‌های دارای عامل $u_x$ باید حذف شوند، در همان جریان محاسبه پدید می‌آید.}.
\end{center}

برای فرم‌های دارای $n$ متغیر، دو نوع بسط به‌صراحت ساخته شده است\srcfn{**)}{نگاه کنید به Gordan, Über Kombinanten. § 2 und 5 (Math. Ann. Bd. V).}:
1) برای فرم‌هایی که هر یک یک ردیف از متغیرهای پادهموردا دارند: $s_x^m u_\tau^\nu$ رشته‌ای پیش‌رونده برحسب توان‌های $u_x$ است که ضرایبش \glqq{}فرم‌های نرمال\grqq{} هستند؛
2) برای فرم‌های دارای دو ردیف از متغیرهای هموردا، یعنی $s_x^m t_x^n$، بسطی برحسب قطبی‌های همورداهای بنیادی (قطبی‌های رشتهٔ فرم‌های $s_x^m t_x^n$) وجود دارد که در حالت سه‌تایی به صورت زیر است:
\[
\tag{I}
s_x^m t_x^{n-\lambda}t_y^\lambda
=\sum_{\varrho}
\frac{\binom{m}{\varrho}\binom{\lambda}{\varrho}}
 {\binom{m+n-\varrho+1}{\varrho}}\cdot
\bigl[(st\widehat{xy})^\varrho s_x^{m-\varrho}t_x^{n-\varrho}\bigr]_{y^{\lambda-\varrho}}.
\]

این دو بسط را با یکدیگر ترکیب می‌کنیم؛ بدین‌صورت که برای فرم‌هایی که هر یک دو ردیف متغیر پادهموردا دارند،
\[
s_x^m t_x^{n-\lambda}t_y^\lambda u_\sigma^\mu u_\tau^{\nu-\varkappa}v_\tau^\varkappa
\]
بسط‌هایی برحسب توان‌های $u_x$ می‌سازیم که ضرایبشان قطبی‌های مرکبِ رشتهٔ فرم‌های $s_x^m t_x^n u_\sigma^\mu u_\tau^\nu$ هستند و نسبت به متغیرهای $x$ و $u$ گرفته می‌شوند.
کافی است رشته را برای هر دو ردیف از نمادهای پادهموردا (یا، به‌ترتیب، متغیرهای پادهموردا) بسازیم؛ زیرا فرم‌های دارای ردیف‌های نماد (یا متغیر) بیشتر را می‌توان با یکی‌کردن هر دو ردیف نماد (یا متغیر) در یک ردیف جدید، به این حالت فروکاست.
برای آنکه همهٔ فرم‌های رشتهٔ فرم‌ها هر کدام فقط دو ردیف نماد یا متغیر داشته باشند، ویژه‌سازی‌های زیر را انجام می‌دهیم:
\[
\begin{array}{ll}
\text{a) } \sigma=\widehat{st}, & \text{a') } y=\widehat{uv},\\[0.3em]
\text{b) } s=\widehat{\sigma\tau}, & \text{b') } v=\widehat{xy},
\end{array}
\]
و بسط را برای ویژه‌سازی‌های a)، a') انجام می‌دهیم.
با انتقال مستقیم بسط I، برای فرم‌های $(stu)^\mu s_x^m t_x^{n-\lambda}t_y^\lambda$ قطبی‌های مرکب به دست می‌آیند؛ اما برای فرم‌های $(stu)^\mu u_\tau^{\nu-\varkappa}v_\tau^\varkappa s_x^m t_x^{n-\lambda}t_y^\lambda$ به‌جای قطبی‌های مرکب، جمله‌هایی پدید می‌آیند که ارزیابی آن‌ها مستلزم واردکردن پیوستهٔ متغیرهای کمکیِ تازه‌ای است که فقط در پایان با یکدیگر برابر گذاشته می‌شوند؛ و این امر محاسبه را بی‌دلیل پیچیده می‌کند.
ازاین‌رو راهی غیرمستقیم در پیش می‌گیریم و قطبی‌های مرکبِ همهٔ فرم‌های رشتهٔ فرم‌ها، یعنی عبارت‌های
\[
\bigl[s_\tau^\varrho(stu)^{\mu-\varrho+r}\bigr]_{y^\lambda}^{v^\varkappa}
\]
را برحسب جمله‌های آغازین این قطبی‌ها بیان می‌کنیم و با وارون‌سازی دستگاه معادلات بازگشتیِ حاصل، بسط مطلوب برحسب قطبی‌ها را به دست می‌آوریم.
بنا بر اصل انتقال، برای قطبیِ مرتبهٔ $\lambda$ از فرم $s_x^m t_x^n:[s_x^m t_x^n]_{y^\lambda}$ ($\lambda\le n$) بسط زیر برقرار است\srcfn{*)}{Gordan, Die Resultante binärer Formen (Kap. I, § 5). Rend. Circ. Mat. di Palermo XXII. 1906.} (حالت $\lambda>n$ با بسط $[s_y^m t_y^n]_{x^{m+n-\lambda}}$ به حالت نخست فروکاسته می‌شود):
\[
\bigl[s_x^m t_x^n\bigr]_{y^\lambda}
=\sum_r (-1)^r\,
\frac{\binom{m}{r}\binom{\lambda}{r}}{\binom{m+n}{r}}\cdot
(st\widehat{xy})^r s_x^{m-r}t_x^{n-\lambda}t_y^{\lambda-r},
\]
و به‌طور متناظر برای قطبیِ مرتبهٔ $\varkappa$ از $u_\sigma^\mu u_\tau^\nu:[u_\sigma^\mu u_\tau^\nu]^{v^\varkappa}$
\[
\bigl[u_\sigma^\mu u_\tau^\nu\bigr]^{v^\varkappa}
=\sum_\varrho (-1)^\varrho\,
\frac{\binom{\mu}{\varrho}\binom{\varkappa}{\varrho}}{\binom{\mu+\nu}{\varrho}}\,
(\sigma\tau\widehat{uv})^\varrho u_\sigma^{\mu-\varrho}u_\tau^{\nu-\varkappa}v_\tau^{\varkappa-\varrho}.
\]

با ضرب و با واردکردن ویژه‌سازی‌های a) و a')، به رابطهٔ زیر می‌رسیم:
\[
\begin{aligned}
\bigl[s_x^m t_x^n(stu)^\mu u_\tau^\nu\bigr]_{\widehat{uv}^{\lambda}}^{v^\varkappa}
={}&
\sum_{r,\varrho}c_{r\varrho}\cdot
(stx\widehat{uv})^r(\widehat{st}\tau\widehat{uv})^\varrho
s_x^{m-r}t_x^{n-\lambda}(tuv)^{\lambda-r}(stu)^{\mu-\varrho}
u_\tau^{\nu-\varkappa}v_\tau^{\varkappa-\varrho},
\end{aligned}
\]
و از آن، با بسط برحسب توان‌های $u_x$ و جداکردن جملهٔ نخست ($\sum'_{r,\varrho}$ یعنی جملهٔ $r=0$، $\varrho=0$ حذف می‌شود) و با درنظرگرفتن بند (a) در ص. 26، به رابطهٔ زیر می‌رسیم:
\[
\tag{II}
\begin{aligned}
&s_x^m t_x^{n-\lambda}(tuv)^\lambda(stu)^\mu u_\tau^{\nu-\varkappa}v_\tau^\varkappa\\
&\quad =
\bigl[s_x^m t_x^n(stu)^\mu u_\tau^\nu\bigr]_{\widehat{uv}^{\lambda}}^{v^\varkappa}\\
&\qquad
-\sum_{r,\varrho}' c_{r\varrho}\cdot s_\tau^\varrho(stu)^{\mu-\varrho+r}
u_\tau^{\nu-\varkappa}v_\tau^{\varkappa-\varrho}
s_x^{m-r}t_x^{n-\lambda}(tuv)^{\lambda-r+\varrho}\cdot v_x^r\\
&\qquad
+u_x\cdot\sum_{\varrho}\cdot\sum_{r=1}^{\lambda} c_{r\varrho}\cdot
s_\tau^\varrho(stu)^{\mu-\varrho+r-1}(stv)
u_\tau^{\nu-\varkappa}v_\tau^{\varkappa-\varrho}
s_x^{m-r}t_x^{n-\lambda}(tuv)^{\lambda-r+\varrho}\cdot v_x^{r-1}\\
&\qquad
\vdots\\
&\qquad
-(-1)^\lambda u_x^\lambda\cdot\sum_\varrho c_{\lambda\varrho}\cdot
s_\tau^\varrho(stu)^{\mu-\varrho}(stv)^\lambda
u_\tau^{\nu-\varkappa}v_\tau^{\varkappa-\varrho}
s_x^{m-\lambda}t_x^{n-\lambda}(tuv)^\varrho,
\end{aligned}
\]
که در آن
\[
c_{r\varrho}=(-1)^{r+\varrho}\cdot
\frac{\binom{m}{r}\binom{\lambda}{r}}{\binom{m+n}{r}}\cdot
\frac{\binom{\mu}{\varrho}\binom{\varkappa}{\varrho}}{\binom{\mu+\nu}{\varrho}}
\qquad
\begin{array}{l}
\text{برای }\lambda\le n,\\
\varkappa\le\nu.
\end{array}
\]
و به‌طور متناظر برای دیگر آرایش‌های اعداد
\[\lambda,n;\varkappa,\nu.\]
و عبارت
\[s_x^m t_x^{n-\lambda}(tuv)^\lambda(stu)^\mu u_\tau^{\nu-\varkappa}v_\tau^\varkappa\]
بدین‌ترتیب به یک قطبی مرکب فروکاسته شده است و با به‌کاربردن همانیِ
\[
(stv)u_\tau=(stu)v_\tau-s_\tau(tuv),
\]
به مجموعی از عبارت‌های هم‌ساخت می‌انجامد که متناظر با فرم‌های بالاترِ رشتهٔ فرم‌ها هستند.
با تکرار این فرایند به بسط زیر می‌رسیم:
\[
\tag{III}
s_x^m t_x^{n-\lambda}(tuv)^\lambda(stu)^\mu u_\tau^{\nu-\varkappa}v_\tau^\varkappa
=
\sum_{p,r,\varrho}u_x^p\cdot C_{pr\varrho}\cdot
\bigl[s_\tau^\varrho(stu)^{\mu-\varrho+r}\bigr]_{\widehat{uv}^{\lambda-r+\varrho}v^{\varkappa-\varrho+p},\,v_x^{r-p}}.
\]
ضرایب عددیِ $C_{pr\varrho}$ به‌طور یکتا و بازگشتی از ضرایب $c_{r\varrho}$ و $c'_{r\varrho}$ دستگاه معادلاتی محاسبه می‌شوند که با تکرار (II.) پدید می‌آید (نگاه کنید به مثال ص. 42). برای دیگر ویژه‌سازی‌ها نیز بسط‌های مشابهی به دست می‌آیند.
\begin{center}
\textbf{§ 3.}\quad \emph{قضیه‌های کاهش.}
\end{center}

اکنون باید قضیه‌های شناخته‌شده دربارهٔ کاهش فرم‌ها و دستگاه‌های فرم‌ها را با بندهای 1 و 2 مرتبط کنیم؛ برای این کار، چند تعریف برگرفته از نظریهٔ فرم‌های دوتایی لازم است.
الف) \emph{دستگاه کامل نسبی} به پیمانهٔ یک رشتهٔ معین از فرم‌ها، دستگاهی از فرم‌هاست با این ویژگی که همهٔ ساخته‌های حاصل از تاخوردگیِ حاصل‌ضربی دلخواه از این فرم‌ها را می‌توان به صورت توابع گویای صحیحِ فرم‌های دستگاه بیان کرد، به‌جز یک جملهٔ جمعی از عبارت‌هایی که از تاخوردگیِ فرم‌های دستگاه با دستگاه پیمانه پدید آمده‌اند\srcfn{*)}{نگاه کنید به Gordan-Kerschensteiner, Vorlesungen über Invariantentheorie. Bd. II. S. 227.}.
ب) فرمی را \emph{کاهش‌پذیر} می‌نامیم هرگاه بتوان آن را با فرم‌هایی بیان کرد که ناورداها را به‌عنوان عامل دارند، یا با \glqq{}فرم‌های بالاتر\grqq{}؛ یعنی با فرم‌های بالاترِ کلِ رشتهٔ فرم‌هایی که این فرم به آن تعلق دارد، یا با فرم‌هایی که نمادهای پیمانه را در مرتبه‌ای بالاتر در بر دارند.
اکنون می‌توان قضیه‌های مربوط به کاهشگرها را در کلی‌ترین صورتشان چنین بیان کرد\srcfn{**)}{نگاه کنید به Gordan, Math. Annalen Bd. XVII. S. 222.}:
\textbf{تعریف:} \emph{کاهشگر، رشته‌ای از فرم‌هاست که کاهش‌پذیر است.}
\textbf{قضیهٔ II.} \emph{اگر فرم آغازینِ یک رشتهٔ فرم‌ها از آن رو کاهش‌پذیر باشد که یکی از جمله‌هایش با تاخوردگی با یک کاهشگر پدید آمده است (کاهشگری را به‌عنوان عامل دارد)، و اگر فرم پایانیِ رشتهٔ فرم‌ها نیز با تاخوردگیِ همان جمله پدید آمده باشد، آنگاه کل رشتهٔ فرم‌ها کاهش‌پذیر است}\srcfn{*)}{ساده‌سازی حاصل از قضیهٔ II را می‌توان با مثال زیر نشان داد: برای فرم ویژهٔ $f=x_1^3x_2+x_2^3x_3+x_3^3x_1$، عبارت $a_y^2a_x^2u_y^2$ کاهشگر است. از این رو، طبق قضیهٔ II، رشتهٔ فرم‌های
\[
\theta_\nu(\vartheta\nu x)\theta_x^3u_\vartheta u_\nu^2
=
a_\nu^2(abu)-\underbrace{a_\nu b_\nu(aba)}_{0},
\]
کاهش می‌یابد، زیرا $\theta_\nu^4=a_\nu^2b_\nu^2(abu)^2$ با تاخوردگی از جملهٔ $a_\nu^2(abu)$ پدید آمده است. اما در Gordan, a. a. O. S. 231، کاهش فرم‌های
\[
\theta_\nu(\vartheta\nu x),\quad \theta_\nu(\vartheta\nu x)^2,\quad
\theta_\nu^2,\quad \theta_\nu^2(\vartheta\nu x),\quad
\theta_\nu^2(\vartheta\nu x)^2
\]
جداگانه انجام می‌شود.}.
\textbf{برهان:} رشتهٔ فرم‌ها از جملهٔ آغازین با تاخوردگیِ آن در خودش پدید می‌آید؛ یعنی از یک سو با تاخوردگیِ بالاتر با کاهشگر و از سوی دیگر با تاخوردگیِ کاهشگر در خودش. در هر دو حالت، طبق بند 2، همهٔ فرم‌های رشته را می‌توان به صورت قطبی‌ها (ترانس‌وکشن‌ها) بر روی رشتهٔ فرم‌های کاهشگر نمایش داد.
در دیگر روش‌های کاهش، دیگر نمی‌توان از فرم آغازین به کل رشتهٔ فرم‌ها نتیجه گرفت. این روش‌ها عبارت‌اند از:
1) به‌اصطلاح \emph{\glqq{}کاهش دوگانه\grqq{}}.
منظور، کاهشِ دوگانهٔ عبارتی است که دو کاهشگرِ مستقل از یکدیگر را به‌عنوان عامل دارد؛ این کاهش از دو راه انجام می‌شود و بدین‌ترتیب میان فرم‌های بالاتری که عبارت به آنها کاهش یافته است رابطه‌هایی پدید می‌آید.
با به‌کارگیریِ نظام‌مندِ \glqq{}کاهش دوگانه\grqq{} باید، از یک سو، همهٔ رابطه‌ها به دست آیند\srcfn{**)}{برای نمونه، با \glqq{}کاهش دوگانه\grqq{} این موارد یافت شدند: کاهش فرم $a_\eta^3$ (بند 10)، تنها فرمی که Maisano بیهوده در دستگاه گنجانده بود؛ و سپس رابطه‌های یادشده در یادداشت مقدمه میان سه هموردا و سه ناوردا.}؛ از سوی دیگر، اگر کاربرد \glqq{}کاهش دوگانه\grqq{} بر عبارت‌های مختلف هیچ رابطهٔ تازه‌ای ندهد، هم معیاری برای استقلال فرم‌های بالاتر داریم و هم کنترلی بر محاسبه.
2) \emph{تاخوردگی با فرم‌های تجزیه‌پذیر.}
منظور از \glqq{}فرم‌های تجزیه‌پذیر\grqq{} --- با تعمیمی اندک از مفهوم متعارف --- فرم‌هایی است که بتوان آنها را با حاصل‌ضربِ فرم‌های درجهٔ پایین‌تر و با \glqq{}فرم‌های بالاتر\grqq{} بیان کرد. حاصل‌ضربِ فرم‌ها پس از یک تاخوردگی به حاصل‌ضرب تبدیل می‌شود؛ همچنین حاصل‌ضربِ فرم‌های \emph{غیرخطی} پس از دو تاخوردگی در ویژه‌سازی‌های a') و b') (بند 2)، بنا بر قضیهٔ حاصل‌ضرب:
\[
a_yb_y a_xb_x=\frac12a_y^2\cdot b_x^2+\frac12b_y^2\cdot a_x^2-\frac12(ab\widehat{yx})^2.
\]
بنابراین قطبی‌های نخست (ترانس‌وکشن‌ها) و برخی قطبی‌های دومِ فرم‌های تجزیه‌پذیر، خود نیز فرم‌های تجزیه‌پذیرند. در نتیجه:
\emph{جمله‌ای} که با یک تاخوردگی (یا به‌ترتیب دو تاخوردگی) با یک فرم تجزیه‌پذیر، در یک ترانس‌وکشنِ یک‌باره (یا به‌ترتیب دوگانه) بر روی یک فرم تجزیه‌پذیر پدید آمده است، خود نیز فرمی تجزیه‌پذیر می‌شود:
الف) اگر فرم‌های بالاترِ بعدی در رشتهٔ فرمِ تجزیه‌پذیر، تجزیه‌پذیر یا کاهش‌پذیر شوند؛ یا نیز اگر در حالت یک تاخوردگی، یکی از ویژه‌سازی‌های $y=\widehat{uv}$ یا $v=\widehat{xy}$ رخ دهد؛
ب) اگر دیگر تاخوردگی‌های یک‌باره به فرم‌های بالاتر بینجامند (نگاه کنید به بند 12، B.، $H_k$ و $H_k(k\eta x)$).
چنین جمله‌ای از یک ترانس‌وکشن ممکن است از راه زیر نیز تجزیه‌پذیر شود:
ج) اگر دیگر تاخوردگی‌های یک‌باره به فرم‌های پایین‌تری بینجامند که مستقل از تاخوردگی با فرم تجزیه‌پذیر، تجزیه‌پذیر می‌شوند؛ یعنی بتوان آنها را با حاصل‌ضرب‌ها و با فرمی که باید کاهش یابد بیان کرد؛ با این همه، هر بار باید نخست با محاسبه اطمینان یافت که ضریب عددیِ جملهٔ مورد نظر به‌طور همانی صفر نمی‌شود. این چیزی جز \glqq{}کاهش دوگانه\grqq{} فرم پایین‌تر نیست (نگاه کنید به بند 23، C.، $H_q(\theta H u)(\theta s u)$).
از همهٔ تاخوردگی‌های یک‌باره با دترمینان‌های ژاکوبین\srcfn{*)}{نگاه کنید به Gordan, a. a. O. § 3.}، و نیز از برخی تاخوردگی‌های بالاتر، فرم‌های تجزیه‌پذیر پدید می‌آیند. داریم:
\[
(st\widehat{yx})s_y=\frac12\,t\cdot s_y^2-\frac12\,s\cdot t_y^2+\frac12(st\widehat{yx})^2,
\]
\[
(st\widehat{yx})^2s_y^2=\frac13\,t\cdot s_y^3-\frac13\,s\cdot t_y^3+s_y(st\widehat{yx})^2-\frac13(st\widehat{yx})^3.
\]
برای فرم‌های دوگان نیز فرمول‌های مشابهی برقرار است
\[
(\sigma\tau\widehat{vu})v_\sigma
\quad\text{و}\quad
(\sigma\tau\widehat{vu})v_\sigma^2.
\]
بنا بر قضیه‌های این بند، برای ساختن همهٔ فرم‌های کاهش‌ناپذیری که از تاخوردگیِ $S$ با $T$ پدید می‌آیند، قاعدهٔ زیر به دست می‌آید:
از پایین‌ترین تاخوردگی‌ها آغاز کنید و در ساخت رشتهٔ فرم‌های $ST$، در مسیری دلخواه ولی از پیش تثبیت‌شده (مثلاً سطر‌به‌سطر یا متقارن نسبت به جملهٔ قطری)، آن‌قدر پیش بروید تا به فرمی برسید که کاهشگری را به‌عنوان عامل دارد. رشتهٔ فرم‌هایی که این فرم تعریف می‌کند، به‌محض آنکه فرم پایانی شرط قضیهٔ II را برآورده کند، باید کنار گذاشته شود. پس از حذف همهٔ رشته‌های فرمِ کاهش‌پذیر، با به‌کارگیری کاهش دوگانه، همهٔ فرم‌های کاهش‌پذیرِ باقی‌مانده و نیز همهٔ فرم‌های تجزیه‌پذیر حذف شوند. جایگاه‌هایی از کل رشتهٔ فرم‌ها که دو رشتهٔ فرمِ کاهش‌پذیر و مستقل از یکدیگر آنها را دو بار می‌پوشانند، موجب کاهش فرم‌های بالاتر می‌شوند (نگاه کنید به بند 11، D).
\begin{center}
\textbf{فصل دوم. \quad دستگاه‌های بنیادیِ کامل نسبی.}
\textbf{§ 4.}\quad\emph{نخستین دستگاه کامل نسبی.}
\end{center}

در ادامه، به‌استثنای قضیهٔ نخست، خود را به فرم سه‌تاییِ درجهٔ چهار محدود می‌کنیم. با این همه، با انتخاب مناسب پیمانه‌ها، می‌توان نتایج بندهای 5 و 6 را به‌آسانی به فرم‌های درجهٔ بالاتر انتقال داد؛ همان‌گونه که این نتایج برای فرم سه‌تاییِ درجهٔ سه نیز برقرارند\srcfn{*)}{برای فرم سه‌تاییِ درجهٔ سه، فرمول‌های متناظر که به‌طور مشابه به دست می‌آیند، چنین‌اند:
\[
(abc)a_x^2b_y^2c_z^2=\frac13(sxy)(sxz)(syz)+\frac23\Delta_x\Delta_y\Delta_z\cdot(xyz),\qquad (\mathrm{IV}.)
\]
\[
\begin{gathered}
s=(abc)(abu)(acu)(bcu),\qquad \Delta=(abc)^2a_xb_xc_x,\\
t=\Delta_\vartheta(\Delta\theta u)^2u_\vartheta=\theta_s^2u_su_\vartheta^2,
\qquad T=a_t^3,
\end{gathered}
\]
\[
\left.
\begin{aligned}
(a\theta u)^2&=\frac13 u_xs,\\
a_\vartheta&=\frac13u_x\cdot\Delta,
& a_\vartheta(a\theta u)&=0,
& a_\vartheta(a\theta u)^2&=s,\\
a_\vartheta^2&=\Delta,
& a_\vartheta^2(a\theta u)&=0.
\end{aligned}
\right\}\tag{1.}
\]
}.
بنا بر قضیهٔ شناخته‌شده، نخستین دستگاه کامل نسبی برای فرم‌هایی با درجهٔ دلخواه به دست می‌آید:
\clearpage
\noindent\textbf{قضیهٔ III.}
\emph{فرم‌های برگرفته از حوزهٔ دوتایی، یعنی فرم‌هایی که از فرم‌های دستگاهِ فرم دوتاییِ متناظر، بنا بر اصل انتقالِ کلِبش و با افزودن سطر و ستون به دترمینان‌ها پدید می‌آیند، دستگاهی کامل نسبی به پیمانهٔ $(abc)$ تشکیل می‌دهند.}
\noindent\textbf{برهان:}
یک رابطه در حوزهٔ فرم‌های دوتایی بیان می‌کند که فرم‌های $Q'_1,\ldots,Q'_n$ دستگاهی کامل برای یک فرم بنیادیِ دوتایی تشکیل می‌دهند:
\[
Q'-F(Q')=0=\sum_\lambda\bigl\{(ab)c_x+(bc)a_x+(ca)b_x\bigr\}\varphi_\lambda\srcfnmark{*)}
\]
\srcfntext{*)}{Gordan--Kerschensteiner, a. a. O. S. 134.}
با افزودن سطر و ستون به دترمینان‌ها، رابطهٔ زیر در حوزهٔ فرم‌های سه‌تایی با این رابطه متناظر می‌شود:
\[
Q-F(Q_i)=\sum_\lambda u_x\cdot(abc)\varphi_\lambda .
\]
اما این همان معادلهٔ تعریف‌کنندهٔ یک دستگاه کامل نسبی به پیمانهٔ $(abc)$ است. برای فرم سه‌تاییِ درجهٔ چهار، دستگاه فرم‌های برگرفته از فرم‌های زیر تشکیل می‌شود:
\[
\begin{gathered}
f=a_x^4,\qquad
\theta=\theta_x^4u_\vartheta^2=(abu)^2a_x^2b_x^2,
\qquad
K=K_x^6u_k^3=(a\theta u)u_\vartheta^2a_x^3\theta_x^3,\\
j=u_j^6=(a\theta u)^4u_\vartheta^2,\qquad
\nu=u_\nu^4=(abu)^4,
\end{gathered}
\]
از میان آنها، چهار فرم نخست دستگاهی کامل نسبی به پیمانهٔ $((abc),\nu)$ تشکیل می‌دهند.
\begin{center}
\textbf{§ 5.}\quad\emph{کاهش پیمانهٔ $(abc)$ به پیمانه‌های $\Delta$ و$\nu$.}
\end{center}

پیمانهٔ $(abc)$ که فقط با یک عامل درون پرانتز داده شده است، مطابق تعریف (§ 3a) باید به فرم‌های دستگاه کاهش یابد. به بیان دیگر: رشتهٔ فرم‌های $(abc)$ باید به‌طور یکتا نرمال‌سازی شود (نگاه کنید به § 1).
\[
a_s^2a_xu_s=\frac13a_s^3\cdot u_x=\frac13S\cdot u_x. \tag{2.}
\]
\[
\left.
\begin{aligned}
(a\Delta u)^2&=a_s-\frac{S}{6}u_x^2,\\
(a\Delta u)^3&=0.
\end{aligned}
\right\} \tag{3.}
\]
\[
\left.
\begin{aligned}
\theta(s)=(ss_1x)^2&=2a_t+\frac13S\cdot\theta-\frac13Tu_x^2,\\
\theta(t)=(tt_1x)^2&=-\frac13S\cdot a_t+\frac23T\cdot a_s+\frac1{18}S^2\cdot\theta-\frac1{18}u_x^2\cdot S\cdot T.
\end{aligned}
\right\} \tag{4.}
\]

رشتهٔ پیمانه‌ها: $(abc)$؛ $\Delta$؛ $s$؛ $t$ (دستگاه پیمانهٔ $t$ تنها از فرم $t$ تشکیل می‌شود).
نشان خواهیم داد که فرم‌ها:
\[
\Delta=(abc)^2a_x^2b_x^2c_x^2,\qquad \nu=(abu)^4
\]
به‌عنوان پیمانه کافی‌اند؛ یعنی فرم‌ها
\[
\begin{gathered}
\Delta=(abc)^2a_x^2b_x^2c_x^2,
\qquad a_\nu=(abc)(bcu)^3a_x^3,\\
a_\nu^2=(abc)^2(bcu)^2a_x^2,
\qquad i=a_\nu^4=(abc)^4
\end{gathered}
\]
رشتهٔ فرم‌های $(abc)$ را تعریف می‌کنند. فرم $a_\nu^3$ در رشتهٔ فرم‌های $a_\nu$ نیست، بنا بر همانی:
\[
a_\nu^3=(abc)^3a_x(bcu)=\frac13u_x\cdot(abc)^4=\frac13i\cdot u_x.
\]

برای کاهش یک پیمانه به پیمانه‌ای بالاتر، بنا بر تعریف باید تاخوردگی‌های با آن پیمانه ــ کلی‌ترینِ آن‌ها، یا همهٔ تاخوردگی‌های ویژهٔ مورد نظر ــ را به تاخوردگی‌های با پیمانهٔ بالاتر کاهش دهیم.
در حالت ما، کلی‌ترین تاخوردگی‌ها از عبارت‌های زیر پدید می‌آیند:
\[
(abc)a_x^3b_y^3c_z^3,
\qquad
(abc)a_x^2b_y^2c_z^2a_xb_xc_x
\]
(برگرفته از فرم‌های درجهٔ سه)،
\[
(abc)a_xb_yc_za_x^2b_x^2c_x^2
\]
(برگرفته از فرم‌های درجهٔ دو) و نیز با گرفتن قطبی‌های این عبارت‌ها.
برای محاسبهٔ این عبارت‌ها از طریق قطبی‌های $\Delta$ و رشتهٔ فرم‌های $a_\nu$، یا به‌عبارت دقیق‌تر جمله‌های آغازین آن، همانند § ۲ راه غیرمستقیم را در پیش می‌گیریم. جمله‌های قطبی را

\[
(abc)^2a_x^2b_y^2c_z^2,
\qquad
 a_\nu(\nu xy)(\nu yz)(\nu zx)a_xa_ya_z=(abc)(bc\widehat{xy})(bc\widehat{yz})(bc\widehat{zx})a_xa_ya_z\quad\text{و غیره}
\]
به‌صورت ترکیب‌های خطی از عبارت‌هایی با عامل نمادی $(abc)$ بسط می‌دهیم و با وارون‌سازی دستگاه معادلات، روابط مطلوب را به‌دست می‌آوریم، و نیز رابطه‌ای میان قطبی‌هایی که از ترکیب‌های گوناگون متغیرها گرفته شده‌اند. برای ارزیابی، قضیهٔ همانی‌ها و قضیهٔ حاصل‌ضرب دترمینان‌ها به‌کار می‌روند.

دستگاه معادلات برای $(abc)a_x^3b_y^3c_z^3$ به‌صورت زیر است:
\begingroup\small
\[
\begin{aligned}
1)\quad
\sum_3a_\nu(\nu yz)^3a_x^3
&=6(abc)a_x^3b_y^3c_z^3-6\sum_3(abc)a_x^3b_y^2b_zc_z^2c_y\srcfnmark{*)},\\[0.4em]
2)\quad
3(abc)^2a_x^2b_y^2c_z^2\cdot(xyz)-\frac12\sum_3a_\nu^2(\nu yz)^2a_x^2\cdot(xyz)
&=2\sum_3(abc)a_x^3b_y^2b_zc_z^2c_y\\
&\quad-4\sum_3(abc)a_xa_ya_zb_y^2b_xc_z^2c_x,\\[0.4em]
3)\quad
a_\nu(\nu xy)(\nu yz)(\nu zx)a_xa_ya_z
&=2\sum_3(abc)a_xa_ya_zb_y^2b_xc_z^2c_x,\\[0.4em]
4)\quad
\sum_3a_\nu(\nu yz)^3a_x^3-\frac32\sum_3a_\nu^2(\nu yz)^2a_x^2\cdot(xyz)+\frac16i\cdot(xyz)^3
&=6\sum_3(abc)a_xa_ya_zb_y^2b_xc_z^2c_x.
\end{aligned}
\]
\endgroup
\srcfntext{*)}{$\sum_3$ به مجموعی بر سه جمله اشاره دارد که از جملهٔ آغازین با جایگشت دوریِ متغیرها پدید می‌آید. معادلات برای هر جملهٔ این مجموع نیز به‌تنهایی برقرارند.}

از این، برای $(abc)a_x^3b_y^3c_z^3$ و با محاسبه‌ای مشابه برای $(abc)a_x^2b_y^2c_z^2a_xb_xc_x$ و$(abc)a_xb_yc_za_x^2b_x^2c_x^2$ نتیجه می‌شود:
\[
\begin{aligned}
(abc)a_x^3b_y^3c_z^3
&=\frac32(abc)^2a_x^2b_y^2c_z^2\cdot(xyz)
+\frac32a_\nu(\nu xy)(\nu yz)(\nu zx)a_xa_ya_z
-\frac1{12}i\cdot(xyz)^3,\\[0.4em]
(\mathrm{IV}.)\quad
(abc)a_x^2b_y^2c_z^2a_xb_xc_x
&=\frac23(abc)^2a_x^2b_yb_xc_zc_x\cdot(xyz)
+\frac13a_\nu(\nu xy)(\nu yz)(\nu zx)a_x^3\\
&\quad-\frac16a_\nu^2(\nu xy)(\nu zx)a_x^2\cdot(xyz),\\[0.4em]
(abc)a_xb_yc_za_x^2b_x^2c_x^2
&=\frac16(abc)^2a_x^2b_x^2c_x^2\cdot(xyz).
\end{aligned}
\]

بدین‌سان دستگاهی کامل نسبی به پیمانهٔ $(\Delta,\nu)$ به دست آورده‌ایم که از چهار فرم $f,\theta,K,j$ تشکیل شده است.

از این نتیجه می‌شود: همهٔ تاخوردگی‌های $f$ بر $\theta$ که از حوزهٔ دوتایی گرفته نشده‌اند، و نیز تاخوردگیِ $(a\theta u)^2$ که در حالت دوتایی به $(ab)^4$ کاهش‌پذیر است، با نمادهای $\Delta$ و$\nu$ بیان‌پذیرند.

فرمول‌های کاهشِ بنیادی برای مطالب بعد را با تخصصی‌کردن بسط IV، یا کوتاه‌تر با محاسبهٔ مستقیم (جابه‌جایی نمادها)، برای $a_\vartheta(a\theta u)^2$ و$a_\vartheta^2(a\theta u)^2$ و$(a\theta u)^2$ بر پایهٔ صورت زیر به دست می‌آوریم:
\[
\begin{aligned}
a_\vartheta(a\theta u)^2
&=(abc)(bcu)(abu)^2-\frac13(abc)(bcu)\{(bcu)a_x-u_x\cdot(abc)\}^2,\\
a_\vartheta^2(a\theta u)^2
&=(abc)^2(abu)^2-\frac13(abc)^2\{(bcu)a_x-u_x\cdot(abc)\}^2.
\end{aligned}
\]
برای $(a\theta u)^2$\srcfn{*)}{Gordan--Kerschensteiner، همان‌جا، ص. 92.}:
\[
\begin{aligned}
(abu)a_y^3b_x^3&=\frac32u_\vartheta(\vartheta yx)\theta_y^2+\frac14(\nu yx)^3,\\
(abu)(acu)^3b_x^3&=\frac32u_\vartheta(\vartheta\widehat{au}x)(\theta au)^2+\frac14(\nu\widehat{au}x)^3.
\end{aligned}
\]
\begingroup\small
\[
\begin{array}{r@{\quad}|l@{\qquad}c|}
& \underline{(a\theta u)^2}=\frac16\nu\cdot f-\frac23u_x\cdot a_\nu+\frac23u_x^2\cdot a_\nu^2-\frac{i}{18}u_x^4
& \bigstar\\[0.55em]
a_\vartheta=\frac13u_x\cdot\Delta
& \underline{a_\vartheta(a\theta u)}=0
& \\[0.55em]
& \underline{a_\vartheta^2}=\Delta
&
\end{array}
\]
\vspace{-0.5em}
\[
\begin{array}{r@{\quad}c@{\quad}|l@{\qquad}|l}
(1.) & \bigstar
& \underline{(a\theta u)^3}=0
& \underline{(a\theta u)^4}=j\\[0.55em]
&
& \underline{a_\vartheta(a\theta u)^2}=-\frac56a_\nu+\frac76u_x\cdot a_\nu^2-\frac{i}{9}u_x^3
& \underline{a_\vartheta(a\theta u)^3}=0\\[0.55em]
&
& \underline{a_\vartheta^2(a\theta u)}=0
& \underline{a_\vartheta^2(a\theta u)^2}=\frac23a_\nu^2-\frac{i}{9}u_x^2
\end{array}
\]
\endgroup

در کنار آن، فرمول زیر را نیز به‌کار می‌گیریم که آن هم از کاهش پیمانهٔ $(abc)$ پدید آمده است:
\[
\begin{array}{r@{\qquad}l}
(2.) & a_\nu^3a_xu_\nu=\frac13i\cdot u_x.
\end{array}
\]

\emph{از فرمول‌های (1.) و(2.) و فرمول‌هایی که به‌طور مشابه برای فرم بنیادیِ $\nu$ ساخته شده‌اند، همهٔ فرمول‌های کاهشِ بعدی با گرفتن قطبی‌ها به دست می‌آیند}، به‌جز روابط مربوط به فرم‌های تجزیه‌پذیر. از فرمول‌های (1.) می‌بینیم:

رشتهٔ فرم‌ها به‌صورت زیر می‌انجامد:
\[
\begin{array}{rcl}
a_\vartheta(a\theta u)^2 &\text{به}& \text{نماد }\nu\text{ به‌تنهایی},\\[0.2em]
a_\vartheta &\text{به}& \text{نمادهای }\Delta\text{ و }\nu,\\[0.2em]
(a\theta u)^2 &\text{به}& \text{نمادهای }j,\Delta\text{ و }\nu,\\[0.2em]
(a\theta u)^2\text{ در تاخوردگی I} &\text{به}& \text{نمادهای }j\text{ و }\nu,\\[0.2em]
(a\theta u)^2\text{ در تاخوردگی II} &\text{به}& \text{نمادهای }\Delta\text{ و }\nu.
\end{array}
\]

پس می‌توانیم جایگزین کنیم:
\begingroup
\renewcommand{\labelenumi}{\arabic{enumi})}
\begin{enumerate}
\item به پیمانهٔ $(\Delta,\nu)$: تاخوردگی‌های با $j$ را با تاخوردگی‌های متناظر با $(a\theta u)^2$ یا $(a\theta u)^3$ جایگزین کنیم؛
\item به پیمانهٔ $(\nu)$: تاخوردگی‌های با $\Delta$ را با تاخوردگی‌های متناظر با $a_\vartheta$ یا $a_\vartheta(a\theta u)$، یا نیز با تاخوردگی‌های ویژهٔ با $(a\theta u)^2$ جایگزین کنیم (که فقط به یک بار تاخوردگی I می‌انجامند).
\end{enumerate}
\endgroup
به‌ویژه به دست می‌آید:
\[
\begin{aligned}
(a\theta u)^2a_y^2\theta_x^2&=\frac13(jyx)^2\modu{\nu},&
(a\theta u)^2a_y\theta_y&=-\frac16(jyx)^2\modu{\nu},\\[0.4em]
(a\theta u)^2v_\vartheta^2&=\frac13(\Delta uv)^2\modu{\nu},&
(a\theta u)(a\theta v)u_\vartheta v_\vartheta&=-\frac16(\Delta uv)^2\modu{\nu},\\[0.4em]
(a\theta u)^2a_yv_\vartheta^2&=\frac13(\Delta uv)^2\Delta_y\modu{\nu}.
\end{aligned}
\]

\begin{center}
\textbf{§ ۶.}\quad\emph{کاهش پیمانهٔ $(\Delta,\nu)$ به پیمانهٔ $(\nu)$.}
\end{center}

فرمول‌های کاهشِ بند پیشین ابزار ساخت مستقیم دستگاه کامل نسبی به پیمانهٔ $\nu$ را در اختیار ما می‌گذارند؛ خواهیم دید که به دستگاه فرم‌های برگرفته‌شده فقط باید فرم‌های $\Delta$ و$(a\Delta u)=N$ را بیفزاییم (همانند فرم درجهٔ سه و در واقع به‌طور کلی معتبر).

برای کاهش پیمانهٔ $\Delta$، همهٔ تاخوردگی‌های ویژهٔ مورد نظر، یعنی تاخوردگی‌های دستگاه کامل نسبی به پیمانهٔ $(\Delta,\nu)$ با دستگاه $\Delta$ را بررسی می‌کنیم و از فرم‌های با کمترین درجه در ضرایب، یعنی تاخوردگی‌های $f$ با $\Delta$ آغاز می‌کنیم.

برای کاهش فرم‌های $(a\Delta u)^2$ و$(a\Delta u)^3$ و$(a\Delta u)^4$، بنا بر § ۵ نمادهای $\Delta$ را با فرم‌های پایین‌ترِ رشتهٔ فرم‌ها جایگزین می‌کنیم؛ یعنی عبارت‌های زیر را بسط می‌دهیم
\[
a_\vartheta b_\vartheta(b\theta u)^2,
\qquad a_\vartheta(a\theta u)b_\vartheta(b\theta u)^2,
\qquad a_\vartheta(a\theta u)b_\vartheta(b\theta u)^3
\]
۱) برحسب قطبی‌های رشتهٔ فرم‌های $a_\vartheta$ (نمادهای $\Delta$ و$\nu$)،

۲) برحسب قطبی‌های رشتهٔ فرم‌های $b_\vartheta(b\theta u)^2$ (نماد $\nu$)

و فرمول‌های کاهش را به دست می‌آوریم (محاسبه را در پایین ببینید):
\[
\begin{aligned}
\text{a)}\quad \underline{(a\Delta u)^2}
&=\frac12(\vartheta \nu x)^2-\frac25f\cdot a_\nu^2
-\frac45u_x\cdot\theta_\nu(\vartheta \nu x)^2
+u_x^2\left\{\frac16i f-\frac1{40}s\right\},\\
(3.)\quad\text{b)}\quad \underline{(a\Delta u)^3}
&=-\frac3{10}\theta_\nu(\vartheta \nu x),\\
\text{c)}\quad \underline{(a\Delta u)^4}
&=\frac{21}{10}\theta_\nu^2-\frac3{10}H
-\frac65u_x\cdot\theta_\nu^3+\frac1{10}u_x^2\cdot\varrho.
\end{aligned}
\]

(کاهش دوگانهٔ عبارت‌های $a_\vartheta^2(b\theta u)^2$ و$a_\vartheta^2(b\theta u)^3$ و$a_\vartheta^2(a\theta u)^2(b\theta u)^2$ و$a_\vartheta^2(a\theta u)(b\theta u)^3$ پس از حذف $a_\vartheta^2$ نیز به همین نتایج می‌انجامد.)
از فرمول‌های (3.) نتیجه می‌شود که $(a\Delta u)^2$ نسبت به پیمانهٔ $\nu$ کاهشگر است. از این رو:

۱) کاهش دستگاه $\Delta$: رشتهٔ فرم‌های $(\Delta\Delta' u)^2\equiv a_\vartheta^2(a\Delta u)^2$ کاهشگر $(a\Delta u)^2$ را به‌عنوان عامل دارد.

۲) کاهش دستگاهی که از تاخوردگی با پیمانهٔ $\Delta$ پدید آمده است:

رشتهٔ فرم‌های $(\theta\Delta u)^2$ کاهشگر $(a\Delta u)^2$ را به‌عنوان عامل دارد.

برای فرم‌های $(\theta\Delta u)$ و$\Delta_\vartheta$ به دست می‌آید:
\[
(a\Delta u)^2(abu)=(\theta\Delta u)-u_x\cdot\Delta_\vartheta(\theta\Delta u),
\]
\[
(a\Delta u)(a\Delta b)=-\frac12\Delta_\vartheta+\frac12u_x\cdot\Delta_\vartheta^2.
\]

اما از کاهشگر $(\theta\Delta u)$، کاهش دستگاهی که از تاخوردگی با پیمانهٔ $\Delta$ پدید می‌آید نتیجه می‌شود.

\bigskip

پس دستگاه کامل نسبی به پیمانهٔ $\nu$ از شش فرم زیر تشکیل می‌شود:
\[
f,\ \theta,\ K,\ j,\ \A,\ N.
\]

به‌عنوان نمونه‌ای از بسط رشته‌ای در § ۲، استخراج فرمول (3.)a را به‌تفصیل می‌آوریم؛ در محاسبات بعدی باید فرمول نهایی III مستقیماً نوشته شود. (قطبی‌های فرم‌های صفرشونده از همان هنگام محاسبه حذف شده‌اند؛ سطر نخست مقادیری را می‌دهد که پس از بسط II وارد شده‌اند، و سطر دوم مقادیر نهایی را که با آغاز از آخرین معادلهٔ دستگاه به‌صورت بازگشتی یافته شده‌اند.)

بنا بر بسط II به دست می‌آید:
\[
\begin{aligned}
\text{1)}\quad&m=3.\quad n=4.\quad \lambda=2.\quad \mu=0.\quad \nu=\varkappa=1,\\[0.25em]
\underline{a_\vartheta b_\vartheta(b\theta u)^2}
&=[a_\vartheta]_{\widehat{ub^2}}b+\frac67\{a_\vartheta b_\vartheta(a\theta u)(b\theta u)-u_x\cdot a_\vartheta b_\vartheta(a\theta b)(b\theta u)\}\\
&\quad-\frac17\{a_\vartheta(a\theta u)^2b_\vartheta-2u_x\cdot a_\vartheta(a\theta u)(a\theta b)b_\vartheta+u_x^2\cdot a_\vartheta(a\theta b)^2b_\vartheta\}\\
&=\frac23(a\A u)^2+\frac15[a_\vartheta(a\theta u)^2]b+\frac1{30}f\cdot[a_\vartheta^2(a\theta u)^2]-\frac25u_x\cdot[a_\vartheta(a\theta u)^2]b^2\\
&\quad-\frac1{15}u_x\cdot[a_\vartheta^2(a\theta u)^2]b+\frac15u_x^2\cdot[a_\vartheta(a\theta u)^2]b^3+\frac1{30}u_x^2\cdot[a_\vartheta^2(a\theta u)^2]b^2;\\[0.7em]
\text{2)}\quad&m=2.\quad n=3.\quad \lambda=1.\quad \mu=1.\quad \nu=\varkappa=1,\\[0.25em]
\underline{a_\vartheta(a\theta u)b_\vartheta(b\theta u)}
&=\frac25\{a_\vartheta(a\theta u)^2b_\vartheta-u_x\cdot a_\vartheta(a\theta u)(a\theta b)b_\vartheta\}+\frac12a_\vartheta^2(b\theta u)^2\\
&\quad-\frac15\{a_\vartheta^2(b\theta u)(a\theta u)-u_x\cdot a_\vartheta^2(a\theta b)(b\theta u)\}\\
&=\frac12(a\A u)^2+\frac25[a_\vartheta(a\theta u)^2]b+\frac1{15}[a_\vartheta^2(a\theta u)^2]\cdot f\\
&\quad-\frac25u_x\cdot[a_\vartheta(a\theta u)^2]b^2-\frac1{10}u_x\cdot[a_\vartheta^2(a\theta u)^2]b+\frac1{30}u_x^2\cdot[a_\vartheta^2(a\theta u)^2]b^2;\\[0.7em]
\text{3)}\quad&m=2.\quad n=3.\quad \lambda=1.\quad \mu=1.\quad \nu=1.\quad \varkappa=2,
\end{aligned}
\]
\[
\begin{aligned}
\underline{a_\vartheta(a\theta b)b_\vartheta(b\theta u)}
&=\frac25\{a_\vartheta(a\theta b)(a\theta u)b_\vartheta-u_x\cdot a_\vartheta(a\theta b)^2b_\vartheta\}\\
&=\frac25[a_\vartheta(a\theta u)^2]b^2+\frac1{30}[a_\vartheta^2(a\theta u)^2]b-\frac25u_x\cdot[a_\vartheta(a\theta u)^2]b^3\\
&\quad-\frac1{30}u_x\cdot[a_\vartheta^2(a\theta u)^2]b^2;\\[0.7em]
\text{4)}\quad&m=1.\quad n=2.\quad \lambda=0.\quad \mu=2.\quad \nu=\varkappa=1,\\[0.25em]
\underline{a_\vartheta(a\theta u)^2b_\vartheta}
&=[a_\vartheta(a\theta u)^2]b+\frac23a_\vartheta^2(a\theta u)(b\theta u)\\
&=[a_\vartheta(a\theta u)^2]b+\frac16[a_\vartheta^2(a\theta u)^2]\cdot f-\frac16u_x\cdot[a_\vartheta^2(a\theta u)^2]b;\\[0.7em]
\text{5)}\quad&m=1.\quad n=2.\quad \lambda=0.\quad \mu=2.\quad \nu=1.\quad \varkappa=2\quad(\varkappa>\nu),\\[0.25em]
\underline{a_\vartheta(a\theta u)(a\theta b)b_\vartheta}
&=[a_\vartheta(a\theta u)^2]b^2+\frac13a_\vartheta^2(a\theta b)(b\theta u)\\
&=[a_\vartheta(a\theta u)^2]b^2+\frac1{12}[a_\vartheta^2(a\theta u)^2]b-\frac1{12}u_x\cdot[a_\vartheta^2(a\theta u)^2]b^2;\\[0.7em]
\text{6)}\quad&m=1.\quad n=2.\quad \lambda=0.\quad \mu=2.\quad \nu=1.\quad \varkappa=3,\\[0.25em]
\underline{a_\vartheta(a\theta b)^2b_\vartheta}&=[a_\vartheta(a\theta u)^2]b^3;\\[0.7em]
\text{7)}\quad&m=2.\quad n=4.\quad \lambda=2.\quad \mu=\nu=\varkappa=0.\quad(\text{پس از بسط I}),\\[0.25em]
\underline{a_\vartheta^2(b\theta u)^2}&=[a_\vartheta^2]_{\widehat{ub^2}}+\frac1{10}\{[a_\vartheta^2(a\theta u)^2]\cdot f-2u_x\cdot[a_\vartheta^2(a\theta u)^2]b+u_x^2\cdot[a_\vartheta^2(a\theta u)^2]b^2\};\\[0.7em]
\text{8)}\quad&m=1.\quad n=3.\quad \lambda=1.\quad \mu=1.\quad \nu=\varkappa=0.\quad(\text{بسط I}),\\[0.25em]
\underline{a_\vartheta^2(a\theta u)(b\theta u)}&=\frac14[a_\vartheta^2(a\theta u)^2]\cdot f-\frac14u_x\cdot[a_\vartheta^2(a\theta u)^2]b;\\[0.7em]
\text{9)}\quad&m=1.\quad n=3.\quad \lambda=1.\quad \mu=\varkappa=1.\quad \nu=0.\quad(\text{بسط I}),\\[0.25em]
\underline{a_\vartheta^2(a\theta b)(b\theta u)}&=\frac14[a_\vartheta^2(a\theta u)^2]b-\frac14u_x\cdot[a_\vartheta^2(a\theta u)^2]b^2.
\end{aligned}
\]

از 1) و4) نتیجه می‌شود:
\[
\begin{aligned}
(a\A u)^2
&=\frac65[a_\vartheta(a\theta u)^2]b+\frac15f\cdot[a_\vartheta^2(a\theta u)^2]+\frac35u_x\cdot[a_\vartheta(a\theta u)^2]b^2-\frac3{20}u_x\cdot[a_\vartheta^2(a\theta u)^2]b\\
&\quad-\frac3{10}u_x^2\cdot[a_\vartheta(a\theta u)^2]b^3-\frac1{20}u_x^2\cdot[a_\vartheta^2(a\theta u)^2]b^2,\\[0.4em]
(a\A u)^2&=-a_\nu b_\nu+\frac35f\cdot a_\nu^2+\frac45u_x\cdot a_\nu^2b_\nu-\frac3{20}u_x^2\cdot a_\nu^2b_\nu^2-\frac1{12}u_x^2\cdot i\cdot f.
\end{aligned}
\]

(برای تبدیل به نمادهای $\theta$، نک. § ۸ و§ ۹.)
فرمول (3.)a را می‌توان کوتاه‌تر با انتقال مستقیم بسط I و با معرفی متغیر کمکی $w=x\widehat{u}b$ محاسبه کرد، در حالی که از فرمول‌های (3.)b و(3.)c به بعد راهی که در پیش گرفته‌ایم ساده‌تر می‌شود؛ در (3.)b بنا بر § ۹ از پیش می‌دانیم که همهٔ جمله‌های دارای عامل $u_x$ باید حذف شوند. برای محاسبهٔ (3.)b و(3.)c به دست می‌آوریم:
\[
\begin{aligned}
(3.)\mathrm{b)}\quad 1)&\quad a_\vartheta(a\theta u)b_\vartheta(b\theta u)^2=\frac12(a\A u)^3+\frac45[a_\vartheta(a\theta u)^2]_{\widehat{ub}}b+\frac{77}{360}[a_\vartheta^2(a\theta u)^2]_{\widehat{ub}},\\
2)&\quad a_\vartheta(a\theta u)b_\vartheta(b\theta u)^2=[a_\vartheta(a\theta u)^2]_{\widehat{ub}}b+\frac{13}{36}[a_\vartheta^2(a\theta u)^2]_{\widehat{ub}};\\[0.5em]
(3.)\mathrm{c)}\quad 1)&\quad a_\vartheta(a\theta u)b_\vartheta(b\theta u)^3=\frac12(a\A u)^4+\frac65[a_\vartheta(a\theta u)^2]_{\widehat{ub^2}}b+\frac{53}{120}[a_\vartheta^2(a\theta u)^2]_{\widehat{ub^2}}\\
&\qquad-\frac65u_x\cdot[a_\vartheta(a\theta u)^2]_{\widehat{ub^2}}b^2-\frac35u_x\cdot[a_\vartheta^2(a\theta u)^2]_{\widehat{ub^2}}b+\frac{19}{120}u_x^2\cdot[a_\vartheta^2(a\theta u)^2]_{\widehat{ub^3}}b^2,\\
2)&\quad a_\vartheta(a\theta u)b_\vartheta(b\theta u)^3=\frac34[a_\vartheta^2(a\theta u)^2]_{\widehat{ub^2}}.
\end{aligned}
\]

یادداشت پایانی: تاخوردگی‌های $f$ بر $j$ که بنا بر § ۵ کاهش‌پذیرند، به‌طور مشابه محاسبه می‌شوند. به دست می‌آید
\[
\tag{4.}
\begin{aligned}
a_j&=-\theta_\nu+u_x\cdot\left\{\frac92\theta_\nu^2-\frac12H\right\}-3u_x^2\cdot\theta_\nu^3+\frac12u_x^3\cdot\varrho,\\
a_j^2&=\frac{21}{10}\theta_\nu^2-\frac3{10}H-2u_x\cdot\theta_\nu^3+\frac12u_x^2\cdot\varrho,\\
a_j^3&=-\frac7{10}\theta_\nu^3+\frac12u_x\cdot\varrho,\\
a_j^4&=\frac35\varrho.
\end{aligned}
\]

از (3.) و(4.)، با انتقال از حوزهٔ دوتایی، نتیجه می‌شود
\[
(\theta\theta'u)^2=\frac13 j\cdot f\modu{\nu}\srcfn{*)}{Gordan--Kerschensteiner، همان‌جا، ص. 181.}.
\]
دیگر فرم‌های رشتهٔ $(\theta\theta'u)^2$ و فرم $\theta'_\vartheta$ به نمادهای $\nu$ کاهش‌پذیرند. پس می‌توانیم جایگزین کنیم:
\[
\modu{\nu}:\quad \text{تاخوردگی‌های با } f\cdot j \text{ با تاخوردگی‌های متناظر با }(\theta\theta'u)^2.
\]
همچنین به دست می‌آید:
\[
(\vartheta\vartheta'x)^2=\frac43 f\cdot\A,\qquad
\theta_{\vartheta'}(\vartheta\vartheta'x)=-N.
\]

\begin{center}
{\Large\bfseries فصل سوم.\quad دستگاه کامل نسبی به پیمانهٔ $(s,\varrho,t)$.}
\end{center}

\begin{center}
\textbf{\S\ ۷.}\quad\emph{مروری بر تشکیل دستگاه.}
\end{center}

برای گذار از دستگاه کامل نسبی به پیمانهٔ $\nu$ به دستگاه کامل نسبی با پیمانهٔ یک مرتبه بالاتر، بنا بر تعریف \S\ ۳ باید دستگاه فرم‌های $\nu$ را نسبت به این پیمانهٔ بالاتر تشکیل داد و تاخوردگی آن را با فرم‌های دستگاه کامل نسبی به پیمانهٔ $\nu$ گرفت. اما $\nu=u_\nu^4$ به‌عنوان پادهموردای درجهٔ چهارمِ فرم $f$، از راه دوگانگی در برابر آن قرار دارد. پس اگر $\nu$ را فرم بنیادی بگیریم، دستگاه کامل نسبی‌ای از شش فرم به پیمانهٔ $\nu(\nu)=(\nu\nu_1x)^4=s_x^4$ به دست می‌آوریم.

پس برای تشکیل دستگاه کامل نسبی به پیمانهٔ $s$ (دستگاه III) باید بر هم نهیم
\[
\begin{array}{ll}
\text{دستگاه I:} & f,\ \theta,\ K,\ j,\ \A,\ N\quad \text{بر}\vphantom{\dfrac11}\\[0.3em]
\text{دستگاه II:} & \nu,\ H=(\nu\nu_1x)^2u_\nu^2u_{\nu_1}^2,\ L=(\nu\eta x)H_x^2u_\nu^3u_\eta^3,\ g=(\nu\eta x)^4H_x^2,\\
& \sigma=H_\eta^2u_\nu^2u_\eta^4,\ (\nu\sigma x).
\end{array}
\]
نشان خواهیم داد (فرمول (13.)) که فرم $g$ کاهش‌پذیر است و بنابراین دستگاه II فقط از پنج فرم تشکیل می‌شود.

در جریان محاسبه، دو فرم درجهٔ دوِ زیر
\[
\varrho=u_\varrho^2=\theta_\nu^4u_\vartheta^2,\qquad t=t_x^2=a_\eta^4H_x^2
\]
به‌عنوان پیمانه‌ها افزوده می‌شوند، به‌طوری که دستگاهی کامل نسبی به پیمانهٔ $(s,\varrho,t)$ حاصل می‌شود؛ در اینجا پیمانهٔ $(\varrho,t)$ بالاتر از پیمانهٔ $s$ به شمار می‌آید.

فرم‌های منفرد برحسب مرتبه نسبت به ضرایب مرتب می‌شوند.

تاخوردگی
\[
\theta_\eta(K_\eta,\theta_\iota,\ldots)
\]
را بالاتر از تاخوردگی
\[
H_\vartheta(H_\kappa,L_\vartheta,\ldots).
\]
قرار می‌دهیم. بدین ترتیب، بنا بر \S\ ۱ و\S\ 3b، ترتیب فرم‌های منفردِ هم‌مرتبه به‌طور یکتا تعیین می‌شود.

ساخت فرم‌های هم‌مرتبه به‌صورت جدولی انجام می‌شود:

I. تعیین اینکه کدام فرم‌های دستگاه‌های I وII باید با یکدیگر تاخورده شوند تا مرتبهٔ معین نسبت به ضرایب حاصل شود.

II. فهرست‌کردن فرم‌های کاهش‌ناپذیر مطابق طرح رشتهٔ فرم‌ها.

III. کاهش رشته‌های فرم‌ها یا فرم‌های کاهش‌پذیر یا تجزیه‌پذیر، یعنی نقاطی که رشتهٔ کامل فرم‌ها در آنها قطع می‌شود، و از این راه اثبات اینکه همهٔ ساخت‌های کاهش‌ناپذیر به موارد ذکرشده منحصرند.

IV. نتایج، شامل: ۱) تعیین کاهشگرهای تازه‌پدید، ۲) تعیین فرم‌هایی که در قالب حاصل‌ضرب با فرم‌های همان دستگاه، در تاخوردگی با فرم‌های دستگاه دیگر وارد نمی‌شوند.

نشان خواهیم داد که فقط موارد زیر پدید می‌آیند:

۱) تاخوردگیِ یک فرم از دستگاه I با یک فرم از دستگاه II.

۲) تاخوردگیِ توان‌های $\theta$ و به‌ترتیب $H$، یا دو فرم از یکی از دستگاه‌ها، با یک فرم از دستگاه دوم.

برای تاخوردگی به طرح زیر می‌رسیم:
\[
\begin{array}{r@{\ }c@{\ }l@{\qquad}c@{\qquad}r@{\ }c@{\ }l}
\multicolumn{3}{c}{\text{دستگاه I}} & \text{تاخورده با} & \multicolumn{3}{c}{\text{دستگاه II}}\\[0.25em]
1. & \text{مرتبه} & f & & 2. & \text{مرتبه} & \nu\\
2. & ,, & \theta & & 4. & ,, & H\\
3. & ,, & K,j,\A & & 6. & ,, & L,\sigma\\
4. & ,, & N,\theta^2,f\cdot j & & 8. & ,, & (\nu\sigma x),H^2\\
5. & ,, & \theta\cdot K & & 10. & ,, & H\cdot L\\
6. & ,, & \theta^3,\A\cdot j & & 12. & ,, & H^3.
\end{array}
\]

برای کنترل محاسبه، افزون بر «کاهش دوگانه»، دو فرم ویژهٔ زیر به کار می‌روند:
\[
\begin{array}{ll}
\text{I.}&\begin{array}{rlrlrl}
f&=\displaystyle\sum_3x_1^3x_2,& a_\nu^2&=\dfrac{i}{6}u_x^2,& s&=\dfrac{i}{3}f,\\
a_\eta^2&=-\dfrac19i\theta,& \A_\nu^2&=\dfrac{i}{6}\theta,& \sigma&=-\dfrac{i}{6}j,\\
g&=-\dfrac23i\A,& \varrho&=0,& t&=0,
\end{array}\\[1.2em]
\text{II.}&\begin{array}{rlrlrl}
f&=\displaystyle\sum_3x_1^4,& s&=\dfrac43if,& a_\eta^2&=\dfrac29i\theta,\\
\A_\nu^2&=\dfrac1{15}i\theta,& \sigma&=\dfrac43ij,& g&=\dfrac43i\A,\\
\varrho&=0,& t&=0.
\end{array}
\end{array}
\]

یادداشت پایانی: برای فرم بنیادیِ $\nu$، فرمول‌های زیر متناظر با فرمول‌های \srcLTR{(1.)}، \srcLTR{(2.)}، \srcLTR{(3.)} و\srcLTR{(4.)} هستند:
\[
\tag{5.}
\begin{array}{@{}c|c|c@{}}
\cdot & \underline{H_\nu}=\dfrac13u_x\cdot\sigma & \cdot\\[0.45em]
\underline{(\nu\eta x)^2}=A & \underline{H_\nu(\nu\eta x)}=0 & \underline{H_\nu^2}=\sigma\\[0.45em]
\underline{(\nu\eta x)^3}=0 & \underline{H_\nu(\nu\eta x)^2}=B & \underline{H_\nu^2(\nu\eta x)}=0\\[0.45em]
\underline{(\nu\eta x)^4}=g & \underline{H_\nu(\nu\eta x)^3}=0 & \underline{H_\nu^2(\nu\eta x)^2}=C,
\end{array}
\]
\[
\begin{aligned}
A&=\frac16s\cdot\nu-\frac23u_x\cdot s_\nu+\frac23u_x^2\cdot s_\nu^2-\frac{J}{18}u_x^4,\\
B&=-\frac56s_\nu+\frac76u_x\cdot s_\nu^2-\frac{J}{9}u_x^3,\\
C&=\frac23s_\nu^2-\frac{J}{9}u_x^2.
\end{aligned}
\]
\[
\tag{6.}
s_\nu^3u_\nu s_x=\frac13u_x\cdot s_\nu^4=\frac13J\cdot u_x.
\]

\[
\tag{7.}
\begin{aligned}
\text{a)}\quad (\nu\sigma x)^2
&=\frac12(Hsu)^2-\frac25\nu\cdot s_\nu^2-\frac45u_x\cdot s_\eta(Hsu)^2+u_x^2\left\{\frac16J\cdot\nu-\frac1{40}(ss'u)^4\right\},\\
\text{b)}\quad (\nu\sigma x)^3&=-\frac3{10}s_\eta(Hsu),\\
\text{c)}\quad (\nu\sigma x)^4&=\frac{21}{10}s_\eta^2-\frac3{10}Z-\frac65u_x\cdot s_\eta^3+\frac1{10}u_x^2\cdot s_\eta^4.
\end{aligned}
\]
\[
\tag{8.}
\begin{aligned}
\text{a)}\quad g'_\nu&=-s_\eta+u_x\left\{\frac92s_\eta^2-\frac12Z\right\}-3u_x^2\cdot s_\eta^3+\frac12u_x^3\cdot s_\eta^4,\\
\text{b)}\quad g_\nu^2&=\frac{21}{10}s_\eta^2-\frac3{10}Z-2u_x\cdot s_\eta^3+\frac12u_x^2\cdot s_\eta^4,\\
\text{c)}\quad g_\nu^3&=-\frac7{10}s_\eta^3+\frac12u_x\cdot s_\eta^4,\\
\text{d)}\quad g_\nu^4&=\frac35s_\eta^4.
\end{aligned}
\]

\begin{center}
\textbf{\S\ ۸.}\quad\emph{فرم‌های مرتبهٔ سوم (دستگاه III).}
\end{center}

\noindent\textbf{تاخوردگیِ $f$ با $\nu$.}

\noindent\emph{فرم‌های کاهش‌ناپذیر:}
\[
\begin{array}{cccc}
a_\nu&\cdot&\cdot&\cdot\\
\cdot&a_\nu^2&\cdot&\cdot\\
\cdot&\cdot&\cdot&\cdot\\
\cdot&\cdot&\cdot&a_\nu^4=i.
\end{array}
\]

\noindent\emph{کاهش:}
\[
a_\nu^3=\frac13iu_x\qquad\text{(فرمول (2.)).}
\]
\noindent\emph{نتایج:} ۱) $a_\nu^3$ کاهشگر است. ۲) $f$ فقط در حاصل‌ضرب‌هایی با فرم‌های پادهموردا $(j)$ در تاخوردگی با $\nu$ وارد می‌شود. همین امر دربارهٔ $\nu$ نسبت به $f$ نیز برقرار است.

\begin{center}
\textbf{\S\ ۹.}\quad\emph{فرم‌های مرتبهٔ چهارم (دستگاه III).}
\end{center}

\noindent\textbf{تاخوردگیِ $\theta$ با $\nu$.}

\noindent\emph{فرم‌های کاهش‌ناپذیر:}
\[
\begin{array}{cccc}
(\vartheta\nu x)&\theta_\nu&\cdot&\cdot\\
(\vartheta\nu x)^2&\theta_\nu(\vartheta\nu x)&\theta_\nu^2&\cdot\\
\cdot&\theta_\nu(\vartheta\nu x)^2&\cdot&\theta_\nu^3.
\end{array}
\]

\noindent\emph{کاهش‌ها:} از کاهشگر $a_\nu^3$، با کاهش دوگانهٔ عبارت‌های
$a_\nu^3(abu)$، $a_\nu^3b_\nu$، $a_\nu^3b_\nu(abu)$، بر اساس فرض زیر:
\[
\begin{aligned}
\theta_\nu^2(\vartheta\nu x)
&=a_\nu^2(\widehat{ab}\nu x)(abu)-\frac13(\nu\nu_1x)^3
=a_\nu b_\nu(\widehat{ab}\nu x)(abu)+\frac16(\nu\nu_1x)^3\\
&=a_\nu^3(abu)-a_\nu^2b_\nu(abu)=2a_\nu^2b_\nu(abu)=\frac23a_\nu^3(abu),\\[.4em]
\theta_\nu^2(\vartheta\nu x)^2
&=a_\nu^2(\widehat{ab}\nu x)^2-\frac13(\nu\nu_1x)^4
=a_\nu b_\nu(\widehat{ab}\nu x)^2+\frac16(\nu\nu_1x)^4\\
&=i f-2a_\nu^3b_\nu+a_\nu^2b_\nu^2-\frac13s
=2a_\nu^3b_\nu-2a_\nu^2b_\nu^2+\frac16s
=\frac23i f-\frac23a_\nu^3b_\nu-\frac16s,\\[.4em]
\theta_\nu^3(\vartheta\nu x)&=a_\nu^2b_\nu(\widehat{ab}\nu x)(abu)=a_\nu^3b_\nu(abu).
\end{aligned}
\]

فرمول‌های زیر به دست می‌آیند:
\[
\tag{9}
\begin{array}{c|c|c}
\underline{\theta_\nu^2(\vartheta\nu x)}=0
&\underline{\theta_\nu^3(\vartheta\nu x)}=0
&\underline{\theta_\nu^4u_\vartheta^2}=u_\varrho^2=\varrho\quad(\text{پیمانهٔ الحاقی، \S\ ۷}),\\[0.6em]
\underline{\theta_\nu^2(\vartheta\nu x)^2}=\dfrac49 i f-\dfrac16s
&&
\end{array}
\]

\noindent\emph{نتایج:} ۱) $\theta_\nu^2(\vartheta\nu x)^2$ کاهشگر است.
۲) $\nu$ فقط در حاصل‌ضرب‌هایی با فرم‌های پادهموردا $(g)$ در تاخوردگی با $\theta$ وارد می‌شود.

\begin{center}
\textbf{\S\ ۱۰.}\quad\emph{فرم‌های مرتبهٔ پنجم (دستگاه III).}
\end{center}

\[
\text{تاخوردگیِ }\left\{\begin{array}{l}
K,j,\A \text{ با }\nu,\\
f\text{ با }H.
\end{array}\right.
\]

\noindent\emph{فرم‌های کاهش‌ناپذیر:}
\[
\begin{array}{c@{\qquad}c@{\qquad}c@{\qquad}c}
\text{A)} & \begin{array}{cccc}
\cdot&K_\nu&\cdot&\cdot\\
(k\nu x)^2&\cdot&K_\nu^2&\cdot\\
(k\nu x)^3&K_\nu(k\nu x)^2&\cdot&\cdot\\
\cdot&K_\nu(k\nu x)^3&\cdot&\cdot
\end{array}
&
\text{B)}\ \begin{array}{c}(j\nu x)\\(j\nu x)^2\\(j\nu x)^3\\ \cdot\end{array}
&
\text{C)}\ \begin{array}{c}\A_\nu\\\cdot\\\cdot\\\cdot\end{array}
\end{array}
\]
\[
\text{D)}\qquad
\begin{array}{cccc}
(aHu)&(aHu)^2&\cdot&\cdot\\
a_\eta&a_\eta(aHu)&a_\eta(aHu)^2&\cdot\\
\cdot&a_\eta^2&\cdot&\cdot
\end{array}
\]

\noindent\emph{کاهش‌ها:}

\noindent A) رشتهٔ فرم‌های $K_\nu^3$: کاهشگر $a_\nu^3$.

فرم $K_\nu^2(k\nu x)^2$: کاهشگر $\theta_\nu^2(\vartheta\nu x)^2$.
\[
(k\nu x),\quad K_\nu(k\nu x),\quad K_\nu^2(k\nu x)
\]
بنا بر \S\ ۳ فرم‌های تجزیه‌پذیرند.

\noindent B) فرم
\[
\begin{aligned}
(j\nu x)^4&=(\widehat{a\theta}\nu x)^4
=a_\nu^4\theta+\theta_\nu^4 f-4a_\nu^3\theta_\nu
-4a_\nu\{a_\nu\theta_x-(\widehat{a\theta}\nu x)\}^3\\
&\quad+6a_\nu^2\{a_\nu\theta_x-(\widehat{a\theta}\nu x)\}^2
=\R f+\frac53 i\theta-6a_\nu^2(\widehat{a\theta}\nu x)^2
+4a_\nu(\widehat{a\theta}\nu x)^3,
\end{aligned}
\]
و ازاین‌رو
\[
(j\nu x)^4=\R f+\frac53 i\theta\modu{(\A,\nu)},\qquad (\text{نک. پایان \S\ ۵}).
\]

\noindent C) فرم $\A_\nu^4$: کاهشگر $\theta_\nu^4$.

فرم‌های $\A_\nu^2$ و$\A_\nu^3$: کاهشگر $a_\nu^3$ یا $\theta_\nu^4$؛ با جایگزین‌کردن فرم $\A$ با فرم‌های پایین‌ترِ رشتهٔ فرم‌ها (پایان \S\ ۵)، یعنی با کاهش دوگانهٔ عبارت‌های
\[
a_\vartheta a_\nu^3,\qquad
a_\vartheta a_\nu^3\theta_\nu,\qquad
a_\vartheta\theta_\nu^4.
\]
محاسبهٔ دوگانهٔ $\A_\nu^3$ زمینهٔ کاهش فرم بالاتر $a_\eta^3$ را فراهم می‌کند.

\noindent\emph{فرمول‌های کاهش:}
\[
\tag{10}
\begin{aligned}
\text{a)}\quad \A_\nu^2
&=\frac35a_\eta^2-\frac3{10}(asu)^2+\frac13 i\theta
+\frac15u_x^2(2a_\R^2-t),\\
\text{b)}\quad \A_\nu^3
&=-\frac3{10}a_\R+\frac35u_x\left(\frac{13}{10}a_\R^2-\frac25t\right),\\
\text{c)}\quad \A_\nu^4&=a_\R^2-\frac25t.
\end{aligned}
\]

\noindent\emph{استخراج فرمول‌ها:}
\[
\begin{aligned}
\text{a)}\quad
 a_\vartheta a_\nu^3
&=\frac13 i\theta
=[a_\vartheta]_{\nu^3}+\frac67[a_\vartheta^2]_{\nu^2}
+\frac65[a_\vartheta(a\theta u)^2]_{\nu^2}\,\widehat{\nu}x^2
+\frac{11}{30}[a_\vartheta^2(a\theta u)^2]_{\nu}\,\widehat{\nu}x^2\\
&\quad-\frac67u_x[a_\vartheta^2]_{\nu^3}
-\frac16u_x[a_\vartheta^2(a\theta u)^2]_{\nu^2}\,\widehat{\nu}x^2,\\
\frac13 i\theta
&=\A_\nu^2-\frac23u_x\A_\nu^3-a_\nu a_{\nu_1}(\nu\nu_1x)^2
+\frac25a_\nu^2(\nu\nu_1x)^2
+\frac15u_xa_\nu^2a_{\nu_1}(\nu\nu_1x)^2;
\end{aligned}
\]
\[
\begin{aligned}
\text{b)}\quad
 a_\vartheta a_\nu^3\theta_\nu
&=0=[a_\vartheta]_{\nu^4}+\frac9{14}[a_\vartheta^2]_{\nu^3}
+\frac35[a_\vartheta(a\theta u)^2]_{\nu^2}\,\widehat{\nu}x^2
+\frac{17}{120}[a_\vartheta^2(a\theta u)^2]_{\nu}\,\widehat{\nu}x^2\\
&\quad-\frac9{14}u_x[a_\vartheta^2]_{\nu^4}
-\frac1{24}u_x[a_\vartheta^2(a\theta u)^2]_{\nu^2}\,\widehat{\nu}x^2,\\
0&=\frac56\A_\nu^3-\frac12u_x\A_\nu^4-\frac14a_\eta^3+\frac1{20}u_xa_\eta^4;
\end{aligned}
\]
\[
\begin{aligned}
\text{b')}\quad
 a_\vartheta\theta_\nu^4
&=a_\R=a_\vartheta\theta_\nu\{a_\nu-(a\theta\widehat{\nu}x)\}^3\\
&=-\frac32[a_\vartheta^2]_{\nu^3}
+\frac35[a_\vartheta(a\theta u)^2]_{\nu^2}\,\widehat{\nu}x^2
-\frac{37}{120}[a_\vartheta^2(a\theta u)^2]_{\nu}\,\widehat{\nu}x^2\\
&\quad+\frac32u_x[a_\vartheta^2]_{\nu^4}
+\frac{49}{120}u_x[a_\vartheta^2(a\theta u)^2]_{\nu^2}\,\widehat{\nu}x^2,\\
a_\R&=-\frac32\A_\nu^3+\frac32u_x\A_\nu^4-\frac{11}{20}a_\eta^3+\frac7{20}u_xa_\eta^4;
\end{aligned}
\]
\[
\begin{aligned}
\text{c)}\quad
 a_\vartheta^2\theta_\nu^4
&=a_\R^2=[a_\vartheta^2]_{\nu^4}+\frac35[a_\vartheta^2(a\theta u)^2]_{\nu^2}\,\widehat{\nu}x^2,\\
a_\R^2&=\A_\nu^4+\frac25a_\eta^4.
\end{aligned}
\]

\noindent D) کاهش $a_\eta^3$ از طریق کاهش دوگانهٔ $\A_\nu^3$ (C).

دیگر کاهش‌ها با محاسبه‌ای به دست می‌آیند که دوگانِ محاسبهٔ \S\ ۹ است.

\noindent\emph{فرمول‌های کاهش:}
\[
\tag{11}
\begin{array}{rl|rl}
\underline{a_\eta^2(aHu)}&=-\dfrac16(asu)^3,&
\underline{a_\eta^2(aHu)^2}&=\dfrac49 i \nu-\dfrac16(asu)^4,\\[0.7em]
\underline{a_\eta^3}&=-a_\R+\dfrac35u_xa_\R^2+\dfrac15u_xt,
&\underline{a_\eta^3(aHu)}&=0,\\[0.7em]
&&\underline{a_\eta^4}&=t\quad(\text{به‌عنوان پیمانه ملحق شده}).
\end{array}
\]

\noindent\emph{نتایج:}
۱) $(j\nu x)^4$ و$\A_\nu^2$ و$a_\eta^2(aHu)$ کاهشگرند.

۲) $\nu$ فقط در حاصل‌ضرب‌هایی با فرم‌های پادهموردا، هنگام تاخوردگی با $K,j,\A$، وارد می‌شود؛ همین حکم دربارهٔ $f$ نسبت به $H$ و دربارهٔ $j$ نسبت به $\nu$ نیز برقرار است، زیرا $\theta_\nu(\vartheta\cdot j,\nu,x)^3$ بنا بر \S\ ۱۱ B کاهش‌پذیر می‌شود.

\begin{center}
\textbf{\S\ ۱۱.\quad فرم‌های مرتبهٔ ششم (دستگاه III).}
\end{center}
\[
\text{تاخوردگیِ }\left\{\begin{array}{l}
\theta^2,\ f\cdot j,\ N \text{ با }\nu,\\
\theta\text{ با }H.
\end{array}\right.
\]


\noindent\emph{فرم‌های کاهش‌ناپذیر:}
\[
\renewcommand{\arraystretch}{1.25}
\begin{array}{@{}r@{\quad}c@{\qquad\qquad}r@{\quad}c@{\qquad\qquad}r@{\quad}c@{}}
\text{A)}&(\vartheta^2\nu x)^3&\text{B)}&a_\nu(j\nu x)&\text{C)}&N_\nu,\\[0.45em]
&\cdot&&\cdot&&\\[-0.1em]
&\theta_\nu(\vartheta^2\nu x)^3&&a_\nu^2(j\nu x)&&
\end{array}
\]

\[
\renewcommand{\arraystretch}{1.25}
\begin{array}{@{}r@{\quad}c|c|c|c@{}}
&\cdot&(\theta Hu)&(\theta Hu)^2&\cdot\\
\text{D)}&(\vartheta\eta x)&\theta_\eta&H_\vartheta(\theta Hu),\ \theta_\eta(\theta Hu)&\theta_\eta(\theta Hu)^2\\
&(\vartheta\eta x)^2&H_\vartheta(\vartheta\eta x),\ \theta_\eta(\vartheta\eta x)&\theta_\eta^2&\theta_\eta^2(\theta Hu)\\
&\cdot&\theta_\eta(\vartheta\eta x)^2&\cdot&\cdot
\end{array}
\]
\noindent\emph{کاهش‌ها:}


\noindent A) عبارت $(\vartheta^2\nu x)^4$ بنا بر فرمول (3.)a تجزیه می‌شود:
\[
(a\A\widehat{\vartheta}x)^2=\frac12(\vartheta^2\nu x)^4-\frac35 f\cdot a_\nu^2(\nu\vartheta x)^2
=\frac13\A^2+f\cdot\A_\vartheta^2.
\]

\noindent B) عبارت $a_\nu(j\nu x)^2$ بنا بر فرمول (9.)a تجزیه می‌شود؛ بدین‌گونه که مطابق نتیجهٔ \S\ ۶، $f\cdot j$ را با $(\theta\theta'u)^2$ جایگزین می‌کنیم و کاهش‌پذیریِ $\theta_\vartheta'$ را در نظر می‌گیریم:
\[
\frac13a_\nu(j\nu x)^2=(\widehat{\theta\theta'}\nu x)^2\theta_\nu
=\theta_\nu^3\cdot\theta-\theta_\nu^2\theta_\nu'
=\theta_\nu^3\cdot\theta+\theta_\nu^2(\widehat{\vartheta}\nu\theta'u)
=\theta_\nu^3\cdot\theta-\frac23\theta_\nu^3\cdot\theta.
\]
فرم‌های $a_\nu(j\nu x)^3$ و$a_\nu^2(j\nu x)^2$؛ و کاهشگرهای $a_j$ و$(j\nu x)^4$:
\[
a_\nu(j\nu x)^3=a_j(j\nu x)^3-(j\nu x)^3(j\nu\widehat{a}u),
\]
\[
a_\nu^2(j\nu x)^2=a_j^2(j\nu x)^2-2a_j(j\nu x)^2(j\nu\widehat{a}u)+(j\nu x)^2(j\nu\widehat{a}u)^2.
\]

\noindent C) رشتهٔ فرم‌های $N_\nu^2$؛ کاهشگر $\A_\nu^2$.
عبارت‌های $(n\nu x)$ و$N_\nu(n\nu x)$ بنا بر \S\ ۳ به‌صورت ترانس‌وکشن روی دترمینان‌های ژاکوبین تجزیه می‌شوند.

\noindent D) \emph{مروری بر کاهش‌ها:}

a) رشتهٔ فرم‌های $\theta_\eta^2H_\vartheta$: کاهشگر $a_\eta^2(aHu)$. (به فرم‌هایی با نمادهای بالاتر می‌انجامد.)

b) رشتهٔ فرم‌های $\theta_\eta H_\vartheta$: کاهشگر $a_\eta^2(aHu)$ با کاهش دوگانه.
(به فرم‌هایی با نمادهای بالاتر و نیز به فرم‌های بالاترِ رشتهٔ کلیِ فرم‌ها، یعنی رشتهٔ فرم‌های $\theta_\eta^2$، می‌انجامد.) به‌جای $\theta_\eta H_\vartheta$ رشتهٔ فرم‌های زیر را در نظر می‌گیریم:
\[
\theta_\eta(\vartheta\eta x)(\theta Hu)=\theta_\eta H_\vartheta+\theta_\eta^2,
\]
و در آن نمادهای $\theta$ را با فرم پایین‌ترِ $(abu)$ جایگزین می‌کنیم؛ بدین‌ترتیب به کاهش دوگانهٔ رشتهٔ فرم‌های
\[
a_\eta^2(aHu)(abu)=\theta_\eta H_\vartheta+\frac32\theta_\eta^2\modu{s}
\]
تا فرم پایانیِ
\[
a_\eta^3b_\eta(bHu)(abu)=\theta_\eta^3H_\vartheta+\frac12\theta_\eta^4.
\]

c) رشتهٔ فرم‌های $\theta_\eta^3$: با کاهش دوگانهٔ فرم‌هایی که هم‌زمان به رشته‌های فرمِ $\theta_\eta H_\vartheta$ و$\theta_\eta^2H_\vartheta$ تعلق دارند ـ یعنی فرم‌های رشتهٔ $\theta_\eta^2H_\vartheta$ ـ از یک سو به فرم‌هایی با نمادهای بالاتر، و از سوی دیگر به فرم‌هایی با نمادهای بالاتر و به رشتهٔ فرم‌های بالاترِ $\theta_\eta^3$ می‌رسیم.

d) فرم‌های $\theta_\eta^2(\vartheta\eta x)$، $\theta_\eta^2(\vartheta\eta x)^2$ و$\theta_\eta^3(\vartheta\eta x)$: با کاهش دوگانه به‌وسیلهٔ کاهشگر $a$، همانند محاسبهٔ \S\ ۹.

e) فرم‌های $\theta_\eta^2(\theta Hu)^2$ و$\theta_\eta^2(\vartheta\eta x)^2$: با کاهش دوگانهٔ عبارت‌های $\A_\nu^2(a\A u)^4$ و$(a\A\widehat{\nu}x)^4$ به‌وسیلهٔ کاهشگرهای $\A_\nu^2$ و$(a\A u)^4$.

f) فرم‌های $H_\vartheta$ و$H_\vartheta^2$: فرم‌هایی تجزیه‌پذیرند. این امر برای $H_\vartheta^2$ از کاهش دوگانهٔ عبارت $\A_\nu^2(a\A u)^2$، یا نیز از محاسبهٔ مستقیم حاصل‌ضرب‌های $(a_\nu^2)^2$ و$a_\nu^2\cdot b_{\nu_1}$ برحسب فرم‌های دستگاه، به دست می‌آید. (محاسبهٔ مشابهِ $(a_\nu)^2$ رابطه‌ای میان فرم‌های تجزیه‌پذیر می‌دهد.)

g) کاهش فرم‌های بالاتر از این واقعیت حاصل می‌شود که فرم‌های رشتهٔ $\theta\cdot H$ دو امکانِ کاهش دارند (به دو رشتهٔ فرمِ کاهش‌پذیر تعلق دارند)، بدین شرح:
\[
\begin{array}{rl}
g&\text{با کاهش دوگانهٔ }\theta_\eta^2(\vartheta\eta x)^2:\text{ کاهش‌های d) وe) را می‌پذیرد؛}\\[0.3em]
s_\vartheta^2(\theta s u)&\text{با کاهش دوگانهٔ }\theta_\eta^3(\vartheta\eta x):\text{ کاهش‌های c) وd) را می‌پذیرد؛}\\[0.3em]
s_\vartheta^2(\theta s u)^2&\text{با کاهش دوگانهٔ }\theta_\eta^2H_\vartheta^2:\text{ کاهش a) را می‌پذیرد و}\\
&\text{کاهشگر }\theta_\nu^2(\vartheta\nu x)^2\text{ را به‌عنوان عامل دارد؛}\\[0.3em]
s_\nu^2&\text{با کاهش دوگانهٔ }\theta_\eta^3H_\vartheta:\text{ کاهش‌های a) وc) را می‌پذیرد.}
\end{array}
\]
استقلال فرم‌های باقی‌مانده با کاهش دوگانهٔ رشتهٔ فرم‌های $\A_\nu^2(a\A u)^2\equiv\theta_\eta^2$ اثبات می‌شود.
\noindent\emph{اجرای کاهش‌ها:}

وجود کاهش‌های a)، b)، c) وd) از پیش روشن است، زیرا بنا بر قانون تشکیلِ بیان‌شده، رابطه‌های حاصل از کاهش دوگانه مستقل از یکدیگرند. دربارهٔ کاهش‌های e)، f) وg) باید با محاسبه نشان داد که هیچ همانی‌ای پدید نمی‌آید.

همچنین فرمول‌های حاصل از کاهش‌های a)، b)، c) وd) را، با پیش‌رفتن به‌طور متقارن نسبت به جملهٔ قطری در رشتهٔ فرم‌های $\theta\cdot H$ تا فرم $\theta_\eta$ و گاه با اشاره‌ای به محاسبه، می‌آوریم؛ زیرا بخشی از آن‌ها در کاهش‌های e)، f) وg) و بخشی بعدتر به کار می‌روند.

\noindent ۱) $H_\vartheta$ و$H_\vartheta^2$ بنا بر f). با قرار دادن\srcfn{*)}{علامت $\sim$ به‌جای علامت تساوی بدین معناست که حاصل‌ضرب‌های $u_x$ در فرم‌هایی بالاتر از فرم‌های حاصل از تاخوردگی‌های III وIV حذف شده‌اند.}
\[
\begin{aligned}
a_\nu\cdot b_{\nu_1}^2
&=\nu\cdot a_\nu b_\nu^2-(aHu)(bHu)b_\eta
\sim \nu\cdot a_\nu b_\nu^2-f\cdot a_\eta(aHu)^2-(abu)(aHu)b_\eta+(abu)^2b_\eta,\\
a_\nu^2\cdot b_{\nu_1}^2
&=b_{\nu_1}^2\{a_{\nu_1}u_\nu-(a_\nu\widehat{\nu}_1u)\}^2
\sim \nu\cdot a_\nu^2b_\nu^2-2a_\eta b_\eta(aHu)(bHu)-\frac16(asu)^2(bsu)^2\\
&\sim \nu\cdot a_\nu^2b_\nu^2-2a_\nu^2(aHu)(abu)-\frac16(asu)^2(bsu)^2+(\vartheta\eta x)^2(\theta Hu)^2.
\end{aligned}
\]
با جای‌گذاریِ مقدار $\theta_\eta H_\vartheta$، فرمول‌های زیر به دست می‌آیند:
\[
\tag{12}
\begin{aligned}
\text{a)}\quad H_\vartheta&\sim 2a_\nu\cdot a_\nu^2+2\nu\cdot\theta_\nu(\vartheta\nu x)^2+2f\cdot a_\eta(aHu)^2-2\theta_\eta,\\
\text{b)}\quad H_\vartheta^2&\sim (a_\nu^2)^2+2\theta_\eta^2+\frac23(\theta su)^2+\frac1{12}s\cdot\nu-\frac19 i f\cdot\nu-\frac16f\cdot(asu)^4.
\end{aligned}
\]

\noindent ۲) $\theta_\eta H_\vartheta$ بنا بر b):
\[
\begin{aligned}
a_\eta^2(aHu)(abu)
&\sim [a_\eta^2(aHu)]_{\widehat{bu}}+\frac12 f\cdot[a_\eta^2(aHu)^2]
=a_\eta b_\eta(aHu)(abu)\\
&\quad+a_\eta(\widehat{ab}\eta x)(abu)(aHu)-\theta_\eta(\vartheta\eta x)(\theta Hu)+\frac12\theta_\eta^2+\frac14(\nu\eta x)^2.
\end{aligned}
\]
از فرمول‌های (5.) و(11.) نتیجه می‌شود:
\[
\theta_\eta H_\vartheta\sim-\frac32\theta_\eta^2-\frac14(\theta su)^2-\frac1{12}s\cdot\nu+\frac29 i f\cdot\nu.
\]

\noindent ۳) $\theta_\eta H_\vartheta(\theta Hu)$ بنا بر b)، همانند $\theta_\eta H_\vartheta$، محاسبه می‌شود:
\[
\theta_\eta H_\vartheta(\theta Hu)\sim-\theta_\eta^2(\theta Hu)-\frac16(\theta su)^3.
\]

\noindent ۴) $\theta_\eta H_\vartheta(\vartheta\eta x)$ بنا بر b)؛ و$\theta_\eta^2(\vartheta\eta x)$ بنا بر d) (مقایسه کنید با \S\ ۹):
\[
\theta_\eta H_\vartheta(\vartheta\eta x)\sim-\theta_\eta^2(\vartheta\eta x)-\frac16s_\vartheta(\theta su)
\sim\frac13(\vartheta\R x)-\frac16s_\vartheta(\theta su),
\]
\[
\theta_\eta^2(\vartheta\eta x)\sim\frac23a_\eta^3(abu)\sim-\frac13(\vartheta\R x).
\]


\noindent ۵)
\[
\begin{array}{ccc}
\theta_\eta H_\vartheta^2\text{ بنا بر b)},&
\theta_\eta^2H_\vartheta\text{ بنا بر a)},&
\theta_\eta^2H_\vartheta\text{ بنا بر b)، و از آن }\theta_\eta^3\text{ بنا بر c)},\\
\text{از }a_\eta^2(aHb)(bHu);&
\text{از }a_\eta^2(Hab)(abu);&
\text{از }a_\eta^3(bHu)(abu):
\end{array}
\]
\[
\theta_\eta H_\vartheta^2\sim-\frac32\theta_\eta^2H_\vartheta-\frac14s_\vartheta(\theta su)^2+\frac16s_\nu-\frac49ia_\nu
\sim-\theta_\varrho-\frac16s_\nu-\frac19ia_\nu,
\]
\[
\theta_\eta^2H_\vartheta\sim\frac23\theta_\varrho-\frac16s_\vartheta(\theta su)^2+\frac29s_\nu-\frac29ia_\nu,
\]
\[
\theta_\eta^3H_\vartheta\sim\frac23\theta_\varrho-\frac23\theta_\eta^3+\frac5{36}s_\nu\quad\text{و از آن}
\]
\[
\theta_\eta^3\sim\frac14s_\vartheta(\theta su)^2-\frac18s_\nu+\frac13ia_\nu.
\]

\noindent ۶) $\theta_\eta^2(\theta Hu)^2$ بنا بر e):
\[
\A_\nu^2(a\A u)^4=[(a\A u)^4]_{\nu^2}=-\frac{12}{5}\theta_\eta^2(\theta Hu)^2-\frac3{10}\sigma-\frac1{20}(\theta su)^4+\nu\cdot\varrho
\quad(\text{بنا بر فرمول (3.) c}),
\]
\[
\A_\nu^2(a\A u)^4=[\A_\nu^2]^4_{\widehat{u a}}=\frac35\theta_\eta^2(\theta Hu)^2+\frac15\sigma-\frac3{10}(\theta su)^4+\frac13ij
\quad(\text{بنا بر فرمول (10.) a}),
\]
\[
\theta_\eta^2(\theta Hu)^2=-\frac16\sigma+\frac1{12}(\theta su)^4-\frac19ij+\frac13\nu\cdot\varrho.
\]

\noindent ۷) $\theta_\eta^2(\vartheta\eta x)^2$ بنا بر d) وe). از آن، کاهشِ فرم بالاترِ $g$ به دست می‌آید.

از محاسبه‌ای همانند محاسبهٔ \S\ ۹ نتیجه می‌شود:
\[
\begin{aligned}
\theta_\eta^2(\vartheta\eta x)^2
&=\frac12t\cdot f-\frac12a_\eta^2b_\eta-\frac1{12}g=\frac23t\cdot f-\frac23a_\eta^3b_\eta-\frac16g\\
&=-\frac16g-\frac13(\vartheta\varrho x)^2+\frac4{15}f\cdot a_\varrho^2+\frac8{15}f\cdot t\quad(\text{بنا بر فرمول (11.)}).
\end{aligned}
\]
از محاسبه‌ای متناظر با محاسبهٔ بالا نتیجه می‌شود:
\[
(a\A\widehat{\nu}x)^4=[(a\A u)^4]_{\widehat{\nu}x^4}=\frac{21}{10}\theta_\eta^2(\vartheta\eta x)^2+\frac7{10}s_\vartheta^2-\frac3{10}g
\quad(\text{بنا بر فرمول (3.) c}),
\]
\[
\begin{aligned}
(a\A\widehat{\nu}x)^4&=-\frac13i\A+6[\A_\nu^2]a^2-4[\A_\nu^3]a+\A_\nu^4\cdot f\\
&=-\frac{36}{5}\theta_\eta^2(\vartheta\eta x)^2-\frac95s_\vartheta^2-\frac35g-\frac35(\vartheta\varrho x)^2\\
&\quad+\frac53i\A+\frac{37}{25}f\cdot a_\varrho^2+\frac{74}{25}f\cdot t\quad(\text{بنا بر فرمول (10.)}).
\end{aligned}
\]
پس
\[
\begin{aligned}
\frac{93}{10}\theta_\eta^2(\vartheta\eta x)^2
&=-\frac3{10}g-\frac52s_\vartheta^2-\frac35(\vartheta\varrho x)^2\\
&\quad+\frac53i\A+\frac{37}{25}f\cdot a_\varrho^2+\frac{74}{25}f\cdot t,
\end{aligned}
\]
و با مقایسهٔ دو مقدارِ $\theta_\eta^2(\vartheta\eta x)^2$ بنا بر g):
\[
\tag{13.)\srcfnmark{*)}}
0=g-2s_\vartheta^2+2(\vartheta\varrho x)^2+\frac43i\A-\frac45f\cdot a_\varrho^2-\frac85f\cdot t.
\]
\srcfntext{*)}{این همان رابطه‌ای است که در حاشیهٔ مقدمه میان سه هموردای دارای مرتبهٔ ۶ و درجهٔ ۶ ذکر شد.}

\noindent ۸) $\theta_\eta^2H_\vartheta(\theta Hu)$ بنا بر a) وb)، و از آن $\theta_\eta^3(\theta Hu)$ بنا بر c) (فرم‌های دارای عامل $u_x$ ناپدید می‌شوند):
\[
\theta_\eta^2H_\vartheta(\theta Hu)=-\frac16s_\vartheta(\theta su)^3-\frac13(\nu\varrho x),
\]
\[
\theta_\eta^2H_\vartheta(\theta Hu)=-\frac13\theta_\eta^3(\theta Hu)-\frac13(\nu\varrho x),
\]
\[
\theta_\eta^3(\theta Hu)=\frac12s_\vartheta(\theta su)^3.
\]

\noindent ۹) $\theta_\eta^2H_\vartheta(\vartheta\eta x)$ بنا بر a) وb)، و از آن $\theta_\eta^3(\vartheta\eta x)$ بنا بر c). نیز $\theta_\eta^3(\vartheta\eta x)$ بنا بر d)، و از آن کاهشِ فرم بالاترِ $s_\vartheta^2(\theta su)$ بنا بر g) (فرم‌های دارای عامل $u_x$ ناپدید می‌شوند):
\[
\theta_\eta^2H_\vartheta(\vartheta\eta x)=\frac13(atu),
\]
\[
\theta_\eta^2H_\vartheta(\vartheta\eta x)=\frac13(atu)+\frac13\theta_\eta^3(\vartheta\eta x)-\frac16s_\vartheta^2(\theta su),
\]
\[
\theta_\eta^3(\vartheta\eta x)=\frac12s_\vartheta^2(\theta su),
\]
\[
\theta_\eta^3(\vartheta\eta x)=\frac25\theta_\varrho(\vartheta\varrho x)+\frac15(atu):
\]
\srcnumdisplay{(14.)}{%
s_\vartheta^2(\theta su)=\frac45\theta_\varrho(\vartheta\varrho x)+\frac25(atu).
}

\noindent ۱۰) $\theta_\eta^2H_\vartheta^2$ بنا بر a) و به‌وسیلهٔ کاهشگر $\theta_\nu^2(\vartheta\nu x)^2$، $\theta_\eta^3H_\vartheta$ بنا بر a) وc)، و$\theta_\eta^4$ بنا بر c)؛ و از آن، کاهشِ فرم‌های بالاترِ $s_\vartheta^2(\theta su)^2$ و$s_\nu^2$ بنا بر g):
\[
\begin{aligned}
\text{بنا بر a)}\quad
\theta_\eta^2H_\vartheta^2&=[a_\eta^2(aHu)^2]b^2+\frac13[a_\eta^4]_{bu^2}-\frac13H_\nu^2(\nu\eta x)^2\\
&=\frac49i\cdot a_\nu^2-\frac16s_\vartheta^2(\theta su)^2-\frac5{18}s_\nu^2+\frac13(atu)^2+\frac1{54}J\cdot u_x^2.
\end{aligned}
\]
\[
\begin{aligned}
\text{به‌وسیلهٔ }\theta_\nu^2(\vartheta\nu x)^2\quad
\theta_\eta^2H_\vartheta^2
&=[\theta_\nu^2(\vartheta\nu x)^2]_{\nu_1}+\frac13[\theta_\nu^4]_{\nu_1}\widehat{\nu}_1x^2-\frac13s_\vartheta^2(\theta su)^2\\
&=\frac49i\cdot a_\nu^2-\frac16s_\nu^2-\frac13s_\vartheta^2(\theta su)^2+\frac13(\nu\varrho x)^2.
\end{aligned}
\]
\[
\begin{aligned}
\text{بنا بر a)}\quad
\theta_\eta^3H_\vartheta&=-[a_\eta^2(aHu)]_{bu}b^2-\frac12[a_\eta^2(aHu)^2]b^2-\frac13[a_\eta^3]b+\frac1{12}[a_\eta^4]_{bu^2}\\
&\quad+\frac12u_x[a_\eta^2(aHu)^2]b^3\\
&=-\frac29ia_\nu^2+\frac1{12}s_\nu^2+\frac4{15}\theta_\varrho^2-\frac7{90}(\nu\varrho x)^2+\frac1{30}(atu)^2+\frac2{27}i^2\cdot u_x^2-\frac1{36}J\cdot u_x^2.
\end{aligned}
\]
\[
\begin{aligned}
\text{بنا بر c)}\quad
\theta_\eta^3H_\vartheta:\quad a_\eta^3(Hab)(bHu)&=\frac12(atu)^2
=\theta_\eta^2H_\vartheta(\theta Hu)(\vartheta\eta x)+\frac14H_\nu^2(\nu\eta x)^2\\
&\quad+\frac12\theta_\eta^2H_\vartheta^2
=\frac32\theta_\eta^2H_\vartheta^2+\theta_\eta^3H_\vartheta-u_x[a_\eta^2(aHu)^2]b^3\\
&\quad+\frac14H_\nu^2(\nu\eta x)^2,
\end{aligned}
\]
و بنابراین
\[
\begin{aligned}
\theta_\eta^3H_\vartheta&=-\frac23i\cdot a_\nu^2+\frac5{12}s_\nu^2+\frac12(\nu\varrho x)^2-\frac12(atu)^2\\
&\quad+\frac4{27}i^2\cdot u_x^2-\frac1{12}J\cdot u_x^2.
\end{aligned}
\]
بنا بر c)
\[
\theta_\eta^4=\frac49ia_\nu^2-\frac16s_\nu^2+\frac{16}{15}\theta_\varrho^2-\frac{14}{45}(\nu\varrho x)^2+\frac2{15}(atu)^2.
\]
بنا بر g)، از دو مقدارِ $\theta_\eta^2H_\vartheta^2$:
\par\smallskip\noindent
\begin{minipage}[c]{3.6em}
\centering (15.)
\end{minipage}%
\begin{minipage}[c]{\dimexpr\linewidth-3.6em\relax}
\centering
\(\displaystyle
\begin{aligned}
s_\vartheta^2(\theta su)^2
&=\frac23s_\nu^2-2(atu)^2+2(\nu\varrho x)^2-\frac19J\cdot u_x^2\\
&=\frac89i\cdot a_\nu^2+\frac8{15}\theta_\varrho^2-\frac{14}{15}(atu)^2+\frac{38}{45}(\nu\varrho x)^2-\frac4{27}i^2\cdot u_x^2.
\end{aligned}
\)
\end{minipage}
\par\smallskip
بنا بر g)، از دو مقدارِ $\theta_\eta^3H_\vartheta$:
\par\smallskip\noindent
\begin{minipage}[c]{3.6em}
\centering (16.)
\end{minipage}%
\begin{minipage}[c]{\dimexpr\linewidth-3.6em\relax}
\centering
\(\displaystyle
s_\nu^2=\frac43i\cdot a_\nu^2+\frac45\theta_\varrho^2+\frac85(atu)^2-\frac{26}{15}(\nu\varrho x)^2+\frac16J\cdot u_x^2-\frac29i^2\cdot u_x^2.
\)
\end{minipage}
\par\smallskip
فرمول (16.) مبنای کاهشِ پیمانهٔ $(ss'u)^4$ و فرمول (13.) مبنای کاهشِ دستگاهِ $s$ را به دست می‌دهد. (\S\ ۱۷.)

\noindent\emph{نتیجه‌ها:}

۱) عبارت‌های $a_\nu(j\nu x)^3$، $\theta_\eta H_\vartheta$، $\theta_\eta^2(\vartheta\eta x)$ و$\theta_\eta^2(\theta Hu)^2$ کاهشگرند؛ و$a_\nu(j\nu x)^2$، $H_\vartheta$ و$H_\vartheta^2$ فرم‌هایی تجزیه‌پذیرند.

۲) فرمِ دستگاه II، یعنی $j(\nu)=g=(\nu\eta x)^4H_x^2$، کاهش‌پذیر است.

۳) $\nu$ در تاخوردگی با دستگاه I اصلاً در توان‌ها یا در حاصل‌ضرب‌ها با فرم‌های همان دستگاه وارد نمی‌شود، زیرا تنها فرم پادهموردا، یعنی $g$، کاهش‌پذیر می‌شود.

$\theta$ تنها هنگامی که پادهموردا در نظر گرفته شود، در حاصل‌ضرب‌ها یا توان‌ها وارد تاخوردگی می‌شود؛ و به‌طور متناظر، $H$ هنگامی که هموردا در نظر گرفته شود (مقایسه کنید با \S\ ۱۳ B).
\begin{center}
\textbf{\S\ 12.}\quad\emph{فرم‌های مرتبهٔ هفتم.}
\end{center}
\[
\text{تاخوردگیِ }\left\{\begin{array}{l}
\theta\cdot K\text{ با }\nu,\\
K,j,\A\text{ با }H,\\
f\text{ با }L,\sigma.
\end{array}\right.
\]


\noindent\emph{فرم‌های کاهش‌ناپذیر:}
\[
\renewcommand{\arraystretch}{1.35}
\begin{array}{@{}r@{\quad}c|c|c|c@{}}
&\cdot&\cdot&(KHu)^2&\cdot\\
&\cdot&K_\eta&\cdot&K_\eta(KHu)^2\\
\text{B)}&(k\eta x)^2&\cdot&K_\eta^2&\cdot\\
&(k\eta x)^3&H_k(k\eta x)^2,\ K_\eta(k\eta x)^2&\cdot&\cdot\\
&\cdot&K_\eta(k\eta x)^3&\cdot&\cdot
\end{array}\, ,
\]
\[
\renewcommand{\arraystretch}{1.35}
\begin{array}{@{}r@{\quad}c@{\qquad}c@{\qquad}c@{\qquad\qquad}r@{\quad}c@{}}
\text{C)}&(j\eta x)&H_j&\cdot&\text{D)}&(\A Hu)\\
&\cdot&H_j(j\eta x)&H_j^2,&&\A_\eta,
\end{array}
\]
\[
\renewcommand{\arraystretch}{1.35}
\begin{array}{@{}r@{\quad}c@{\qquad}c@{\qquad}c@{\qquad}c@{}}
&\cdot&(aLu)^2&(aLu)^3&\cdot\\
\text{E)}&a_l&\cdot&a_l(aLu)^2&a_l(aLu)^3\\
&\cdot&a_l^2&\cdot&\cdot
\end{array}
\]
\noindent\emph{کاهش‌ها:}

\noindent A) عبارت $(\vartheta\cdot k,\nu,x)^4$ تجزیه می‌شود، زیرا از یک بار ترانس‌وکشن روی فرم تجزیه‌پذیرِ $(\vartheta^2\nu x)^4$ حاصل شده است.

\noindent B) رشتهٔ فرم‌های $H_kK_\eta$: کاهشگر $H_\vartheta\theta_\eta$.

رشتهٔ فرم‌های $K_\eta^2(KHu)$: کاهشگر $a_\eta^2(aHu)$.

رشتهٔ فرم‌های $K_\eta^2(k\eta x)$: کاهشگر $\theta_\eta^2(\vartheta\eta x)$.

کاهش دوگانهٔ رشتهٔ فرم‌های $K_\eta^3$ که از این راه حاصل می‌شود، به کاهش رشتهٔ فرم‌های بالاترِ $(j\eta x)^3$ می‌انجامد.

عبارت‌های $(KHu)$، $(k\eta x)$ و$K_\eta(k\eta x)$ بنا بر \S\ 3 فرم‌هایی تجزیه‌پذیرند.

عبارت‌های $H_k(KHu)$ و$K_\eta(KHu)$ تجزیه می‌شوند، زیرا از یک بار تاخوردگی با فرم تجزیه‌پذیرِ $K_\nu(k\nu x)$ و با دترمینان ژاکوبینِ زیر حاصل شده‌اند:
\[
K_\nu^2\equiv\theta_\nu^2(a\theta u)\equiv a_\nu^2(a\theta u)\modu{\nu^2}.
\]
با نمایش دوگانهٔ $K_\nu^2$ رابطه‌ای میان فرم‌های تجزیه‌پذیر متناظر است. داریم:
\[
H_k(KHu)=2K_{\nu_1}(\nu\nu_1k)=2\theta_{\nu_1}(\nu\nu_1\widehat{a\theta})
=2\theta_\nu^2\cdot a_\nu-a_\nu^2\cdot\theta_\nu+f\cdot\theta_\eta(\theta Hu)^2-f\cdot \nu\cdot\theta_\nu^3\quad(\text{قضیهٔ حاصل‌ضرب})\modu{(\nu^3)},
\]
\[
\begin{aligned}
K_\eta(KHu)&=-K_\nu^2(\nu\nu_1x)=-a_\nu^2(a\theta u)(\nu\nu_1x)=-\theta_\nu^2(a\theta u)(\nu\nu_1x)\\
&=a_\eta(aHu)^2\cdot\theta+a_\nu^2\cdot\theta_\nu
=-\theta_\eta(\theta Hu)^2\cdot f-\theta_\nu^2\cdot a_\nu.
\end{aligned}
\]
\par\smallskip\noindent
\begin{minipage}[c]{4.8em}
\centering (17.)\srcfnmark{*)}
\end{minipage}%
\begin{minipage}[c]{\dimexpr\linewidth-4.8em\relax}
\centering
\(\displaystyle
0=a_\nu^2\cdot\theta_{\nu_1}+\theta_\nu^2\cdot a_{\nu_1}+f\cdot\theta_\eta(\theta Hu)^2+\theta\cdot a_\eta(aHu)^2\modu{(\nu^3)}.
\)
\end{minipage}
\par\smallskip
\srcfntext{*)}{حاصل‌ضربِ فرم‌هایی که یکی از آن‌ها پادهمورداست، مانند $f\cdot\theta_\nu^3\cdot \nu$ در این مورد، در همهٔ روابط میان فرم‌های تجزیه‌پذیر حذف شده است؛ زیرا هنگام به‌کاربردن فرمول‌ها در دستگاه $s_{IV}$ نقشی ندارد.}
فرم‌های $H_k$، $H_k(k\eta x)$، $H_k^2$ و$H_k^2(k\eta x)$ تجزیه می‌شوند، زیرا از یک بار تاخوردگی با فرم‌های تجزیه‌پذیرِ $H_\vartheta$ و$H_\vartheta^2$ حاصل شده‌اند (مقایسه کنید با \S\ 3. 2)، a) وb)).

داریم
\[
H_k=H_\vartheta(a\theta u)\sim[H_\vartheta]_{\widehat{ua}}-f\cdot H_\vartheta(\theta Hu)\qquad(\text{تخصیصِ }y=\widehat{uv}\text{ از 2) a}).
\]
\[
H_k(k\eta x)=H_\vartheta a_\eta-H_\vartheta\theta_\eta\cdot f,
\qquad
H_\vartheta a_\eta=\frac54[H_\vartheta]a-\frac14H_\vartheta a_\vartheta
\quad(\S\ 3. 2)b),
\]
\[
H_\vartheta^2(a\theta u)=[H_\vartheta^2]_{\widehat{ua}},
\qquad
H_\vartheta^2a_\eta=[H_\vartheta^2]a=H_k^2(k\eta x)+H_\vartheta^2\theta_\eta\cdot f.
\]

\noindent C) رشتهٔ فرم‌های $(j\eta x)^3$: به‌وسیلهٔ کاهش دوگانهٔ رشتهٔ فرم‌های $K_\eta^3$ بنا بر B) کاهش‌پذیر است؛ با آغاز از عبارتِ زیر:
\[
\begin{aligned}
0&=a_\eta^3(a\theta u)=\theta_\eta^3(a\theta u)
=\{\theta_\eta+(\widehat{a\theta}\eta x)\}^3(a\theta u)-\theta_\eta^3(a\theta u)\\
&=\frac12(j\eta x)^3\modu{(\nu^3,s,\R,t,i,J)}.\qquad(\text{با نتیجهٔ \S\ 5 مقایسه کنید.})
\end{aligned}
\]
عبارتِ متناظر تا فرم پایانیِ رشتهٔ فرم‌ها برقرار است:
\[
0=H_\vartheta^2(\widehat{a\theta}\eta x)a_\eta^3
=H_\vartheta^2(\widehat{a\theta}\eta x)\theta_\eta^3=\frac12H_j^2(j\eta x)^4\modu{(\nu^3,s,\R,t,i,J)}.
\]
فرم $(j\eta x)^2$: تجزیه می‌شود 1) با دو بار قطبی‌سازیِ رابطهٔ فرم تجزیه‌پذیرِ $H_\vartheta^2$ ((12.) b)؛

2) با محاسبهٔ مستقیمِ حاصل‌ضربِ $a_\nu\cdot\theta_\nu^3$؛ از این راه فرم بالاترِ $(\A Hu)^2$ نیز تجزیه می‌شود. داریم:
\par\smallskip\noindent
\begin{minipage}[c]{4.8em}
\centering (18.)
\end{minipage}%
\begin{minipage}[c]{\dimexpr\linewidth-4.8em\relax}
\centering
\(\displaystyle
\left.\begin{array}{ll}
\text{a)}&\dfrac13(\A Hu)^2+\dfrac13(j\eta x)^2=\theta_\nu^2\cdot a_{\nu_1}^2,\\[0.6em]
\text{b)}&(j\eta x)^2=\dfrac43\theta_\nu^3\cdot a_{\nu_1},
\end{array}\right\}\modu{(\nu^3,s,\R,t,i,J)}
\)
\end{minipage}
\par\smallskip
با آغاز از عبارتِ زیر:
\[
\begin{aligned}
H_\vartheta^2(a\theta u)^2=\frac13(\A Hu)^2
&=2\theta_\eta^2(a\theta u)(aHu)+\theta_\nu^2\cdot a_{\nu_1}^2\quad(\text{قضیهٔ حاصل‌ضرب})\\
&=(a\theta u)(aHu)\{a_\eta-(\widehat{a\theta}\eta x)\}^2+\theta_\nu^2\cdot a_{\nu_1}^2,
\end{aligned}
\]
\[
\begin{aligned}
a_{\nu_1}\cdot\theta_\nu^3
&=-(a\widehat{\nu\nu_1}u)\theta_\nu^2(\theta\widehat{\nu\nu_1}u)-\frac12(a\widehat{\nu\nu_1}u)\theta_\nu\theta_{\nu_1}(\theta\widehat{\nu\nu_1}u)\\
&=-\frac32\theta_\eta^2(\theta Hu)(aHu)=-\frac32\theta_\eta^2(\theta Hu)(a\theta u)\\
&=-\frac32\{a_\eta-(\widehat{a\theta}\eta x)\}^2\{(aHu)-(a\theta u)\}(a\theta u).
\end{aligned}
\]

\noindent D) رشتهٔ فرم‌های $\A_\eta(\A Hu)$: کاهشگر $\A_\nu^2$.

عبارت $(\A Hu)^2$ بنا بر C) فرمی تجزیه‌پذیر است.

\noindent E) رشتهٔ فرم‌های $a_l^2(aLu)$: کاهشگر $a_\eta^2(aHu)$.
فرم‌های $(aLu)$ و$a_l(aLu)$: تجزیه می‌شوند، زیرا بنا بر \S\ 3 از ترانس‌وکشن روی دترمینان‌های ژاکوبین حاصل شده‌اند.

\noindent F) رشتهٔ فرم‌های $a_\sigma^3$: کاهشگر $a_\eta^3$.

فرم‌های $a_\sigma^2$ و$a_\sigma^3$: کاهشگرهای $a_\nu^3$ و$a_\eta^3$، به‌وسیلهٔ کاهش دوگانهٔ عبارت‌های متناظرِ
\[
H_\nu a_\nu^3\quad\text{و}\quad H_\nu a_\eta^3,
\qquad
H_\nu a_\nu^3a_\eta\quad\text{و}\quad H_\nu a_\eta^4
\]
(مقایسه کنید با \S\ 10. C)، کاهشِ $\A_\nu^2$ و$\A_\eta^3$).

از این راه، با در نظر گرفتن (16.)، کاهشِ فرم‌های بالاترِ $a_\nu s_\nu(asu)^2$ و$a_\nu^2s_\nu(asu)^2$ و نیز فرم‌ها بر حسب نمادهای $\R,t(a_\nu,a_\eta^2)$ به دست می‌آید:
\par\smallskip\noindent
\begin{minipage}[c]{4.8em}
\centering (19.)
\end{minipage}%
\begin{minipage}[c]{\dimexpr\linewidth-4.8em\relax}
\centering
\(\displaystyle
\left.\begin{array}{rcl}
a_\nu s_\nu(asu)^2&=&\dfrac13 i\theta_\nu^2-\dfrac1{18}iH,\\[0.5em]
a_\nu^2s_\nu(asu)^2&=&0,
\end{array}\right\}\modu{(\R,t)}
\)
\end{minipage}
\par\smallskip
فرم $a_\sigma$: با قطبی‌سازیِ (12.)b یا، به‌طور کوتاه‌تر، با بسط مستقیمِ حاصل‌ضربِ $a_\nu^2\cdot\theta_\nu^3$ تجزیه می‌شود؛ با آغاز از عبارتِ زیر:
\[
a_\nu^2\cdot\theta_{\nu_1}^3=a_\nu^2\theta_{\nu_1}\{a_{\nu_1}-(\widehat{a\theta}\nu_1x)\}^2
=-2a_\eta(aHu)(a\theta H)(\widehat{a\theta}\eta x)+(\widehat{a\theta}\nu_1x)^2a_\nu^2\theta_{\nu_1}
\]
و با در نظر گرفتن کاهشِ $H_j(j\eta x)^2$ (C):
\par\smallskip\noindent
\begin{minipage}[c]{4.8em}
\centering (20.)
\end{minipage}%
\begin{minipage}[c]{\dimexpr\linewidth-4.8em\relax}
\centering
\(\displaystyle a_\sigma=2a_\nu^2\cdot\theta_{\nu_1}^3\modu{(s,\R,t,i)}.\)
\end{minipage}
\par\smallskip

\noindent\emph{نتیجه‌ها:}

1) عبارت‌های $(j\eta x)^3$، $\A_\eta(\A Hu)$ و$a_\sigma^2$ کاهشگرند؛ و عبارت‌های $(j\eta x)^2$، $(\A Hu)^2$ و$a_\sigma$ فرم‌هایی تجزیه‌پذیرند.

2) فرم $f$ تنها در حاصل‌ضرب با فرم‌های پادهموردا وارد تاخوردگی با $L$ می‌شود؛ $f$ و بنابراین $K$ و$N$ نیز اصلاً وارد تاخوردگی با $\sigma$ و در نتیجه با $(\nu\sigma x)$ نمی‌شوند. فرم $j$ تنها در حاصل‌ضرب با فرم‌های پادهموردا، و $\A$ اصلاً نه در حاصل‌ضرب‌ها، وارد تاخوردگی با $H$ و$L$ می‌شوند (این امر از $\A_\eta(\A Hu)$ و$(\A Hu)^2$ نتیجه می‌شود). همچنین $\sigma$ و در نتیجه $(\nu\sigma x)$ در حاصل‌ضرب‌ها وارد تاخوردگی با هیچ فرمی از دستگاه I نمی‌شوند. با در نظر گرفتن نتیجه‌های \S\ 11، از حاصل‌ضرب‌های دستگاه II تنها توان‌های $H$ و حاصل‌ضرب‌های $H$ با $L$ باقی می‌مانند.

\begin{center}
\textbf{\S\ 13.}\quad\emph{فرم‌های مرتبهٔ هشتم (دستگاه III).}
\end{center}
\[
\text{تاخوردگیِ }\left\{\begin{array}{l}
\A\cdot j\text{ با }\nu,\\
\theta^2,
 f\cdot j,
 N\text{ با }H,\\
\theta\text{ با }L,\sigma.
\end{array}\right.
\]


\noindent\emph{فرم‌های کاهش‌ناپذیر:}
\[
\renewcommand{\arraystretch}{1.28}
\begin{array}{@{}r@{\quad}c@{\qquad\qquad}r@{\quad}c@{\qquad}c@{}}
\text{A)}&\A_\nu(j\nu x),&\text{B)}&(\vartheta^2\eta x)^3&\cdot\\
&&&(\vartheta^2\eta x)^4&H_\vartheta(\vartheta^2\eta x)^3,\ \theta_\eta(\vartheta^2\eta x)^3,
\end{array}
\]
\[
\renewcommand{\arraystretch}{1.28}
\begin{array}{@{}r@{\quad}c@{\qquad}c@{\qquad\qquad}r@{\quad}c@{}}
\text{C)}&(aHu)H_j&\cdot&\text{D)}&H_\eta;\ N_\eta,\\
&\cdot&a_\eta(aHu)H_j&&
\end{array}
\]
\[
\renewcommand{\arraystretch}{1.30}
\begin{array}{@{}r@{\quad}c@{\qquad}c@{\qquad}c@{\qquad}c@{\qquad}c@{}}
&\cdot&\cdot&(\theta Lu)^2&(\theta Lu)^3&\cdot\\
&\cdot&\theta_l&\cdot&L_\vartheta(\theta Lu)^2,\ \theta_l(\theta Lu)^2&\theta_l(\theta Lu)^3\\
\text{E)}&(\vartheta lx)^2&\cdot&\theta_l^2&\cdot&\cdot\\
&\cdot&\theta_l(\vartheta lx)^2&\cdot&\cdot&\cdot
\end{array}
\]
\noindent\emph{کاهش‌ها:}

\noindent A) $\A_\nu(j\nu x)^3$: کاهشگر $a_\nu(j\nu x)^3$.

فرم $\A_\nu(j\nu x)^2$ به‌وسیلهٔ فرم تجزیه‌پذیرِ $a_\nu(j\nu x)^2$ بنا بر \S\ 3. 2)b و با در نظر گرفتن کاهشگرِ $\theta_{\vartheta j}$ کاهش‌پذیر است. زیرا بنا بر \S\ 11 B) داریم:
\[
\frac3{10}\A_\nu(j\nu x)^2+\frac35a_\nu(j\nu\vartheta)a_\vartheta(j\nu x)
+\frac1{10}a_\nu(j\nu\vartheta)^2
=\frac35\theta_\nu^3\cdot\theta_{\vartheta\nu}^{\prime2}+\frac25\theta_\nu^3\theta_{\vartheta\nu}\theta_{\vartheta\nu}'.
\]
(کاهشگرها: $a_j$ و $\theta_{\vartheta j}$).

\noindent B) فرم‌های $H_\vartheta(\vartheta'\eta x)^2$، $H_\vartheta^2(\vartheta'\eta x)$ و$H_\vartheta^2(\vartheta'\eta x)^2$ تجزیه می‌شوند، زیرا از تاخوردگی‌های متوالی، هر بار یک‌مرتبه، با فرم‌های تجزیه‌پذیرِ $H_\vartheta$ و$H_\vartheta^2$، یا به‌ترتیب $H_\vartheta a_\eta$ و$H_\vartheta^2a_\eta$ حاصل شده‌اند؛ بنا بر \S\ 3. 2)b) و رابطهٔ زیر:
\[
H_\vartheta(\vartheta'\eta x)^2=2H_\vartheta a_\eta^2\cdot f-2H_\vartheta a_\eta b_\eta.
\]

\noindent C) رشتهٔ فرم‌های $a_\eta(j\eta x)^2$: کاهشگرها: $a_j$ و $(j\eta x)^3$.

فرم $a_\eta(j\eta x)$ تجزیه می‌شود: کاهشگرِ $a_j$ و یک بار تاخوردگی با فرم تجزیه‌پذیرِ $(j\eta x)^2$ بنا بر \S\ 3. 2)a (تخصیصِ $y=uv$).

فرم $a_\eta H_j$ تجزیه می‌شود: 1) با یک بار تاخوردگی با فرم تجزیه‌پذیرِ $a_\nu(j\nu x)^2$؛

2) با دو بار تاخوردگیِ ویژه با فرم تجزیه‌پذیرِ $(j\eta x)^2$؛ از این راه رابطه‌ای میان فرم‌های تجزیه‌پذیر به دست می‌آید.

از $a_\nu(j\nu x)^2$:
\[
0=\theta_\nu^3\cdot\theta_\nu+2a_\eta H_j\modu{(s,\R,t,i,J)}.
\]
از $(j\eta x)^2$:
\[
0=\theta_\nu^3\cdot\theta_\nu-\frac32a_\eta H_j\modu{(s,\R,t,i,J)}
\]
(حاصل‌ضرب‌های شامل پادهمورداها حذف شده‌اند)، با آغاز از عبارتِ زیر:
\[
0=\frac12a_\nu(abu)^2\theta_\nu^3
+\frac12a_\nu(abu)\theta_{\nu_1}^3(\theta bu)
-\frac34(\widehat{j\eta}bu)^2+\frac34H_j(\widehat{j\eta}bu)
-\frac18f\cdot H_j^2
\]
و با جابه‌جاییِ نمادها در $a_\nu(abu)(\theta bu)\theta_{\nu_1}^3$.

\noindent D) رشتهٔ فرم‌های $N_\eta(NHu)$: کاهشگر $\A_\nu^2$.

عبارت‌های $(NHu)$، $(n\eta x)$، $N_\eta(n\eta x)$ و$H_n(NHu)$ فرم‌هایی تجزیه‌پذیرند (مقایسه کنید با \S\ 12 B)؛ فرم $(NHu)^2$ با یک بار تاخوردگی با فرم $(\A Hu)^2$ تجزیه می‌شود.

\noindent E) رشته‌های فرم‌های $\theta_lL_\vartheta$, $\theta_l^2(\theta Lu)^2$, $\theta_l^2(\vartheta lx)$:

کاهشگرها: $\theta_\eta H_\vartheta$, $\theta_\eta^2(\theta Hu)^2$, $\theta_\eta^2(\vartheta\eta x)$.

فرم‌های $(\theta Lu)$، $(\vartheta lx)$، $L_\vartheta$، $L_\vartheta(\theta Lu)$، $\theta_l(\theta Lu)$، $L_\vartheta(\vartheta lx)$، $\theta_l(\vartheta lx)$، $L_\vartheta^2$، $L_\vartheta^2(\theta Lu)$ و$\theta_l^2(\theta Lu)$ تجزیه می‌شوند. (این‌ها به‌واسطهٔ دوگانی در برابر فرم‌های تجزیه‌پذیرِ \S\ 12 B) قرار می‌گیرند.)

\noindent F) رشتهٔ فرم‌های $\theta_\sigma(\vartheta\sigma x)$: کاهشگر $a_\sigma^2$.

فرم‌های $(\vartheta\sigma x)$ و$(\vartheta\sigma x)^2$ تجزیه می‌شوند، زیرا از یک بار تاخوردگی با فرم تجزیه‌پذیرِ $a_\sigma$ حاصل شده‌اند؛ فرم $\theta_\sigma$ نیز که از دو بار تاخوردگی حاصل شده است تجزیه می‌شود، زیرا رشتهٔ فرم‌ها رابطهٔ زیر را دارد:
\[
a_\nu^2(a\theta u)\theta_{\nu_1}^3\equiv H_j^2(j\eta x)\equiv0\;\text{به پیمانهٔ }(s,\R,t,i,J):
\]
\[
\theta_\sigma=a_\sigma(abu)^2=a_\nu^2(abu)^2\cdot\theta_\nu^3.
\]

\noindent\emph{نتیجه‌ها:}

عبارت $\theta^3$ یا $\theta^2K$ وارد تاخوردگی با $H$ نمی‌شود؛ $\theta$ تنها با در نظر گرفتن آن به‌عنوان پادهموردا، در توان‌ها یا حاصل‌ضرب‌ها، وارد تاخوردگی با $L$ می‌شود و اصلاً وارد تاخوردگی با $\sigma$ و$(\nu\sigma x)$ نمی‌شود.

فرم $N$ در حاصل‌ضرب‌ها وارد تاخوردگی با هیچ فرمی از دستگاه II نمی‌شود؛ همچنین با توان‌ها یا حاصل‌ضرب‌های فرم‌های دستگاه II نیز وارد تاخوردگی نمی‌شود. با در نظر گرفتن نتیجه‌های بندهای اخیر، از حاصل‌ضرب‌های دستگاه I تنها حاصل‌ضرب‌های $f\cdot j$ و$\A\cdot j$، توان‌های $\theta$ و حاصل‌ضرب‌های $\theta\cdot K$ باقی می‌مانند.

\begin{center}
\textbf{\S\ 14.}\quad\emph{فرم‌های مرتبهٔ نهم (دستگاه III).}
\end{center}
\[
\text{تاخوردگیِ }\left\{\begin{array}{l}
\theta\cdot K\text{ با }H,\\
K,j,\A\text{ با }L,\\
j,\A\text{ با }\sigma,\\
f\text{ با }H^2.
\end{array}\right.
\]


\noindent\emph{فرم‌های کاهش‌ناپذیر:}
\[
\renewcommand{\arraystretch}{1.35}
\begin{array}{@{}r@{\quad}c@{\qquad}c@{\qquad\qquad}r@{\quad}c@{\qquad}c@{\qquad}c@{\qquad}c@{}}
&\cdot&\cdot&&\cdot&\cdot&(KLu)^3&\cdot\\
\text{A)}&(\vartheta\cdot k,\eta,x)^4&\cdot&\text{B)}&\cdot&\cdot&\cdot&K_l(KLu)^3\\
&\cdot&H_k(\vartheta\cdot k,\eta,x)^4&&\cdot&K_l^2&\cdot&\cdot\\
&&&& (klx)^3&\cdot&\cdot&\cdot\\
&&&&\cdot&K_l(Klx)^3&\cdot&\cdot
\end{array}
\]
\[
\renewcommand{\arraystretch}{1.35}
\begin{array}{@{}r@{\quad}c@{\qquad}c@{\qquad}r@{\quad}c@{\qquad}r@{\quad}c@{\qquad}r@{\quad}c@{\qquad}c@{}}
\text{C)}&L_j&\cdot&\text{D)}&\A_l&\text{E)}&(j\sigma x),&\text{G)}&(aH^2u)^3&(aH^2u)^4\\
&\cdot&L_j^2&&&&&&\cdot&a_\eta(aH^2u)^3
\end{array}
\]
\noindent\emph{کاهش‌ها:}

A) فرم‌های $H_\vartheta(k\eta x)^3$، $H_\vartheta^2(k\eta x)^2$ و$H_\vartheta^2(k\eta x)^3$ تجزیه می‌شوند، زیرا از تاخوردگی‌های متوالی، هر بار یک‌مرتبه، با فرم‌های تجزیه‌پذیر حاصل شده‌اند.

B) رشته‌های فرم‌های $K_lL_k$, $K_l^2(KLu)$, $K_l^2(klx)$:

کاهشگرها: $\theta_\eta H_\vartheta$, $a_\eta^2(aHu)$, $\theta_\eta^2(\vartheta\eta x)$.

فرم‌های $K_l$، $K_l(KLu)$، $K_l(klx)$، $(KLu)^2$ و$(klx)^2$ تجزیه می‌شوند، زیرا ترانس‌وکشن‌هایی بر دترمینان‌های ژاکوبین‌اند (\S\ 3).

فرم‌های $L_k$, $L_k(KLu)$, $L_k(KLu)^2$, $L_k(klx)$, $L_k(klx)^2$, $L_k^2$, $L_k^2(KLu)$, $L_k^2(klx)$, $L_k^3$, $K_l(KLu)^2$, $K_l(klx)^2$ تجزیه می‌شوند، زیرا از یک بار تاخوردگی با فرم‌های تجزیه‌پذیرِ زیر حاصل شده‌اند:
\[
L_\vartheta,
\ H_k(KHu),
\ L_\vartheta(\vartheta lx),
\ H_k^2,
\ L_\vartheta^2,
\ L_\vartheta^2(\theta Lu),
\ K_\eta(KHu),
\ \theta_l(\vartheta lx)
\]
بنا بر \S\ 3. 2), a) و b).

C) رشتهٔ فرم‌های $(jlx)^3$: کاهشگرِ $(j\eta x)^3$.

فرم‌های $(jlx)^2$ و $L_j(jlx)$: از یک بار تاخوردگی با فرم تجزیه‌پذیرِ $(j\eta x)^2$ حاصل شده‌اند.

D) رشتهٔ فرم‌های $\A_l(\A Lu)$: کاهشگرِ $\A_\nu^2$.

فرم‌های $(\A Lu)^2$، $(\A Lu)^3$: یک بار تاخوردگی با فرم تجزیه‌پذیرِ $(\A Hu)^2$.

E) رشتهٔ فرم‌های $(j\sigma x)^2$: کاهشگرِ $a_\sigma^2$ با جایگزین کردن $j$ با فرم‌های پایین‌ترِ رشتهٔ فرم‌ها (\S\ 5):
\[
a_\sigma^2(a\theta u)^2=\frac13(j\sigma x)^2,
\qquad
 a_\sigma^2(a\theta\sigma x)^2=\frac13(j\sigma x)^4.
\]
F) رشتهٔ فرم‌های $\A_\sigma^2$: کاهشگرِ $a_\sigma^2$.

فرم $\A_\sigma$: کاهشگرِ $a_\sigma^2$ با جایگزین کردن $\A$ با فرم‌های پایین‌ترِ رشتهٔ فرم‌ها:
\[
a_\vartheta a_\sigma^2=\frac23\A_\sigma\modu{\nu}.
\]

\noindent\emph{نتیجه‌ها:} تنها فرم $j$ در تاخوردگی با $\sigma$ و $(\nu\sigma x)$ وارد می‌شود.

$N$ در تاخوردگی با $L$ وارد نمی‌شود، زیرا فرم متناظر با $\A_l$، یعنی $N_l$، تجزیه می‌شود.

\begin{center}
\textbf{\S\ 15.}\quad\emph{فرم‌های مرتبهٔ دهم (دستگاه III).}
\end{center}
\[
\text{تاخوردگیِ }\left\{\begin{array}{l}
\theta^2,
\ f\cdot j\text{ با }L,\\
\theta\text{ با }H^2.
\end{array}\right.
\]
\noindent\emph{فرم‌های کاهش‌ناپذیر:}
\[
\begin{array}{c@{\quad}l@{\qquad}c@{\quad}l}
\text{A)}&(\vartheta^2lx)^4,
\ \theta_l(\vartheta^2lx)^4&
\text{B)}&(aLu)^2L_j,
\ a_l(aLu)^2L_j,\\[0.8em]
\text{C)}&(\theta H^2u)^3,
\ (\theta H^2u)^4,
\ H_\vartheta(\theta H^2u)^3,
\ \theta_\eta(\theta H^2u)^3.
\end{array}
\]

\noindent\emph{کاهش‌ها:}

A) و B): تجزیهٔ دیگر فرم‌ها با یک بار تاخوردگی با فرم‌های تجزیه‌پذیر.

C) تجزیهٔ دیگر فرم‌ها بنا بر \S\ 13 B) با بررسی از سوی مقابل بر اساس دوگانی.

\noindent\emph{نتیجه‌ها:}

$\theta^2\cdot K$ در تاخوردگی با $L$ وارد نمی‌شود؛ به همین ترتیب، $H^2L$ در تاخوردگی با $K$ وارد نمی‌شود. $H^3$ و $H^2L$ در تاخوردگی با $\theta$ وارد نمی‌شوند. پس طرح تاخوردگیِ \S\ 7 اثبات شده است.

\begin{center}
\textbf{\S\ 16.}\quad\emph{فرم‌های مرتبه‌های 11، 12، 13 و 15 (دستگاه III).}
\end{center}

\noindent\emph{مرتبهٔ 11.}
\[
\text{تاخوردگیِ }\left\{\begin{array}{l}
\theta\cdot K\text{ با }L,\\
K,j\text{ با }H^2,\\
j\text{ با }(\nu\sigma x),\\
f\text{ با }H\cdot L.
\end{array}\right.
\]
\noindent\emph{فرم‌های کاهش‌ناپذیر:}
\[
\begin{array}{c@{\quad}l@{\qquad}c@{\quad}l}
\text{A)}&(\vartheta\cdot k,l,x)^5,
\ \theta_l(\vartheta\cdot k,l,x)^5&
\text{B)}&(KH^2u)^4,
\ K_\eta(KH^2u)^4,\\[0.8em]
\text{C)}&(\nu\sigma j),&
\text{D)}&(a,H\cdot L,u)^4.
\end{array}
\]

\noindent\emph{مرتبهٔ 12.}
\[
\text{تاخوردگیِ }\left\{\begin{array}{l}
\theta^3\text{ با }L,\\
\theta\text{ با }H\cdot L.
\end{array}\right.
\]
\noindent\emph{فرم‌های کاهش‌ناپذیر:}
\[
\text{E)}\ (\vartheta^3lx)^6,
\qquad
\text{F)}\ (\theta,H\cdot L,u)^4,
\qquad L_\vartheta(\theta,H\cdot L,u^4).
\]

\noindent\emph{مرتبهٔ 13.}
\[
\text{تاخوردگیِ }K\text{ با }H\cdot L.
\]
\noindent\emph{فرم‌های کاهش‌ناپذیر:}
\[
\text{G)}\ (K,H\cdot L,u)^5,
\qquad K_\eta(K,H\cdot L,u)^5.
\]

\noindent\emph{مرتبهٔ 15.}
\[
\text{تاخوردگیِ }K\text{ با }H_3.
\]
\noindent\emph{فرم کاهش‌ناپذیر:}
\[
\text{H)}\ (KH^3u)^6.
\]

\noindent\emph{کاهش‌ها:}

فرم‌های $H_{j^3}^2$, $H_{j^4}^2$. کاهشگرِ $L_j^3$ از رابطهٔ زیر:
\[
(\nu lx)L_x^3=-\frac12H^2\modu{(\sigma,s)}.
\]
تجزیه یا کاهش‌پذیریِ دیگر فرم‌ها بر پایهٔ بررسی‌های بندهای پیشین.

\noindent\emph{نتیجه‌ها:}

بر اساس طرح تاخوردگیِ \S\ 7 و نتیجه‌های \S\S\ 8 تا 15، بدین‌ترتیب همهٔ فرم‌های کاهش‌ناپذیرِ دستگاه III (117 فرم) به دست آمده‌اند.



\clearpage

\begin{center}
{\Large\bfseries فصل IV.\quad دستگاه کامل نسبی به پیمانهٔ $(\varrho,t)$.}
\end{center}

\begin{center}
\textbf{\S\ 17.}\quad\emph{کاهش پیمانهٔ $(ss'u)^4$ و دستگاهِ $s$.}
\end{center}

برای ساختن دستگاه کامل نسبیِ بالاترِ بعدی، باید به قیاسِ \S\ 7، دستگاهِ $s$ را نسبت به پیمانهٔ بعدی
\[
\nu(s)=(ss'u)^4
\]
ساخت و آن را با فرم‌های دستگاه کامل نسبی به پیمانهٔ $s$ تاخوردگی داد. نشان خواهیم داد که با وارد کردن پیمانهٔ $(\varrho,t)$ به‌عنوان پیمانه‌ای بالاتر،
\[
\begin{array}{ll}
\text{A)}&\text{پیمانهٔ }(ss'u)^4\text{ کاهش‌پذیر می‌شود،}\vphantom{\dfrac11}\\
\text{B)}&\text{دستگاهِ }s\text{ تنها از فرمِ }s\text{ تشکیل می‌شود.}
\end{array}
\]

A) کاهش پیمانهٔ $(ss'u)^4$ بی‌درنگ از کاهشگرِ $s_\nu^2$ به دست می‌آید که با فرمول (16.) در \S\ 11 یافته شد. از
\[
 s_\nu^2=\frac43 i a_\nu^2+\frac45\theta_\varrho^2+\frac85(atu)^2-
 \frac{26}{15}(\nu\varrho x)^2+\frac16J\cdot u_x^2-
 \frac29i^2\cdot u_x^2
\]
با گرفتن قطبی‌ها از $x$ به $\nu_1$ نتیجه می‌شود:
\srcnumdisplay{(21.)\srcfnmark{*)}}{%
\begin{aligned}
(ss'u)^4&=-\frac23J\cdot \nu+6s_\nu^2s_{\nu_1}^2\\
&=\frac{24}{5}\theta_\nu^2\theta_\varrho^2+
\frac{48}{5}a_\nu^2(atu)^2-
\frac{52}{5}H_\varrho^2+
\frac43i\cdot(asu)^4+\frac13J\cdot \nu-
\frac49i^2\cdot \nu,
\end{aligned}
}
و بدین‌ترتیب کاهش پیمانهٔ $(ss'u)^4$ حاصل می‌شود.

\srcfntext{*)}{با به‌کارگیری فرمول (10.)c، از فرمول (21.) رابطهٔ میان سه ناوردای مرتبهٔ 9 که در یادداشتِ مقدمه ذکر شد، نتیجه می‌شود:
\srcnumdisplay{(22.)}{%
\begin{aligned}
(ass')^4&=\frac{24}{5}\theta_\nu^2\theta_\varrho^2a_\nu^2a_\vartheta^2+
\frac{48}{5}a_\nu^2b_\nu^2(tab)^2-
\frac{52}{5}H_\varrho^2a_\eta^4+
\frac53J\cdot i-\frac49i^3,\\
(ass')^4&=\frac{24}{5}a_\varrho^2a_{\varrho'}^2-
\frac{12}{5}t_\varrho^2+\frac53J\cdot i-\frac49i^3.
\end{aligned}
}
}

B) کاهشِ
\[
Z=(ss'u)^2=\theta(s)
\]
و در نتیجه کاهش دستگاهِ $s$، با جایگزینیِ فرم $Z$ با فرم پایین‌ترِ $g_\nu^2$ طبق فرمول (8.)b در \S\ 7 و با در نظر گرفتن رابطهٔ زیر حاصل می‌شود:
\[
 s_\eta^2=s_\nu^2(\nu\nu_1x)^2-\frac13Z.
\]
فرم $g_\nu^2$ طبق فرمول (13.) در \S\ 11، کاهشگرِ $s_\nu^2$ را به‌عنوان عامل دارد. بنا بر فرمول‌های (8.) و (16.) داریم:
\[
 g_\nu^2=-Z+\frac{14}{5}ia_\eta^2+\frac{14}{15}i(asu)^2
 \modu{(\varrho,t)},
\]
بنا بر فرمول‌های (13.)، (14.)، (15.) و (16.):
\[
\begin{aligned}
g_\nu^2&=2[s_\vartheta^2]_{\nu^2}-\frac43 i\A_\nu^2
      =2s_\vartheta^2s_\nu^2-\frac65[s_\vartheta^2(\theta su)^2]\widehat{\nu x}^{2}
        -\frac43i\A_\nu^2\\
&=\frac45i\cdot a_\eta^2-\frac25i\cdot(asu)^2+\frac13J\cdot\theta
\modu{(\varrho,t)}.
\end{aligned}
\]
از این نتیجه می‌شود
\srcnumdisplay{(23.)}{%
Z=2i\cdot a_\eta^2+\frac43i\cdot(asu)^2-\frac13J\cdot\theta
 \modu{(\varrho,t)}.
}

پس دستگاه کامل نسبی به پیمانهٔ $((ss'u)^4,\varrho,t)$ به دستگاهی کامل نسبی به پیمانهٔ $(\varrho,t)$ تبدیل می‌شود، و از این دستگاه، با ترانس‌وکشن بر دستگاه شناخته‌شدهٔ دو فرم درجهٔ دو، دستگاه کامل مطلق حاصل می‌شود. برای ساختن دستگاه کامل نسبی به پیمانهٔ $(\varrho,t)$، از آنجا که طبق فرمول (23.) دستگاهِ $s$ به فرم $s$ کاهش می‌یابد، باید دستگاه کامل نسبی به پیمانهٔ $(s,\varrho,t)$ را بر $s$ و توان‌های $s$ ترانس‌وکشن دهیم. خواهیم دید که توان‌های $s$ تنها با دستگاه I وارد تاخوردگی می‌شوند.

در میان فرم‌های هم‌مرتبه و هم‌درجه، \glqq{}فرم‌های بالاتر\grqq{} را فرم‌هایی تعریف می‌کنیم که از فرم‌های بالاترِ دستگاه کامل نسبی به پیمانهٔ $(s,\varrho,t)$ با تاخوردگی با $s$ حاصل شده‌اند؛ در اینجا تاخوردگی با $s$ را تنها گرفتن قطبی‌های فرم‌های اصلی در نظر می‌گیریم. بدین‌ترتیب ترتیب تک‌تک فرم‌ها به‌طور یکتا تعیین می‌شود (نگاه کنید به \S\ 1. رشته‌های فرم‌ها 2). برای مثال:
\[
(\theta_\eta(\theta Hu),s,u)^2\text{ بالاتر از }s_\eta((\theta Hu)^2,s,u)\srcfnmark{*)},
\]
\[
(\theta_\eta(\theta Hu)^2,s,u)\text{ بالاتر از }(\theta_\eta(\theta Hu),s,u)^2.
\]

دستگاه حاصل را به چهار زیردستگاه تقسیم می‌کنیم:
\[
\begin{array}{ll}
\text{دستگاه }s_I:&\text{حاصل از تاخوردگیِ }s\text{ و توان‌های }s\text{ با دستگاه I،}\\
\text{دستگاه }s_{II}:&\text{حاصل از تاخوردگیِ }s\text{ با دستگاه II،}\\
\text{دستگاه }s_{III}:&\text{حاصل از تاخوردگیِ }s\text{ با تک‌تک فرم‌های}\\
&\text{(بدون حاصل‌ضرب‌ها) از دستگاه III و با حاصل‌ضرب‌های}\\
&\text{یک فرم از هر یک از دستگاه‌های I و II،}\\
\text{دستگاه }s_{IV}:&\text{حاصل از تاخوردگیِ }s\text{ با حاصل‌ضرب‌های یک فرم از هر یک از جفت دستگاه‌های}\\
&\text{I و III، II و III، و III و III.}
\end{array}
\]
همچنین توجه شود که از این پس همهٔ فرمول‌ها به پیمانهٔ $(\varrho,t)$، و از فرمول (27.) به بعد به پیمانهٔ $(\varrho,t,i,J,\text{ فرم‌های بالاتر})$ فهمیده می‌شوند، بی‌آنکه هر بار این امر صریحا ذکر شود.

برای کاهش، فرمول‌های به‌دست‌آمده در گذشته را گرد می‌آوریم:
\srcnumdisplay{(24.)}{%
\begin{aligned}
s_\nu^2&=\frac43i\,a_\nu^2+\frac16J\cdot u_x^2-\frac29i^2\cdot u_x^2,\\
s_\nu^3&=\frac13J\cdot u_x,
& s_\nu^4&=J,\\
g&=2s_\vartheta^2-\frac43i\Delta,\\
s_\vartheta^2(\theta su)&=0,\\
s_\vartheta^2(\theta su)^2&=\frac89i\cdot a_\nu^2-\frac4{27}i^2\cdot u_x^2,\\
Z&=2i\cdot a_\eta^2+\frac43i(asu)^2-\frac13J\cdot\theta,\\
g_\nu&=-s_\eta+u_x\left\{2ia_\eta^2-\frac23i(asu)^2+\frac23J\cdot\theta\right\},\\
g_\nu^2&=\frac45ia_\eta^2-\frac25i(asu)^2+\frac13J\cdot\theta,
&g_\nu^3&=0,\quad g_\nu^4=0.
\end{aligned}
}
\srcfntext{*)}{برای اختصار در نگارش، جملهٔ ترانس‌وکشن را به‌جای ترانس‌وکشن کاملِ $[(\theta Hu)^2]_{\widehat{su}}s$ به‌عنوان فرمی از دستگاه وارد می‌کنیم. هنگام تنظیم فرمول‌ها، همهٔ ترانس‌وکشن‌ها با جمله‌های آغازینشان جایگزین می‌شوند؛ برای مثال $(\theta_\eta(\theta Hu),s,u)^2$ با $\theta_\eta(\theta Hu)(\theta su)^2$.}

\noindent\emph{نتیجه‌ها:}
\[
 s_\nu^2,\quad s_\vartheta^2(\theta su),\quad s_\vartheta^2s_\nu,\quad s_\vartheta^2\theta_\nu
\]
کاهشگر هستند.

\begin{center}
\textbf{\S\ 18.}\quad\emph{دستگاه $s_I$.}
\end{center}

\[
\text{تاخوردگیِ }\left\{\begin{array}{l}
s\text{ با } f;\ \theta;\ K,j,\Delta;\ f\cdot j,N;\ \Delta\cdot j,\\
s^2\text{ با }\theta,K.
\end{array}\right.
\]

\noindent\emph{فرم‌های کاهش‌ناپذیر.}

\noindent مرتبهٔ ۵:
\[
(asu);\qquad (asu)^2;\qquad (asu)^3;\qquad (asu)^4.
\]

\noindent مرتبهٔ ۶:
\[
\begin{array}{cccc}
(\theta su)&(\theta su)^2&(\theta su)^3&(\theta su)^4\\
s_\vartheta&s_\vartheta(\theta su)&s_\vartheta(\theta su)^2&s_\vartheta(\theta su)^3\\
\cdot&s_\vartheta^2&\cdot&\cdot
\end{array}
\]

\noindent مرتبهٔ ۷:
\[
\begin{array}{c@{\quad}c@{\quad}c@{\quad}c@{\qquad}l}
\cdot&(Ksu)^2&(Ksu)^3&(Ksu)^4&\\
\text{A)}\ s_k&s_k(Ksu)&s_k(Ksu)^2&s_k(Ksu)^3&
\text{B)}\ s_j,\\
\cdot&s_k^2&\cdot&\cdot&
\text{C)}\ (\Delta su),\ (\Delta su)^2.
\end{array}
\]

\noindent مرتبهٔ ۸:
\[
\begin{array}{ll}
\text{A)}\ s_j(asu),\ s_j(asu)^2,\ s_j(asu)^3,
&
\text{B)}\ \begin{array}[t]{cc}
\cdot&(Nsu)^2\\
s_n&s_n(Nsu).
\end{array}
\end{array}
\]

\noindent مرتبهٔ ۱۰:
\[
\text{A)}\ s_j(\Delta su),\qquad \text{B)}\ s_\vartheta(\theta s'u)^4.
\]

\noindent مرتبهٔ ۱۱:
\[
\begin{array}{cc}
(Ks^2u)^6&\cdot\\
s_k(Ks^2u)^5&s_k(Ks^2u)^6.
\end{array}
\]

\noindent\emph{کاهش‌ها:}

رشته‌های فرم‌ها
\[
 s_\vartheta^2(\theta su),\quad s_k^2(Ksu),\quad s_j^2,\quad (\Delta su)^3,\quad (Nsu)^3
\]
کاهشگرِ $s_\vartheta^2(\theta su)$ را به‌عنوان عامل دارند؛ $s_j^2$ و$(\Delta su)^3$ نیز با بازگرداندن $j$ و$\Delta$ به فرم‌های پایین‌ترِ رشتهٔ فرم‌ها به دست می‌آیند:
\[
\begin{aligned}
s_\vartheta^2(\theta su)(a\theta u)^3&=-\frac12s_j^2,\\
s_\vartheta^2(\theta su)(a\theta u)^2&=\frac13(\Delta su)^3,\\
s_\vartheta^2(\theta su)^2(a\theta u)^2&=\frac13(\Delta su)^4+\frac13s_j^2.
\end{aligned}
\]
فرم‌های $s_\vartheta^2s_{\vartheta'}$ و$s_\vartheta^2s_{\vartheta'}^2$: کاهشگرهای $s_\vartheta^2(\theta su)$ و$\theta_{\vartheta'}$:
\[
\begin{aligned}
s_\vartheta^2s_{\vartheta'}&=s_\vartheta^2\theta_{\vartheta'}-s_\vartheta^2(\theta s\widehat{\vartheta'}x),\\
s_\vartheta^2s_{\vartheta'}^2&=s_\vartheta^2s_{\vartheta'}\theta_{\vartheta'}-s_\vartheta^2s_{\vartheta'}(\theta s\widehat{\vartheta'}x).
\end{aligned}
\]
رشتهٔ فرم‌های $s_j(\Delta su)^2$: کاهشگرهای $\Delta_j$ و$(\Delta su)^3$.

\noindent\emph{نتیجه‌ها:}

$s$ با توان‌هایش وارد تاخوردگی با $f,j,\Delta,N$ نمی‌شود. $f,j,\Delta,N$ فقط در حاصل‌ضرب با فرم‌های پادهموردا وارد تاخوردگی می‌شوند؛ اما این فرم‌ها در حاصل‌ضرب با $f$ می‌توانند هنوز درجهٔ دو نسبت به $x$، و در حاصل‌ضرب با $\Delta$ هنوز درجهٔ یک نسبت به $x$ داشته باشند؛ یعنی باید حاصل‌ضرب‌های زیر از فرم‌ها را در نظر گرفت:
\[
 f\cdot(0,\lambda),\quad f\cdot(1,\lambda),\quad f\cdot(2,\lambda),\quad
 j\cdot(\mu,0),\quad N,\quad
 \Delta\cdot(0,\lambda),\quad \Delta\cdot(1,\lambda)\quad(\lambda>0),
\]
که در آن $(\mu,\lambda)$ فرمی با درجهٔ $\mu$ نسبت به $x$ و درجهٔ $\lambda$ نسبت به $u$ است.

$\theta$ یا $K$، در توان‌ها یا حاصل‌ضرب‌ها با فرم‌های دلخواه، وارد تاخوردگی با $s$ یا $s^2$ نمی‌شوند. برای حاصل‌ضرب‌ها با دترمینان ژاکوبین $K$، یعنی $K\cdot\varphi_x$ که $\varphi\ne f$ یا $\theta$، رابطهٔ زیر میان فرم‌های تجزیه‌پذیر برقرار است:
\[
K\cdot\varphi_x=(a\theta u)\varphi_x=f\cdot(\varphi\theta u)-\theta\cdot(\varphi a u).
\]
بدین‌ترتیب طرح بالا برای تاخوردگیِ دستگاه $s_I$ اثبات شد.

\begin{center}
\textbf{\S\ 19.}\quad\emph{دستگاه $s_{II}$.}
\end{center}

تاخوردگیِ $s$ با $\nu;\ H;\ L;\ H^2$.

\noindent\emph{فرم‌های کاهش‌ناپذیر:}
\[
\begin{array}{c@{\qquad}c@{\qquad}c@{\qquad}c}
\text{مرتبهٔ ۶} & \text{مرتبهٔ ۸} & \text{مرتبهٔ ۱۰} & \text{مرتبهٔ ۱۲}\\[0.3em]
s_\nu & (Hsu),\ (Hsu)^2 & (Lsu)^2,\ (Lsu)^3 & (H^2su)^3\\
\cdot & s_\eta & s_l & \cdot\\
\cdot & \cdot & \cdot & \cdot\\
s_\nu^4=J & \cdot & \cdot & \cdot
\end{array}
\]

\noindent\emph{کاهش‌ها:}

رشته‌های فرمِ $s_\eta(Hsu)$ و$s_l(Lsu)$: کاهشگر $s_\nu^2$.

رشتهٔ فرم‌های $s_\sigma$: کاهشگر $s_\nu^2$ حاصل از $H$، که $s_\nu^2=\dfrac23s_\sigma$.

تجزیه‌پذیریِ $(H^2su)^4$ از فرمول (7.)a نتیجه می‌شود (نگاه کنید به \S\ 11. A). از این، تجزیه‌پذیریِ $(H\cdot L,s,u)^4$ نتیجه می‌شود.

\noindent\emph{نتیجه‌ها:}

$s^2$ وارد تاخوردگی با دستگاه II نمی‌شود.

$\nu$ فقط در حاصل‌ضرب با فرم‌های پادهموردا وارد تاخوردگی می‌شود؛ $H$ نیز، هنگامی که هموردا در نظر گرفته شود، فقط در حاصل‌ضرب با فرم‌های $(1,\lambda)$، $(2,\lambda)$ و$(3,\lambda)$، که در آن $\lambda$ دلخواه است، وارد تاخوردگی می‌شود.

$\sigma$ و$(\nu\sigma x)$ به‌هیچ‌وجه وارد تاخوردگی نمی‌شوند، و$L$ نیز، به‌عنوان دترمینان ژاکوبین، در حاصل‌ضرب‌ها وارد تاخوردگی نمی‌شود (زیرا ترانس‌وکشن‌های کامل بر همورداهای دستگاه III کاهش‌پذیر می‌شوند).

از حاصل‌ضرب‌های دستگاه‌های I وII موارد زیر باقی می‌مانند:
\[
f\cdot \nu,\qquad \Delta\cdot \nu.
\]

\begin{center}
\textbf{\S\ 20.}\quad\emph{فرم‌های مرتبهٔ ۷ (دستگاه $s_{III}$).}
\end{center}

تاخوردگیِ $s$ با $f\cdot \nu$، $a_\nu$ و$a_\nu^2$.

\noindent\emph{فرم‌های کاهش‌ناپذیر:}
\[
\begin{gathered}
s_\nu(asu),\quad a_\nu(asu),\quad s_\nu(asu)^2,\quad a_\nu(asu)^2,\quad
s_\nu(asu)^3,\quad a_\nu(asu)^3,\\
a_\nu s_\nu,\qquad a_\nu s_\nu(asu),\qquad a_\nu^2(asu),\qquad a_\nu^2(asu)^2.
\end{gathered}
\]

\noindent\emph{کاهش‌ها:}

کاهش‌های بدیهیِ حاصل از تاخوردگی مستقیم با کاهشگرهای $s_\nu^2$ و$s_\eta(Hsu)$ ذکر نمی‌شوند؛ همچنین تجزیه‌پذیریِ ترانس‌وکشن‌ها با دترمینان‌های ژاکوبین نیز ذکر نمی‌شود.

فرم‌های $a_\nu s_\nu(asu)^2$ و$a_\nu^2s_\nu(asu)^2$: فرمول (19.) و \S\ 12 را ببینید؛
\[
a_\nu s_\nu(asu)^3=a_\nu(a\widehat{\nu_1}\nu_2u)^3(\nu_1\nu_2\nu)=0.
\]
رشتهٔ فرم‌های $a_\nu^2s_\nu$: کاهشگر $s_\vartheta^2(\theta su)$ با بازگرداندن $a_\nu^2$ به فرم‌های پایین‌تر، یعنی با کاهش دوگانهٔ عبارت‌های
\[
s_\vartheta^2(a\theta s)(a\theta u),\qquad
s_\vartheta^2(a\theta s)(a\theta u)^2
\]
بنا بر صورتِ مفروض زیر:
\[
\begin{aligned}
s_\vartheta^2(a\theta s)(a\theta u)
&=[(a\theta u)^2]s^3+\frac15[a_\vartheta(a\theta u)^2]s^2+\frac1{60}[a_\vartheta^2(a\theta u)^2]s\\
&=\frac12a_\nu^2s_\nu-\frac23a_\nu s_\nu^2,\\
s_\vartheta^2(a\theta s)(a\theta u)^2
&=\frac45[a_\vartheta(a\theta u)^2]_{\widehat{us}}s^2+\frac{23}{180}[a_\vartheta^2(a\theta u)^2]_{\widehat{us}}s\\
&=-\frac12a_\nu^2s_\nu(asu).
\end{aligned}
\]
در این‌جا $a_\nu^2b_\nu(abu)=0$ برقرار است.

فرمول‌های زیر به دست می‌آیند:
\srcnumdisplay{(25.)}{%
\begin{array}{rcl@{\qquad}rcl}
\underline{a_\nu^2s_\nu}&=&\dfrac43i\cdot a_\nu^2b_\nu,
& \underline{a_\nu s_\nu(asu)^2}&=&\dfrac13i\cdot\theta_\nu^2-\dfrac1{18}iH,\\[0.7em]
\underline{a_\nu s_\nu(asu)^3}&=&0,
& \underline{a_\nu^2s_\nu(asu)}&=&0,\\[0.7em]
\underline{a_\nu^2s_\nu(asu)^2}&=&0.
\end{array}
}

\noindent\emph{نتیجه‌ها:}

1) $a_\nu s_\nu(asu)^2$ و$a_\nu^2s_\nu$ کاهشگرند.

2) مقادیر فرم‌های $(\Delta su)^3$ و$(\Delta su)^4$ که در \S\ 19 کاهش‌پذیر شناخته شدند، به دست می‌آیند؛ همین‌طور مقادیر رشتهٔ فرم‌های متناظر $(agu)^2$ که در تاخوردگی با $\nu$، به‌وسیلهٔ کاهشگر $g_\nu$، موجب کاهش دوگانه می‌شود.
\srcnumdisplay{(26.)}{%
\begin{aligned}
\text{a)}\quad (\Delta su)^3&=-\frac65a_\nu^2(asu)+3a_\nu s_\nu(asu)+2i\theta_\nu(\vartheta \nu x),\\
\text{b)}\quad (\Delta su)^4&=-\frac25a_\nu^2(asu)^2+\frac43i\theta_\nu^2+\frac19iH.
\end{aligned}
}

از فرمول‌های (24.) و(26.)، به پیمانهٔ $(i,J)$ (مقایسه کنید با ص. ۷۱):
\srcnumdisplay{(27.)}{%
\begin{aligned}
\text{a)}\quad (agu)^2&=\frac23(\Delta su)^2-\frac43a_\nu s_\nu+\frac35s\cdot a_\nu^2,\\
\text{b)}\quad (agu)^3&=2a_\nu s_\nu(asu)-a_\nu^2(asu)^2,\\
\text{c)}\quad (agu)^4&=-a_\nu^2(asu)^2.
\end{aligned}
}


\begin{center}
\textbf{\S\ 21.}\quad\emph{فرم‌های مرتبهٔ ۸ (دستگاه $s_{III}$).}
\end{center}

تاخوردگیِ $s$ با فرم‌های مرتبهٔ ۴ (دستگاه III). (\S\ 9.)

\noindent\emph{فرم‌های کاهش‌ناپذیر:}
\[
\begin{array}{c@{\quad}c|c|c}
\cdot
& (\vartheta \nu x)s_\nu,\ ((\vartheta \nu x)^2,s,u)
& ((\vartheta \nu x),s,u)^2,\ (\theta_\nu su)
& \bigstar\\[0.45em]
& (\vartheta \nu x)^2s_\nu
& ((\vartheta \nu x)^2,s,u)^2,\ \theta s_\nu
& \\[0.45em]
&& ((\vartheta \nu x)^2,s,u)s_\nu,\ (\theta_\nu(\vartheta \nu x)^2,s,u)
&
\end{array}
\]
\[
\begin{array}{c|c|c|c}
\bigstar
& ((\vartheta \nu x),s,u)^3,\ (\theta_\nu su)^2
& ((\vartheta \nu x),s,u)^4,\ (\theta_\nu su)^3
& \cdot\\[0.45em]
& ((\vartheta \nu x)^2,s,u)^3,\ (\theta_\nu(\vartheta \nu x),s,u)^2,\ (\theta_\nu^2su)
& (\theta_\nu(\vartheta \nu x),s,u)^3,\ (\theta_\nu^2su)^2
& (\theta_\nu(\vartheta \nu x),s,u)^4\\[0.45em]
& \cdot
& (\theta_\nu^3su)
& \cdot
\end{array}
\]

\noindent\emph{کاهش‌ها:}

$((\vartheta \nu x),s,u)s_\nu$ بنا بر \S\ 3 به‌صورت ترانس‌وکشنی بر دترمینان‌های ژاکوبین تجزیه‌پذیر است.

رشتهٔ فرم‌های $((\vartheta \nu x),s,u)^2s_\nu$: کاهشگرهای $s_\nu^2$ و$s_\vartheta^2\theta_\nu$:
\[
(\widehat{\vartheta \nu su})s_\nu(\theta su)=s_\nu s_\vartheta(\theta su)=-\frac12(\widehat{\vartheta \nu su})^2(\theta su),
\]
\[
\theta_\nu(\widehat{\vartheta \nu su})s_\nu=\theta_\nu s_\nu s_\vartheta=-\frac12\theta_\nu(\widehat{\vartheta \nu su})^2.
\]

\noindent(قاعدهٔ ضرب.)

رشتهٔ فرم‌های $\theta_\nu s_\nu(\theta su)$: کاهشگر $a_\nu s_\nu(asu)^2$ با کاهش دوگانهٔ عبارت‌های
\[
a_\nu s_\nu(asu)^2(abu)\quad\text{و}\quad a_\nu s_\nu(asu)^2(abu)(sbu);
\]
$\theta_\nu s_\nu(\theta su)^3$ نیز با کاهش دوگانهٔ $(agu)^3(abu)(bgu)^3$ کاهش می‌یابد.

رشتهٔ فرم‌های $\theta_\nu(\vartheta \nu x)s_\nu$: کاهشگر $a_\nu^2s_\nu$ از رابطهٔ $\theta_\nu(\vartheta \nu x)s_\nu=a_\nu^2(abu)s_\nu$.

رشتهٔ فرم‌های $(\theta_\nu(\vartheta \nu x)^2,s,u)$: کاهش دوگانهٔ فرم‌های رشتهٔ $\theta_\nu(\widehat{\vartheta \nu su})s_\nu$ که هم‌زمان به رشته‌های فرمِ $((\vartheta \nu x)su)^2s_\nu$ و$\theta_\nu(\vartheta \nu x)s_\nu$ تعلق دارند.

فرم $(\theta_\nu(\vartheta \nu x),s,u)$ به‌صورت یک ترانس‌وکشن منفرد بر دترمینان ژاکوبین $\theta_\nu(\vartheta \nu x)=a_\nu^2(abu)$ تجزیه‌پذیر است.

فرم $((\vartheta \nu x)^2,s,u)^4$: کاهش دوگانهٔ عبارت $(a\A u)^2(\A su)^4$، یا نیز $(\theta gu)^4$ بر پایهٔ فرمول‌های (27.) و به قیاس با کاهش $(j\eta x)^4$ (\S\ 12 C).

از کاهش دوگانه فرمول‌های زیر به دست می‌آیند:
\srcnumdisplay{(28.)}{%
\begin{aligned}
0&=\theta_\nu s_\nu(\theta su)-\frac16(\widehat{\vartheta \nu su})^2(\theta su)-\frac1{12}(Hsu),\\
0&=\theta_\nu s(\theta su)^2-\frac14(\widehat{\vartheta \nu su})^2(\theta su)^2-\frac1{24}(Hsu)^2,\\
0&=(\widehat{\vartheta \nu su})^2(\theta su)^2+2\theta_\nu(\widehat{\vartheta \nu su})(\theta su)^2
  +3\theta_\nu^2(\theta su)^2-\frac13H(su)^2,\\
0&=\theta_\nu s_\nu(\theta su)^3+\frac12\theta_\nu(\widehat{\vartheta \nu su})(\theta su)^3,\\
0&=\theta_\nu s_\nu s_\vartheta-\frac14s_\eta,
\quad 0=\theta_\nu s_\nu s_\vartheta(\theta su),
\quad 0=\theta_\nu s_\nu s_\vartheta(\theta su)^2.
\end{aligned}
}
مورد اخیر متناظر است با $((\theta_\nu(\vartheta \nu x)^2,s,u)^2,\ldots)$.

\noindent\emph{نتیجه‌ها:}

1) $\theta_\nu(\vartheta \nu x)s_\nu$، $(\widehat{\vartheta \nu su})(\theta su)s_\nu$، $\theta_\nu(\widehat{\vartheta \nu su})^2$ و$(\widehat{\vartheta \nu su})^2(\theta su)^2$ کاهشگرند.

2) فرم‌های مرتبهٔ چهارمِ $(5,\lambda)$، $(6,\lambda)$ و مانند آن‌ها وارد تاخوردگی با $s^2$ نمی‌شوند.

3) برای رشتهٔ فرم‌های $(\theta gu)^2$، بنا بر (27.) و با در نظر گرفتن کاهش‌های این بخش، فرمول‌های زیر به دست می‌آیند:
\srcnumdisplay{(29.)}{%
\begin{aligned}
\text{a)}\quad (\theta gu)^2&=\frac43\theta_\nu s_\nu+\frac23s_\nu s_\vartheta-\frac13s\cdot\theta_\nu^2-\frac19s\cdot H-\frac13f\cdot a_\nu^2(asu)^2+\frac13a_\nu^2\cdot(asu)^2,\\
\text{b)}\quad (\theta gu)^3&=-\frac13(\widehat{\vartheta \nu su})^2(\theta su)-\theta_\nu(\widehat{\vartheta \nu su})(\theta su)-\theta_\nu^2(\theta su),\\
\text{c)}\quad (\theta gu)^4&=2\theta_\nu^2(\theta su)^2-\frac13(Hsu)^2,\\
g\vartheta(\theta gu)&=(\widehat{\vartheta \nu su})(\vartheta \nu x)s_\nu-\theta_\nu(\vartheta \nu x)(\widehat{\vartheta \nu su}),\\
g\vartheta(\theta gu)^2&=\frac23s\cdot\theta_\nu^3,
\quad g\vartheta(\theta gu)^3=\frac13\theta_\nu^3(\theta su),
\quad g\vartheta(\theta gu)^4=0.
\end{aligned}
}

\begin{center}
\textbf{\S\ 22.}\quad\emph{فرم‌های مرتبهٔ نهم (دستگاه $s_{III}$).}
\end{center}

تاخوردگیِ $s$ با فرم‌های مرتبهٔ پنجم (دستگاه III) و با حاصل‌ضرب $\A\cdot \nu$ (\S\ 10).

\noindent\emph{فرم‌های کاهش‌ناپذیر:}
\[
\begin{array}{c@{\quad}c|c|c|c}
\text{A)}& \cdot & \cdot & (K_\nu su)^2 & \bigstar\\[0.45em]
& \cdot & ((k\nu x)^2,s,u)^2,\ K_\nu s_\nu & ((k\nu x)^2,s,u)^3 & \\[0.45em]
& (k\nu x)^2s_\nu,\ ((k\nu x)^3,s,u) & ((k\nu x)^3,s,u)^2 & \cdot & \cdot\\[0.45em]
& (k\nu x)^3s_\nu & ((k\nu x)^3,s,u)s_\nu,\ (K_\nu(k\nu x)^3,s,u) & \cdot & \cdot\\[0.45em]
& \bigstar\ (K_\nu su)^3 & (K_\nu su)^4 & \cdot & (K_\nu^2su)^4
\end{array}
\]
\[
\begin{array}{c@{\quad}c|c@{\qquad}c@{\quad}c}
\text{B)}& (j\nu x)s_\nu,\ ((j\nu x)^2,s,u) & ((j\nu x)^2,s,u)^2 && \\
& ((j\nu x)^3,s,u) & \cdot && \text{C)}\quad s_\nu(\A su),\ (\A_\nu su),
\end{array}
\]
\[
\begin{array}{c@{\quad}c|c|c}
\text{D)}& \cdot & ((aHu),s,u)^2,\ ((aHu)^2,s,u) & \bigstar\\[0.45em]
& (aHu)s_\eta,\ (a_\eta su) & (aHu)^2s_\eta,\ (a_\eta su)^2 & \\[0.45em]
& \cdot & (a_\eta^2su) & \\
\bigstar& ((aHu),s,u)^3,\ ((aHu)^2,s,u)^2 & ((aHu),s,u)^4 & \\
& (a_\eta su)^3,\ (a_\eta(aHu),s,u)^2,\ (a_\eta(aHu)^2,s,u) & (a_\eta(aHu),s,u)^3 &
\end{array}
\]

\noindent\emph{کاهش‌ها:}

A) کاهش‌ها مستقیماً با یک تاخوردگی از کاهش‌های \S\ 21 به دست می‌آیند. رشتهٔ فرم‌های $K_\nu s_\nu(Ksu)^2$ کاهشگرهای $a_\nu s_\nu(asu)^2$ و$\theta_\nu s_\nu(\theta su)^2$ را به‌عنوان عامل دارد. از این‌رو فرم‌های مرتبهٔ بالاترِ $K_\nu^2(Ksu)^2$ و$K_\nu^2(Ksu)^3$ بر پایهٔ فرض زیر تجزیه‌پذیرند:
\[
0=a_\nu s_\nu(asu)^2(a\theta u)=K_\nu s_\nu(Ksu)^2=\theta_\nu s_\nu(\theta su)^2(a\theta u).
\]

B) رشتهٔ فرم‌های $(j\nu x)^2s_\nu$: کاهشگر $a_\nu^2s_\nu$ از
\[
(j\nu x)^2s_\nu=3a_\nu^2s_\nu(a\theta u)^2.
\]
برای $((j\nu x)^3,s,u)^2$: کاهشگر $a_\nu^2s_\nu$ از
\[
((j\nu x)^3,s,u)^2=-6a_\nu^2s_\nu(a\theta u)^2(\theta su).
\]

C) فرم‌های $\A_\nu(\A su)^2$ و$s_\nu(\A su)^2$: کاهشگر $g_\nu$ از فرمول (27.)a.

فرم $\A_\nu s_\nu$: 1) کاهشگر $a_\nu^2s_\nu$ از $a_ga_\nu^2s_\nu=\frac23\A_\nu s_\nu$؛

2) بنا بر فرمول (27.)a با کاهش دوگانهٔ عبارت $(\A s\widehat{\nu}x)^2$ تجزیه‌پذیر است.

از این‌رو: فرم مرتبهٔ بالاترِ $a_\eta s_\eta$ تجزیه‌پذیر است.

D) رشتهٔ فرم‌های $(aHu)^2s_\eta(asu)$: کاهشگر $a_\nu s_\nu(asu)^2$ با کاهش دوگانهٔ رشتهٔ فرم‌های $[a_\nu s_\nu(asu)^2]_{\widehat{\nu_1}x}$؛ و کاهشِ $(aHu)^2s_\eta(asu)^2$ از $a_\nu s_\nu(asu)^2$ و$a_\nu s_\nu(asu)^3$. از این‌رو $(a_\eta,s,u)^4$ کاهش می‌یابد.

رشتهٔ فرم‌های $a_\eta(aHu)s_\eta$: کاهشگر $a_\nu^2s_\nu$؛ و $a_\eta^2s_\eta$ با کاهش دوگانهٔ عبارت $(\A s\widehat{\nu}x)^3$ کاهش می‌یابد، زیرا $a_\eta^2s_\eta$ با تاخوردگی از جملهٔ کاهش‌پذیرِ $a_\nu^2s_\nu(\nu\nu_1x)$ در فرم $a_\eta(aHu)s_\eta$ پدید نمی‌آید.

رشتهٔ فرم‌های $(a_\eta^2su)^2$: کاهشگر $a_\nu^2s_\nu$ از
\[
a_\nu^2s_\nu s_{\nu_1}=-\frac12a_\eta^2(Hsu)^2.
\]
فرم $(a_\eta(aHu),s,u)$ به‌صورت ترانس‌وکشنی بر یک دترمینان ژاکوبین تجزیه‌پذیر است.

فرمول‌های کاهش زیر به دست می‌آیند:
\srcnumdisplay{(30.)}{%
\begin{aligned}
\text{a)}\quad 0&=(aHu)^2(asu)s_\eta-a_\eta(asu)(Hsu)^2-a_\eta(aHu)(asu)(Hsu),\\
\text{b)}\quad 0&=(aHu)^2(asu)^2s_\eta-a_\eta(aHu)(asu)^2(Hsu),\\
\text{c)}\quad 0&=a_\eta(asu)^2(Hsu)^2,\\
0&=a_\eta s_\eta+\frac14a_\nu^2\cdot g-\frac14s\cdot a_\eta^2,\\
0&=a_\eta(aHu)s_\eta-\frac12a_\eta^2(Hsu),
\quad 0=a_\eta^2(Hsu)^2,\ldots,
\quad 0=a_\eta^2s_\eta.
\end{aligned}
}

\noindent\emph{نتیجه‌ها:}

1) فرم‌های $(j\nu x)^2s_\nu$، $\A_\nu s_\nu$، $\A_\nu(\A su)^2$ و در نتیجه $N_\nu(Nsu)^2$، $(aHu)^2(asu)s_\eta$، $a_\eta(aHu)s_\eta$ و$a_\eta^2(Hsu)^2$ کاهشگرند؛ $a_\eta s_\eta$ فرمی تجزیه‌پذیر است که تحت همهٔ تاخوردگی‌های یک‌باره و دوباره به فرم‌های تجزیه‌پذیر تبدیل می‌شود.

2) فرم‌های مرتبهٔ پنجمِ $(5,\lambda)$، $(6,\lambda)$ و مانند آن‌ها وارد تاخوردگی با $s^2$ نمی‌شوند.


\begin{center}\textbf{\S\ 23.}\quad\emph{فرم‌های مرتبهٔ دهم (دستگاه $s_{III}$).}\end{center}
تاخوردگیِ $s$ با فرم‌های مرتبهٔ ششم (دستگاه III) (\S\ 11).

\noindent\emph{فرم‌های کاهش‌ناپذیر:}
\[
\begin{array}{r@{\quad}c}
\text{A)}&
\begin{array}{c|c}
(\vartheta^2\nu x)^3s_\nu&\cdot\\[0.45em]
\cdot&(\vartheta^2\nu x)^3s_\nu s_\vartheta,\ \theta_\nu(\vartheta^2\nu x)^3s_\vartheta
\end{array}
\end{array}
\]
\[
\begin{array}{r@{\quad}c}
\text{B)}&
\begin{array}{c|ccc}
\cdot&(a_\nu(j\nu x),s,u)^2&(a_\nu(j\nu x),s,u)^3&(a_\nu(j\nu x),s,u)^4\\[0.45em]
a_\nu(j\nu x)s_\nu&\cdot&(a_\nu^2(j\nu x),s,u)^2&(a_\nu^2(j\nu x),s,u)^3
\end{array}
\end{array}
\]
\[
\begin{array}{r@{\quad}c}
\text{C)}&
\begin{array}{c|c|c|c}
\cdot&\cdot&((\theta Hu),s,u)^2,\ ((\theta Hu)^2,s,u)&\bigstar\\[0.35em]
\cdot&((\vartheta\eta x),s,u)^2,\ (\theta Hu)s_\eta,\ (\theta_\eta su)&((\vartheta\eta x),s,u)^3,\ (\theta_\eta su)^2&\\[0.35em]
((\vartheta\eta x)^2,s,u)&\cdot&(\theta_\eta^2su)&\\[0.35em]
\cdot&(\theta_\eta(\vartheta\eta x)^2,s,u)&\cdot&
\end{array}
\end{array}
\]
\[
\begin{array}{c@{\quad}|c|c}
\bigstar&((\theta Hu),s,u)^3,\ ((\theta Hu)^2,s,u)^2&((\theta Hu),s,u)^4\\[0.45em]
&(H_\vartheta(\theta Hu),s,u)^2,\ (\theta_\eta(\theta Hu),s,u)^2,\ (\theta_\eta(\theta Hu)^2,s,u)&(\theta_\eta su)^4,\ (\theta_\eta(\theta Hu),s,u)^3\\[0.45em]
&\cdot&(\theta_\eta^2(\theta Hu),s,u)^2
\end{array}
\]

\noindent\emph{کاهش‌ها:}

A) فرم‌های $((\vartheta^2\nu x)^3,s,u)^2$ و~$((\vartheta^2\nu x)^3,s,u)^3$ تجزیه‌پذیرند:
\[
(\vartheta \nu x)(\widehat{\vartheta \nu su})(\widehat{\vartheta'\nu su})
=(\vartheta \nu x)s_\vartheta s_{\vartheta'}-(\vartheta \nu x)s_\nu s_\vartheta
=-2\theta(\vartheta \nu x)s_\nu s_\vartheta+(\vartheta \nu x)s_\vartheta^2.
\]
فرم‌های باقی‌مانده: یا کاهشگرهایی را به‌عنوان عامل دارند، یا از یک تاخوردگی با فرم‌های تجزیه‌پذیر به دست می‌آیند.

B) کاهشگرهای $a_\nu^2s_\nu$ و~$(\widehat{j\nu su})s_\nu$.

C) 1. رشتهٔ فرم‌های $(\theta Hu)^2s_\eta$: a) کاهشگر $(aHu)^2(asu)s_\eta$؛ b) کاهش دوگانهٔ رشته‌های فرمِ $(\theta gu)^3g_\nu$ و~$g_\nu(agu)^3(bgu)^2(abu)$.

از این‌رو فرم‌های مرتبهٔ بالاترِ $(\theta_\eta su)^3$، $(H_\vartheta(\theta Hu),s,u)^3$ و~$(a_\nu\cdot b_{\nu_1}^2,s,u)^4$ کاهش می‌یابند [فرم اخیر به‌جای $(H_\vartheta su)^4$].

2. $((\vartheta\eta x),s,u)^4$: کاهشگر $(a_\eta su)^4$.

3. $((\vartheta\eta x)^2,s,u)^3$: کاهشگر $s_\vartheta^2s_\nu$ از $(\widehat{\vartheta\eta}su)^2(Hsu)$.
فرم‌های $(\vartheta\eta x)s_\eta$، $(\vartheta\eta x)^2s_\eta$، $((\vartheta\eta x)^2,s,u)^2$ و~$\theta_\eta s_\eta$ با تاخوردگی با فرم تجزیه‌پذیرِ $a_\eta s_\eta$ تجزیه‌پذیرند.

4. رشته‌های فرمِ $H_\vartheta(\theta Hu)s_\eta$، $\theta_\eta(\theta Hu)s_\eta$، $H_\vartheta(\vartheta\eta x)s_\eta$ و~$\theta_\eta(\vartheta\eta x)s_\eta$:
کاهشگرهای $a_\eta(aHu)s_\eta$ و~$a_\eta^2s_\eta$.

رشتهٔ فرم‌های $(\theta_\eta(\vartheta\eta x),s,u)^2$: کاهشگر $a_\eta^2(Hsu)^2$ [به‌استثنای $(\theta_\eta^2(\theta Hu),s,u)^2$ که با تاخوردگی از جملهٔ $a_\eta^2(bsu)(Hsu)$ پدید می‌آید].

5. $(H_\vartheta(\vartheta\eta x),s,u)$ و~$(H_\vartheta(\vartheta\eta x),s,u)^2$: کاهش دوگانهٔ $a_\eta(aHu)s_\eta(abu)$ و~$g_\vartheta(\theta gu)g_\nu$.

رشتهٔ فرم‌های $(H_\vartheta(\vartheta\eta x),s,u)^3$: کاهشگر $((\vartheta \nu x)^2,s,u)^2s_\nu$.

6. $(H_\vartheta(\theta Hu),s,u)$ بنا بر \S\ 3. 2)، c) با یک تاخوردگی با فرم تجزیه‌پذیرِ $(\theta_\nu(\vartheta \nu x),s,u)$ تجزیه‌پذیر است؛ یعنی با کاهش دوگانهٔ فرم تجزیه‌پذیرِ مرتبهٔ پایین‌ترِ $(H_\vartheta su)^2$:\srcfn{*)}{علامت $\equiv$ بدین معناست که حاصل‌ضرب‌های فرم‌های درجات پایین‌تر حذف شده‌اند.}
\[
\begin{aligned}
0&\equiv \theta_\nu(\vartheta\nu\nu_1)(\theta su)+\frac23\theta_\nu^2(\vartheta \nu_1 x)(\theta su)+\theta_\nu(\vartheta \nu x)(\theta su)s_{\nu_1}\\
&\equiv -\frac12H_\vartheta(\theta Hu)(\theta su)+\theta_\eta(\theta su)(Hsu)+\frac12H_\vartheta(\theta su)(Hsu)\\
&\equiv -\frac12H_\vartheta(\theta Hu)(\theta su)+\theta_\eta(\theta su)(Hsu)+\frac12[H_\vartheta]_{\widehat{su}^2}-\frac12\cdot\frac35H_\vartheta(\theta Hu)(\theta su)\\
&\equiv -\frac35H_\vartheta(\theta Hu)(\theta su).
\end{aligned}
\]
با کاهش دوگانه، بنا بر 1. و~5.، فرمول‌های زیر به دست می‌آیند:
\[
\begin{array}{@{}r@{\quad}l@{}}
&\displaystyle
\begin{aligned}
0={}&\theta_\eta(\theta su)(Hsu)^2+\frac14H_\vartheta(\theta Hu)(\theta su)^2-\frac34\theta_\eta(\theta Hu)(\theta su)(Hsu)\\
&-\frac54\theta_\eta(\theta Hu)^2(\theta su)-(asu)^3\cdot b_\eta(bHu)^2-a_\nu(asu)^3\cdot b_{\nu_1}^2,
\end{aligned}\\[0.55em]
(31.)&\displaystyle 0=H_\vartheta(\theta Hu)(\theta su)^3-7\theta_\eta(\theta Hu)(\theta su)^2(Hsu),\\[0.55em]
&\displaystyle 0=a_\nu(asu)^2b_{\nu_1}^{2}(bsu)^2-\theta_\eta(\theta su)^2(Hsu)^2+6\theta_\eta(\theta Hu)(\theta su)^2(Hsu),\\[0.55em]
&\displaystyle 0\equiv (H_\vartheta(\vartheta\eta x),s,u),\qquad
0\equiv H_\vartheta(\widehat{\vartheta\eta su})(Hsu)-4\theta_\eta^2(Hsu).
\end{array}
\]

\noindent\emph{نتیجه‌ها:}

1) فرم‌های $(\theta Hu)^2s_\eta$، $H_\vartheta(\theta Hu)s_\eta$، $\theta_\eta(\theta Hu)s_\eta$، $H_\vartheta(\vartheta\eta x)s_\eta$، $\theta_\eta(\vartheta\eta x)s_\eta$، $(H_\vartheta(\vartheta\eta x),s,u)$ و~$(\theta_\eta(\vartheta\eta x),s,u)^2$ کاهشگرند.

2) $s^2$ وارد تاخوردگی با فرم‌های مرتبهٔ ششمِ $(5,\lambda)$ و مانند آن‌ها نمی‌شود.

\begin{center}\textbf{\S\ 24.}\quad\emph{فرم‌های مرتبهٔ یازدهم (دستگاه $s_{III}$).}\end{center}
تاخوردگیِ $s$ با فرم‌های مرتبهٔ هفتم (دستگاه III) (\S\ 12).

\noindent\emph{فرم‌های کاهش‌ناپذیر:}
\[
\text{A)}\quad
\begin{array}{c|c|c|c|c}
\cdot&\cdot&\cdot&((KHu)^2su)^2&\bigstar\\[0.35em]
\cdot&\cdot&(K_\eta su)^2&\cdot&\\[0.35em]
\cdot&\cdot&\cdot&\cdot&\\[0.35em]
((k\eta x)^3,s,u)&\cdot&\cdot&\cdot&\\[0.35em]
\cdot&(K_\eta(k\eta x)^3,s,u)&\cdot&\cdot&
\end{array}
\]
\[
\begin{array}{c@{\quad}|c|c}
\bigstar&((KHu)^2,s,u)^3&((KHu)^2,s,u)^4\\[0.45em]
&\cdot&(K_\eta(KHu)^2,s,u)^3
\end{array}
\]
\[
\text{B)}\quad ((j\eta x),s,u)^2,\ (H_jsu)\ \big|\ ((j\eta x),s,u)^3,
\qquad \text{C)}\quad (\D_\eta su),
\]
\[
\text{D)}\quad
\begin{array}{c|c|c}
\cdot&((aLu)^2,s,u)^2,\ ((aLu)^3,s,u)&((aLu)^2,s,u)^3\\[0.5em]
(aLu)^2s_\eta,\ (a_lsu)^2&(a_lsu)^3&(a_l(aLu)^2,s,u)^2
\end{array}
\]

\noindent\emph{کاهش‌ها:}

A) کاهش یا تجزیه‌شدن بر اثر یک تاخوردگی با فرم‌های متناظرِ دارای نمادهای $\theta\cdot H$.

فرم‌های $(K_\eta(KHu)^2,s,u)$ و~$(K_\eta(KHu)^2,s,u)^2$ بنا بر \S\ 3. 2)، a)، با تاخوردگی با $K_\nu^2(Ksu)$ و~$K_\nu^2(Ksu)^2$ تجزیه‌پذیرند.

تجزیه‌پذیریِ فرم $(K_\eta(KHu),s,u)^4\equiv(a_\nu^2\cdot\theta_{\nu_1},s,u)^4$ بنا بر \S\ 3. 2)، b)، از فرم تجزیه‌پذیرِ $K_\nu^2(Ksu)^3$ نتیجه می‌شود.

B) رشتهٔ فرم‌های $(j\eta x)s_\eta$: کاهشگر $(j\nu x)^2s_\nu$.

فرم $H_j(\widehat{j\eta su})(Hsu)$: کاهشگر $(j\nu x)(\widehat{j\nu su})^2$.

فرم $H_j(j\eta x)(Hsu)$: با کاهشِ فرم تجزیه‌پذیرِ $((j\eta x)^2,s,u)^2$ به‌وسیلهٔ کاهشگر $(j\nu x)^2s_\nu$ تجزیه‌پذیر می‌شود (فرمول (18.)a):
\[
0=-2(j\nu x)^2s_\nu s_{\nu_1}=[(j\eta x)^2]_{\widehat{su}^2}+H_j(j\eta x)(Hsu).
\]
C) کاهشگرهای $\D_\nu s_\nu$ و~$\D_\nu(\D su)^2$.

D) کاهش و تجزیه‌شدن بر اثر یک تاخوردگی با فرم‌های متناظرِ دارای نمادهای $f\cdot H$.

E) از (30.)c، کاهشِ رشتهٔ فرم‌های تجزیه‌پذیرِ $(a_\sigma su)^2$ به دست می‌آید:
\[
\begin{aligned}
0&=a_\eta(Hsu)^2(as\widehat{\nu}x)^2
=a_\eta(Hsu)^2(a\widehat{\nu}\eta u)^2+2a_\eta(a\widehat{\nu}\eta u)(\widehat{\nu}\eta su)(Hsu)^2
=\frac13a_\sigma(asu)^2,\\
0&=a_\eta a_\nu(Hsu)^2(as\widehat{\nu}x)(asu)
=a_\eta(a\widehat{\nu}\eta u)^2(Hsu)^2(asu)+a_\eta(a\widehat{\nu}\eta u)(\widehat{\nu}\eta su)(Hsu)^2(asu)
=\frac13a_\sigma(asu)^3.
\end{aligned}
\]
\noindent\emph{نتیجه:} $s^2$ وارد تاخوردگی با فرم‌های مرتبهٔ هفتم نمی‌شود.


\begin{center}\textbf{\S\ 25.}\quad\emph{فرم‌های مرتبهٔ دوازدهم، سیزدهم، چهاردهم و پانزدهم (دستگاه $s_{III}$).}\end{center}
تاخوردگیِ $s$ با فرم‌های مرتبهٔ هشتم، نهم، دهم و یازدهم (دستگاه III) (\S\S\ 13--16).

\noindent\emph{فرم‌های کاهش‌ناپذیر:}

\noindent مرتبهٔ دوازدهم.
\[
\begin{adjustbox}{max width=0.98\linewidth}
$\displaystyle
\begin{array}{c@{\qquad}c@{\qquad}c}
\text{A)}&
\begin{array}{c|c}
((\vartheta^2\eta x)^4,s,u)&\cdot\\[0.35em]
\cdot&\theta_\eta(\vartheta^2\eta x)^3s_\vartheta
\end{array}
&
\text{B)}\quad ((aHu)H_j,s,u)^2\;\big|\;((aHu)H_j,s,u)^3
\\[1.0em]
\text{C)}&
\begin{array}{c|c|c}
\cdot&((\theta Lu)^2,s,u)^2,\ ((\theta Lu)^3,s,u)&((\theta Lu)^2,s,u)^3\\[0.35em]
(\theta_lsu)^2&\cdot&(\theta_l(\theta Lu)^2,s,u)^2
\end{array}
&
\end{array}$
\end{adjustbox}
\]
\noindent\emph{کاهش‌ها:} کاهش‌های حاصل از تاخوردگیِ مستقیم با کاهشگرها یا فرم‌های تجزیه‌پذیر، دیگر ذکر نمی‌شوند.

B) تجزیهٔ $((aHu)H_j,s,u)$ و~$(a_\eta(aHu)H_j,s,u)$ به‌صورت ترانس‌وکشنی بر دترمینان‌های ژاکوبین.

تجزیهٔ $(a_\eta(aHu)H_j,s,u)^2$: کاهش دوگانهٔ فرم تجزیه‌پذیرِ $[\theta_\nu\cdot\theta_{\nu_1}^{\prime 3}]_{\widehat{su^3}}$ (\S\ 13. C)
\begin{align*}
0&\equiv[\theta_{\nu_1}^{\prime 3}\cdot\theta_\nu]_{\widehat{su^3}}
\equiv-\frac34\theta_\nu(\theta su)^2(\theta\theta'u)\theta_{\nu_1}^{\prime 3}\\
&\equiv-\frac98(\widehat{\theta\theta'}\widehat{\eta su})(\theta\theta'u)(\theta Hu)(\theta'Hu)\theta_\eta'(\theta su)
\equiv-\frac38a_\eta(aHu)H_j(asu)^2.
\end{align*}

C) از فرمول‌های (31.)، کاهشِ فرم‌های تجزیه‌پذیرِ زیر به دست می‌آید:
\[
 [L_\vartheta(\theta Lu)]_{\widehat{su^3}}\quad\text{و}\quad [\theta_l(\theta Lu)]_{\widehat{su^3}},
\]
یعنی حاصل‌ضرب‌های $[\theta_\nu^2\cdot H]_{\widehat{su^3}}$ و~$[a_\nu^2\cdot(bHu)^2]_{\widehat{su^3}}$:
\[
\begin{array}{@{}r@{\quad}l@{}}
&\displaystyle 0\equiv\theta_\nu^2(\theta su)^2(Hsu),\qquad
0\equiv [L_\vartheta(\theta Lu)]_{\widehat{su^4}}-7[\theta_l(\theta Lu)]_{\widehat{su^4}},\\[0.55em]
&\displaystyle 0\equiv a_\nu^2(asu)^2(bHu)^2(bsu)+6\theta_l(\theta Lu)^2(\theta su)(Lsu),\\[0.55em]
(32.)&\displaystyle 0\equiv\theta_\nu^3(\theta su)(Hsu)^2
\quad\text{از }0\equiv\theta_l^2(\theta Lu)(\theta su)(Lsu)^2\\[-0.05em]
&\displaystyle\hfill(\text{کاهشگر }a_\eta^2(Hsu)^2).
\end{array}
\]

\noindent\emph{فرم‌های کاهش‌ناپذیر.}

\noindent مرتبهٔ سیزدهم.
\[
\begin{array}{rl}
\text{A)}&((KLu)^3,s,u)^2;\ ((KLu)^3,s,u)^3,\\[0.25em]
\text{B)}&(L_jsu)^2,\\[0.25em]
\text{C)}&((aH^2u)^3,s,u)^2.
\end{array}
\]
\noindent مرتبهٔ چهاردهم.
\[
\begin{array}{rl}
\text{A)}&((aLu)^2L_j,s,u)^2,\\[0.25em]
\text{B)}&((\theta H^2u)^3,s,u)^2.
\end{array}
\]
\noindent مرتبهٔ پانزدهم.
\[
 ((KH^2u)^4,s,u)^2.
\]
\noindent\emph{کاهش‌ها:} فرم‌های $(K_l(KLu)^3,s,u)^2$، $(H_\vartheta(\theta H^2u)^3,s,u)$ و~$((K,H\cdot L,u)^5,s,u)$ بنا بر \S\ 3. 2)، c) تجزیه می‌شوند؛ یعنی با بررسیِ هم‌زمانِ فرم‌های تجزیه‌پذیرِ $(K_l(KLu)^2,s,u)^3$، $(H_\vartheta(\theta H'u)^2,s,u)^2$ و~$((K,H\cdot L,u)^4,s,u)^2$.

\noindent\emph{نتیجه:} $s^2$ با هیچ‌یک از فرم‌های منفردِ دستگاه III وارد تاخوردگی نمی‌شود.

\begin{center}\textbf{\S\ 26.}\quad\emph{دستگاه $s_{IV}$.}\end{center}
\noindent 1) \emph{تاخوردگیِ $s$ با دستگاه I$\cdot$III.}

\noindent\emph{کاهش‌ها:} بنا بر کاهش‌های بندهای اخیر ($(a_\nu^2s_\nu,\ (aHu)^2(asu)s_\eta$ و غیره)، فرم‌های $(0,\lambda)$، $(1,\lambda)$ و~$(2,\lambda)$ تاخوردگی‌های I و III را هم‌زمان نمی‌پذیرند؛ ازاین‌رو، بنا بر \S\ 18، $f,\D,N$ برای تشکیل دستگاه $s_{IV}$ در حاصل‌ضرب‌ها وارد تاخوردگی نمی‌شوند.

$j$ وارد تاخوردگی نمی‌شود، زیرا بنا بر همان کاهش‌ها، تاخوردگی‌های پادهموردا نسبت به $j$
\[
(K_\nu(k\nu x)^3,s,u),\qquad ((\vartheta^2\eta x)^4,s,u)\quad\text{و غیره}
\]
با تاخوردگی‌های هموردا نسبت به $j$ متناظرند:
\[
K_\nu(k\nu x)^2s_k,
\qquad (\vartheta^2\eta x)^3s_\vartheta\quad\text{و غیره}
\]

\noindent 2) \emph{تاخوردگیِ $s$ با دستگاه II$\cdot$III}
یعنی تاخوردگیِ $s$ با $H\cdot(1,\lambda)$، $H\cdot(2,\lambda)$ و~$H\cdot(3,\lambda)$.

\noindent\emph{فرم‌های کاهش‌ناپذیر:}
\[
\begin{array}{rl}
\text{A)}&(a_\nu\cdot H,s,u)^4,
((aHu)^2\cdot H,s,u)^3,
((aHu)^2\cdot H,s,u)^4,\\[0.25em]
&((aLu)^2\cdot H,s,u)^4,
((aLu)^3\cdot H,s,u)^3,
((aH^2u)^3\cdot H,s,u)^4,\\[0.45em]
\text{B)}&(a_\nu^2\cdot H,s,u)^3,
(a_\eta(aHu)^2\cdot H,s,u)^3,
(a_l(aLu)^2\cdot H,s,u)^4,\\[0.45em]
\text{C)}&(\theta_\nu\cdot H,s,u)^4,
((\theta Hu)^2\cdot H,s,u)^3,
((\theta Hu)^2\cdot H,s,u)^4,\\[0.25em]
&((\theta Lu)^2\cdot H,s,u)^4,
((\theta Lu)^3\cdot H,s,u)^3,
((\theta H^2u)^3\cdot H,s,u)^4,\\[0.45em]
\text{D)}&(\theta_\eta(\theta Hu)^2\cdot H,s,u)^3,
(\theta_l(\theta Lu)^2\cdot H,s,u)^4,\\[0.45em]
\text{E)}&((KLu)^3\cdot H,s,u)^4,
((KH^2u)^4\cdot H,s,u)^4,\\[0.45em]
\text{F)}&((j\nu x)^2\cdot H,s,u)^3,
((j\nu x)^2\cdot H,s,u)^4,
((j\eta x)\cdot H,s,u)^4,\\[0.25em]
&(H_j\cdot H,s,u)^3,
(L_j\cdot H,s,u)^4.
\end{array}
\]
\noindent\emph{کاهش‌ها:} $\nu$ نیز همانند $j$ وارد تاخوردگی نمی‌شود.

B) و D). $(a_\nu^2\cdot H,s,u)^4$ و~$(\theta_\nu^2\cdot H,s,u)^4$ بنا بر فرمول‌های (27.)c) و (29.)c)، با در نظر گرفتن $(gHu)^2=-\frac13s\cdot\sigma$، تجزیه می‌شوند:
\[
(a_\nu^2\cdot H,s,u)^4\equiv\frac13(asu)^4\cdot\sigma,
\qquad
(\theta_\nu^2\cdot H,s,u)^4\equiv-\frac16(\theta su)^4\cdot\sigma;
\]
از اینجا تجزیهٔ $(a_\eta(aHu)\cdot H,s,u)^4$ و~$(\theta_\eta(\theta Hu)\cdot H,s,u)^4$ به دست می‌آید.

$(\theta_\nu^2\cdot H,s,u)^3$، $(\theta_\nu^3\cdot H,s,u)^3$
و در نتیجهٔ آن‌ها $(\theta_\eta^2(\theta Hu)\cdot H,s,u)^4$، بنا بر مورد C) از فرم‌های مرتبهٔ دوازدهم در \S\ 25، تجزیه می‌شوند.
\noindent 3) \emph{تاخوردگیِ $s$ با دستگاه III$\cdot$III}، یعنی تاخوردگیِ $s$ با فرم‌های دستگاه III:
\[
(3,\lambda)\cdot(3,\lambda),\quad (3,\lambda)\cdot(2,\lambda),\quad (3,\lambda)\cdot(1,\lambda),\quad
(2,\lambda)\cdot(2,\lambda),\quad (2,\lambda)\cdot(1,\lambda),\quad
(1,\lambda)\cdot(1,\lambda).
\]
\noindent\emph{فرم‌های کاهش‌ناپذیر:}
\[
\begin{array}{rl}
\text{A)}&(a_\nu^2\cdot a_\nu^2,s,u)^3,
(a_\nu^2\cdot a_\nu^2,s,u)^4,
(a_\nu^2\cdot a_\eta(aHu)^2,s,u)^3,\\[0.45em]
\text{B)}&(a_\nu\cdot H_j,s,u)^4,
((aHu)^2\cdot H_j,s,u)^3,
((aLu)^2\cdot H_j,s,u)^4,\\[0.25em]
&((aLu)^3\cdot H_j,s,u)^2,
((aH^2u)^3\cdot H_j,s,u)^3.
\end{array}
\]
\noindent\emph{کاهش‌ها:}

حاصل‌ضرب‌های II$\cdot$III، در صورت یکسان‌بودن مرتبه نسبت به نمادهای $\nu$ و $f$، باید بالاتر از حاصل‌ضرب‌های III$\cdot$III به شمار آیند.

کاهش‌ها بخشی از روابط میان حاصل‌ضرب‌ها (سیزیجی‌ها)، و بخشی از جایگزین‌کردن حاصل‌ضرب‌ها با فرم‌های تجزیه‌پذیرِ متناظر و کاهشِ تاخوردگیِ این فرم‌ها با $s$، به دست می‌آیند. (حاصل‌ضرب‌هایی که پادهموردا دربر دارند همواره حذف شده‌اند، زیرا به حاصل‌ضرب‌ها تبدیل می‌شوند.)

A) مرتبهٔ دو نسبت به نماد $f$؛ نمادهای $\nu$ با مرتبهٔ دلخواه. بنا بر فرمول‌های (12.) در \S\ 11 و با محاسبهٔ مشابه، داریم:
\begin{align*}
1)\quad 0&=a_\nu\cdot b_{\nu_1}+\frac12f\cdot(aHu)^2-\frac14\theta\cdot H+\frac12u_xH_\vartheta-\frac14u_x^2\cdot H_\vartheta^2,\\
2)\quad 0&=a_\nu\cdot b_{\nu_1}^2+f\cdot a_\eta(aHu)^2-\theta_\eta-\frac12H_\vartheta,\\
3)\quad 0&=a_\nu^2\cdot b_{\nu_1}^2-H_\vartheta^2+2\theta_\eta^2.
\end{align*}
از اینجا نتیجه می‌شود:
\begin{align*}
4)\quad 0&=a_\nu^2\cdot b_{\nu_1}^2(j\nu_1x)\quad(\text{از }3)),\\
5)\quad L_\vartheta(\theta Lu)&=-a_\nu\cdot b_\eta(bHu)^2+\frac34a_\nu^2\cdot(bHu)^2+\frac34\theta_\nu^2\cdot H\quad(\text{از }2)),\\
6)\quad L_\vartheta(\theta Lu)&=-6a_\nu\cdot b_\eta(bHu)^2+2a_\nu^2\cdot(bHu)^2-\frac12\theta_\nu^2\cdot H\quad(\text{از }1)),\\
7)\quad 0&=a_\nu\cdot(bHu)^2+f\cdot(aLu)^3+\theta_\nu\cdot H\quad(\text{از }1)),\\
\end{align*}
\[
\left.
\begin{aligned}
8)\quad 0&=a_\nu\cdot(bLu)^3+\frac34(\theta Hu)^2\cdot H\\
9)\quad 0&=(aHu)^2\cdot(bHu)^2+2(\theta Hu)^2\cdot H
\end{aligned}
\right\}\quad\text{(از 7)),}
\]
\begin{align*}
10)\quad 0&=a_\eta(aHu)^2\cdot b_\eta(bHu)^2\quad(\text{از }5)\text{ و }6)),\\
11)\quad 0&= [a_\nu^2\cdot b_\eta(bHu)]_{\widehat{su}^4}\quad(\text{از فرمول (30.)b)}).
\end{align*}
کاهش‌های
\[
(H_\vartheta su)^4,
\qquad (L_\vartheta(\theta Lu),s,u)^3,
\qquad (L_\vartheta(\theta Lu),s,u)^4
\]
(فرمول‌های (31.) و (32.)) همراه با فرمول‌های 1) تا 11)، کاهشِ همهٔ تاخوردگی‌های $s$ با حاصل‌ضرب‌هایی را به دست می‌دهند که نسبت به نماد $f$ دارای مرتبهٔ دو و نسبت به $\nu$ دارای مرتبهٔ دلخواه‌اند.

B) حاصل‌ضرب‌های دوتاییِ نمادهای $f,\theta,j$؛ $\nu$ با مرتبهٔ دلخواه. بنا بر فرمول‌های (17.)، (18.) و (20.) و با محاسبهٔ مستقیم، روابط زیر برقرارند:
\[
\left.
\begin{aligned}
1)\quad 0&=a_\nu\cdot\theta_{\nu_1}+\frac14\theta\cdot(aHu)^2+\frac14f\cdot(\theta Hu)^2,\\
2)\quad 0&=\theta_\nu\cdot\theta_{\nu_1}+\frac12\theta\cdot(\theta Hu)^2
\end{aligned}
\right\}\quad\text{(محاسبهٔ مستقیم مشابه A 1)).}
\]
از اینجا نتیجه می‌شود:
\[
\left.
\begin{aligned}
3)\quad 0&=3a_\nu\cdot(\theta Hu)^2+\theta\cdot(aLu)^3+2f\cdot(\theta Lu)^3\\
4)\quad 0&=3\theta_\nu\cdot(aHu)^2+f\cdot(\theta Lu)^3+2\theta\cdot(aLu)^3
\end{aligned}
\right\}\quad\text{(از 1) و A 6)),}
\]
\begin{align*}
5)\quad 0&=a_\nu\cdot(\theta Lu)^3,\quad 0=\theta_\nu\cdot(aLu)^3,\quad 0=(aHu)^2\cdot(\theta Hu)^2
\quad\text{(از 3), 4) و A 8)),}\\
6)\quad 0&=a_\nu\cdot\theta_\nu^2+\theta_\nu\cdot a_\nu^2+f\cdot\theta_\eta(\theta Hu)^2+\theta\cdot a_\eta(aHu)^2\quad(\text{فرمول (17.)}),\\
7)\quad 0&=\theta_\nu\cdot\theta_\nu^2+\theta\cdot\theta_\eta(\theta Hu)^2+\frac16f\cdot H_j\\
&\qquad [\text{استنتاجی مشابهِ استنتاجِ }6)\text{ از }\theta_\nu^2(\theta\theta'\widehat{\nu}_1x)],\\
8)\quad 0&=a_\nu\cdot(j\nu x)^2+f\cdot H_j\quad(\text{از }6)),\\
9)\quad 0&=(aHu)^2\cdot(j\nu x)^2-2a_\nu\cdot H_j\quad(\text{از }8)),\\
10)\quad 0&=\theta_\nu\cdot(j\nu x)^2,\quad 0=\theta\cdot H_j\quad(\text{از }7)\text{ و }8)),
\end{align*}
\[
\left.
\begin{aligned}
11)\quad 0&=a_\nu\cdot\theta_\nu^3-\frac34(j\eta x)^2,\\
12)\quad 0&=a_\nu^2\cdot\theta_\nu^2-\frac13(j\eta x)^2-\frac13(\D Hu)^2
\end{aligned}
\right\}\quad\text{(فرمول (18.)),}
\]
\begin{align*}
13)\quad 0&=a_\nu^2\cdot\theta_\nu^3-\frac12a_\sigma\quad(\text{فرمول (20.)}),\\
14)\quad 0&=\theta_\nu\cdot\theta_\nu^3,
&0&=a_\eta H_j\quad(\S\ 13. C)).
\end{align*}
فرمول‌ها تا آنجا پیش برده شده‌اند که حاصل‌ضرب‌ها قطبی‌پذیر شوند؛ یعنی بر اثر تاخوردگی تنها \emph{یک} حاصل‌ضرب جدید پدید آید، بدین‌گونه که
\begin{center}
\begin{tabular}{rl}
a)&تاخوردگیِ حاصل‌ضرب در خودش کاهش‌پذیر شود،\\
b)&دیگر حاصل‌ضرب‌های پدیدآمده از تاخوردگی کاهش‌پذیر شوند
\end{tabular}
\end{center}
یا به پادهمورداها بینجامند (مقایسه کنید با \S\ 3. 2)، a) و b)).

فرمول‌های 1) تا 5)، کاهشِ همهٔ تاخوردگی‌های $s$ با حاصل‌ضرب‌هایی را به دست می‌دهند که از تاخوردگیِ $a_\nu\cdot\theta_\nu$ پدید می‌آیند،

و فرمول‌های 6) تا 10) نیز کاهشِ حاصل‌ضرب‌هایی را به دست می‌دهند که از تاخوردگیِ $a_\nu\cdot\theta_\nu^2$ پدید می‌آیند.

بنا بر \S\ 24 A)، با در نظر گرفتن فرمول‌های 6) تا 10)، تاخوردگی‌های $s$ با حاصل‌ضرب‌هایی که از $a_\nu^2\cdot\theta_\nu$ پدید می‌آیند، تجزیه می‌شوند.

داریم:
\begin{align*}
0&\equiv a_\nu^2(asu)^2\theta_{\nu_1}(\theta su)^2,
&0&\equiv a_\nu^2(asu)\theta_{\nu_1}(\theta su)^3-K_\eta(KHu)^2(Ksu)^3,\\
0&\equiv a_\nu^2(asu)^2\theta_{\nu_1}(\theta su),
&0&\equiv a_\nu^2(asu)\theta_{\nu_1}(\theta su)^2,\\
0&\equiv a_\nu^2(asu)\theta_{\nu_1}(\theta su).
\end{align*}
کاهشگرها:
\[
 ((j\eta x)^2,s,u)^3
 \ [\text{از }(\vartheta\eta\widehat{su})^2(Hsu).\ \S\ 23.\ C\ 3)],
\]
\[
(H_j(j\eta x),s,u)^2\ [\S\ 24,\ B],
\quad ((\D Hu)^2,s,u)^2\ [\S\ 24,\ C],
\]
\[
(a_\sigma su)^2\ [\S\ 24,\ E]
\]
همراه با فرمول‌های 11) تا 14)، کاهشِ همهٔ حاصل‌ضرب‌های پدیدآمده از
$a_\nu\cdot\theta_{\nu_1}^3$، $a_\nu^2\cdot\theta_{\nu_1}^2$ و $a_\nu^2\cdot\theta_{\nu_1}^3$ را به دست می‌دهند.

محاسبات پیشین ما را می‌توان در
\emph{نتیجهٔ زیر} خلاصه کرد:

دستگاه کامل نسبی به پیمانهٔ $\modu{(\R,t)}$ از \srcLTR{331} فرم و زیردستگاهِ زیر تشکیل می‌شود؛ آن‌ها را در ادامه نیز به‌ترتیب در جدول می‌آوریم:
\[
\left.
\begin{array}{lr}
u_x&1\ \text{فرم}\\
\text{دستگاه }\srcLTR{I}&6\ \text{فرم}\\
\text{دستگاه }\srcLTR{II}&5\ \text{فرم}\\
\text{دستگاه }\srcLTR{III}&117\ \text{فرم}
\end{array}
\right\}
\quad
\begin{gathered}
\text{دستگاه کامل نسبی به پیمانهٔ }(s,\varrho,t)\\
[\text{جدول }\srcLTR{I}\text{ را ببینید}].
\end{gathered}
\]
\[
\left.
\begin{array}{lr}
s&1\ \text{فرم}\\
\text{دستگاه }s_I&35\ \text{فرم}\\
\text{دستگاه }s_{II}&9\ \text{فرم}\\
\text{دستگاه }s_{III}&125\ \text{فرم}\\
\text{دستگاه }s_{IV}&32\ \text{فرم}
\end{array}
\right\}
\quad [\text{جدول }\srcLTR{II}\text{ را ببینید}].
\]

\clearpage
\begin{NoetherWideTable}{297mm}{420mm}
\thispagestyle{plain}
\begin{otherlanguage}{persian}
\begin{center}\bfseries جدول \srcLTR{I} (ص. \srcLTR{90} را ببینید).\end{center}
\end{otherlanguage}
\vspace{-0.6em}
\begingroup
\setlength{\tabcolsep}{2pt}
\renewcommand{\arraystretch}{1.35}
\scriptsize
\begin{center}
\begin{adjustbox}{max width=0.98\linewidth,max totalheight=0.86\textheight,center}
\begin{tabular}{|c|c|c|c|c|c|c|c|c|c|c|c|c|c|c|c|c|c|c|c|c|c|c|}
\hline
{} & 0 & 1 & 2 & 3 & 4 & 5 & 6 & 7 & 8 & 9 & 10 & 11 & 12 & 13 & 14 & 15 & 16 & 17 & 18 & 19 & 20 & 21 \\ \hline
0 & \mcell{i} &  &  &  & \mcell{f} &  & \mcell{\A} &  & \mcell{K_{\nu}(k\nu x)^3} &  & \mcell{K_{\eta}(k\eta x)^3} &  &  &  & \mcell{(\vartheta^2\eta x)^4} & \mcell{H_k(k\cdot\vartheta,\eta x)^4} &  & \mcell{\theta_l(\vartheta\cdot k,lx)^5} &  &  &  & \mcell{(\vartheta^3lx)^6} \\ \hline
1 &  & \mcell{u_x} &  &  &  & \mcell{\theta_\nu(\vartheta\nu x)^2} &  & \mcell{\theta_\eta(\vartheta\eta x)^2} & \mcell{N} & \mcell{(k\nu x)^3} & \mcell{\theta_\nu(\vartheta^2\nu x)^3} & \mcell{(k\eta x)^3} & \mcell{H_{\vartheta}(\vartheta^2\eta x)^3\\ \theta_\eta(\vartheta^2\eta x)^3} &  & \mcell{\theta_l(\vartheta^2lx)^4} &  & \mcell{(\vartheta k,\eta,x)^4} &  & \mcell{(\vartheta k,l,x)^5} &  &  &  \\ \hline
2 &  &  & \mcell{a_\nu^2} &  & \mcell{\theta\\ a_\eta^2} &  & \mcell{(\vartheta\nu x)^2} & \mcell{K_\nu(k\nu x)^2} & \mcell{(\vartheta\eta x)^2} & \mcell{H_k(k\eta x)^2\\K_\eta(k\eta x)^2} &  & \mcell{(\vartheta^2\nu x)^3\\ K_l(klx)^3} &  & \mcell{(\vartheta^2\eta x)^3} &  & \mcell{(\vartheta^2lx)^4} &  &  &  &  &  &  \\ \hline
3 &  & \mcell{\theta_\nu^3} &  & \mcell{a_\nu} & \mcell{\theta_\nu(\vartheta\nu x)} & \mcell{\A_\nu\\a_\eta} & \mcell{K\\H_{\vartheta}(\vartheta\eta x)\\\theta_\eta(\vartheta\eta x)} & \mcell{\A_\eta} & \mcell{(k\nu x)^2\\\theta_l(\vartheta lx)^2} &  & \mcell{(k\eta x)^2} &  & \mcell{(klx)^3} &  &  &  &  &  &  &  &  &  \\ \hline
4 & \mcell{\nu} &  & \mcell{H\\\theta_\nu^2} & \mcell{(j\nu x)^3\\a_\eta(aHu)} & \mcell{\theta_\eta^2} & \mcell{(\vartheta\nu x)\\a_l^2} &  & \mcell{N_\nu\\(\vartheta\eta x)} &  & \mcell{H_n\\N_\eta\\(\vartheta lx)^2} &  &  &  &  &  &  &  &  &  &  &  &  \\ \hline
5 &  & \mcell{a_\eta(aHu)^2} & \mcell{\theta_\eta^2(\theta Hu)} & \mcell{\theta_\nu} & \mcell{K_\nu^2\\(aHu)} & \mcell{\theta_\eta} & \mcell{K_\eta^2\\(\A Hu)\\a_l} &  & \mcell{\A_l} &  &  &  &  &  &  &  &  &  &  &  &  &  \\ \hline
6 & \mcell{j\\\sigma} &  & \mcell{(j\nu x)^2\\(aHu)^2} & \mcell{L\\a_\nu^2(j\nu x)\\H_{\vartheta}(\theta Hu)\\\theta_\eta(\theta Hu)} &  & \mcell{K_\nu\\\theta_l^2} &  & \mcell{K_\eta} &  &  &  &  &  &  &  &  &  &  &  &  &  &  \\ \hline
7 &  & \mcell{\theta_\eta(\theta Hu)^2} & \mcell{H_j(j\eta x)\\a_l(aLu)^2} &  & \mcell{a_\nu(j\nu x)\\(\theta Hu)} &  & \mcell{\A_\nu(j\nu x)\\\theta_l} & \mcell{K_l^2} &  &  &  &  &  &  &  &  &  &  &  &  &  &  \\ \hline
8 & \mcell{H_j^2\\a_l(aLu)^3} & \mcell{(\nu\sigma x)\\(j\nu x)} & \mcell{(\theta Hu)^2} & \mcell{K_\eta(KHu)^2\\(j\eta x)\\(aLu)^2} &  &  &  &  &  &  &  &  &  &  &  &  &  &  &  &  &  &  \\ \hline
9 &  & \mcell{H_j\\(aLu)^3} & \mcell{a_\eta(aHu)H_j\\L_{\vartheta}(\theta Lu)^2\\\theta_l(\theta Lu)^2} &  & \mcell{(KHu)} &  &  &  &  &  &  &  &  &  &  &  &  &  &  &  &  &  \\ \hline
10 & \mcell{\theta_l(\theta Lu)^3} & \mcell{L_j^2\\(j\sigma x)\\a_\eta(aH^2u)^3} &  & \mcell{H_j(aHu)\\(\theta Lu)^2} &  &  &  &  &  &  &  &  &  &  &  &  &  &  &  &  &  &  \\ \hline
11 &  & \mcell{(\theta Lu)^3} & \mcell{K_l(KLu)^3\\L_j\\(aH^2u)^3} &  &  &  &  &  &  &  &  &  &  &  &  &  &  &  &  &  &  &  \\ \hline
12 & \mcell{(aH^2u)^4} & \mcell{a_l(aLu)^2L_j\\H_{\vartheta}(\theta H^2u)^3\\\theta_\eta(\theta H^2u)^3} &  & \mcell{(KLu)^3} &  &  &  &  &  &  &  &  &  &  &  &  &  &  &  &  &  &  \\ \hline
13 & \mcell{(\nu\sigma j)} &  & \mcell{(aLu)^2L_j\\(\theta H^2u)^3} &  &  &  &  &  &  &  &  &  &  &  &  &  &  &  &  &  &  &  \\ \hline
14 & \mcell{(\theta H^2u)^4} & \mcell{K_\eta(KH^2u)^4\\(a,H\!\cdot L,u)^4} &  &  &  &  &  &  &  &  &  &  &  &  &  &  &  &  &  &  &  &  \\ \hline
15 & \mcell{L_\vartheta(\theta,H\!\cdot L,u)^4} &  & \mcell{(KH^2u)^4} &  &  &  &  &  &  &  &  &  &  &  &  &  &  &  &  &  &  &  \\ \hline
16 &  & \mcell{(\theta,H\!\cdot L,u)^4} &  &  &  &  &  &  &  &  &  &  &  &  &  &  &  &  &  &  &  &  \\ \hline
17 & \mcell{K_\eta(K,H\!\cdot L,u)^5} &  &  &  &  &  &  &  &  &  &  &  &  &  &  &  &  &  &  &  &  &  \\ \hline
18 &  & \mcell{(K,H\!\cdot L,u)^5} &  &  &  &  &  &  &  &  &  &  &  &  &  &  &  &  &  &  &  &  \\ \hline
19 &  &  &  &  &  &  &  &  &  &  &  &  &  &  &  &  &  &  &  &  &  &  \\ \hline
20 &  &  &  &  &  &  &  &  &  &  &  &  &  &  &  &  &  &  &  &  &  &  \\ \hline
21 & \mcell{(KH^3u)^6} &  &  &  &  &  &  &  &  &  &  &  &  &  &  &  &  &  &  &  &  &  \\ \hline
\end{tabular}
\end{adjustbox}
\end{center}
\vspace{-0.2em}
\begin{otherlanguage}{persian}
\begin{center}\footnotesize (اعدادِ ردیف افقیِ بالایی درجه نسبت به متغیرهای $x$ و اعدادِ ستون عمودیِ نخست درجه نسبت به متغیرهای $u$ را نشان می‌دهند.)\end{center}
\end{otherlanguage}
\endgroup
\end{NoetherWideTable}
\clearpage
\begin{NoetherWideTable}{420mm}{594mm}
\thispagestyle{plain}
\begin{otherlanguage}{persian}
\begin{center}\bfseries جدول \srcLTR{II} (ص. \srcLTR{90} را ببینید).\end{center}
\end{otherlanguage}
\vspace{-0.6em}
\begingroup
\setlength{\tabcolsep}{2pt}
\renewcommand{\arraystretch}{1.35}
\scriptsize
\begin{center}
\begin{adjustbox}{max width=0.98\linewidth,max totalheight=0.86\textheight,center}
\begin{tabular}{|c|c|c|c|c|c|c|c|c|c|c|c|c|c|c|c|c|c|}
\hline
{} & 0 & 1 & 2 & 3 & 4 & 5 & 6 & 7 & 8 & 9 & 10 & 11 & 12 & 13 & 14 & 15 & 16 \\ \hline
0 & \mcell{J} &  &  &  & \mcell{s} &  & \mcell{s_\vartheta^2} &  &  &  &  & \mcell{s_n} & \mcell{(k\nu x)^3s_\nu} & \mcell{(\vartheta^3\nu x)^3s_\nu s_\vartheta\\\theta_\nu(\vartheta^2\nu x)^3s_\vartheta} &  & \mcell{\theta_\eta(\vartheta^2\eta x)^3s_\vartheta} &  \\ \hline
1 &  &  &  &  &  &  & \mcell{(asu)} & \mcell{s_\vartheta} & \mcell{s_k^2\\(\Delta su)} & \mcell{s_n(Nsu)\\(\vartheta\nu x)^2s_\nu} & \mcell{((k\nu x)^3,s,u)s_\nu\\(K_\nu(k\nu x)^3,s,u)} &  & \mcell{(K_\eta(k\eta x)^3,s,u)} &  & \mcell{(\vartheta^2\nu x)^3s_\nu} &  & \mcell{((\vartheta^2\eta x)^4,s,u)} \\ \hline
2 &  &  &  &  & \mcell{(asu)^2} & \mcell{s_\vartheta(\theta su)} & \mcell{(\Delta su)^2\\a_\nu s_\nu} & \mcell{((\vartheta\nu x)^2,s,u)s_\nu\\(\theta_\nu(\vartheta\nu x)^2,s,u)} &  & \mcell{s_k\\(\theta_\eta(\vartheta\eta x)^2,s,u)} &  & \mcell{(k\nu x)^2s_\nu\\((k\nu x)^3,s,u)} &  & \mcell{((k\eta x)^3,s,u)} &  &  &  \\ \hline
3 &  &  & \mcell{(asu)^3} & \mcell{s_\vartheta(\theta su)^2\\s_\nu} & \mcell{a_\nu s_\nu(asu)\\a_\nu^2(asu)} & \mcell{s_\eta} & \mcell{(\theta su)\\(a_\eta^2su)} & \mcell{s_k(Ksu)} & \mcell{(Nsu)^2\\(\vartheta\nu x)s_\nu\\((\vartheta\nu x)^2,s,u)} & \mcell{((k\nu x)^3,s,u)^2} & \mcell{((\vartheta\eta x)^2,s,u)} &  &  &  &  &  &  \\ \hline
4 & \mcell{(asu)^4} & \mcell{s_\vartheta(\theta su)^3} & \mcell{a_\nu^2(asu)^2} & \mcell{(\theta_\nu^3su)} & \mcell{(\theta su)^2} & \mcell{s_k(Ksu)^2\\a_\nu(asu)\\s_\nu(asu)} & \mcell{((\vartheta\nu x)^2,s,u)^2\\\theta_\nu s_\nu} & \mcell{s_\nu(\Delta su)\\(\Delta_\nu su)\\(aHu)s_\eta\\(a_\eta su)} &  & \mcell{(\Delta_\eta su)} &  &  &  &  &  &  &  \\ \hline
5 &  &  & \mcell{(\theta su)^3} & \mcell{s_k(Ksu)^3,s_j,\\s_\vartheta(\theta s^\prime u)^4\\s_\nu(asu)^2\\a_\nu(asu)^2} & \mcell{(Hsu)\\((\vartheta\nu x)^2,s,u)^3\\(\theta_\nu(\vartheta\nu x),s,u)^2\\(\theta_\nu^2su)} & \mcell{((j\nu x)^3,s,u)\\(aHu)^3s_\eta\\(a_\eta su)^2} & \mcell{(Ksu)^2\\s_l\\(\theta_\eta^2su)} &  & \mcell{((k\nu x)^2,s,u)^2\\K_\nu s_\nu} &  &  &  &  &  &  &  &  \\ \hline
6 & \mcell{(\theta su)^4} & \mcell{s_\nu(asu)^3\\a_\nu(asu)^3} & \mcell{(Hsu)^2\\(\theta_\nu(\vartheta\nu x),s,u)^3\\(\theta_\nu^2su)^2} & \mcell{(a_\eta su)^3\\(a_\eta(aHu),s,u)^2\\(a_\eta(aHu)^2,s,u)} & \mcell{(Ksu)^3} & \mcell{s_j(asu)\\((\vartheta\nu x),s,u)^2\\(\theta_\nu su)} & \mcell{((k\nu x)^2,s,u)^3} & \mcell{s_j(\Delta su)\\a_\nu(j\nu x)s_\nu\\((\vartheta\eta x),s,u)^2\\(\theta Hu)s_\eta\\(\theta_\eta su)} &  &  &  &  &  &  &  &  &  \\ \hline
7 & \mcell{(\theta_\nu(\vartheta\nu x),s,u)^4} & \mcell{(a_\eta(aHu),s,u)^3} & \mcell{(Ksu)^4\\(\theta_\eta^2(\theta Hu),s,u)^2\\(a_\nu^2\cdot a_\nu^2,s,u)^3} & \mcell{s_j(asu)^2\\s_k(Ks^2u)^5\\((\vartheta\nu x),s,u)^3\\(\theta_\nu su)^2} & \mcell{(j\nu x)s_\nu\\((j\nu x)^2,s,u)\\((aHu),s,u)^2\\((aHu)^2,s,u)} & \mcell{((\vartheta\eta x),s,u)^3\\(\theta_\eta su)^2} & \mcell{(aLu)^2s_l\\(a_lsu)^2} &  &  &  &  &  &  &  &  &  &  \\ \hline

8 & \mcell{(a_\nu^2\cdot a_\nu^2,s,u)^4} & \mcell{s_j(asu)^3\\s_k(Ks^2u)^6\\((\vartheta\nu x),s,u)^4\\(\theta_\nu su)^3} & \mcell{((j\nu x)^2,s,u)^2\\((aHu),s,u)^3\\((aHu)^2,s,u)^2} & \mcell{(Lsu)^2\\(a_\nu^2(j\nu x),s,u)^2\\(H_\vartheta(\theta Hu),s,u)^2\\(\theta_\eta(\theta Hu),s,u)^2\\(\theta_\eta(\theta Hu)^2,s,u)} & \mcell{(a_lsu)^3} & \mcell{(K_\nu su)^2} &  & \mcell{(K_\eta su)^2} &  &  &  &  &  &  &  &  &  \\ \hline
9 & \mcell{(K_\nu^2su)^4\\((aHu),s,u)^4} & \mcell{(Lsu)^3\\(a_\nu^2(j\nu x),s,u)^3\\(\theta_\eta su)^4\\(\theta_\eta(\theta Hu),s,u)^3} & \mcell{(Ks^2u)^6\\(a_l(aLu)^2,s,u)^2\\(a_\nu^2\cdot H,s,u)^3} & \mcell{(K_\nu su)^3} & \mcell{(a_\nu(j\nu x),s,u)^2\\((\theta Hu),s,u)^2\\((\theta Hu)^2,s,u)} &  & \mcell{(\theta_lsu)^3} &  &  &  &  &  &  &  &  &  &  \\ \hline
10 &  & \mcell{(K_\nu su)^4\\(a_\nu^2\cdot a_\eta(aHu)^2,s,u)^3} & \mcell{(a_\nu(j\nu x),s,u)^3\\((\theta Hu),s,u)^3\\((\theta Hu),s,u)^2} & \mcell{((j\eta x),s,u)^2\\(H_jsu)\\((aLu)^2,s,u)^2\\((aLu)^3,s,u)} &  &  &  &  &  &  &  &  &  &  &  &  &  \\ \hline
11 & \mcell{(a_\nu(j\nu x),s,u)^4\\((\theta Hu),s,u)^4} & \mcell{(K_\eta(KHu)^2,s,u)^3\\((j\eta x),s,u)^3\\((aLu)^2,s,u)^3\\(a_\nu\cdot H,s,u)^4} & \mcell{(H^2su)^3\\(\theta_l(\theta Lu)^2,s,u)^2} &  & \mcell{((KHu)^2,s,u)^2} &  &  &  &  &  &  &  &  &  &  &  &  \\ \hline
12 &  & \mcell{(a_\eta(aHu)^2\cdot H,s,u)^3} & \mcell{((KHu)^2,s,u)^3} & \mcell{((aHu)H_j,s,u)^2\\((\theta Lu)^2,s,u)^2\\((\theta Lu)^3,s,u)} &  &  &  &  &  &  &  &  &  &  &  &  &  \\ \hline
13 & \mcell{((KHu)^2,s,u)^4} & \mcell{((aHu)H_j,s,u)^3\\((\theta Lu)^2,s,u)^3\\(\theta_\nu\cdot H,s,u)^4} & \mcell{(L_jsu)^2\\((aH^2u)^3,s,u)^2\\((j\nu x)^2\cdot H,s,u)^3\\((aHu)^2\cdot H,s,u)^3} &  &  &  &  &  &  &  &  &  &  &  &  &  &  \\ \hline
14 & \mcell{((j\nu x)^2\cdot H,s,u)^4\\((aHu)^2\cdot H,s,u)^4} & \mcell{(\theta_\eta(\theta Hu)^2\cdot H,s,u)^3} &  & \mcell{((KLu)^3,s,u)^2} &  &  &  &  &  &  &  &  &  &  &  &  &  \\ \hline
15 & \mcell{(a_l(aLu)^2\cdot H,s,u)^4} & \mcell{((KLu)^3,s,u)^3} & \mcell{((aLu)^2L_j,s,u)^2\\((\theta H^2u)^3,s,u)^2\\((\theta Hu)^2\cdot H,s,u)^3} &  &  &  &  &  &  &  &  &  &  &  &  &  &  \\ \hline
16 & \mcell{((\theta Hu)^2\cdot H,s,u)^4\\(a_\nu\cdot H_j,s,u)^4} & \mcell{((j\eta x)\cdot H,s,u)^4\\(H_j\cdot H,s,u)^3\\((aLu)^2\cdot H,s,u)^4\\((aLu)^3\cdot H,s,u)^3} &  &  &  &  &  &  &  &  &  &  &  &  &  &  &  \\ \hline
17 & \mcell{(\theta_l(\theta Lu)^2\cdot H,s,u)^4} &  & \mcell{((KH^2u)^4,s,u)^3} &  &  &  &  &  &  &  &  &  &  &  &  &  &  \\ \hline
18 &  & \mcell{((\theta Lu)^2\cdot H,s,u)^4\\((\theta Lu)^3\cdot H,s,u)^3\\((aHu)^2\cdot H_j,s,u)^3} &  &  &  &  &  &  &  &  &  &  &  &  &  &  &  \\ \hline
19 & \mcell{((aH^2u)^3\cdot H,s,u)^4\\(L_j\cdot H,s,u)^4} &  &  &  &  &  &  &  &  &  &  &  &  &  &  &  &  \\ \hline
20 &  & \mcell{((KLu)^3\cdot H,s,u)^4} & \mcell{((aLu)^3\cdot H_j,s,u)^2} &  &  &  &  &  &  &  &  &  &  &  &  &  &  \\ \hline
21 & \mcell{((\theta H^2u)^3\cdot H,s,u)^4\\((aLu)^2\cdot H_j,s,u)^4} &  &  &  &  &  &  &  &  &  &  &  &  &  &  &  &  \\ \hline
22 &  &  &  &  &  &  &  &  &  &  &  &  &  &  &  &  &  \\ \hline
23 & \mcell{((KH^2u)^4\cdot H,s,u)^4} & \mcell{((aH^2u)^3\cdot H,s,u)^3} &  &  &  &  &  &  &  &  &  &  &  &  &  &  &  \\ \hline
\end{tabular}
\end{adjustbox}
\end{center}
\vspace{-0.2em}
\begin{otherlanguage}{persian}
\begin{center}\footnotesize (اعدادِ ردیف افقیِ بالایی درجه نسبت به متغیرهای $x$ و اعدادِ ستون عمودیِ نخست درجه نسبت به متغیرهای $u$ را نشان می‌دهند.)\end{center}
\end{otherlanguage}
\endgroup
\end{NoetherWideTable}

\setcounter{footnote}{0}


\begin{center}
{\Large\bfseries 3. دربارهٔ نظریهٔ ناورداهای فرم‌های دارای $n$ متغیر.\srcfn{1)}{سخنرانی‌ای که در گردهمایی طبیعت‌پژوهان زالتسبورگ در سال 1909 ایراد شد.}}\par
\vspace{1.2em}
نوشتهٔ \srcspaced{امی نوتر} در ارلانگن.\par
\medskip
\srcLTR{Jahresbericht d. DMV 19 (1910), S. 101--104}
\end{center}

\vspace{1em}
در نظریهٔ ناورداهای پروژکتیوِ فرم‌های دارای $n$ متغیر، مسئلهٔ اصلی حل شده است و متناهی‌بودن دستگاه فرم‌ها به اثبات رسیده است. بااین‌حال، شماری از پرسش‌های دیگر بررسی نشده‌اند، یعنی پرسش‌هایی که به روش‌های پژوهش در روابط میان فرم‌ها مربوط می‌شوند. مایلم نتایجی را که دربارهٔ چنین پرسش‌هایی به دست آورده‌ام به‌اختصار بیان کنم؛ اما برای روشنی، نخست خودِ پرسش‌ها را در حوزهٔ سه‌تایی، که در آن شناخته‌شده‌اند، به‌اجمال طرح می‌کنم.

نخست، دو قضیهٔ بنیادیِ روش نمادی را در نظر می‌گیرم: یکی اینکه همهٔ ساخته‌های ناوردا را می‌توان به‌صورت نمادی نمایش داد، و دیگری اینکه همهٔ روابط موجود میان ناورداها با به‌کارگیریِ پیاپیِ شمار اندکی از همانی‌های نمادی به دست می‌آیند؛ بنابراین روش نمادی همهٔ روابط را بی‌آنکه از حوزهٔ ناورداها بیرون رود به دست می‌دهد.\srcfn{1)}{Study: \emph{Methoden zur Theorie der ternären Formen}. II. \S\ 6. (Leipzig, Teubner 1889.)} بر این مبنا، سپس فرایندهای ساخت فرم‌ها، یعنی فرایند نمادیِ تاخوردگی و فرایندهای غیرنمادیِ مشتق‌گیری، و نیز نظریهٔ بسط‌های سری و مانند آن‌ها مطرح می‌شوند. در همین پیوند، مایلم قضیه‌ای را نیز یاد کنم که در رسالهٔ دکتری‌ام\srcfn{2)}{\emph{Über die Bildung des Formensystems der ternären biquadratischen Form}. \S\ 1. (Crelle's Journal Bd. 134. 1908.)} ثابت کرده‌ام: همهٔ تاخوردگی‌ها را می‌توان با به‌کارگیریِ پیاپیِ دو تاخوردگیِ ممکن که \glqq{}هموردا\grqq{} نامیده می‌شوند ساخت؛ مرادم از آن‌ها، برای حاصل‌ضرب نمادی $s_x t_x u_\sigma u_\tau$، دو تاخوردگیِ $(stu)$ و $(\sigma\tau x)$ است. بر پایهٔ این قضیه می‌توان نظریهٔ کاهشگرها و کاهش دستگاه‌های فرم را، همانند حوزهٔ دوتایی، بنا کرد؛ چنان‌که در رسالهٔ دکتری‌ام نشان داده‌ام.

اگر اکنون به حوزهٔ فرم‌های دارای $n$ متغیر بگذریم، حتی قضیهٔ نخست روش نمادی نیز فقط به‌طور مشروط اعتبار دارد. چنان‌که می‌دانیم، حتی در فرم‌های دارای چهار متغیر، افزون بر مختصات نقطه‌ها و صفحه‌ها، یعنی به‌ترتیب $x$ و $u$، مختصات خط‌ها $p_{ik}$ نیز ظاهر می‌شوند؛ این‌ها دترمینان‌های فرعیِ حاصل از دو ردیف مختصات نقطه‌اند. به‌همین‌سان، در حوزهٔ فرم‌های دارای $n$ متغیر ردیف‌های متغیر $p_\rho$ را داریم که عناصر منفردشان دترمینان‌های فرعیِ حاصل از $\rho$ ردیف $x$ هستند، و ردیف‌های نماد $S^\rho$ را داریم که عناصرشان مانند دترمینان‌های فرعیِ حاصل از $\rho$ ردیف $s$ که با $u$ هموردا هستند، رفتار می‌کنند. چنان‌که می‌دانیم، نمایش نمادی با دترمینان‌ها فقط در صورتی معتبر است که ردیف‌های نماد $S^\rho$ به ردیف‌های $s$ تفکیک و این ردیف‌ها میان دترمینان‌های مختلف توزیع شوند؛ آنگاه فقط ترکیب‌هایی از دترمینان‌ها که به $S^\rho$ وابسته‌اند معنای واقعی دارند. اما نمی‌توان تشخیص داد کدام ترکیب‌های دترمینانی این شرط را برآورده می‌کنند؛ ازاین‌رو همهٔ قضایای دیگری که در حوزهٔ سه‌تایی نام بردیم از دست می‌روند.

اکنون می‌خواهم اصلی ساده بیان کنم تا \emph{نمایش نمادیِ صریح برحسب ردیف‌های $S^\rho$} به دست آید و بدین‌وسیله دیگر قضایا نیز بازیافته شوند. این کاربردی از قضیهٔ شناخته‌شده دربارهٔ ماتریس‌های متناظر است: \glqq{}با هر ردیف متغیر $p_\rho$ می‌توان ردیف دیگری $q_{n-\rho}$ را به‌طور یک‌به‌یک متناظر کرد که عناصر منفردش دترمینان‌های فرعیِ حاصل از $n-\rho$ ردیف $u$ هستند؛ این ردیف‌ها با $\rho$ ردیف $x$ به‌وسیلهٔ معادله‌های $(u\mid x)=0$ مرتبط‌اند.\grqq{} به‌طور متناظر، به هر ردیف نماد $S^\rho$ ردیف دیگری $S_{n-\rho}$ نسبت داده می‌شود که چنان رفتار می‌کند که گویی از $n-\rho$ ردیفِ هموردا با $x$ تشکیل شده است.

این قضیه صورت‌بندیِ دومی نیز دارد که برای کاربردها مناسب است: \glqq{}به هر حاصل‌ضرب ماتریسی $(v^{(1)}\cdots v^{(\rho)}\mid y^{(1)}\cdots y^{(\rho)})$،\srcfn{1)}{یعنی مجموع حاصل‌ضرب‌های دترمینان‌های متناظر که یکی از ماتریسِ متشکل از $\rho$ ردیف $v$ و دیگری از ماتریسِ متشکل از $\rho$ ردیف $y$ گرفته می‌شود.} که ردیف‌های ماتریس دوم نسبت به ردیف‌های ماتریس نخست پادهموردا هستند، یک دترمینان نسبت داده می‌شود؛ و برعکس، هر دترمینان را می‌توان به حاصل‌ضربی ماتریسی با $\rho$ ردیف، که شمارشان از پیش تعیین شده است، تبدیل کرد.\grqq{} بدین‌ترتیب، دترمینان و حاصل‌ضرب ماتریسی کاملا هم‌ارز می‌شوند و قضیهٔ نخست روش نمادی به صورت زیر درمی‌آید: \glqq{}همهٔ ساخته‌های ناوردا را می‌توان با حاصل‌ضرب‌های ماتریسی به‌صورت نمادی نمایش داد.\grqq{} اکنون با کمک ردیف‌های متغیر می‌توان از دو ردیف نماد $S^\sigma$ و $T^\tau$ به‌آسانی حاصل‌ضربی ماتریسی ساخت که این ردیف‌های نماد را صریحا دربردارد و در نتیجه یک ساختهٔ ناوردا به‌صورت $(S^\sigma\mid p_\sigma)(T^\tau\mid p_\tau)$ را نمایش می‌دهد، یعنی حاصل‌ضرب زیر:
\srcnumdisplay[3.2em]{(1)}{%
  \pair{q_{n-\sigma-\lambda}S^\sigma}{T_{n-\tau}p_{\tau-\lambda}},
  \qquad (\lambda=1,2,\ldots,\tau\ \text{یا}\ n-\sigma).
}

مهم‌تر، عکسِ این گزاره است: \emph{همهٔ ساخته‌های ناوردا را می‌توان به‌طور صریح با حاصل‌ضرب‌های ماتریسی نمایش داد}؛ یعنی فرایند بالا تعمیمِ فرایند تاخوردگی در حوزه‌های دوتایی و سه‌تایی است. برای اثبات این عکس لازم است قضیهٔ دوم روش نمادی را به حوزهٔ فرم‌های دارای $n$ متغیر منتقل کنیم، یعنی مجموعهٔ کامل همانی‌های مستقلِ تجزیه‌ناپذیر را صورت‌بندی کنیم که در آن‌ها همهٔ ردیف‌های $S^\rho$ و $p_\rho$ خود تجزیه‌ناپذیر در نظر گرفته می‌شوند. این همانی‌ها از همانی‌های شناخته‌شده‌ای که فقط ردیف‌های $x$ و $u$ را دربر دارند،\srcfn{2)}{E. Pascal in \emph{Memorie d. R. Acc. d. Lincei}. (4) 5 (1888), p. 375.} به دست می‌آیند؛ بدین‌گونه که قضیهٔ حاصل‌ضرب دترمینان‌ها با دستگاهی از قضایای حاصل‌ضرب ماتریس‌ها جایگزین می‌شود و چند حاصل‌ضرب ماتریسی نیز با همین قضایای حاصل‌ضرب در یک حاصل‌ضرب ماتریسیِ واحد ادغام می‌شوند. بدین‌سان به دو فرمول دوگان می‌رسیم که در اینجا فقط یکی از آن‌ها آورده می‌شود:
\srcnumdisplay[3.2em]{(2)}{%
\begin{aligned}
 \pair{S_1^{\rho_1}S_2^{\rho_2}\cdots S_k^{\rho_k}}{x p_{\rho-1}}
 &=\sum_{i=1}^{k}\eps\,\bigl(S_1^{\rho_1}\cdots S_{i-1}^{\rho_{i-1}}S_{i+1}^{\rho_{i+1}}\cdots S_k^{\rho_k}q_{n-\rho+1},\,S_{n-\rho_i}x\bigr);\\
 \rho&=\sum_{i=1}^{k}\rho_i\le n;\qquad \eps=\pm1.
\end{aligned}
}

تنها ویژه‌سازی‌ای که این فرمول‌ها دربر دارند، یعنی واردشدن یک ردیف $x$، اجتناب‌ناپذیر است؛ در حالت کاملا عمومیِ ردیف‌های نماد و ردیف‌های متغیر، به یک \glqq{}همانیِ تجزیه\grqq{} می‌رسیم، یعنی همانی‌ای که در آن یک ردیف $p_\rho$ به ردیف‌های منفرد $x$ تجزیه می‌شود. کلِبش پیش‌تر ساده‌ترین حالت ویژهٔ همانیِ تجزیه را در اثر خود، \glqq{}Fundamentalaufgabe der Invariantentheorie\grqq{}، برای استخراج همورداهای بنیادی در بسط‌های سری به کار برده بود. اکنون با کمک این همانی‌ها، که میانجیِ گذار از عبارت‌های دترمینانیِ غیرصریح به حاصل‌ضرب ماتریسی هستند، می‌توان در واقع نشان داد که نمایش ماتریسی همان نمایش نمادیِ صریحِ مطلوب را فراهم می‌کند.

برای فرایند تاخوردگی نیز قضیه‌ای کاملا همانند قضیهٔ حوزهٔ سه‌تایی برقرار است، و قضایای کاهش را نیز می‌توان به همان شیوه بر آن بنا کرد. اگر در (1) عدد $\lambda$ را نقص تاخوردگی بنامیم و تاخوردگی‌هایی را که برایشان $\sigma=\tau$ است هموردا بنامیم، \emph{همهٔ تاخوردگی‌ها را می‌توان با به‌کارگیریِ پیاپیِ تاخوردگی‌های هموردای دارای نقص 1 ساخت}. برهان این قضیه اساسا بر آن استوار است که می‌توان عمومی‌ترین فرم‌ها را با \glqq{}فرم‌های نرمال\grqq{} جایگزین کرد، یعنی با فرم‌هایی که با اعمال همهٔ عملیات دیفرانسیلیِ ناوردا عینا صفر می‌شوند. در حوزهٔ فرم‌های دارای چهار متغیر، آقای \emph{مرتنس}\srcfn{1)}{Über invariante Gebilde quaternärer Formen. Wiener Berichte. Bd. 98. 1889.} این واقعیت اخیر را اثبات کرده است؛ شناخت ما از عمومی‌ترین عملیات دیفرانسیلی ــ که چنان‌که می‌دانیم با جایگزین‌کردن ردیف‌های نماد با نمادهای دیفرانسیلی از فرایندهای تاخوردگی پدید می‌آیند ــ به ما امکان می‌دهد با به‌کارگیریِ همانیِ تجزیه، برهان مرتنس را تقریبا واژه‌به‌واژه به حوزهٔ فرم‌های دارای $n$ متغیر منتقل کنیم.

\clearpage
\setcounter{footnote}{0}
\begin{center}
{\Large\bfseries دربارهٔ نظریهٔ ناورداهای\\
فرم‌های دارای $n$ متغیر\footnote{نگاه کنید به اطلاعیهٔ کوتاهی برای گردهمایی طبیعت‌پژوهان زالتسبورگ: \srcLTR{Jahresbericht d. deutsch. Math.-Vereinigung. Bd.\ 19, S.\ 101.}}}\par
\vspace{0.5em}
نوشتهٔ دوشیزه \emph{امی نوتر} در ارلانگن.\par
\medskip
\rule{0.22\linewidth}{0.4pt}\par
\vspace{0.7em}
\srcLTR{Journal f.\ d.\ reine u.\ angew.\ Math.\ 139 (1911), S.\ 118--154}
\end{center}

\begin{center}
\textbf{مقدمه.}
\end{center}
در نظریهٔ ناورداهای پروژکتیوِ فرم‌های دارای $n$ متغیر، اگر از حوزهٔ دوتایی و سه‌تایی فراتر رویم، پرسش‌هایی که به مسئلهٔ اصلی ــ اثبات متناهی‌بودن دستگاه فرم‌ها ــ می‌پیوندند به‌ندرت بررسی شده‌اند. این‌ها پرسش دربارهٔ وابستگی متقابل فرم‌ها و روش‌های مهار این پیوند میان فرم‌ها هستند. این روش‌ها اساساً ــ چنان‌که در حوزهٔ دوتایی و سه‌تایی به‌طور کامل ثابت شده است ــ بر دو قضیه استوارند که آقای اشتودی آن‌ها را \glqq{}قضایای بنیادیِ روش نمادی\grqq{} نامیده است، یعنی قضیهٔ نمایش‌پذیری نمادیِ فرم‌ها و قضیهٔ همانی‌های نمادی.\footnote{\srcLTR{E. Study, \emph{Methoden zur Theorie der ternären Formen} (Leipzig, Teubner, 1889). II.} \S\ 6. -- برای حضور این قضایای بنیادی در نظریهٔ ناورداهای گروه‌های تبدیل ویژه نیز نک. \srcLTR{Study, \emph{Das Apollonische Problem} (Math. Ann. Bd. 49 [1897])} و \srcLTR{\emph{Zur Differentialgeometrie der analytischen Kurven} (Transactions of the American Math. Society, Bd. 1 [1909]).}} قضیهٔ اخیر می‌گوید همهٔ روابط موجود میان ساخته‌های ناوردای چندجمله‌ای با به‌کارگیریِ پیاپیِ شمار اندکی از همانی‌های نمادی به دست می‌آیند؛ بنابراین روش نمادی، و تنها همین روش، شناخت پیوند میان فرم‌ها را بی‌آنکه از حوزهٔ ناورداها بیرون رود فراهم می‌کند.

آقای گوردان پیش‌تر در برنامهٔ ارلانگن خود\footnote{\srcLTR{P. Gordan, \emph{Über das Formensystem binärer Formen}. S. 46. (Leipzig, Teubner, 1875).}} به علت دشواریِ گذار به $n$ متغیر اشاره کرده است؛ علت در این است که قضیهٔ نخستِ روش نمادی در صورت‌بندیِ عام خود دیگر برقرار نیست. چنان‌که می‌دانیم، در حوزهٔ $n$-تایی ردیف‌های متغیرِ ردهٔ بالاتر $p_\rho$\footnote{برای نمادگذاری نک. \S\ 1.} و، به‌همین قیاس، ردیف‌های نمادِ ردهٔ بالاتر $S^\rho$ پدید می‌آیند؛ این ردیف‌ها را هم هنگامی به دست می‌آوریم که، همان‌گونه که کلِبش\footnote{\srcLTR{A. Clebsch, über eine Fundamentalaufgabe der Invariantentheorie. (Abh. der Ges. d. Wiss. zu Göttingen, Bd. XVII, 1872).}} نخستین‌بار چنین کرد، از فرم‌هایی با تعداد دلخواهی ردیف $p_\rho$ آغاز کنیم، و هم هنگامی که، به پیروی از ف. درویتس و آ. کاپلی، از فرم‌هایی آغاز کنیم که تنها تعداد دلخواهی ردیف هموردای $x$ دارند.\footnote{نک. \srcLTR{A. Capelli, \emph{Lezioni sulla teoria delle forme algebriche} (Napoli 1902),} \S\ 22--24، و منابع یادشده در آنجا.} اکنون نمایش نمادی‌ای که ردیف‌های $p_\rho$ و $S^\rho$ را به‌طور صریح\footnote{تعریف دقیق‌ترِ \glqq{}نمایش صریح\grqq{} در \S\ 2 آمده است.} دربرداشته باشد وجود ندارد؛ لازم است $p_\rho$ و $S^\rho$ را به ردیف‌های منفرد $x$ و ردیف‌های $s$ پادهموردای آن‌ها تفکیک کنیم و این ردیف‌ها را میان دترمینان‌های گوناگون توزیع کنیم. البته می‌توان به‌کمک شرط‌های دیفرانسیلی (نک. فرمول (19.)) بررسی کرد که یک مجموع دترمینانیِ معین این ویژگی را دارد که فقط به ردیف‌های $p_\rho$ و $S^\rho$ وابسته باشد؛ اما از این راه نمی‌توان بر کلِ ساخته‌های ناوردا و فرایندهای ساخت آن‌ها احاطه یافت.\footnote{گسترشی از نمادپردازی که آقای اشتودی عرضه کرده و از نظر محاسباتی بسیار ارزشمند است (به نقل از \srcLTR{F. Engel, Ber. d. K. Sächs. Ges. d. Wiss., math.-phys. Kl., 1900, S. 126}، و برای حوزهٔ چهارتایی از \srcLTR{E. Study, \emph{Geometrie der Dynamen}, S. 135}) نیز در معنای ما نمایش صریحی نیست. بعید است مثلاً به فرایندهای مشتق‌گیریِ غیرنمادی برای ساخت فرم‌ها بینجامد.}

اکنون در \S\S\ 2 و 6 نشان می‌دهیم که چگونه با کاربردِ \glqq{}قضیهٔ ماتریس‌های متناظر\grqq{}\footnote{گذشته از \srcLTR{Graßmanns Ausdehnungslehre} سال 1862، شمارهٔ 112، این قضیه نخست نزد \srcLTR{A. Brill, über zwei Berührungsprobleme,} \S\ 2 \srcLTR{(Math. Ann. Bd. 4. 1871)} آمده است.} قضیهٔ بنیادیِ نخست را در کلیت خود بازیابی می‌کنیم؛ بدین‌گونه که نمایش با دترمینان‌ها را با نمایشی از طریقِ \glqq{}حاصل‌ضرب‌های ماتریسی\grqq{} جایگزین می‌کنیم. در \S\ 2 نشان داده می‌شود که هر حاصل‌ضرب ماتریسی که ردیف‌های $S^\rho,p_\rho$ را به‌طور صریح دربردارد، ساخته‌ای ناوردا از فرم بنیادی است؛ عکسِ این گزاره، که در \S\ 6 می‌آید و می‌گوید هر ساختهٔ ناوردا را می‌توان به‌طور صریح با حاصل‌ضرب‌های ماتریسی نمایش داد، تنها با انتقالِ قضیهٔ بنیادیِ دوم (\S\S\ 3--5) به دست می‌آید.

این قضیه برای فرم‌هایی که فقط ردیف‌های متغیر $x$ را دربر دارند شناخته شده است.\footnote{\srcLTR{E. Pascal, Giorn. di mat. 26 (1888), S. 33, 102; Rendic. d. Accad. d. Linc. (4) 4 (1888), S. 119; ib. Mem. (4) 5 (1888), S. 375.}} مسئلهٔ عامِ صورت‌بندیِ کلِ همانی‌های تجزیه‌ناپذیری که در آن‌ها همهٔ ردیف‌های $S^\rho,p_\rho$ تجزیه‌ناپذیر در نظر گرفته می‌شوند، از سوی آقای اشتودی طرح شده است؛\footnote{\srcLTR{``Methoden \dots'', S. 204, Anmerkung 17).}} او هم‌زمان شمار همانی‌های پدیدارشونده را معیاری برای دشواریِ روش‌ها در حوزهٔ $n$-تایی می‌داند. اکنون که می‌توان کلِ همانی‌ها را در یک فرمول واحد (36.) گرد آورد، این دشواری دست‌کم به کمینه کاهش یافته است. بااین‌حال، هنگام کاربرد غالباً بعضی حالت‌های ویژهٔ برجستهٔ (36.) پدید می‌آیند، مانند (20.) و (21.) که ویژگی‌شان این است که بی‌جایگذاریِ ردیف‌های تازه نمایش صریح را ممکن می‌کنند؛ یا فرمول (31.) که \glqq{}همانیِ تجزیه\grqq{} نامیده می‌شود، زیرا در آن یک ردیف $p_\rho$ به‌ظاهر به ردیف‌های $y$ تجزیه‌شده پدیدار می‌شود.

اکنون می‌توان ادامهٔ نظریه را به‌آسانی بر دو قضیهٔ بنیادی بنا کرد. قضیهٔ نخست مستقیماً فرایندهای ساخت فرم‌ها، یعنی فرایند نمادیِ تاخوردگی و فرایندهای غیرنمادیِ مشتق‌گیری، را به دست می‌دهد،\footnote{در مقالهٔ \srcLTR{``Invariante Gebilde von Nullsystemen'' (Sitzungsber. d. Ak. d. Wiss. Wien, Bd. 97, Abt. IIa, 1888)}، آقای مرتنس از راهی دیگر که مستقیماً تعمیم‌پذیر نیست به فرایند تاخوردگیِ چهارتایی رسیده است. تفاوتِ ناوردای متناوبِ شش ردیف $S^2$ که در آنجا پدید می‌آید با ناوردایی که نزد ما با یک‌بار به‌کارگیریِ جایگذاریِ (11.) حاصل می‌شود، مجموعی از ناورداهای زوج است. در یک حالت ویژه، فرایند تاخوردگیِ چهارتایی نزد آقای پ. گوردان نیز آمده است: \srcLTR{\emph{Das simultane System von zwei quadratischen quaternären Formen} (Math. Ann. Bd. 56. 1903).}} و همانیِ تجزیه در میان این فرایندها همهٔ فرایندهای هم‌ارز را حذف می‌کند (\S\ 6). شناخت فرایندهای مشتق‌گیری، افزون بر این، انتقالِ مفهوم \glqq{}فرم نرمال\grqq{} به حوزهٔ $n$-تایی و، به‌کمک روش اثباتی که آقای مرتنس\footnote{\srcLTR{F. Mertens, Über invariante Gebilde quaternärer Formen. (Sitzber. d. Ak. d. Wiss. Wien. Bd. 98. Abt. IIa. 1889.)}} در حوزهٔ چهارتایی عرضه کرده است، این قضیه را به دست می‌دهد که می‌توان یک دستگاه از ساخته‌های ناوردا را با دستگاهی هم‌ارز از فرم‌های نرمال جایگزین کرد (\S\ 7). با فرضِ اخیر، در رسالهٔ دکتری‌ام\footnote{\srcLTR{Über die Bildung des Formensystems der ternären biquadratischen Form. (Dieses Journal, Bd. 134. 1908.)}} نشان داده‌ام که فرایند تاخوردگیِ سه‌تایی را می‌توان با به‌کارگیریِ پیاپیِ دو تاخوردگیِ بنیادی ساخت. سپس بر پایهٔ این قضیه و نظریهٔ بسط‌های سری، با معرفیِ مفهوم \glqq{}رشتهٔ فرم‌ها\grqq{}، مهم‌ترین قضایای کاهش را برای دستگاه‌های فرم بسط داده‌ام. کاملاً به‌همین قیاس، در حوزهٔ $n$-تایی می‌توان فرایند تاخوردگی را از $n-1$ تاخوردگیِ بنیادی ساخت (\S\ 8) و قضایای کاهش را صورت‌بندی کرد (\S\ 9).

\paragraph{یادداشت پایانی.} در پی اطلاعیهٔ زالتسبورگ، آقای پروفسور ا. مولر از وین توجه مرا به این نکته جلب کرد که محاسبه با حاصل‌ضرب‌های ماتریسی، که بنیاد این مقاله است، در اصل با حسابِ نظریهٔ امتدادِ گراسمن یکسان است.\footnote{\srcLTR{Hermann Graßmanns gesammelte mathematische und physikalische Werke. Ersten Bandes zweiter Teil: Die Ausdehnungslehre von 1862. In Gemeinschaft mit H. Graßmann d. Jüngeren herausgegeben v. Fr. Engel (Leipzig, Teubner 1896).}} در واقع، می‌توان \glqq{}حاصل‌ضرب ترکیبیاتی\grqq{} گراسمن $[S^{\rho_1}S^{\rho_2}\cdots S^{\rho_i}]$ را ماتریسی دانست که این ردیف‌ها تعیین می‌کنند (نک. \S\ 2): \glqq{}حاصل‌ضرب پیش‌رونده\grqq{} ماتریسِ خودِ ردیف‌هاست و \glqq{}حاصل‌ضرب پس‌رونده\grqq{} ماتریسِ ردیف‌های متناظر $S_{n-\rho_1},S_{n-\rho_2},\ldots,S_{n-\rho_i}$ است؛ درحالی‌که ماتریسِ متناظر با \glqq{}حاصل‌ضرب آمیخته\grqq{} از آن راه پدید می‌آید که چند ردیف $S$ را با جایگذاریِ (9.) در ردیف‌های تازهٔ $P,Q$ گرد آوریم. در این تناظر، \glqq{}ردیف متناظر\grqq{} ما با \glqq{}متمم\grqq{} گراسمن و \glqq{}حاصل‌ضرب ماتریسی\grqq{} ما با \glqq{}حاصل‌ضرب آمیخته با عددِ ردهٔ 0\grqq{} متناظر است.

بر پایهٔ این تناظر، بسط داده‌شده در نظریهٔ امتدادِ 1862، شمارهٔ 113، با فرمولی دوگان نسبت به فرمول (32.) ما یکسان می‌شود. در حالت ویژهٔ $\rho=1$، (32.) به (17.) و فرمول دوگان به (16.) بدل می‌شود. آقای مولر\footnote{\srcLTR{E. Müller, Beiträge zur Graßmannschen Ausdehnungslehre. 1. Mitteilung: Einige allgemeine Sätze. (Sitzungsber. d. Ak. d. Wiss. Wien; Bd. 118. Abt. IIa. Juli 1909),} قضیه‌های 16 و 18.} اخیراً همانی‌های (20.) و (21.) یادشده در بالا را از این حالت‌های ویژهٔ بسطِ گراسمن استخراج کرده است.

استنتاج ما، برخلاف استنتاج گراسمن--مولر، هم‌زمان اثباتِ کامل‌بودنِ دستگاه همانی‌ها را نیز به دست می‌دهد ــ امری که برای پرسش‌های موردنظر در اینجا اساسی‌ترین نکته به نظر می‌رسد ــ و از (20.) و (21.) فراتر می‌رود و به عام‌ترین همانی‌های تجزیه‌ناپذیر (36.) می‌رسد که ظاهراً تاکنون مورد توجه قرار نگرفته‌اند.

\section*{\S\ 1. نمادگذاری و تعریف‌ها. خلاصهٔ نتایج شناخته‌شده.}
مطابق معمول، ردیف متغیرهای $x$ را ردیف عناصر $x_1,x_2,\ldots,x_n$ می‌نامیم و، به همین قیاس، ردیف متغیرهای $p_\rho$ ($\rho=1,2,\ldots,n-1$) را ردیفِ $\binom n\rho$ عنصرِ $p_{i_1i_2\ldots i_\rho}$؛ که در آن، $p_{i_1i_2\ldots i_\rho}$ دترمینان‌هایی هستند که از ماتریسِ دارای $\rho$ سطرِ $x^{(1)},x^{(2)},\ldots,x^{(\rho)}$ ساخته شده‌اند، تعدادشان $\binom n\rho$ است، و به صورت
\[
\left|\begin{array}{cccc}
 x_{i_1}^{(1)}&x_{i_2}^{(1)}&\cdots&x_{i_\rho}^{(1)}\\
 x_{i_1}^{(2)}&x_{i_2}^{(2)}&\cdots&x_{i_\rho}^{(2)}\\
 \vdots&\vdots&&\vdots\\
 x_{i_1}^{(\rho)}&x_{i_2}^{(\rho)}&\cdots&x_{i_\rho}^{(\rho)}
\end{array}\right|,
\qquad (i_1<i_2<\cdots<i_\rho=1,2,\ldots,n)
\]
نوشته می‌شوند. بنا بر قضیهٔ ماتریس‌های متناظر\footnote{نک. برای نمونه: \srcLTR{Gordan-Kerschensteiner, \emph{Vorlesungen über Invariantentheorie}, Bd. I, S. 99.}}، با هر ردیف $p_\rho$ ردیف دومی، $q_{n-\rho}$، به‌واسطهٔ رابطهٔ زیر در تناظری یک‌به‌یک قرار می‌گیرد؛ رابطه‌ای که برای هر عنصر ردیف برقرار است:
\begin{equation}
 p_{i_1i_2\ldots i_\rho}=(-1)^{\frac{\rho(\rho+1)}{2}+i_1+i_2+\cdots+i_\rho}q_{j_1j_2\ldots j_{n-\rho}},
\tag{1}
\end{equation}
در این رابطه، اندیس‌های $i$ و $j$، هر دسته در میان خود، به‌ترتیب صعودی‌اند و در کنار هم یک جایگشت کامل از اعداد 1 تا $n$ را تشکیل می‌دهند. ردیف $q_{n-\rho}$ از $\binom n\rho$ دترمینان درجهٔ $(n-\rho)$، یعنی $q_{j_1j_2\ldots j_{n-\rho}}$، تشکیل می‌شود که از ماتریسِ دارای $n-\rho$ سطرِ پادهموردای $x$، یعنی $u^{(1)},u^{(2)},\ldots,u^{(n-\rho)}$، ساخته شده‌اند. این سطرها $\rho(n-\rho)$ معادلهٔ شرطی زیر را برآورده می‌کنند:
\begin{equation}
 (u^{(k)}\mid x^{(l)})=\sum_{\alpha=1}^{n}u_\alpha^{(k)}x_\alpha^{(l)}=0;
 \qquad (k=1,2,\ldots,n-\rho;\ l=1,2,\ldots,\rho).
\tag{2}
\end{equation}
با پیروی از کلِبش\footnote{کلِبش، اثرِ پیش‌گفته، \S\ 5.}، این سطرها را چنان نرمال‌شده فرض می‌کنیم که ضریب تناسبی که در غیر این صورت در (1.) وارد می‌شد، برابر یک گرفته شود.

عام‌ترین فرم‌ها را، بنا بر کلِبش\footnote{همان، \S\ 12.} (و نیز بنا بر پژوهش‌های بعدیِ ف. درویتس و آ. کاپلی؛ نک. مقدمه)، می‌توان به فرم‌هایی فروکاست که از هر یک از $n-1$ ردیفِ $p_\rho$ حداکثر یکی را دربردارند؛ بنابراین می‌توان آن‌ها را به‌شکل نمادی زیر نمایش داد:
\begin{equation}
 (S^1\mid p_1)^{m_1}(S^2\mid p_2)^{m_2}\cdots(S^{n-1}\mid p_{n-1})^{m_{n-1}},
\tag{3}
\end{equation}
که در آن، به قیاس با (2.)، چنین تعریف کرده‌ایم:
\[
 (S^\rho\mid p_\rho)=\sum S^{i_1i_2\ldots i_\rho}p_{i_1i_2\ldots i_\rho}.
\]
ردیف‌های نمادیِ $S^\rho$ را، به‌سبب اتحادهای غیرخطی‌ای که میان عناصر یک ردیف $p_\rho$ برقرارند، می‌توان چنان نرمال کرد که دقیقاً همان اتحادها را برآورده سازند؛ بنابراین همانند دترمینان‌های یک ماتریسِ $\rho$-سطری رفتار می‌کنند که سطرهای $s^{(1)},s^{(2)},\ldots,s^{(\rho)}$ آن نسبت به $x$ پادهموردا هستند.\footnote{\srcLTR{Clebsch,} همان، \S\ 4.} پس می‌توان به هر ردیف $S^\rho$ ردیف دیگری، $S_{n-\rho}$، را از طریق رابطهٔ زیر به‌طور یک‌به‌یک نسبت داد:
\begin{equation}
 S_{i_1i_2\ldots i_{n-\rho}}=(-1)^{\frac{(n-\rho)(n-\rho+1)}2+i_1+i_2+\cdots+i_{n-\rho}}S^{j_1j_2\ldots j_\rho},
\tag{4}
\end{equation}
در اینجا نیز اندیس‌های $i$ و $j$، هر دسته در میان خود، به‌ترتیب صعودی‌اند و روی‌هم دقیقاً همهٔ اعداد 1 تا $n$ را دربرمی‌گیرند. عناصر $S_{i_1i_2\ldots i_{n-\rho}}$ از ردیف $S_{n-\rho}$ همچون دترمینان‌هایی رفتار می‌کنند که از $(n-\rho)$ سطرِ هموردای $x$، یعنی $\sigma^{(1)},\sigma^{(2)},\ldots,\sigma^{(n-\rho)}$، ساخته شده‌اند؛ و سطرهای $s$ و $\sigma$ به‌واسطهٔ $\rho(n-\rho)$ شرط زیر به هم پیوسته‌اند:
\begin{equation}
 (s^{(k)}\mid \sigma^{(l)})=0,\qquad(k=1,2,\ldots,\rho;\ l=1,2,\ldots,n-\rho).
\tag{5}
\end{equation}
در نتیجهٔ روابط (1.) و (4.)، هر عاملِ (3.) نمایش نمادی چهارگانهٔ هم‌ارز زیر را می‌پذیرد:
\begin{equation}
 (S^\rho\mid p_\rho)=(S^\rho q_{n-\rho})=(S_{n-\rho}p_\rho)=(-1)^{\rho(n-\rho)}(q_{n-\rho}\mid S_{n-\rho})\footnotemark,
\tag{6}
\end{equation}
\footnotetext{نک. \srcLTR{Clebsch,} همان، \S\ 5.}
که در اینجا و همواره در ادامه، $(S^\rho q_{n-\rho})$ و $(S_{n-\rho}p_\rho)$ به‌اختصار به‌جای دترمینان‌های $(s^{(1)}\cdots s^{(\rho)}u^{(1)}\cdots u^{(n-\rho)})$ و $(\sigma^{(1)}\cdots\sigma^{(n-\rho)}x^{(1)}\cdots x^{(\rho)})$ به کار می‌روند؛ بسط لاپلاس نشان می‌دهد که این دترمینان‌ها به‌ترتیب فقط به ردیف‌های $S^\rho,q_{n-\rho}$ و $S_{n-\rho},p_\rho$ وابسته‌اند، و نگارش بالا نیز برای اشاره به همین امر است. عبارت‌های $(S^\rho\mid p_\rho)$ و $(q_{n-\rho}\mid S_{n-\rho})$، بنا بر توضیحات بالا، به‌ترتیب حاصل‌ضربِ ماتریس‌های تشکیل‌شده از سطرهای $s$ و $x$ و حاصل‌ضربِ ماتریس‌های تشکیل‌شده از سطرهای $u$ و $\sigma$ را نشان می‌دهند؛ مراد از حاصل‌ضرب ماتریسی، مجموعِ حاصل‌ضرب‌های دترمینان‌های متناظر، یعنی مقدار دترمینانِ ماتریسِ حاصل‌ضرب، است.

عدد مشخصهٔ $\rho(n-\rho)$ را که به یک ردیف نمادی یا ردیف متغیرها تعلق دارد، به پیروی از کلِبش وزن ردیف می‌نامیم،\footnote{همان، \S\ 11.} عدد $\rho$ را که شمار سطرهای $s$ را مشخص می‌کند اندیس وزنی بالا، و عدد $n-\rho$ را که شمار سطرهای $\sigma$ را مشخص می‌کند اندیس وزنی پایین می‌نامیم. از روابط (4.) برمی‌آید که یک ردیف نمادی می‌تواند در یک رابطهٔ واحد، یک بار با اندیس وزنی بالا و بار دیگر با اندیس وزنی پایین ظاهر شود؛ همین حکم دربارهٔ حضور ردیف‌های متناظرِ $p_\rho$ و $q_{n-\rho}$ نیز برقرار است. همچنین یادآور می‌شویم که به‌طور کلی نمادگذاری زیر را به کار می‌بریم: ردیف‌های متغیرِ مرتبهٔ بالاتر را که ماتریسِ مربوط به آن‌ها از سطرهای $x$ تشکیل شده است با $p$، و آن‌هایی را که ماتریسشان از سطرهای $u$ تشکیل شده است با $q$ نشان می‌دهیم؛ ردیف‌های نمادیِ هموردای $x$ را با حروف کوچک یونانی، یعنی $\sigma,\tau,\alpha,\beta$؛ ردیف‌های نمادیِ هموردای $u$ را با حروف کوچک لاتین، یعنی $s,t,a,b$؛ و همهٔ ردیف‌های نمادی دیگر را با حروف بزرگ لاتینِ دارای اندیس وزنی، یعنی $S^\sigma$، $T^\tau$، $A^\rho$، یا به‌ترتیب $S_{n-\sigma}$، $T_{n-\tau}$، $A_{n-\rho}$. برحسب اینکه اندیس وزنی بالا باشد یا پایین، تک‌تک عناصر یک ردیف را نیز با اندیس‌های بالا یا پایین می‌نویسیم؛ چنان‌که مثلاً در (4.) آمده است.

\section*{\S\ 2. نمایش به‌وسیلهٔ حاصل‌ضرب‌های ماتریسی.}
در ادامه، از ``حاصل‌ضرب ماتریسی'' همواره مجموع حاصل‌ضرب‌های دترمینان‌های متناظرِ دو ماتریس با تعداد سطرهای برابر را مراد می‌کنیم، که در آن سطرهای ماتریس دوم نسبت به سطرهای متناظرِ ماتریس نخست پادهموردا هستند. در اینجا سطرهای هر ماتریس می‌توانند: 1. تک‌تک داده شده باشند؛ 2. چنان‌که در (6.) آمده است، در یک ردیف $S^\rho$ گرد آمده باشند؛ 3. یا در ردیف‌های گوناگون $S^{\rho_1},S^{\rho_2}\ldots$ گروه‌بندی شده باشند. حاصل‌ضرب ماتریسی را با نماد $(\ \mid\ )$ به‌کاررفته در (6.) نشان می‌دهیم. نمایش (6.) کاربرد ویژه‌ای از صورت‌بندی دوم قضیهٔ ماتریس‌های متناظر است: ``به هر حاصل‌ضرب ماتریسی با $\rho$ سطر، $(v^{(1)}v^{(2)}\cdots v^{(\rho)}\mid y^{(1)}\cdots y^{(\rho)})$، یک دترمینان نسبت داده می‌شود؛ و برعکس، می‌توان به هر دترمینان یک حاصل‌ضرب ماتریسی با تعداد سطرِ ازپیش‌مقررِ $\rho$ ($\rho=1,2,\ldots,n-1$) نسبت داد''.
\footnote{از قرار معلوم، برای $\rho>n$ حاصل‌ضرب ماتریسی صفر می‌شود؛ برای $\rho=n$ به حاصل‌ضرب دو دترمینان تبدیل می‌شود.}
در واقع، کافی است ردیف‌های $y$ را در یک ردیف $p_\rho$ گرد آوریم (یا ردیف‌های $v$ را در یک ردیف $q_\rho$) و جایگذاری‌های داده‌شده در (1.) را انجام دهیم تا از یک حاصل‌ضرب ماتریسیِ داده‌شده به دترمینان برسیم و برعکس. بااین‌حال، درایه‌های منفرد دترمینان به‌تنهایی مقدار آن را به دست نمی‌دهند، بلکه این مقدار تنها از بسط لاپلاس حاصل می‌شود. پس، چون دترمینان و حاصل‌ضرب ماتریسی هم‌ارز از کار درمی‌آیند، قضیهٔ نمایش‌پذیریِ نمادیِ فرم‌ها (قضیهٔ بنیادی نخست) را نیز می‌توان چنین بیان کرد:

\emph{قضیهٔ \srcLTR{I}:} ``همهٔ ساخته‌های ناوردا را می‌توان به‌طور نمادی با حاصل‌ضرب‌های ماتریسی نمایش داد؛ و برعکس، همهٔ حاصل‌ضرب‌های ماتریسی ساخته‌های ناوردا هستند''.
\footnote{البته ساخته‌های ناوردای فرم‌های بنیادیِ ازپیش‌مقرر تنها آن ترکیب‌هایی از حاصل‌ضرب‌های ماتریسی‌اند که به ردیف‌های نماد $S^\rho\ldots$ وابسته‌اند.}

نمایش با حاصل‌ضرب‌های ماتریسی امکان می‌دهد بر دشواری اساسیِ نمایش غیرصریح، که از نمایش دترمینانی ناشی می‌شود، غلبه کنیم.
\footnote{نک. مقدمه.}
در اینجا از ``نمایش صریح'' نمایشی را مراد می‌کنیم که در آن تنها خودِ ردیف‌های $S^\rho,T^\tau,\ldots$، یا ردیف‌های $R^\rho$ که به‌طور یکتا از آن‌ها به دست می‌آیند، وارد می‌شوند و دیگر ردیف‌های جزءِ منفرد $s,t,\ldots$ در آن ظاهر نمی‌شوند. در \S\ 6، برای همهٔ ساخته‌های ناوردایی که تنها به ردیف‌های نماد $S^\rho,T^\tau,\ldots$ وابسته‌اند، نمایشی به دست خواهیم آورد که نسبت به همین ردیف‌های نماد صریح است؛ و از این راه، فرایندهای تولیدکننده، یعنی فرایند تاخوردگی و فرایندهای مشتق‌گیری، را به حوزهٔ $n$-تایی منتقل خواهیم کرد. امکان نمایش صریح از اینجا روشن می‌شود که تنها ساده‌ترین حاصل‌ضرب‌های ماتریسی دقیقاً با دترمینان‌هایی متناظر می‌شوند که در نمایش دترمینانی پدیدارند؛ بنابراین همهٔ حاصل‌ضرب‌های ماتریسی دیگر، ترکیب چند دترمینان گوناگون در یک عبارت واحدند. زیرا اگر ردیف‌های نماد و متغیر را به ردیف‌های منفرد $s,u,x$ تجزیه کنیم، بنا بر قضیهٔ بنیادی نخست در صورت معمولش، به‌ضرورت مجموعی از دترمینان‌ها پدید می‌آید. اثبات اینکه، برعکس، همهٔ مجموع‌های دترمینانی که ساخته‌های ناوردای فرم‌های بنیادیِ ازپیش‌مقرر را نمایش می‌دهند می‌توانند در حاصل‌ضرب‌های ماتریسیِ صریح گرد آورده شوند، با استفاده از قضیهٔ بنیادی دوم (\S\ 6) به دست می‌آید.

اکنون برای آنکه از دو ردیف نماد $S^\sigma,T^\tau$، با کمک ردیف‌های متغیر، ساخته‌ای ناوردا از فرم بنیادیِ $(S^\sigma\mid p_\sigma)(T^\tau\mid p_\tau)$ به دست آوریم که همهٔ ردیف‌ها را به‌صراحت دربر گیرد (فرایند تاخوردگی)، کافی است مطابق تعریف داده‌شده در آغاز این بند، حاصل‌ضرب ماتریسیِ نوع سوم زیر را تشکیل دهیم:
\begin{equation}
 \pair{q_{n-\sigma-\lambda}S^\sigma}{T_{n-\tau}p_{\tau-\lambda}},
 \qquad (\lambda=1,2,\ldots,\tau,n-\sigma\footnotemark).
\tag{7}
\end{equation}
\footnotetext{عدد $\lambda$ همواره تا کوچک‌ترِ دو عددی که در پایان ذکر شده‌اند پیش می‌رود.}
با بسط (7.) داریم:
\begingroup
\setlength{\arraycolsep}{2pt}
\begin{equation}
\begin{aligned}
&\pair{v^{(1)}\cdots v^{(n-\sigma-\lambda)}s^{(1)}\cdots s^{(\sigma)}}
       {\alpha^{(1)}\cdots\alpha^{(n-\tau)}y^{(1)}\cdots y^{(\tau-\lambda)}} \\
&=\sum
\left|\begin{array}{ccc}
 v_{i_1}^{(1)}&\cdots&v_{i_{n-\lambda}}^{(1)}\\
 \vdots&&\vdots\\
 v_{i_1}^{(n-\sigma-\lambda)}&\cdots&v_{i_{n-\lambda}}^{(n-\sigma-\lambda)}\\
 s_{i_1}^{(1)}&\cdots&s_{i_{n-\lambda}}^{(1)}\\
 \vdots&&\vdots\\
 s_{i_1}^{(\sigma)}&\cdots&s_{i_{n-\lambda}}^{(\sigma)}
\end{array}\right|
\left|\begin{array}{ccc}
 \alpha_{i_1}^{(1)}&\cdots&\alpha_{i_{n-\lambda}}^{(1)}\\
 \vdots&&\vdots\\
 \alpha_{i_1}^{(n-\tau)}&\cdots&\alpha_{i_{n-\lambda}}^{(n-\tau)}\\
 y_{i_1}^{(1)}&\cdots&y_{i_{n-\lambda}}^{(1)}\\
 \vdots&&\vdots\\
 y_{i_1}^{(\tau-\lambda)}&\cdots&y_{i_{n-\lambda}}^{(\tau-\lambda)}
\end{array}\right|\\
&=\sum_{i,k,l}\eps\,q_{k_1\ldots k_{n-\sigma-\lambda}}
 S^{k_{n-\sigma-\lambda+1}\ldots k_{n-\lambda}}
 T_{l_1\ldots l_{n-\tau}}
 p^{l_{n-\tau+1}\ldots l_{n-\lambda}},
\end{aligned}
\tag{8}
\end{equation}
\endgroup
یعنی حاصل، عبارتی است که تنها به ردیف‌های $S,T,q,p$ وابسته است. در اینجا $\eps=\pm1$\footnote{این نمادگذاری را در سراسر ادامه حفظ خواهیم کرد.}، و مجموع را باید چنین فهمید که اندیس‌های هر یک از ردیف‌های $S\ldots$ در میان خود به‌ترتیب صعودی‌اند و، برای هر ترکیب $i_1,i_2,\ldots,i_{n-\lambda}$، اندیس‌های $k$ و اندیس‌های $l$، هریک جداگانه، همهٔ جایگشت‌های هنوز ممکنِ اعداد $i$ را طی می‌کنند. عدد $\lambda$ را که در (7.) وارد می‌شود، در حالت $\sigma\geqq\tau$ نقصِ تاخوردگی می‌نامیم؛ علت این قید اخیر در \S\ 6 روشن خواهد شد.

می‌توان دترمینان‌های دو ماتریسِ واردشده در (7.) را عناصر ردیف‌های تازهٔ $Q^{n-\lambda},P_{n-\lambda}$ دانست، و ردیف‌های متناظرِ آن‌ها را $Q_\lambda,P^\lambda$ گرفت. بنا بر (8.)، ردیف‌های $Q,P$ را می‌توان به‌ترتیب برحسب ردیف‌های اصلیِ $S$ و $q$، و $T$ و $p$ بیان کرد. این مطلب را با معادلات زیر نشان می‌دهیم:
\begin{equation}
\begin{aligned}
 q_{n-\sigma-\lambda}S^\sigma&=Q^{n-\lambda}\sim Q_\lambda, &\quad\text{یا به عبارت دیگر}\quad Q_\lambda&=[q_{n-\sigma-\lambda}S^\sigma]_\lambda,\\
 T_{n-\tau}p_{\tau-\lambda}&=P_{n-\lambda}\sim P^\lambda, &\quad\text{یا به عبارت دیگر}\quad P^\lambda&=[T_{n-\tau}p_{\tau-\lambda}]^\lambda.
\end{aligned}
\tag{9}
\end{equation}
با درنظرگرفتن (4.)، حاصل‌ضرب ماتریسیِ (7.)، همانند (6.)، نمایش نمادیِ چهارگانهٔ هم‌ارز زیر را دارد:
\begin{equation}
\pair{q_{n-\sigma-\lambda}S^\sigma}{T_{n-\tau}p_{\tau-\lambda}}
=(q_{n-\sigma-\lambda}S^\sigma P^\lambda)=(Q_\lambda T_{n-\tau}p_{\tau-\lambda})=(-1)^{\lambda(n-\lambda)}\pair{P^\lambda}{Q_\lambda}.
\tag{10}
\end{equation}
جایگذاری دیگری از ردیف‌های تازه را نیز می‌توان در (7.) انجام داد، بدین‌گونه که بگذاریم:
\begin{equation}
 \pair{q_{n-\sigma-\lambda}S^\sigma}{T_{n-\tau}p_{\tau-\lambda}}
 =\pair{R^{\sigma+\lambda}}{p_{\sigma+\lambda}}(R^{\tau-\lambda}p_{\tau-\lambda}).
\tag{11}
\end{equation}
در اینجا نیز (8.) نشان می‌دهد که حاصل‌ضرب‌های عناصر ردیف‌های $R$ را می‌توان به‌طور یکتا برحسب ردیف‌های $S$ و$T$ بیان کرد. جایگذاری‌های (9.) و (11.) به‌ویژه هنگامی به کار می‌روند که بخواهیم با استفاده از (7.) تاخوردگی دیگری با ردیف‌های دیگر انجام دهیم.

نمایش ماتریسی ما را قادر می‌سازد روابط درجهٔ دو میان عناصر یک ردیف $p_\rho$ را ــ که چنان‌که می‌دانیم همهٔ روابط دیگر از آن‌ها نتیجه می‌شوند ــ به‌صورت ناوردا بیان کنیم؛ یعنی، بنا بر \S\ 1، همان روابطی را که در ضرایب فرم بنیادی خطی‌اند و ردیف‌های $S^\rho$ آن‌ها را برآورده می‌کنند. این‌ها همانی‌های زیرند:
\begin{equation}
 \pair{q_{n-\rho-\lambda}S^\rho}{S_{n-\rho}p_{\rho-\lambda}}=0\footnotemark,
 \qquad (\lambda=1,2,\ldots,\rho,n-\rho),
\tag{12}
\end{equation}
\footnotetext{نک. افزوده‌ای از \srcLTR{Study} در \srcLTR{Graßmanns Werke, Bd. I, Teil 2, S. 510} (شرطِ ساده‌بودنِ کمیت $S^\rho$).}
که در آن $p$ و$q$ را باید کمیت‌هایی دلخواه دانست. در \S\ 5 خواهیم دید که همانی‌های دارای عددِ نقصِ فردِ $\lambda$ از همانی‌های دارای نقص بالاتر نتیجه می‌شوند. روابط (12.) نتیجهٔ مستقیم روابط (5.) هستند؛ زیرا داریم:
\[
 \pair{q_{n-\rho-\lambda}S^\rho}{S_{n-\rho}p_{\rho-\lambda}}
 =(v^{(1)}\cdots v^{(n-\rho-\lambda)}s^{(1)}\cdots s^{(\rho)}\mid \sigma^{(1)}\cdots\sigma^{(n-\rho)}y^{(1)}\cdots y^{(\rho-\lambda)}),
\]
و به‌کارگیریِ قضیهٔ حاصل‌ضرب، در نتیجهٔ (5.)، دترمینانی با $(n-\lambda)$ سطر می‌دهد که صفر است. همانی‌های (12.) می‌گویند برای یافتن همهٔ انواع تاخوردگی کافی است تاخوردگی را بر حاصل‌ضرب فرم‌هایی اعمال کنیم که در هر ردیف متغیر خطی‌اند؛ یا، به بیان هم‌ارز بنا بر \S\ 6، فرم نرمال‌شدهٔ $(S^\rho\mid p_\rho)^m$ هیچ ساختهٔ ناوردای خطی در ضرایب ندارد.

شرط‌های (12.) برای مشخصه‌سازی ردیف‌های $p_\rho$ کافی نیز هستند. زیرا ویژگیِ دترمینان‌های یک ماتریس بودنِ عناصر منفردِ $p_\rho$ ناورداست، و ازاین‌رو باید بتوان آن را نیز به‌صورت ناوردا بیان کرد؛ اما شرط‌های (12.) بنا بر قضیهٔ \srcLTR{VI} بزرگ‌ترین قید ناوردای ممکن بر $p_\rho$ را به دست می‌دهند، یعنی اینکه فرم خطیِ ساخته‌شده از ردیف $p_\rho$ به‌عنوان ضرایب، اصلاً هیچ ساختهٔ ناوردایی ندارد.

\section*{\S\ 3. همانی‌های نمادی\footnote{برای صورت‌بندی مسئله و منابع آن، نک. مقدمه.}.}
اکنون باید قضیهٔ بنیادی دومِ روش نمادی را انتقال دهیم؛ یعنی کلِ همانی‌های کاهش‌ناپذیر و تجزیه‌ناپذیر را تعیین کنیم که در آن‌ها همهٔ ردیف‌های $S^\rho,p_\rho$ تجزیه‌ناپذیر در نظر گرفته می‌شوند. در این میان قرارداد زیر را نیز می‌گذاریم: با توجه به معنای ردیف‌های نمادی و ردیف‌های متغیرها، همانی‌هایی را مرکب می‌دانیم که از تخصص‌دادنِ ردیف متغیرهای $p_\rho$ به $p_{\rho-1}y$ و سپس اعمال یکی از همانی‌های دستگاه بر عبارت‌های حاصل پدید می‌آیند؛ در مقابل، همانی‌هایی که $k$ ردیف نمادی دارند تجزیه‌ناپذیر به‌شمار می‌آیند، حتی اگر با فرایند بالا از همانی‌هایی با ردیف‌های نمادی کمتر به دست آمده باشند. لازم نیست ردیف‌های $S^\rho,p_\rho$ به‌طور صریح در همانی‌ها وارد شوند؛ برای آنکه آن‌ها را کمیت‌هایی تجزیه‌ناپذیر بدانیم، کافی است نشان دهیم حاصل‌ضرب‌های ماتریسیِ پدیدارشونده تنها به این ردیف‌ها وابسته‌اند. بااین‌حال نشان خواهیم داد که در این حالت همواره می‌توان، با به‌کارگیریِ جایگذاریِ (11.) در بعضی مواضع، به نمایشی صریح رسید. همانی‌های عمومی را با اصل نمایش ماتریسی از همانی‌های ویژه‌ای به دست می‌آوریم که آقای پاسکال داده است؛ همانی‌هایی که تنها ردیف‌های $x$ و$u$ را دربر دارند و انتقالی مستقیم از همانی‌هایی‌اند که آقای اشتودی برای حوزهٔ سه‌تایی داده است:
\begin{align}
\sum_{i=1}^n(-1)^{i+1}\pair{a^{(i)}}{x}(a^{(1)}\cdots a^{(i-1)}a^{(i+1)}\cdots a^{(n)}u)&=(-1)^{n+1}\pair{u}{x}(a^{(1)}\cdots a^{(n)}),\tag{13}\\
\sum\pm\pair{a^{(1)}}{x^{(1)}}\pair{a^{(2)}}{x^{(2)}}\cdots\pair{a^{(n)}}{x^{(n)}}&=(a^{(1)}\cdots a^{(n)})(x^{(1)}\cdots x^{(n)}),\tag{14}\\
\sum_{i=1}^n(-1)^{i+1}\pair{u}{a^{(i)}}(a^{(1)}\cdots a^{(i-1)}a^{(i+1)}\cdots a^{(n)}x)&=(-1)^{n+1}\pair{u}{x}(a^{(1)}\cdots a^{(n)}).\tag{15}
\end{align}
دو همانی دیگر به‌ترتیب از (13.) و(15.) با جایگذاری‌های $(a^{(i)}\mid x)=(a^{(i)}q_{n-1})$ و$(u\mid a^{(i)})=(p_{n-1}a^{(i)})$ به دست می‌آیند و بنابراین، از دیدگاه ما، اساساً متفاوت به‌شمار نمی‌روند. رابطهٔ (14.) قضیهٔ حاصل‌ضربِ دترمینان‌ها را بیان می‌کند؛ نشان خواهیم داد که چگونه می‌توان (13.) و(14.) ــ و به‌طور دوگان (15.) و(14.) ــ را با دستگاهِ قضیه‌های حاصل‌ضربِ ماتریس‌ها که به شیوه‌ای مناسب ساخته شده‌اند جایگزین کرد. قضیهٔ حاصل‌ضرب، یک بار برای ماتریس‌های $\rho$ سطری و یک بار برای ماتریس‌های $(\rho-1)$ سطری، چنین است:
\[
\pair{a^{(1)}a^{(2)}\cdots a^{(\rho)}}{xp_{\rho-1}}
=(a^{(1)}a^{(2)}\cdots a^{(\rho)}\mid xy^{(2)}\cdots y^{(\rho)})
=\sum\pm\pair{a^{(1)}}{x}\pair{a^{(2)}}{y^{(2)}}\cdots\pair{a^{(\rho)}}{y^{(\rho)}},
\]
\[
\sum\pm\pair{a^{(2)}}{y^{(2)}}\cdots\pair{a^{(\rho)}}{y^{(\rho)}}
=(a^{(2)}\cdots a^{(\rho)}\mid y^{(2)}\cdots y^{(\rho)})
=(a^{(2)}\cdots a^{(\rho)}p_{\rho-1})=(a^{(2)}\cdots a^{(\rho)}q_{n-\rho+1});
\]
از این دو، دستگاه قضیه‌های حاصل‌ضربِ زیر نتیجه می‌شود، که در آن $\rho=2,3,\ldots,n$:
\begin{equation}
\begin{aligned}
\pair{a^{(1)}a^{(2)}\cdots a^{(\rho)}}{xp_{\rho-1}}
 &=\sum_{i=1}^{\rho}(-1)^{i+1}\pair{a^{(i)}}{x}(a^{(1)}\cdots a^{(i-1)}a^{(i+1)}\cdots a^{(\rho)}\mid p_{\rho-1})\\
 &=\sum_{i=1}^{\rho}(-1)^{i+1}\pair{a^{(i)}}{x}(a^{(1)}\cdots a^{(i-1)}a^{(i+1)}\cdots a^{(\rho)}q_{n-\rho+1}).\footnotemark
\end{aligned}
\tag{16}
\end{equation}
\footnotetext{همانی‌های (16.) را همچنین می‌توان همانی‌های (13.) دانست که از حوزهٔ $\rho$ متغیر انتقال یافته‌اند.}
برای $\rho=n$، رابطهٔ (16.) به (13.) تبدیل می‌شود؛ اگر در ادامه $p_{n-1}=y^{(2)}p_{n-2}$، و$p_{n-2}=y^{(3)}p_{n-3}$، و به همین ترتیب، تخصص دهیم و همانی‌های (16.) را بر عبارت‌های $(a^{(1)}\cdots a^{(i-1)}a^{(i+1)}\cdots a^{(n)}\mid y^{(2)}p_{n-2})$، و$(a^{(1)}\cdots a^{(i-1)}a^{(i+1)}\cdots a^{(k-1)}a^{(k+1)}\cdots a^{(n)}\mid y^{(3)}p_{n-3})$، و به همین ترتیب، اعمال کنیم، از (13.) به (14.) می‌رسیم. بنابراین، طبق تعریف ما، همانی (14.) از همانی‌های (16.) مرکب است؛ یا به بیان دیگر، دستگاه (16.) با همانی‌های (13.) و(14.) هم‌ارز است. دستگاه (16.) یکی از شرط‌های خواسته‌شده برای همانی‌های عمومی را برآورده می‌کند؛ این دستگاه ردیف $p_{\rho-1}$ را هم به‌عنوان کمیتی تجزیه‌ناپذیر و هم به‌صورت صریح دربردارد. این تنها دستگاهی است که همین شرط و فقط همین شرط را برآورده می‌کند. زیرا از ردیف‌های $a,x,p$ نمی‌توانیم دترمینان یا حاصل‌ضرب ماتریسیِ دیگری بسازیم؛ میان این‌ها نیز هیچ رابطه‌ای مستقل از (16.) نمی‌تواند برقرار باشد، زیرا در غیر این صورت با حذفِ $(a^{(1)}\cdots a^{(\rho)}\mid xp_{\rho-1})$ رابطه‌ای میان $\rho$ عبارت، با شرطِ $\rho\leqq n$، از نوعِ $(a^{(i)}\mid x)(a^{(1)}\cdots a^{(i-1)}a^{(i+1)}a^{(\rho)}q_{n-\rho+1})$ پدید می‌آمد، حال‌آنکه بنا بر (13.) دست‌کم $n+1$ عبارت از این نوع باید در یک رابطه وارد شوند. از ملاحظات مشابه، دستگاهِ زیر که با (15.) و(14.) هم‌ارز است به دست می‌آید:
\begin{equation}
\begin{aligned}
\pair{u q_{\rho-1}}{a^{(1)}\cdots a^{(\rho)}}
 &=\sum_{i=1}^{\rho}(-1)^{i+1}\pair{u}{a^{(i)}}(q_{\rho-1}\mid a^{(1)}\cdots a^{(i-1)}a^{(i+1)}\cdots a^{(\rho)})\\
 &=\sum_{i=1}^{\rho}(-1)^{i+1}\pair{u}{a^{(i)}}(p_{n-\rho+1}a^{(1)}\cdots a^{(i-1)}a^{(i+1)}\cdots a^{(\rho)}).
\end{aligned}
\tag{17}
\end{equation}
برای رسیدن از (16.) به همانی‌هایی میان ردیف‌های دلخواهِ $A^{\rho_1},B^{\rho_2},\ldots,x,p$، توجه می‌کنیم که این همانی‌ها باید به‌محض آنکه آن‌ها را به‌صورت $A^{\rho_1}=a^{(1)}a^{(2)}\cdots a^{(\rho_1)}$، و$B^{\rho_2}=b^{(1)}b^{(2)}\cdots b^{(\rho_2)}$، و به همین ترتیب، تجزیه کنیم، با (16.) یکسان شوند؛ و هرگاه عبارت‌های منفرد از نوعِ عبارت‌های واردشده در (16.) باشند، تنها با عملیاتی که (16.) تعیین می‌کند می‌توان آن‌ها را در عبارت‌های نهایی گرد آورد. پس اگر همچنین بگذاریم
\begin{equation}
 \rho=\rho_1+\rho_2+\cdots+\rho_k\leqq n
\tag{18}
\end{equation}
به دست می‌آوریم:
\begin{align*}
&\pair{a^{(1)}\cdots a^{(\rho_1)}b^{(1)}\cdots b^{(\rho_2)}\cdots k^{(1)}\cdots k^{(\rho_k)}}{xp_{\rho-1}}
  =\pair{A^{\rho_1}B^{\rho_2}\cdots K^{\rho_k}}{xp_{\rho-1}} \\
&\quad=\sum_{i=1}^{\rho_1}(-1)^{i+1}\pair{a^{(i)}}{x}(a^{(1)}\cdots a^{(i-1)}a^{(i+1)}\cdots a^{(\rho_1)}B^{\rho_2}\cdots K^{\rho_k}q_{n-\rho+1})\\
&\qquad+(-1)^{\rho_1\rho_2}\sum_{i=1}^{\rho_2}(-1)^{i+1}\pair{b^{(i)}}{x}(b^{(1)}\cdots b^{(i-1)}b^{(i+1)}\cdots b^{(\rho_2)}A^{\rho_1}\cdots K^{\rho_k}q_{n-\rho+1})\\
&\qquad+\cdots\\
&\qquad+(-1)^{(\rho_1+\cdots+\rho_{k-1})\rho_k}\sum_{i=1}^{\rho_k}(-1)^{i+1}\pair{k^{(i)}}{x}(k^{(1)}\cdots k^{(i-1)}k^{(i+1)}\cdots k^{(\rho_k)}A^{\rho_1}\cdots (K-1)^{\rho_{k-1}}q_{n-\rho+1}).
\end{align*}
اکنون هر یک از این مجموع‌ها (اما نه بخشی از یک مجموع) واقعاً ردیف‌های $A^{\rho_1},B^{\rho_2}\ldots$ را به‌عنوان کمیت‌هایی تجزیه‌ناپذیر دربردارد؛ زیرا شرط‌های لازم و کافیِ زیر را برای این منظور برآورده می‌کند:\footnote{\srcLTR{Capelli}، در همان مأخذ، \S\ 23، شرط‌های کافیِ $\bigl(a^{(k)}\mid \frac{\partial}{\partial a^{(i)}}\bigr)=0$، $(k>i,\ i=1,2,\ldots,\rho_1-1)$ را به دست می‌دهد. شرط‌های لازم و کافی از این‌ها به ساده‌ترین راه، با به‌کارگیریِ فرایند براکت پواسون بنا بر \srcLTR{Wellstein, Von den Differentialgleichungen der projektiven Invarianten, Math. Ann. 67 (1909), \S\ 3}، نتیجه می‌شوند.}
\begin{equation}
 \left(a^{(i+1)}\middle|\frac{\partial}{\partial a^{(i)}}\right)=0;\ \ldots\ ;
 \left(k^{(j+1)}\middle|\frac{\partial}{\partial k^{(j)}}\right)=0.
 \quad (i=1,2,\ldots,\rho_1-1;\ j=1,2,\ldots,\rho_k-1)
\tag{19}
\end{equation}
نشان می‌دهیم که این مجموع همهٔ ردیف‌ها را نیز به‌طور صریح دربردارد. زیرا اگر $B^{\rho_2}\cdots K^{\rho_k}q_{n-\rho+1}=q_{n-\rho_1+1}$ بگذاریم، بنا بر (16.) و(10.) داریم:
\begin{align*}
&\sum(-1)^{i+1}\pair{a^{(i)}}{x}(a^{(1)}\cdots a^{(i-1)}a^{(i+1)}\cdots a^{(\rho_1)}q_{n-\rho+1})\\
&=\pair{A^{\rho_1}}{xp_{\rho_1-1}}\\
&=(A_{n-\rho_1}xp_{\rho_1-1})=(-1)^{(\rho_1+1)(n+1)}\pair{q_{n-\rho_1+1}}{A_{n-\rho_1}x}\\
&=(-1)^{(\rho_1+1)(n+1)}\pair{B^{\rho_2}\cdots K^{\rho_k}q_{n-\rho+1}}{A_{n-\rho_1}x}.
\end{align*}
اگر مجموع‌های دیگر را نیز به‌طور مشابه محاسبه کنیم و سپس، چون ردیف‌ها از پیش به‌وسیلهٔ اندیس‌های وزنی از یکدیگر متمایز شده‌اند، $B^{\rho_2}=A^{\rho_2}$، و$C^{\rho_3}=A^{\rho_3},\ldots$ بگذاریم، همانی‌های مطلوب را به دست می‌آوریم:
\begin{equation}
\begin{aligned}
\pair{A^{\rho_1}A^{\rho_2}\cdots A^{\rho_k}}{xp_{\rho-1}}
&=\sum_{i=1}^{k}(-1)^{\delta_i}\pair{A^{\rho_1}\cdots A^{\rho_{i-1}}A^{\rho_{i+1}}\cdots A^{\rho_k}q_{n-\rho+1}}{A_{n-\rho_i}x},\\
\delta_i&=(\rho_1+\rho_2+\cdots+\rho_{i-1})\rho_i+(\rho_i+1)(n+1).
\end{aligned}
\tag{20}
\end{equation}
و از (17.) همانی‌های دوگانِ زیر را به دست می‌آوریم:
\begin{equation}
\begin{aligned}
\pair{u q_{\sigma-1}}{A_{\sigma_1}A_{\sigma_2}\cdots A_{\sigma_k}}
&=\sum_{i=1}^{k}(-1)^{\eps_i}\pair{u A^{n-\sigma_i}}{p_{n-\sigma+1}A_{\sigma_1}\cdots A_{\sigma_{i-1}}A_{\sigma_{i+1}}\cdots A_{\sigma_k}},\\
\eps_i&=(\sigma_1+\sigma_2+\cdots+\sigma_{i-1})\sigma_i+(\sigma_i+1)(n+\sigma).
\end{aligned}
\tag{21}
\end{equation}
\emph{قضیهٔ \srcLTR{II}.} همانی‌های (20.) و(21.)، به‌ترتیب، کلِ همانی‌های تجزیه‌ناپذیر میان ردیف‌های $A^{\rho_i},x,p$ و میان ردیف‌های $A_{\sigma_i},u,q$ را دربر می‌گیرند؛ و هم‌زمان کلِ آن همانی‌های تجزیه‌ناپذیری را نیز دربر می‌گیرند که بدون انجام جایگذاریِ (11.) نمایش صریح می‌پذیرند.

بخش اول قضیه از ملاحظات مربوط به اشتقاق (20.) و(21.) به دست می‌آید؛ بخش دوم از ناممکن‌بودنِ وجودِ همانی‌ای میان دو ردیف نماد:
\[
\pair{q_\rho A^\sigma}{B_{\rho+\sigma-\tau}p_\tau}
=\eps_1\pair{q_\rho B^{n-\rho-\sigma+\tau}}{A_{n-\sigma}p_\tau}
+\eps_2\pair{q_\rho q_{n-\tau}}{B_{\rho+\sigma-\tau}A_{n-\sigma}},
\]
که در آن $\rho\leqq\tau\leqq\sigma$ و هیچ‌یک از ردیف‌ها وزن $1\cdot(n-1)$ ندارد. (فرض‌های دیگر دربارهٔ $\rho,\tau,\sigma$ را می‌توان با جابه‌جا کردنِ نام‌های ردیف‌ها به این فرض‌ها بازگرداند.) این ناممکنی، برای مثال، از بسط (8.) در تخصص $q_\rho=lM^{\rho-1}$، و$q_{n-\tau}=lN^{n-\tau-1}$ نتیجه می‌شود، اگر $l_1=1$، و$M^{2,3,\ldots,\rho}=1$، و$A^{\rho+1,\ldots,\rho+\sigma}=1$ بگذاریم و همهٔ مؤلفه‌های دیگر ردیف‌های $l,M,A$ را صفر قرار دهیم، در حالی که $N$ و$B$ دلخواه‌اند. از آنجا که اکنون همانی‌های میان ردیف‌های دلخواه، با شکستن ردیف $p_\tau$ به $y^{(1)}y^{(2)}\cdots y^{(\tau)}$، به همانی‌هایی تبدیل می‌شوند که از (20.) مرکب‌اند، بخش دوم ادعا به‌محض آنکه بتوان (20.) را از همانی‌هایی با فقط دو ردیف نماد تولید کرد، کاملا اثبات شده است. در واقع، از
\begin{equation}
\pair{A^{\rho_1}A^{\rho_2}}{xp_{\rho-1}}
=\eps\pair{A^{\rho_1}q_{n-\rho+1}}{A_{n-\rho_2}x}
+\pair{A^{\rho_1}}{xp_{\rho_1-1}},
\quad \text{که در آن }p_{\rho_1-1}\sim A^{\rho_2}q_{n-\rho+1},
\tag{20a}
\end{equation}
با تخصص‌دادنِ $A^{\rho_1}=B^{\sigma_1}B^{\sigma_2}$، یعنی با
\[
\pair{B^{\sigma_1}B^{\sigma_2}}{xp_{\rho_1-1}}
=\eps\pair{B^{\sigma_1}q_{n-\rho_1+1}}{B_{n-\sigma_2}x}
+\pair{B^{\sigma_1}}{xp_{\sigma_1-1}},
\quad \text{که در آن } p_{\sigma_1-1}\sim B^{\sigma_2}q_{n-\rho_1+1},
\]
یک همانی از نوع (20.) با سه ردیف نماد به دست می‌آید؛ و با تخصص‌دادن‌های بیشترِ $B^{\sigma_1}=C^{\tau_1}C^{\tau_2}$، و به همین ترتیب، همانی‌های عمومی (20.) به دست می‌آیند. بنابراین از نبودِ نمایش صریحِ یک همانی میان دو ردیف نماد و دو ردیف متغیر دلخواه می‌توان نبودِ چنین نمایشی را برای هر تعداد دلخواه ردیف نماد نتیجه گرفت، مشروط بر اینکه هیچ ردیف $x$ یا $u$ در میان ردیف‌های متغیر نباشد.

\clearpage
\section*{\S\ 4. همانی‌های تجزیه.}
اکنون عام‌ترین حالتِ همانی‌های میان ردیف‌های دلخواهِ نماد و متغیر را بررسی می‌کنیم؛ همانی‌هایی که در عبارت‌های نهاییِ آن‌ها، بی‌آنکه جایگذاریِ (11.) انجام شود، ردیف $p_\tau$ به ردیف‌های منفردِ $y^{(1)}\cdots y^{(\tau)}$ تجزیه‌شده ظاهر می‌شود، و این همانی‌ها را \glqq{}همانی‌های تجزیه\grqq{} می‌نامیم. در پیِ ملاحظهٔ پایانِ بند پیشین، فعلاً آن‌ها را برای حالتی تعیین می‌کنیم که تنها دو ردیف نماد دارد؛ یعنی عبارت $\pair{q_\rho A^\sigma}{B_{\rho+\sigma-\tau}p_\tau}$ را بسط می‌دهیم. در اینجا می‌گذاریم:
\begin{equation}
\begin{aligned}
q_{\rho-\alpha}^{(i_1\ldots i_\alpha)}&\sim p_{n-\rho}\,y^{(i_1)}\cdots y^{(i_\alpha)},\qquad
p_{\tau-\alpha}^{(i_1\ldots i_\alpha)}=y^{(i_{\alpha+1})}\cdots y^{(i_\tau)},\qquad
p_\tau=y^{(1)}\cdots y^{(\tau)}.
\end{aligned}
\tag{22}
\end{equation}
در ارتباط با اعداد $\rho,\sigma,\tau$ باید چهار حالت را از یکدیگر متمایز کرد:
\[
\begin{array}{llll}
1)&\rho+\sigma\leqq n,& \tau\leqq\rho,& \tau\leqq\sigma;\\[2pt]
2)&\rho+\sigma\leqq n,& \tau\geqq\rho,& \tau\leqq\sigma;\\[2pt]
3)&\rho+\sigma\leqq n,& \tau\leqq\rho,& \tau\geqq\sigma;\\[2pt]
4)&\rho+\sigma\leqq n,& \tau\geqq\rho,& \tau\geqq\sigma,
\end{array}
\]
از میان این حالت‌ها، حالت نخست را برمی‌گزینیم تا سپس سایر حالت‌ها را به‌اختصار مشخص کنیم. بنا بر (20.)، (10.) و (9.)، اگر ردیف $y^{(1)}y^{(2)}\cdots y^{(\tau)}B_{\rho+\sigma-\tau}$ را به $y^{(1)}$ و $y^{(2)}\cdots y^{(\tau)}B_{\rho+\sigma-\tau}$ تجزیه کنیم، داریم:
\begin{equation}
\begin{aligned}
\pair{q_\rho A^\sigma}{y^{(1)}y^{(2)}\cdots y^{(\tau)}B_{\rho+\sigma-\tau}}
&=\pair{q_{\rho-1}^{(1)}A^\sigma}{y^{(2)}\cdots y^{(\tau)}B_{\rho+\sigma-\tau}}\\
&\quad+(-1)^\rho\pair{q_\rho[A_{n-\sigma}y^{(1)}]^{\sigma-1}}{y^{(2)}\cdots y^{(\tau)}B_{\rho+\sigma-\tau}}.
\end{aligned}
\tag{23}
\end{equation}
اکنون جملهٔ آخرِ (23.) دقیقاً صورتِ عبارت اولیه را دارد ــ فقط $\sigma$ با $\sigma-1$ و $\tau$ با $\tau-1$ جایگزین شده‌اند ــ، و بنابراین معادلهٔ (23.) را دوباره می‌توان بر آن اعمال کرد. پس چون $\tau\leqq\sigma$، می‌توانیم عملِ (23.) را $\tau$ بار اعمال کنیم و با در نظر گرفتنِ (10.) به رابطهٔ زیر برسیم:
\begin{equation}
\begin{aligned}
\pair{q_\rho A^\sigma}{p_\tau B_{\rho+\sigma-\tau}}
&=(-1)^{\rho\sigma+(\sigma+\tau)(n+1)}
  \pair{q_\rho B^{n-\rho-\sigma+\tau}}{A_{n-\sigma}p_\tau}\\
&\quad+\sum_{i_1=1}^{\tau}(-1)^{\rho(i_1-1)}
 \pair{q_{\rho-1}^{(i_1)}[A_{n-\sigma}y^{(1)}\cdots y^{(i_1-1)}]^{\sigma-i_1+1}}
 {y^{(i_1+1)}\cdots y^{(\tau)}B_{\rho+\sigma-\tau}}.
\end{aligned}
\tag{24}
\end{equation}
اکنون هر جملهٔ زیرِ علامت جمع دوباره از نوعِ عبارت اولیه است؛ $\rho$ با $\rho-1$، $\sigma$ با $\sigma-i_1+1$ و $\tau$ با $\tau-i_1$ جایگزین شده‌اند، پس شرط 1) حتی به‌نحو قوی‌تری برقرار است. ازاین‌رو، چون $\tau\leqq\rho$، می‌توانیم عملِ (24.) را $\tau$ بار اعمال کنیم و همانیِ تجزیهٔ مطلوب را به دست آوریم:
\begin{equation}
\begin{aligned}
\pair{q_\rho A^\sigma}{p_\tau B_{\rho+\sigma-\tau}}
&=\eps_{i_0}\pair{q_\rho B^{n-\rho-\sigma+\tau}}{A_{n-\sigma}p_\tau}\\
&\quad+\sum_{i_1=1}^{\tau}\eps_{i_1}
 \pair{q_{\rho-1}^{(i_1)}B^{n-\rho-\sigma+\tau}}{A_{n-\sigma}p_{\tau-1}^{(i_1)}}\\
&\quad+\sum_{i_2=i_1+1}^{\tau}\eps_{i_2}
 \pair{q_{\rho-2}^{(i_1i_2)}B^{n-\rho-\sigma+\tau}}{A_{n-\sigma}p_{\tau-2}^{(i_1i_2)}}+\cdots\\
&\quad+\eps_{i_\tau}\pair{q_{\rho-\tau}^{(i_1\ldots i_\tau)}B^{n-\rho-\sigma+\tau}}{A_{n-\sigma}}\\
&=\sum_{\alpha=0}^{\tau}\sum_{i_\alpha=i_{\alpha-1}+1}^{\tau}\eps_{i_\alpha}
 \pair{q_{\rho-\alpha}^{(i_1\ldots i_\alpha)}B^{n-\rho-\sigma+\tau}}{A_{n-\sigma}p_{\tau-\alpha}^{(i_1\ldots i_\alpha)}}.
\end{aligned}
\tag{25}
\end{equation}
عاملِ علامتِ $\eps$ از (24.) چنین به دست می‌آید:
\begin{equation}
\begin{aligned}
\eps_{i_\alpha}&=(-1)^{\delta_{i_\alpha}},\qquad
\delta_{i_\alpha}=\sum_{\beta=1}^{\alpha}(\rho-\beta+1)(i_{\beta-1}+i_\beta+1)\\
&\quad +(\rho-\alpha)(\sigma+i_\alpha+\alpha)+(\sigma+\tau+\alpha)(n+1),
\end{aligned}
\tag{26}
\end{equation}
که در آن باید $i_0=0$ گذاشت. علامت جمع در (25.) را باید چنین فهمید که به هر ترکیب $i_1i_2\ldots i_{\alpha-1}$، که در آن $i_1<i_2<\cdots<i_{\alpha-1}$، همهٔ مقادیر $i_\alpha$ که $i_{\alpha-1}<i_\alpha\leqq\tau$ را برآورده می‌کنند افزوده می‌شوند. هریک از مجموع‌های (25.)، چنان‌که با در نظر گرفتنِ (26.) به‌آسانی می‌توان دید، معادلات دیفرانسیلِ همانندِ (19.) را برآورده می‌کند:
\begin{equation}
 \left(\frac{\partial}{\partial y^{(i)}}\middle|y^{(i+1)}\right)=0,
 \qquad (i=1,2,\ldots,\tau-1),
\tag{27}
\end{equation}
و بنابراین هر یک فقط به ردیف $p_\tau$ وابسته است. پس همانیِ تجزیه یک همانیِ تجزیه‌ناپذیر میان ردیف‌های $A,B,q,p$ است؛ در بند بعد به نمایش صریحِ آن بازمی‌گردیم.

در حالت 2)، عملِ (23.) را نیز می‌توان $\tau$ بار اعمال کرد، اما عملِ (24.) پس از $\rho$ بار متوقف می‌شود. بنابراین در عبارت نهاییِ (25.) باید نخستین علامت جمع را با $\sum_{\alpha=0}^{\rho}$ جایگزین کرد. در حالت‌های 3) و 4)، عملِ (23.) را فقط می‌توان $\sigma$ بار اجرا کرد؛ در این صورت به‌جای (24.) رابطهٔ زیر قرار می‌گیرد:
\begin{equation}
\begin{aligned}
\pair{q_\rho A^\sigma}{p_\tau B_{\rho+\sigma-\tau}}
&=(-1)^{\rho\sigma}
 \pair{q_{\rho-(\tau-\sigma)}^{(\sigma+1\ldots\tau)}B^{n-\rho+\tau-\sigma}}{A_{n-\sigma}p_{\tau-(\tau-\sigma)}^{(\sigma+1\ldots\tau)}}\\
&\quad+\sum_{i_1=1}^{\sigma}(-1)^{\rho(i_1-1)}
 \pair{q_{\rho-1}^{(i_1)}[A_{n-\sigma}y^{(1)}\cdots y^{(i_1-1)}]^{\sigma-i_1+1}}
 {y^{(i_1+1)}\cdots y^{(\tau)}B_{\rho+\sigma-\tau}};
\end{aligned}
\tag{28}
\end{equation}
تکرارِ (28.) به تعدادِ $(\tau-\sigma)$ بار، که در آن علامت جمع به $\sum_{i_\beta=i_{\beta-1}+1}^{\sigma+\beta-1}$ تبدیل می‌شود ــ چنان‌که اعمالِ (23.) به تعدادِ $(\sigma-i_1+1)$ بار بر جمله‌های زیرِ علامت جمع در (28.) نشان می‌دهد ــ، همهٔ جمله‌های آن مجموع در (25.) را که با اندیس $\alpha=\tau-\sigma$ متناظر است، با علامت متناظر به دست می‌دهد. اما پس از آن، حالت‌های 3) و 4) به حالت‌های 1) و 2) تبدیل می‌شوند، زیرا اعداد متناظر با $\sigma,\tau$ چنین می‌شوند: $\sigma'=\sigma-i_{\tau-\sigma}+\tau-\sigma$، $\tau'=\tau-i_{\tau-\sigma}=\sigma'$؛ و با ادامهٔ روش، $\tau'<\sigma'$ خواهد شد. بنابراین در عبارت نهاییِ (25.) باید نخستین علامت جمع را با $\sum_{\alpha=\tau-\sigma}^{\tau,\rho}$ جایگزین کرد.

به همین قیاس، از (21.) همانیِ تجزیهٔ دوگان به دست می‌آید که در آن ردیف $q_\rho$ به ردیف‌های منفردِ $v^{(1)},\ldots,v^{(\rho)}$ تجزیه می‌شود. اگر بگذاریم
\begin{equation}
 \dot p_{\tau-\beta}^{(i_1\ldots i_\beta)}\sim v^{(i_1)}\cdots v^{(i_\beta)}q_{n-\tau};\qquad
 \dot q_{\rho-\beta}^{(i_1\ldots i_\beta)}=v^{(i_{\beta+1})}\cdots v^{(i_\rho)};\qquad
 q_\rho=v^{(1)}\cdots v^{(\rho)},
\tag{29}
\end{equation}
آنگاه داریم:
\begin{equation}
 \pair{q_\rho A^\sigma}{p_\tau B_{\rho+\sigma-\tau}}
 =\sum_{\beta=\substack{0,\tau-\sigma}}^{\substack{\tau,\rho}}
   \sum_{i_\beta=i_{\beta-1}+1}^{\rho}\eps
\pair{\dot q_{\rho-\beta}^{(i_1\ldots i_\beta)}B^{n-\rho-\sigma+\tau}}{A_{n-\sigma}\dot p_{\tau-\beta}^{(i_1\ldots i_\beta)}} ,
\tag{30}
\end{equation}
که در آن اندیس جمعِ $\beta$ از بزرگ‌ترِ دو مقدار پایین تا کوچک‌ترِ دو مقدار بالا می‌رود. همهٔ جمله‌های هر مجموع منفرد در (25.) و(30.) قطبی‌های یک فرم واحدند که با دو عملیاتِ گرفتن قطبی‌های زیر پدید آمده‌اند:
\[
\left(\frac{\partial}{\partial p_{\tau-\alpha}}\middle|p_{\tau-\alpha}^{(i_1\ldots i_\alpha)}\right),\qquad
\left(q_{\rho-\alpha}^{(i_1\ldots i_\alpha)}\middle|\frac{\partial}{\partial q_{\rho-\alpha}}\right).
\]
برای بیان این امر، $\Pi$ را ترکیبی از چنین عملیاتِ گرفتن قطبی‌ها، ضرب‌شده در ثابت‌ها، می‌گیریم؛ آنگاه به‌جای (25.) و(30.) رابطهٔ زیر را به دست می‌آوریم:
\begin{equation}
 \pair{q_\rho A^\sigma}{p_\tau B_{\rho+\sigma-\tau}}
 =\sum_{\alpha=\substack{0,\tau-\sigma}}^{\substack{\tau,\rho}}
 \Pi\pair{q_{\rho-\alpha}B^{n-\rho-\sigma+\tau}}{A_{n-\sigma}p_{\tau-\alpha}}.
\tag{31}
\end{equation}

\section*{\S\ 5. همانی‌های تجزیه به صورت صریح.}
اگر در (25.)، در حالت 4)، $\tau$ را به $\sigma+\rho$ تخصص دهیم، داریم:
\begin{equation}
 \pair{q_\rho A^\sigma}{y^{(1)}\cdots y^{(\rho+\sigma)}}
 =\sum_{i_\rho=i_{\rho-1}+1}^{\tau}
 \eps_{i_\rho}\pair{q_\rho}{y^{(i_1)}\cdots y^{(i_\rho)}}
 \pair{A^\sigma}{y^{(i_{\rho+1})}\cdots y^{(i_{\rho+\sigma})}},
\tag{32}
\end{equation}
که در آن $\eps_{i_\rho}=(-1)^{\delta_{i_\rho}}$ و$\delta_{i_\rho}=\rho(\rho+1)/2+i_1+\cdots+i_\rho$ می‌شود. این صورت، صورتی عام‌تر از قضیهٔ حاصل‌ضرب ماتریس‌ها در (16.) و(17.) است و از بسط لاپلاسِ دترمینانِ ماتریس حاصل‌ضرب به دست می‌آید:
\[
 \sum\pm(v^{(1)}y^{(1)})\cdots(v^{(\rho)}y^{(\rho)})(a^{(1)}y^{(\rho+1)})\cdots(a^{(\sigma)}y^{(\rho+\sigma)}).
\]
پس قضیهٔ حاصل‌ضربِ عام‌تر، برحسب مورد، از همانی‌های (20.) و(21.) ترکیب شده است و بدین‌ترتیب استخراجِ صورت‌های ویژهٔ (16.) و(17.) توضیح داده می‌شود.

اکنون رابطهٔ (32.) وسیلهٔ نمایش صریحِ همانی‌های تجزیه را فراهم می‌کند. برای این منظور، فرم‌های ظاهرشده در (31.) را فرم‌های بنیادی می‌گیریم؛ یعنی جایگذاریِ (11.) را انجام می‌دهیم. این جایگذاری در نمایش صریح برحسب ردیف‌های $A,B$ تغییری ایجاد نمی‌کند، زیرا بنا بر (8.) ردیف‌های تازه‌واردِ $R$ را می‌توان برحسب خودِ ردیف‌های $A,B$، و نه برحسب ردیف‌های جزءِ آن‌ها، یعنی $a^{(1)}a^{(2)}\ldots$ و$b^{(1)}b^{(2)}\ldots$، بیان کرد. پس اگر
\begin{equation}
 \pair{q_{\rho-\alpha}B^{n-\rho-\sigma+\tau}}{A_{n-\sigma}p_{\tau-\alpha}}
 =\bigl(R_{\rho-\alpha}p_{n-\rho+\alpha}\bigr)
  \pair{R^{\tau-\alpha}}{p_{\tau-\alpha}},
\tag{33}
\end{equation}
آنگاه مجموع منفردِ (25.) که با اندیس $\alpha$ متناظر است چنین می‌شود:
\begin{equation}
\begin{aligned}
&\sum \eps_{i_\alpha}\bigl(R_{\rho-\alpha}p_{n-\rho}y^{(i_1)}\cdots y^{(i_\alpha)}\bigr)
 \pair{R^{\tau-\alpha}}{y^{(i_{\alpha+1})}\cdots y^{(i_\tau)}}\\
&\quad=\sum \eps_{i_\alpha}\bigl([R_{\rho-\alpha}p_{n-\rho}]^\alpha y^{(i_1)}\cdots y^{(i_\alpha)}\bigr)
 \pair{R^{\tau-\alpha}}{y^{(i_{\alpha+1})}\cdots y^{(i_\tau)}}\\
&\quad=\eps\pair{[R_{\rho-\alpha}p_{n-\rho}]^\alpha R^{\tau-\alpha}}{p_\tau}
 =\eps\pair{R^{\tau-\alpha}q_{n-\tau}}{R_{\rho-\alpha}p_{n-\rho}}.
\end{aligned}
\tag{34}
\end{equation}
بدین‌ترتیب اکنون واقعاً یک عبارت نهاییِ صریح به دست آمده است؛ و صورت متقارنی که $q_\rho$ و$p_\tau$ در آن وارد می‌شوند نشان می‌دهد که (30.) نیز به همان عبارت نهایی می‌انجامد، پس این عبارت از تجزیهٔ اولیه مستقل است. تفاوت میان (25.) و(30.) تنها تا زمانی برقرار است که ردیف‌های $y$ و$v$ خود معنای مستقلی داشته باشند، یعنی همانیِ تجزیه، بنا بر تعریف ما، به همانی‌ای تجزیه‌پذیر تبدیل شود. برای رسیدن به همانی‌هایی با ردیف‌های نمادِ بیشتر، بنا بر پایانِ \S\ 3 کافی است ردیف $q_\rho$ را به $q_{\rho_1}C^{\sigma_1}$، یا ردیف $p_\tau$ را به $p_{\tau_1}C_{\tau-\tau_1}$ تجزیه کنیم. عبارت‌های حاصل:
\begin{equation}
 \pair{q_{\rho_1}C^{\sigma_1}}{[R^{\tau-\alpha}q_{n-\tau}]_\alpha R_{\rho_1+\sigma_1-\alpha}},
 \qquad
 \pair{[R_{\rho-\alpha}p_{n-\rho}]^\alpha R^{\tau-\alpha}}{p_{\tau_1}C_{\tau-\tau_1}},
\tag{35}
\end{equation}
اکنون دقیقاً از نوعِ عبارت‌هایی هستند که با دو ردیف نماد پدید می‌آیند؛ پس با جایگذاری‌ای متناظر با (33.) دوباره نمایش صریح می‌پذیرند، و تکرار این کار برای هر تعداد دلخواه از ردیف‌ها نمایش صریح به دست می‌دهد. همچنین باید یادآور شد که، برحسب مورد، نخستین و آخرین مجموع در (25.) و(30.) نمایش صریح را بی‌آنکه (33.) به کار رود، فقط با فرمول (32.) به دست می‌دهند؛ یا اصلاً به یک جملهٔ واحد محدود می‌شوند که همهٔ ردیف‌ها را به‌صراحت دربر دارد. رابطه‌ای که از تجزیهٔ $q_\rho$ به $q_{\rho_1}C^{\sigma_1}$ و مانند آن در (32.) پدید می‌آید، می‌تواند قضیهٔ حاصل‌ضرب در عام‌ترین صورت آن تلقی شود.

در جمع‌بندی نشان خواهیم داد:

\emph{قضیهٔ \srcLTR{III}:} همانی‌های
\begin{equation}
 \pair{q_\rho A^\sigma}{p_\tau B_{\rho+\sigma-\tau}}
 =\sum_{\alpha=\substack{0,\tau-\sigma}}^{\substack{\tau,\rho}}
 \eps\pair{R^{\tau-\alpha}q_{n-\tau}}{R_{\rho-\alpha}p_{n-\rho}},
\tag{36}
\end{equation}
که در آن ردیف‌های $R$ با (33.) تعیین می‌شوند، همراه با همانی‌هایی که با تجزیهٔ $q_\rho,p_\tau$ به $q_{\rho_1}C^{\sigma_1}$ و مانند آن از این‌ها پدید می‌آیند، کلِ همانی‌های تجزیه‌ناپذیر را دربر می‌گیرند.

در واقع، همانی‌های حاصل عام‌ترین همانی‌های نوع خودند، زیرا ردیف‌های منفرد دیگر از حیث وزن و تعداد هیچ محدودیتی ندارند؛ برخلاف آنچه هنوز در (20.) و(21.) برقرار بود. هیچ همانیِ دیگری نیز میان این ردیف‌ها وجود ندارد؛ زیرا با تجزیهٔ یک ردیف $p_\tau$ به $y^{(1)}\cdots y^{(\tau)}$، عام‌ترین همانی‌ها از (20.)، و بنابراین بنا بر پایانِ \S\ 3 از (20a.)، ترکیب می‌شوند؛ و ترکیبی که پس از (16.) بررسی شد دقیقاً همان ترکیبی است که در \S\ 4 انجام گرفت. اما گذار به هر تعداد دلخواه از ردیف‌ها، چنان‌که بحثِ (20a.) نشان می‌دهد، با به‌کارگیریِ همواره همان عملیات بر عبارت‌هایی حاصل می‌شود که از تجزیهٔ $q_\rho$ به $q_{\rho_1}C^{\sigma_1}$ و مانند آن پدید می‌آیند. از اینجا همچنین نتیجه می‌شود که گذار مستقیم از (20.) به‌جای (20a.) چیز تازه‌ای به دست نمی‌دهد؛ این امر را می‌توان به‌آسانی با محاسبه تأیید کرد: یک بار ردیف $q_\rho$ را در (25.) تخصص می‌دهیم و بار دیگر، به روشِ \S\ 4، گذار از (20.) را انجام می‌دهیم. هر دو بار همان مجموع‌ها پدید می‌آیند و این مجموع‌ها با به‌کارگیریِ قضیهٔ حاصل‌ضرب عمومی (نگاه کنید به بالا) به همانی‌های حاصل از (36.) تبدیل می‌شوند. همانی‌های (20.) و(21.) با حالت‌های ویژهٔ $\tau=1$ و$\rho=1$ متناظرند؛ در این صورت در هر مورد فقط دو مجموع منفرد ظاهر می‌شود و بنابراین، مطابق ملاحظهٔ بالا، نمایش صریح بی‌آنکه جایگذاریِ (33.) انجام شود به دست می‌آید، و این با نتیجهٔ پیشین سازگار است.

در پایان، دو حالت ویژهٔ دیگرِ همانیِ تجزیه را به‌اختصار یاد می‌کنیم. اگر در (31.) بگذاریم $\tau=\sigma$ و$\rho=n-\sigma$، و ردیف $p_\tau$ را $p'_\sigma\sim q'_{n-\sigma}$ بنامیم، به دست می‌آید:
\begin{equation}
\begin{aligned}
\pair{A^\sigma}{p_\sigma}\pair{B^\sigma}{p'_\sigma}
-\pair{A^\sigma}{p'_\sigma}\pair{B^\sigma}{p_\sigma}
=\sum_{\alpha=1}^{\sigma,n-\sigma}
\Pi\pair{q_{n-\sigma-\alpha}B^\sigma}{A_{n-\sigma}p_{\sigma-\alpha}}.
\end{aligned}
\tag{37}
\end{equation}
این فرمول کلِبش\footnote{کلِبش، همان اثر، صص. 25--32. کلِبش ثابت می‌کند که جمله‌های منفردِ طرف‌های راست فقط به ردیف‌های $A,B$ وابسته‌اند؛ و نمایش صریح از به‌کارگیریِ (32.) بر عبارت‌های $\Phi_h$ او به دست می‌آید.} برای بسطِ فرمی دارای دو ردیف هموردای $p_\sigma,p'_\sigma$ برحسب همورداهای بنیادی است. در این صورت، همورداهای بنیادیِ متعلق به فرم $\pair{A^\sigma}{p_\sigma}^m\pair{B^\sigma}{p'_\sigma}^n$ چنین‌اند:
\begin{equation}
 \left\{\sum_{\alpha=1}^{\sigma,n-\sigma}
 \pair{q_{n-\sigma-\alpha}B^\sigma}{A_{n-\sigma}p_{\sigma-\alpha}}
 \right\}^{\rho}
 \pair{A^\sigma}{p_\sigma}^{m-\rho}\pair{B^\sigma}{p'_\sigma}^{n-\rho},
 \qquad (\rho=0,1,\ldots,m,n).
\tag{38}
\end{equation}
این‌ها از فرم بنیادی، با آنچه به‌اختصار \glqq{}تاخوردگیِ هموردای دارای نقص $\alpha$\grqq{} خواهیم نامید، پدید می‌آیند (نک. \S\ 2).

دوم، اگر بگذاریم $\tau=\sigma-\lambda$ و$\rho=n-\sigma-\lambda$؛ و $B^\sigma=A^\sigma$، به دست می‌آید:
\begin{equation}
\begin{aligned}
\pair{q_{n-\sigma-\lambda}A^\sigma}{A_{n-\sigma}p_{\sigma-\lambda}}
&=(-1)^\lambda\pair{q_{n-\sigma-\lambda}A^\sigma}{A_{n-\sigma}p_{\sigma-\lambda}}\\
&\quad+(-1)^{(n-\sigma)(\sigma-\lambda)}
 \sum_{\alpha=1}^{\sigma-\lambda,n-\sigma-\lambda}
 \Pi\pair{q_{n-\sigma-\lambda-\alpha}A^\sigma}{A_{n-\sigma}p_{\sigma-\lambda}}.
\end{aligned}
\tag{39}
\end{equation}
در نتیجه، چنان‌که در \S\ 2 یادآوری شد، شرط‌های (12.) که ردیف‌های نماد را مشخص می‌کنند، هنگامی که نقص $\lambda$ فرد است، از شرط‌های دارای نقص بالاتر نتیجه می‌شوند. برای فرم خطی $\pair{A^\sigma}{p_\sigma}\equiv\pair{B^\sigma}{p_\sigma}$، از (39.) نتیجه می‌شود که ساخته‌های ناوردای کاهش‌ناپذیرِ آن که نسبت به ضرایب درجهٔ دو دارند، به ساخته‌هایی با نقص زوج فروکاسته می‌شوند؛ این با واقعیت شناخته‌شده‌ای سازگار است که یک فرم خطی در حوزه‌های دوتایی و سه‌تایی هیچ ساختهٔ ناوردایی ندارد.

یادآور می‌شویم که همانی‌های عمومی در آغاز با این فرض استخراج شدند که ردیف‌های منفرد شرط‌های (12.) را برآورده می‌کنند. اما چون این ردیف‌های منفرد فقط به‌صورت خطی در همانی‌ها وارد می‌شوند، این همانی‌ها برای ردیف‌های عام‌تری که شرط‌های (12.) را برآورده نمی‌کنند نیز برقرارند.

\section*{\S\ 6. نمایش صریحِ ساخته‌های ناوردا.\\فرایندهای تاخوردگی و فرایندهای مشتق‌گیری.}
نتایج بندهای پیشین به ما اجازه می‌دهند قضیهٔ نمایش‌پذیریِ نمادی (قضیهٔ بنیادی نخست) را که در \S\ 2 بیان شد، با افزودن نمایش‌پذیریِ صریح کامل کنیم؛ بدین‌سان که نشان می‌دهیم:

\emph{قضیهٔ \srcLTR{IV}:} ``همهٔ ساخته‌های ناوردای فرم‌های بنیادیِ دلخواه را می‌توان به‌طور صریح با حاصل‌ضرب‌های ماتریسی نمایش داد؛ یعنی در عبارت‌های نهایی، به‌جز ردیف‌های متغیر $p_\rho$، فقط ردیف‌های $S^\sigma,\ldots,T^\tau$ از فرم‌های اصلی یا ردیف‌های $P,Q,R$ که با جایگذاریِ (9.) یا (11.) از آن‌ها به دست آمده‌اند، وارد می‌شوند.''

برای اثبات، فرض می‌کنیم که یک ساختهٔ ناوردا، افزون بر ردیف‌های متغیر، به ردیف‌های $S^\sigma,\ldots,T^\tau$ وابسته باشد؛ و این ردیف‌ها به‌اندازهٔ کافی به ردیف‌های منفرد $S^\sigma=S^{\sigma_1}\cdots S^{\sigma_i},\ldots,T^\tau=T^{\tau_1}\cdots T^{\tau_k}$ تفکیک شده باشند تا نمایش به‌وسیلهٔ دترمینان‌ها ممکن شود. ردیف‌ها را در ترتیب $S^\sigma,\ldots,T^\tau$ می‌گذاریم که فی‌نفسه دلخواه، اما به‌طور معین تثبیت شده است. اکنون یکی از چنین ترکیب‌های دترمینانی را برمی‌گزینیم که به نخستین ردیف کل $S^\sigma=S^{\sigma_1}\cdots S^{\sigma_i}$ وابسته باشد، نه فقط به برخی از ردیف‌های جزءِ آن. اما بنا بر \S\ 3--\S\ 5، این وابستگی نتیجهٔ همانی‌های عمومی است. اگر اکنون ردیف متغیر $p_\rho$ را به ردیف‌های $y,v$، یا یکی از ردیف‌های بعدیِ $T$ را به ردیف‌های منفرد $t$ یا $\tau$ تفکیک کنیم، عام‌ترین همانی‌ها از (20.) و(21.) ترکیب می‌شوند. این همانی‌های اخیر همواره به عبارتی می‌انجامند که ردیف $S^\sigma=S^{\sigma_1}\cdots S^{\sigma_i}$ را به‌طور صریح دربر دارد، یعنی به طرف چپِ (20.) و(21.)؛ زیرا از (19.) روشن می‌شود که تنها مجموع کاملِ طرف راست -- و نه بخشی از آن که با جمله‌ای از (20a.) متناظر است -- این خاصیت را دارد که فقط به $S^\sigma$ وابسته باشد. با ادامهٔ روش، هر ردیفی که یک بار به‌طور صریح وارد شده است همچنان صریح باقی می‌ماند؛ کاربرد همانی‌ها فقط ممکن است یکی‌کردنِ چند ردیف در یک ردیف، یعنی جایگذاریِ (9.)، را ایجاب کند. اگر در پایان فقط یک ردیف $T^\tau$ مانده باشد که باید به صورت صریح درآید، دو حالت ممکن است. ممکن است هنوز یک ردیف متغیر آزاد، یعنی ردیفی که با جایگذاریِ (9.) به ردیف‌های دیگر پیوند ندارد، وجود داشته باشد؛ با تفکیک آن به ردیف‌های جزءِ $y$ یا $v$ می‌توان روش را به پایان رساند، و چنان‌که در (31.)، قطبی‌های یک یا چند فرم پدید می‌آیند که همهٔ ردیف‌ها را به‌طور صریح دربر دارند. اگر دیگر هیچ ردیف متغیر آزادی وجود نداشته باشد، از شیوهٔ پیدایش عبارت‌ها درمی‌یابیم که مجموعهٔ عبارت‌هایی که ردیف‌های منفرد $t$ یا $\tau$ را دربر دارند اما به ردیف $T^\tau$ وابسته‌اند، جمله‌های یک همانیِ تجزیه‌اند؛ و این جمله‌ها بنا بر \S\ 5 و با کاربرد جایگذاریِ (11.) همواره نمایش صریح در همهٔ ردیف‌ها را ممکن می‌سازند. بدین‌ترتیب قضیهٔ بالا در حالت کلی ثابت شد. همچنین یادآور شویم که اگر فقط دو ردیف نماد وجود داشته باشد، جایگذاریِ (11.) هرگز رخ نمی‌دهد؛ زیرا -- به‌سبب $\sigma,n-\sigma,\tau,n-\tau<n$ -- تنها هنگامی هیچ ردیف متغیری وارد نمی‌شود که دو ردیف نماد در یک دترمینان گرد آمده و بنابراین به‌طور صریح ظاهر شده باشند؛ اما به‌محض ورود ردیف‌های متغیر، (34.) و(33.) نشان می‌دهند که جایگذاریِ (11.) زائد است.

از آنچه گفته شد نتیجه می‌شود که برای تولید همهٔ ساخته‌های ناوردا کافی است ردیف‌های نماد را، با افزودن ردیف‌های متغیر، برای تشکیل همهٔ حاصل‌ضرب‌های ماتریسیِ ممکن با هم پیوند دهیم. عام‌ترین حاصل‌ضرب ماتریسی اکنون به صورت زیر است:
\begin{equation}
 \pair{P^{\mu_1}\cdots P^{\mu_i}}{Q_{\nu_1}\cdots Q_{\nu_k}};\qquad \sum\mu=\sum\nu,
\tag{40}
\end{equation}
که در آن $P$ و$Q$ می‌توانند ردیف‌های فرم‌های اصلی باشند، یا از رهگذر توالی‌ای از جایگذاری‌های (9.) یا (11.) از ردیف‌های اصلی پدید آمده باشند، یا سرانجام ردیف‌های متغیر باشند. اما عبارت‌های (40.) را می‌توان با پیوند پیاپی هر دو ردیف نماد در یک حاصل‌ضرب ماتریسی، به یاری ردیف‌های متغیر، تولید کرد؛ زیرا از
\[
 \pair{q_{n-\mu-\lambda}P^\mu}{Q_{\nu_1}p_{n-\nu_1-\lambda}}
 =\bigl(R_\lambda Q_{\nu_1}p_{n-\nu_1-\lambda}\bigr)
\]
با پیوند دادن آن با ردیف $Q_{\nu_2}$، عبارتی به دست می‌آید:
\[
 \pair{[R_\lambda Q_{\nu_1}]^{n-\lambda-\nu_1}}{Q_{\nu_2}p_{n-\nu_1-\nu_2-\lambda}}
 =\pair{q_{n-\mu-\lambda}P^\mu}{Q_{\nu_1}Q_{\nu_2}p_{n-\nu_1-\nu_2-\lambda}},
\]
و بنابراین با ادامهٔ روش، عبارتی از صورت (40.) حاصل می‌شود. اگر $P$ و$Q$ ردیف‌های فرم‌های اصلی نباشند، (9.) و(11.) همراه با آنچه همین اکنون نشان داده شد ثابت می‌کنند که آن‌ها نیز از توالی‌ای از پیوندهای هر دو ردیف نماد پدید آمده‌اند. از این نتیجه می‌شود:

\emph{قضیهٔ \srcLTR{V}:} فرایند تولید همهٔ ساخته‌های ناوردا، یعنی فرایند تاخوردگی، عبارت است از پیوند پیاپی هر دو ردیف نماد در یک حاصل‌ضرب ماتریسی، به یاری ردیف‌های متغیر.

از دو ردیف نماد $S^\sigma,T^\tau$، که در آن $\sigma\geqq\tau$، اکنون می‌توان حاصل‌ضرب‌های ماتریسی زیر را ساخت:
\begin{align}
&\pair{q_{n-\sigma-\lambda}S^\sigma}{T_{n-\tau}p_{\tau-\lambda}}, &&(\lambda=1,2,\ldots,\tau,n-\sigma),\tag{41}\\
&\pair{q_{n-\tau-\mu}T^\tau}{S_{n-\sigma}p_{\sigma-\mu}}, &&(\mu=\sigma-\tau+\lambda),\tag{42}\\
&\pair{q_{n-\tau-\nu}T^\tau}{S_{n-\sigma}p_{\sigma-\nu}}, &&(\nu=1,2,\ldots,\sigma-\tau).
\tag{43}
\end{align}
اما از (31.) برای (42.) و(43.) مقدارهای زیر به دست می‌آیند:
\begin{align}
&\pair{q_{n-\sigma-\lambda}S^\sigma}{T_{n-\tau}p_{\tau-\lambda}}
 +\sum_{\alpha}\Pi\pair{q_{n-\sigma-\lambda-\alpha}S^\sigma}{T_{n-\tau}p_{\tau-\lambda-\alpha}},\tag{42a}\\
&\Pi(q_{n-\sigma}S^\sigma)(T_{n-\tau}p_\tau)
 +\sum_{\beta}\Pi\pair{q_{n-\sigma-\beta}S^\sigma}{T_{n-\tau}p_{\tau-\beta}}.
\tag{43a}
\end{align}
پس می‌توان فرم‌های (42.) را به‌طور خطی بر حسب (41.) و قطبی‌های (41.)، و فرم‌های (43.) را بر حسب قطبی‌های فرم بنیادی و فرم‌های (41.) بیان کرد. همچنین از (31.) وابستگی خطی فرم‌های (41.) به (42.) و قطبی‌های (42.) به دست می‌آید. بنابراین، مطابق تعریف کلبش\footnote{\srcLTR{Clebsch,} همان، \S\ 7.}، فرم‌های (42.) با فرم‌های (41.) هم‌ارزند، در حالی که (43.) به فرم بنیادی فروکاسته می‌شود. پس فرایند تولید را می‌توان به‌طور هم‌ارز با (41.) یا (42.) به دست داد؛ فرایند (41.) را، که از حیث عدد $\lambda$ ساده‌تر است، فرایند تاخوردگی تعریف می‌کنیم. عدد $\lambda$، یعنی نقص تاخوردگی (نک. \S\ 2)، شمار ردیف‌هایی را مشخص می‌کند که برای کامل شدن حاصل‌ضرب دترمینان‌ها کم‌اند؛ آن هم در حاصل‌ضرب ماتریسی‌ای که برای یک تاخوردگی معین بیشترین شمار ردیف‌ها را دارد. چون $\lambda$ فقط تا کوچک‌ترِ دو مقدار $\tau,n-\sigma$ پیش می‌رود، و $\sigma\geqq\tau$، داریم:
\begin{equation}
 \lambda\leqq\frac n2,\ n\text{ زوج};\qquad
 \lambda\leqq\frac{n-1}{2},\ n\text{ فرد}.
\tag{44}
\end{equation}
فروکاستن (42.) و(43.) به (41.) حتی هنگامی نیز معتبر می‌ماند که ردیف‌های $p,q$ بر اثر تاخوردگیِ بیشتر به ردیف‌های نماد تبدیل شده باشند؛ زیرا مطابق (36.) به دست می‌آید:
\[
 \pair{A^{n-\tau-\varkappa}T^\tau}{S_{n-\sigma}B_{\sigma-\varkappa}}
 =\sum_{\alpha=0,\varkappa-\sigma+\tau}^{\tau,n-\sigma}
 \eps\pair{R^{\tau-\alpha}B^{n-\sigma+\varkappa}}{A_{\tau+\varkappa}R_{n-\sigma-\alpha}},
\]
که در آن ردیف‌های $R$ مطابق (33.) از $S$ و$T$ پدید آمده‌اند. بنابراین این فرم‌ها را که فقط شامل ردیف‌های نمادند، از طریق تاخوردن $\pair{q_\tau A^{n-\tau-\varkappa}}{B_{\sigma-\varkappa}p_{n-\sigma}}$ یا، به‌ترتیب، $\pair{q_{\tau-\alpha}B^{n-\sigma+\varkappa}}{A_{\tau+\varkappa}p_{n-\sigma-\alpha}}$ ــ بسته به اینکه $\sigma-\tau\gtreqless 2\varkappa$ ــ با فرم بنیادی و فرم‌های (41.) به دست می‌آوریم.

از فرایند نمادیِ تاخوردگی، به روش شناخته‌شده، فرایندهای مشتق‌گیریِ ناوردا به دست می‌آیند؛ کافی است ردیف‌های نماد $S^\sigma,T_{n-\tau}$ را با ردیف‌های نمادهای مشتق $\frac{\partial}{\partial p_\sigma}$، $\frac{\partial}{\partial q_{n-\tau}}$ جایگزین کنیم. چنان‌که معلوم است، هر حاصل‌ضرب $\frac{\partial}{\partial p_{i_1\ldots i_\sigma}}\cdot\frac{\partial}{\partial q_{j_1\ldots j_{n-\tau}}}$ معنای واقعیِ $\frac{\partial^2}{\partial p_{i_1\ldots i_\sigma}\partial q_{j_1\ldots j_{n-\tau}}}$ را دارد. چون ردیف‌های $\frac{\partial}{\partial p_\sigma}$، $\frac{\partial}{\partial q_{n-\tau}}$ با ردیف‌های $S^\sigma,T_{n-\tau}$ هموردا هستند، از همان همانی‌های آن‌ها پیروی می‌کنند؛ به‌ویژه از همانی‌های میان (42.) و(42a.)، و میان (43.) و(43a.). در جمع‌بندی به دست می‌آوریم:

\emph{قضیهٔ \srcLTR{VI}:} همهٔ ساخته‌های ناوردای فرم‌های بنیادی دلخواه، با کاربرد مکرر فرایند نمادیِ تاخوردگیِ (41.)، یا نیز فرایندهای مشتق‌گیریِ ناوردا، از همین فرم‌ها پدید می‌آیند:
\begin{equation}
 \left(q_{n-\sigma-\lambda}\frac{\partial}{\partial p_\sigma}\middle|\frac{\partial}{\partial q_{n-\tau}}p_{\tau-\lambda}\right),
 \qquad(\sigma\geqq\tau,\ \lambda=1,2,\ldots,\tau,n-\sigma).
\tag{45}
\end{equation}
همچنین باید یادآور شد که در هر تاخوردگیِ (41.) وزن کل فرم، یعنی مجموع وزن‌های ردیف‌های منفردِ متغیر (نک. \S\ 1)، به اندازهٔ عدد مثبت $2(\lambda^2+\lambda(\sigma-\tau))$ کاهش می‌یابد. از این، مستقل از اثبات عمومیِ متناهی‌بودن، متناهی‌بودن رشتهٔ فرم‌ها، یعنی کل ساخته‌های ناوردایی که نسبت به ضرایب یک فرم بنیادی معین خطی‌اند، نتیجه می‌شود.\footnote{نک. \srcLTR{Clebsch,} همان، \S\ 11 و12: متناهی‌بودن دستگاه همورداهای بنیادی.}

همچنین یادآور می‌شویم که قضیهٔ \srcLTR{VI} برای فرم‌هایی که شرایط (12.) را برآورده نمی‌کنند نیز همچنان برقرار است. زیرا می‌توان هر عاملِ (3.)، یعنی $\pair{S^\rho}{p_\rho}^{m_\rho}$، را با حاصل‌ضربِ $m_\rho$ عامل خطی $\pair{A^\rho}{p_\rho}\pair{B^\rho}{p_\rho}\cdots\pair{M^\rho}{p_\rho}$ جایگزین کرد؛ بدین‌ترتیب ساخته‌های ناوردا به ساخته‌های دستگاهی از فرم‌های خطی تبدیل می‌شوند که فقط بعداً لازم است با یکدیگر برابر قرار داده شوند. اما برای هر فرم خطیِ منفرد، شرایط (12.)، که با وارد کردن نمادهای دیفرانسیلی فقط مشتق‌های مرتبهٔ دوم را دربر می‌گیرند، همواره برقرارند.

\section*{\S\ 7. بسط فرم‌های عمومی برحسب فرم‌های نرمال.}
در ادامه (\S\ 8) یکی از ویژگی‌های اصلیِ فرایند تاخوردگیِ (41.)، یعنی ترکیب‌پذیریِ آن از تاخوردگی‌های بنیادی، را استخراج خواهیم کرد؛ اما برای این کار باید قضیهٔ مهمی را که آقای مرتنس\footnote{\srcLTR{F. Mertens, \emph{Über invariante Gebilde quaternärer Formen} (نگاه کنید به مقدمه)، که از آن با نام ``Mertens`` یاد می‌شود.}} در حوزهٔ چهارتایی به دست داده است به حوزهٔ $n$-تایی منتقل کنیم. این قضیه می‌گوید هر دستگاه متناهیِ دلخواه از فرم‌ها را می‌توان با دستگاهی هم‌ارز از فرم‌های نرمال جایگزین کرد. در اینجا، با تعمیم تعریف دوتایی و سه‌تایی،\footnote{\srcLTR{Study, a. a. O. S. 54.}} از «فرم نرمال» فرمی را مراد می‌کنیم که همهٔ ساخته‌های ناوردای آن که نسبت به ضرایب خطی‌اند، به‌استثنای خودِ فرم بنیادی، همانی‌وار صفر شوند. پس بنا بر قضیهٔ ششم (\S\ 6)، یک فرم نرمالِ $\varphi$ دستگاه معادلات دیفرانسیل زیر را ارضا می‌کند:
\begin{equation}
 \left(q_{n-\sigma-\lambda}\frac{\partial}{\partial p_\sigma}\middle|\frac{\partial}{\partial q_{n-\tau}}p_{\tau-\lambda}\right)\varphi=0,
 \qquad(\sigma\geqq\tau,\ \lambda=1,2,\ldots,\tau,n-\sigma),
\tag{46}
\end{equation}
که در آن باید ردیف‌های $q_{n-\sigma-\lambda},p_{\tau-\lambda}$ را کمیت‌های دلخواه در نظر گرفت. این معادلات دیفرانسیل در حوزهٔ چهارتایی ــ با درنظرگرفتنِ (39.) برای $\sigma=\tau=2$ ــ دقیقاً با ده معادلهٔ دیفرانسیلی که آقای مرتنس بدون معرفیِ ردیف‌های دلخواهِ $p,q$ داده است (همان‌جا، فرمول 28)) منطبق‌اند؛ شناخت آن‌ها همراه با قضیهٔ ششم به ما اجازه می‌دهد روش اثباتِ آنجا را تقریباً لفظ‌به‌لفظ به $n$ متغیر منتقل کنیم. تنها چیزی که باید افزوده شود، اثباتِ این است که فرم‌های ساخته‌شده فرم‌های نرمال‌اند؛ این اثبات با همانیِ تجزیهٔ (30.) به دست می‌آید، حال آنکه در حوزهٔ چهارتایی بی‌هیچ تبدیل دیگری حاصل می‌شود. ازاین‌رو بیان کوتاهی از اثبات نیز کفایت می‌کند؛ و برای سادگی، آن را به تنها حالتی محدود می‌کنیم که در ادامه مطرح است، یعنی حالتِ یک فرم بنیادیِ $f$ و رشتهٔ فرم‌های آن (مقایسه کنید با پایان \S\ 6)، زیرا برای ساخته‌های ناوردای با درجهٔ دلخواه نسبت به ضرایبِ هر تعداد فرم بنیادی، اثبات بدون تغییر می‌ماند.

اندیشهٔ بنیادیِ اثبات این است که با معرفیِ ضرایب تبدیل‌یافته، یک فرم با $n-1$ ردیفِ $p_\rho$ را با فرم‌هایی دارای $n$ ردیفِ هموردا با $x$، یعنی ردیف‌های $\xi$ از کمیت‌های تبدیل، جایگزین کنیم. فرم‌های نرمالِ مطلوب مستقیماً با فرایندهای مشتق‌گیریِ ناوردا از این فرم‌ها استخراج می‌شوند. بسط فرم‌های عمومی برحسب فرم‌های نرمال با بسطی سری برای فرم‌هایی دارای $\rho$ ($\rho\leqq n$) ردیفِ $\xi$ هموردا انجام می‌شود؛ فرایندهای مشتق‌گیریِ ناوردا دوباره به فرم‌های دارای ردیف‌های اصلیِ $p_\rho$ بازمی‌گردانند.

فرض کنیم $\xi^{(1)},\ldots,\xi^{(n)}$ ردیف‌های کمیت‌های تبدیل باشند، چنان‌که:
\begin{equation}
\begin{aligned}
x_i&=\xi_i^{(1)}x'_1+\xi_i^{(2)}x'_2+\cdots+\xi_i^{(n)}x'_n,\\
p_{i_1\ldots i_\rho}&=\sum_j(\xi_{i_1}^{(j_1)}\cdots \xi_{i_\rho}^{(j_\rho)})p'_{j_1\ldots j_\rho}.
\end{aligned}
\tag{47}
\end{equation}
ضرایب تبدیل‌یافتهٔ فرم‌های $f,\varphi,\ldots$ را با $(f),(\varphi),\ldots$ نشان می‌دهیم و فرض می‌کنیم
\[
 f\equiv\pair{S^1}{p_1}^{m_1}\pair{S^2}{p_2}^{m_2}\cdots\pair{S^{n-1}}{p_{n-1}}^{m_{n-1}}.
\]
آنگاه:
\begin{equation}
\begin{aligned}
(f)&=\pair{S^1}{\xi^{(1)}}^{m_{11}}\cdots\pair{S^1}{\xi^{(n)}}^{m_{1n}}
\pair{S^2}{\xi^{(1)}\xi^{(2)}}^{m_{21}}\cdots
\pair{S^{n-1}}{\xi^{(2)}\cdots\xi^{(n)}}^{m_{n-1,n}}\\
&=\pair{s^{(1)}}{\xi^{(1)}}^{k_1}\pair{s^{(2)}}{\xi^{(2)}}^{k_2}\cdots\pair{s^{(n)}}{\xi^{(n)}}^{k_n},
\end{aligned}
\tag{48}
\end{equation}
که در آن $m_{11}+\cdots+m_{1n}=m_1$ و به همین ترتیب. از هر ضریبِ $(f)$ که برای آن
\begin{equation}
 k_1\geqq k_2\geqq k_3\geqq\cdots\geqq k_n,
\tag{49}
\end{equation}
برقرار باشد، فرایند مشتق‌گیریِ زیر که ردیف‌های $\xi$ را دقیقاً به‌تمامی مصرف می‌کند:
\begin{equation}
\left(\frac{\partial}{\partial\xi^{(1)}}\middle|p_1\right)^{k_1-k_2}
\left(\frac{\partial}{\partial\xi^{(1)}}\frac{\partial}{\partial\xi^{(2)}}\middle|p_2\right)^{k_2-k_3}
\cdots
\left(\frac{\partial}{\partial\xi^{(1)}}\cdots\frac{\partial}{\partial\xi^{(n-1)}}\middle|p_{n-1}\right)^{k_{n-1}-k_n}
\left(\frac{\partial}{\partial\xi^{(1)}}\cdots\frac{\partial}{\partial\xi^{(n)}}\right)^{k_n}
\tag{50}
\end{equation}
فرمِ
\begin{equation}
 \varphi=\bigl(s^{(1)}\mid p_1\bigr)^{k_1-k_2}
 \bigl(s^{(1)}s^{(2)}\mid p_2\bigr)^{k_2-k_3}\cdots
 \bigl(s^{(1)}\cdots s^{(n-1)}\mid p_{n-1}\bigr)^{k_{n-1}-k_n}
 \bigl(s^{(1)}s^{(2)}\cdots s^{(n)}\bigr)^{k_n}.
\tag{51}
\end{equation}
به دست می‌آید. فرم‌های $\varphi$ همان فرم‌های نرمالِ مطلوب‌اند؛ زیرا دو ویژگیِ تعریف‌کنندهٔ فرم‌های نرمال را دارند:

1. آن‌ها ساخته‌های ناوردای فرم بنیادیِ $f$ هستند که نسبت به ضرایب خطی‌اند. این امر از آنجا نتیجه می‌شود که آن‌ها از $f$، با فرایندهای مشتق‌گیریِ ناوردای $\left(\frac{\partial}{\partial p_\rho}\middle|\xi^{(i_1)}\xi^{(i_2)}\cdots\xi^{(i_\rho)}\right)$ و (50.)، پدید می‌آیند. بنابراین (بنا بر پایان \S\ 6) می‌توان آن‌ها را به‌صورت خطی برحسب قطبی‌های رشتهٔ متناهیِ فرم‌های $f$ بیان کرد. این قطبی‌ها را چنین به دست می‌آوریم که متغیرهای واردشده در $\varphi$، یعنی $p_\rho$، را ضرایب فرمی با ردیف‌های متغیرِ $\mathfrak{p}_\sigma$ در نظر بگیریم که به فرم‌های خطی تجزیه می‌شود:
\begin{equation}
 H=(q_{n-1}\mid \mathfrak{p}_{n-1})^{k_1-k_2}(q_{n-2}\mid \mathfrak{p}_{n-2})^{k_2-k_3}\cdots(q_1\mid \mathfrak{p}_1)^{k_{n-1}-k_n}.
\tag{52}
\end{equation}
ناورداهای هم‌زمانِ رشته‌های فرمِ $f$ و $H$ همهٔ قطبی‌های موردنظر را به دست می‌دهند، زیرا درجهٔ این قطبی‌ها نسبت به متغیرهای هر ردیف منفردِ $p_\rho$ باید با درجهٔ خودِ $\varphi$ نسبت به همان متغیرها برابر باشد.

2. آن‌ها معادلات دیفرانسیلِ (46.) را ارضا می‌کنند. در واقع، اگر فرایند مشتق‌گیریِ (45.) را بر $\varphi$ اعمال کنیم، به دست می‌آید:
\[
 \varphi'=\bigl(q_{n-\sigma-\lambda}s^{(1)}\cdots s^{(\sigma)}\mid [s^{(1)}\cdots s^{(\tau)}]_{n-\tau}p_{\tau-\lambda}\bigr),\qquad (\sigma\geqq\tau).
\]
از همانیِ تجزیهٔ (30.) نتیجه می‌شود، اگر ردیف‌های ساخته‌شده از ردیف‌های $s$، یعنی $q_\sigma=s^{(1)}\cdots s^{(\sigma)}$ و $p_{n-\tau}=[s^{(1)}\cdots s^{(\tau)}]_{n-\tau}$، را با $q_\rho,p_\tau$ آنجا یکی بگیریم، که:
\[
 \varphi'=\sum_{\beta=\sigma-\tau+\lambda}^{n-\tau,\sigma}\sum_{i_\beta}\eps\,
 \pair{\dot{q}_{\sigma-\beta}^{(i_1\ldots i_\beta)}q_{n-\tau+\lambda}}{p_{\sigma+\lambda}\dot{p}_{n-\tau-\beta}^{(i_1\ldots i_\beta)}}.
\]
در اینجا:
\[
 \dot{p}_{n-\tau-\beta}^{(i_1\ldots i_\beta)}\sim s^{(1)}\cdots s^{(\tau)}s^{(i_1)}\cdots s^{(i_\beta)},
\]
که در آن $i_1,\ldots,i_\beta$ هر $\beta$ عددِ متمایز از یکدیگرند که از میان اعداد 1 تا $\sigma$ برگزیده شده‌اند. چون $\beta>\sigma-\tau$، ناگزیر در میان اعداد $i_1,\ldots,i_\beta$ عددی بین 1 تا $\tau$ وجود دارد؛ در نتیجه $\dot{p}_{n-\tau-\beta}^{(i_1\ldots i_\beta)}$ همانی‌وار صفر می‌شود، پس تک‌تک حاصل‌ضرب‌های ماتریسی و همراه با آن‌ها $\varphi'$ نیز صفر می‌شوند.

برای تکمیل هم‌ارزیِ دستگاه فرم‌های نرمال با رشتهٔ فرم‌های $f$، اکنون فقط باید نشان دهیم که همهٔ فرم‌های این رشته را می‌توان، به معنای تعیین‌شده در 1، برحسب قطبی‌های فرم‌های نرمال بسط داد. فرض کنیم $\Pi$ ترکیبی از عملیات قطبی باشد
\srcnumdisplay[4.2em]{(53.)}{%
 \left(\frac{\partial}{\partial\xi^{(i)}}\middle|\xi^{(k)}\right),
 \qquad (i\ne k;\ i=1,2,\ldots,\rho;\ k=1,2,\ldots,\rho)}
و $\Pi^\sigma$ ترکیبی از عملیات ویژهٔ (53.) باشد که در آن‌ها $i=\sigma$ و$k<\sigma$. اگر $F$ فرمی در $\rho$ ردیفِ $\xi^{(1)},\ldots,\xi^{(\rho)}$ باشد که در ردیف $\xi^{(\rho)}$ درجهٔ $k_\rho$ دارد، آنگاه بسط زیر برقرار است\srcfn{*)}{\srcLTR{F. Mertens, \emph{Über eine Formel der Determinantentheorie}. Sitzungsber. d. Ak. d. Wiss., Wien, math.-naturw. Kl., Bd. 91. 2. Abt. (1885). Formel 20).}}:
\srcnumdisplay[4.2em]{(54.)}{%
 F=\sum_{\alpha=0}^{k_\rho}\Pi F_\alpha,
 \qquad(\rho\leqq n)}
که در آن $F_\alpha$ در ردیف آخرِ $\xi^{(\rho)}$ درجهٔ $\alpha$ دارد و از $F$ با عملیات مشتق‌گیریِ ناوردای
\srcnumdisplay[4.2em]{(55.)}{%
 \left(\frac{\partial}{\partial\xi^{(1)}}\cdots\frac{\partial}{\partial\xi^{(\rho)}}\middle|\xi^{(1)}\cdots\xi^{(\rho)}\right)^\alpha}
و $\Pi^\rho$ پدید می‌آید. اگر هنگام ساختن $F_\alpha$، فرایند (55.) را با
\srcnumdisplay[4.2em]{(56.)}{%
 \left(\frac{\partial}{\partial\xi^{(1)}}\cdots\frac{\partial}{\partial\xi^{(\rho)}}\middle|p_\rho\right)^\alpha,}
جایگزین کنیم، فرمی چون $F_\alpha(p_\rho)$ با فقط $\rho-1$ ردیفِ $\xi^{(1)},\ldots,\xi^{(\rho-1)}$ پدید می‌آید که (54.) را دوباره می‌توان بر آن اعمال کرد. ادامهٔ این روش به فرم‌های $G(p_\rho,p_{\rho-1},\ldots,p_1)$ می‌انجامد. برای به‌دست‌آوردن بسط $F$ مطابق (54.) از این فرم‌ها، باید هر ردیف $p_\sigma$ را با $\xi^{(1)},\ldots,\xi^{(\sigma)}$ جایگزین کرد و فرایند قطبیِ $\Pi$ در (53.) را بر فرم‌های حاصل اعمال کرد؛ این کار (مقایسه کنید با (48.)) جمعی از ضرایب تبدیل‌یافتهٔ $(G)$ به دست می‌دهد. برای $\rho=n$، در گام نخست فرایند (55.) حفظ می‌شود. اگر $c$ ثابت‌های عددی باشند و برای اختصار نیز بگذاریم
\srcnumdisplay[4.2em]{(57.)}{%
 (\xi^{(1)}\xi^{(2)}\cdots\xi^{(n)})=\Delta;
 \qquad
 \left(\frac{\partial}{\partial\xi^{(1)}}\frac{\partial}{\partial\xi^{(2)}}\cdots\frac{\partial}{\partial\xi^{(n)}}\right)=\nabla,}
در حالت $\rho=n$ به دست می‌آید:
\srcnumdisplay[4.2em]{(58.)}{%
 F=\sum c\Delta^\alpha(G),}
درحالی‌که برای $\rho<n$ فقط $\alpha=0$، یعنی $\sum c(G)$، در طرف راست ظاهر می‌شود.

اگر (58.) را بر یک ضریب تبدیل‌یافتهٔ $(f)$ در (48.) اعمال کنیم، از جابجاپذیریِ عملیات قطبیِ $\Pi^\sigma$ با همهٔ عملیات (56.) که در آن‌ها $\rho\geqq\sigma$ است، و از اینکه قطبی‌های ضرایب تبدیل‌یافته دوباره ضرایب تبدیل‌یافته می‌دهند، نتیجه می‌شود:
\[
\begin{aligned}
 G={}&\left(\frac{\partial}{\partial\xi^{(1)}}\middle|p_1\right)^{\beta_1}
 \left(\frac{\partial}{\partial\xi^{(1)}}\frac{\partial}{\partial\xi^{(2)}}\middle|p_2\right)^{\beta_2}
 \cdots
 \left(\frac{\partial}{\partial\xi^{(1)}}\cdots\frac{\partial}{\partial\xi^{(n-1)}}\middle|p_{n-1}\right)^{\beta_{n-1}}
 \nabla^{\beta_n}\Pi(f)\\
 ={}&\sum c\varphi
\end{aligned}
\]
و از (58.) نتیجه می‌شود:
\srcnumdisplay[4.2em]{(59.)}{%
 (f)=\sum c\Delta^\alpha(\varphi)\qquad\text{(مرتنس، فرمول 23)).}}
برای گذر از (59.) به بسط خودِ $f$، دوباره متغیرهای $p_\rho$ را ضرایب فرم $H$ در (52.) با اعداد درجهٔ متناظر با $f$، یعنی $m_1,\ldots,m_{n-1}$، در نظر می‌گیریم، و فرض می‌کنیم $m=m_1+\cdots+m_{n-1}$. آنگاه، بنا بر تعریف ناورداها:
\[
 f\cdot\Delta^m=\sum(f)(H)=\sum c\Delta^\alpha(\varphi)(H),
\]
و اعمال فرایندی که ردیف‌های $\xi$ را دقیقاً به‌تمامی مصرف می‌کند، یعنی $\nabla^m$، به دست می‌دهد:
\srcnumdisplay[4.2em]{(60.)}{%
 f=\sum c\Phi,}
که در آن فرم‌های $\Phi$ از $\varphi$ و$H$ با فرایندهای مشتق‌گیریِ ناوردا نسبت به متغیرها مشتق شده‌اند و بنابراین ناورداهای هم‌زمانِ رشته‌های فرمِ $\varphi$ و$H$ هستند\srcfn{*)}{مرتنس در \S\ 7، به‌جای استنتاج مستقیم، رشته‌ای از فرایندهای مشتق‌گیریِ بیشتر را معرفی می‌کند که اساساً برای ساختن رشتهٔ فرم‌های $H$ به کار می‌روند؛ کاری که قضیهٔ ششم ما انجام می‌دهد.}. اما چون $\varphi$ فرم نرمال است، رشتهٔ فرم‌های $\varphi$ به خودِ فرم $\varphi$ محدود می‌شود؛ رشتهٔ فرم‌های $H$ همهٔ ترکیب‌های ناوردای ممکنِ متغیرها را دربر دارد؛ پس ناورداهای هم‌زمانِ $\varphi$ و$H$ همان قطبی‌های $\varphi$ به معنای گفته‌شده در 1 هستند؛ یعنی (60.) بسط $f$ برحسب قطبی‌های فرم‌های نرمال را نمایش می‌دهد. اکنون باید همین بسط را برای یک فرم دلخواهِ $\theta$ از رشتهٔ فرم‌های $f$ نیز ثابت کنیم. فرض کنیم $\Gamma$ فرم‌های نرمالی باشند که از $(\theta)$ به‌وسیلهٔ (50.) مشتق شده‌اند، به‌طوری‌که:
\srcnumdisplay[4.2em]{(61.)}{%
 (\theta)=\sum c\Delta^\beta(\Gamma).}
فرم‌های $\theta$ با رابطهٔ زیر مشخص می‌شوند که برای هر ضریب تبدیل‌یافته برقرار است:
\[
 (\theta)\Delta^\sigma=\sum c(f)\qquad\text{(مرتنس، فرمول 9)).}
\]
اما چون هر فرم $F$ در $n$ ردیفِ $\xi$ که عامل $\Delta^\sigma$ را دربردارد، عامل دوم، یعنی $\nabla^\sigma\Pi F$، را نیز دربردارد (مرتنس، فرمول 20))، نتیجه می‌شود:
\[
 (\theta)=\sum c\nabla^\sigma(f),\qquad\text{پس}\qquad \Gamma=\sum c\varphi,
\]
و از این، بنا بر (61.)، رابطهٔ متناظر با (59.) به دست می‌آید:
\[
 (\theta)=\sum c\Delta^\beta(\varphi),
\]
که برای $\theta$، بسط (60.) برحسب قطبی‌های فرم‌های نرمالِ $\varphi$ را در پی دارد. در جمع‌بندی داریم:

\emph{قضیهٔ هفتم. هر رشتهٔ فرم‌های $f$ را می‌توان با دستگاهی هم‌ارز از فرم‌های نرمال جایگزین کرد.}

\begin{center}
\emph{\S\ 8. فروکاستن فرایند تاخوردگی به تاخوردگی‌های پایه.}
\end{center}
در پی قضیهٔ هفتم، اکنون فروکاستن فرایند تاخوردگی به تاخوردگی‌های پایه را که در آغاز \S\ 7 یاد شد، انجام می‌دهیم. در اینجا (نک. \S\ 5) هر تاخوردگیِ دو ردیف هموردای $S^\rho,T^\rho$ را «\emph{تاخوردگیِ هموردا}» و، به‌ویژه، تاخوردگیِ هموردای دارای نقص 1 دو ردیف $S^\rho,T^\rho$ را «\emph{تاخوردگیِ پایه از نوع $\rho$}» می‌نامیم. سپس، در تکمیل قضیهٔ ششم، قضیهٔ زیر را ثابت می‌کنیم:

\emph{قضیهٔ هشتم: فرایند عمومیِ تاخوردگی (41.) را می‌توان از $(n-1)$ تاخوردگیِ پایه از نوع $\rho$ ترکیب کرد، که در آن $\rho=1,2,\ldots,n-1$. به بیان دیگر: برای تولید همهٔ ساخته‌های ناوردا، کاربرد پیاپیِ $n-1$ تاخوردگیِ پایه کافی است.}

در حوزهٔ دوتایی، فرایند تاخوردگیِ (41.) مستقیماً با تنها تاخوردگیِ پایهٔ ممکن یکی است؛ در حوزهٔ سه‌تایی، قضیهٔ هشتم را در \S\ 1 رسالهٔ دکتری خود (نک. مقدمه) ثابت کرده‌ام. در آنجا، به‌سبب (44.)، تاخوردگی‌های هموردای عمومی هنوز با تاخوردگی‌های پایه یکسان‌اند؛ و شایان توجه است که قضیهٔ هشتم برای $n$ متغیر، فروکاست به تاخوردگی‌های پایه را به دست می‌دهد، نه فقط فروکاستِ بسیار ضعیف‌تر به تاخوردگی‌های هموردا را. برهان در اساس از حالت سه‌تایی الگوبرداری شده است: کلی‌ترین تاخوردگی‌ها با به‌کارگیری همانی‌های نمادی از تاخوردگی‌های پایه ساخته می‌شوند. خواهیم ساخت:
\begin{enumerate}
\renewcommand{\labelenumi}{\arabic{enumi})}
\item تاخوردگی‌های دلخواهِ دارای نقص 1 به‌وسیلهٔ تاخوردگی‌های هموردای دارای نقص 1، یعنی به‌وسیلهٔ تاخوردگی‌های پایه (نظیر حالت سه‌تایی)،
\item تاخوردگی‌های دلخواهِ دارای نقص $\lambda$ به‌وسیلهٔ تاخوردگی‌های هموردای دارای نقص $\lambda$ و به‌وسیلهٔ تاخوردگی‌های دلخواهِ دارای نقص 1،
\item تاخوردگی‌های هموردای دارای نقص $\lambda$ به‌وسیلهٔ تاخوردگی‌های هموردای دارای نقص $\lambda-1$ و به‌وسیلهٔ تاخوردگی‌های دلخواهِ دارای نقص 1.
\end{enumerate}
در اینجا، و این نکته اساسی است، بنابر قضیهٔ هفتم فرض می‌کنیم دو فرم $S$ و$T$ که باید بر آن‌ها تاخوردگی انجام شود، فرم‌های نرمال‌اند. بنابراین، بر پایهٔ (46.)، روابط زیر برقرارند:
\srcnumdisplay[4.2em]{\raisebox{-1.7ex}{(62.)}}{%
\begin{aligned}
\pair{q_{n-\sigma_1-\lambda}S^{\sigma_1}}{S_{n-\sigma_2}p_{\sigma_2-\lambda}}&=0,
&& (\sigma_1\geqq\sigma_2,\ \lambda=1,2,\ldots,\sigma_2,n-\sigma_1),\\
\pair{q_{n-\tau_1-\lambda}T^{\tau_1}}{T_{n-\tau_2}p_{\tau_2-\lambda}}&=0,
&& (\tau_1\geqq\tau_2,\ \lambda=1,2,\ldots,\tau_2,n-\tau_1).
\end{aligned}}
همچنین، بنا بر پایان \S\ 2، کافی است فرم‌های $S$ و$T$ را در همهٔ ردیف‌های متغیر خطی فرض کنیم؛ یعنی فرم‌های ساخته‌شده به رشتهٔ فرم‌های زیر تعلق دارند:
\[
 f=\pair{S^1}{p_1}\pair{S^2}{p_2}\cdots\pair{S^{n-1}}{p_{n-1}}
   \pair{T^1}{p_1}\pair{T^2}{p_2}\cdots\pair{T^{n-1}}{p_{n-1}}.
\]
\emph{1) تولید $\pair{q_{n-\sigma-1}S^\sigma}{T_{n-\tau}p_{\tau-1}}$، که در آن $\sigma>\tau$، از تاخوردگی‌های پایه از نوع $\tau+\alpha$؛ $\alpha=0,1,\ldots,\sigma-\tau$.} از تاخوردگیِ پایه از نوع $\tau$
\[
 \pair{q_{n-\tau-1}S^\tau}{T_{n-\tau}p_{\tau-1}}
\]
با انجام جایگذاریِ زیر به دست می‌آوریم
\[
 w\sim T_{n-\tau}p_{\tau-1}
\]
فرمی از رشتهٔ فرم‌های $f$ با سه ردیف که اندیس وزنیِ بالای هر یک (نک. \S\ 1) برابر $\tau+1$ است:
\srcnumdisplay[4.2em]{(63.)}{%
 \pair{wS^\tau}{p_{\tau+1}}\pair{S^{\tau+1}}{p_{\tau+1}}\pair{T^{\tau+1}}{p_{\tau+1}}.}
به‌کاربردن تاخوردگیِ پایه‌ای از نوع $\tau+1$ بر (63.) حاصل زیر را می‌دهد:
\[
 \pair{q_{n-\tau-2}wS^\tau}{S_{n-\tau-1}p_\tau};
\]
از این، بر پایهٔ (31.) و با جایگذاریِ $q_{n-\tau-2}w=q'_{n-\tau-1}$، مقدار زیر به دست می‌آید:
\[
 \eps\bigl(q_{n-\tau-2}wS^{\tau+1}\bigr)\pair{S^\tau}{p_\tau}
 +\sum_{\alpha=1}^{\tau,n-\tau-1}\Pi\pair{q'_{n-\tau-1-\alpha}S^{\tau+1}}{S_{n-\tau}p_{\tau-\alpha}}.
\]
عبارتِ زیر علامت جمع بنا بر (62.) همانی‌وار صفر می‌شود؛ پس به فرم زیر رسیده‌ایم:
\[
 \pair{wS^{\tau+1}}{p_{\tau+2}}\pair{S^{\tau+2}}{p_{\tau+2}}\pair{T^{\tau+2}}{p_{\tau+2}},
\]
که با تبدیلِ $\tau$ به $\tau+1$ از (63.) به دست می‌آید. بنابراین با $(\sigma-\tau)$ بار تکرار فرایند، به تاخوردگیِ مطلوب می‌رسیم:
\[
 \pair{wS^\sigma}{p_{\sigma+1}}=\eps\pair{q_{n-\sigma-1}S^\sigma}{T_{n-\tau}p_{\tau-1}}.
\]
فرایند انجام‌شده را به اختصار با فرمول زیر نشان می‌دهیم:
\srcnumdisplay[4.2em]{(64.)}{%
 T^\tau S^\tau,\quad S^{\tau+1},\quad S^{\tau+2},\ldots,S^\sigma;}
می‌بینیم که تنها خاصیت فرم نرمالِ $S$ به کار می‌رود، نه خاصیت $T$. فرایندی دوگانِ (64.) را می‌توان، تنها با استفاده از خاصیت فرم نرمالِ $T$، با ترتیب معکوس اجرا کرد؛ یعنی فرایند زیر:
\srcnumdisplay[4.2em]{(65.)}{%
 S^\sigma T^\sigma,\quad T^{\sigma-1},\quad T^{\sigma-2},\ldots,T^\tau}
بر پایهٔ روابط زیر:
\begin{align*}
\pair{q_{n-\sigma-1}S^\sigma}{T_{n-\sigma}p_{\sigma-1}}&=\eps\bigl(T_{n-\sigma}y p_{\sigma-1}\bigr),\quad \text{که در آن }y\sim q_{n-\sigma-1}S^\sigma,\\
\pair{q_{n-\sigma}T^{\sigma-1}}{T_{n-\sigma}yp_{\sigma-2}}&=\eps\bigl(T_{n-\sigma+1}y p_{\sigma-2}\bigr),\\
&\vdots\\
\pair{q_{n-\tau-1}T^\tau}{T_{n-\tau-1}yp_{\tau-1}}&=\eps\pair{q_{n-\sigma-1}S^\sigma}{T_{n-\tau}p_{\tau-1}}.
\end{align*}
همچنین شایسته است به ارتباط با کاهش وزن کلِ یادشده در پایان \S\ 6 اشاره کنیم؛ مقدار آن برای تاخوردگی‌های دارای نقص 1 برابر $2(1+(\sigma-\tau))$ و برای تاخوردگی‌های پایه برابر $2\cdot1$ است. بنابراین اگر ترکیب‌پذیری از تاخوردگی‌های پایه اصولاً ممکن باشد، به‌ضرورت $\sigma-\tau+1$ تاخوردگی از این نوع لازم است، همان‌گونه که در بالا انجام شد.

\emph{2) تولید تاخوردگی‌های عمومی از تاخوردگی‌های هموردا و پایه.} کاهش وزن برای یک تاخوردگی عمومیِ (41.) برابر است با $2(\lambda^2+\lambda(\sigma-\tau))=2[\lambda^2+(\sigma-\tau)(1+(\lambda-1))]$؛ پس امکان تجزیهٔ آن را به یک تاخوردگی هموردای دارای نقص $\lambda$ (کاهش وزن $2\lambda^2$) و $\sigma-\tau$ تاخوردگیِ دارای نقص 1 و تفاضل اندیس‌های وزنیِ $\lambda-1$ (کاهش وزن $2\{1+(\lambda-1)\}$) نشان می‌دهد. این تجزیه در حالت $\lambda=1$ با تجزیهٔ بیان‌شده در 1) سازگار است؛ نشان می‌دهیم که واقعاً می‌توان آن را تحقق بخشید.

از تاخوردگی هموردای دارای نقص $\lambda$
\[
 \pair{q_{n-\tau-\lambda}S^\tau}{T_{n-\tau}p_{\tau-\lambda}}
\]
با انجام جایگذاریِ $R^\lambda\sim T_{n-\tau}p_{\tau-\lambda}$، فرمی با دو ردیف به دست می‌آوریم که اندیس‌های وزنیِ بالای آن‌ها به‌ترتیب $\tau+\lambda$ و$\tau+1$ است:
\srcnumdisplay[4.2em]{(66.)}{%
 \pair{R^\lambda S^\tau}{p_{\tau+\lambda}}\pair{S^{\tau+1}}{p_{\tau+1}}.}
ردیفِ دارای اندیس وزنیِ کوچک‌ترِ $\tau+1$ به یک فرم نرمال تعلق دارد؛ بنابراین (66.) تاخوردگیِ دارای نقص 1 و مرکب از $\lambda$ تاخوردگیِ پایه را، با ترتیب (65.)، می‌پذیرد:
\[
 (R^\lambda S^\tau)S^{\tau+\lambda},\quad S^{\tau+\lambda-1},\quad S^{\tau+\lambda-2},\ldots,S^{\tau+1}.
\]
این، با درنظرگرفتن (31.) و(62.)، حاصل زیر را می‌دهد:
\[
 \pair{q_{n-\tau-\lambda-1}R^\lambda S^\tau}{S_{n-\tau-1}p_\tau}=\eps\pair{R^\lambda S^{\tau+1}}{p_{\tau+\lambda+1}},
\]
یعنی فرمی که با تبدیلِ $\tau$ به $\tau+1$ از (66.) پدید می‌آید، همانند آنچه در 1) دربارهٔ (63.) روی داد. $(\sigma-\tau)$ بار تکرار این فرایند به تاخوردگیِ مطلوب می‌انجامد:
\[
 \pair{R^\lambda S^\sigma}{p_{\sigma+\lambda}}=\eps\pair{q_{n-\sigma-\lambda}S^\sigma}{T_{n-\tau}p_{\tau-\lambda}}.
\]
در سراسر این فرایند فقط خاصیت فرم نرمالِ $S$ به کار رفت؛ درنظرگرفتنِ $T$ به‌عنوان فرم نرمال، با وارونگی کامل ترتیب، به فرایند دوگان بر پایهٔ فرمول‌های زیر می‌انجامد:
\begin{align*}
\pair{q_{n-\sigma-\lambda}S^\sigma}{T_{n-\sigma}p_{\sigma-\lambda}}&=\eps\bigl(T_{n-\sigma}R_\lambda p_{\sigma-\lambda}\bigr),\quad \text{که در آن }R_\lambda\sim q_{n-\sigma-\lambda}S^\sigma,\\
\pair{q_{n-\sigma}T^{\sigma-1}}{T_{n-\sigma}R_\lambda p_{\sigma-\lambda-1}}&=\eps\bigl(T_{n-\sigma+1}R_\lambda p_{\sigma-\lambda-1}\bigr),\\
&\vdots\\
\pair{q_{n-\tau-1}T^\tau}{T_{n-\tau-1}R_\lambda p_{\tau-\lambda}}&=\eps\pair{q_{n-\sigma-\lambda}S^\sigma}{T_{n-\tau}p_{\tau-\lambda}}.
\end{align*}
\emph{3) تولید تاخوردگی‌های هموردای دارای نقص $\lambda$ از تاخوردگی‌های هموردای دارای نقص $\lambda-1$ و تاخوردگی‌های پایه.} کاهش وزن برای تاخوردگی‌های هموردای دارای نقص $\lambda$ برابر است با $2\lambda^2=2((\lambda-1)^2+2\lambda-1)$؛ پس امکان تجزیهٔ چنین تاخوردگی‌ای را به یک تاخوردگی هموردای دارای نقص $\lambda-1$ و $2\lambda-1$ تاخوردگیِ پایه فراهم می‌کند. این تجزیه را به ساده‌ترین صورت هنگامی می‌یابیم که ابتدا حالت خاصِ $n=2\lambda$ را در نظر بگیریم. در این صورت تنها تاخوردگی هموردای ممکنِ دارای نقص $\lambda$ چنین می‌شود: $\pair{S^{n/2}}{T_{n/2}}=\pair{S^\lambda}{T_\lambda}$؛ این تاخوردگی یک ردیف $u$ و یک ردیف $x$ کمتر از تاخوردگیِ دارای نقص $\lambda-1$ برای همان ردیف‌ها، یعنی $\pair{uS^\lambda}{T_\lambda x}$، دارد. برعکس، تاخوردگیِ دارای نقص $\lambda$ را از تاخوردگی‌های پایه‌ای ترکیب می‌کنیم که یک ردیف $u$ و یک ردیف $x$ را ناپدید می‌کنند (کاهش وزن $2(n-1)=2(2\lambda-1)$) و هم‌زمان ردیف نمادِ تازهٔ $R^\lambda$ را می‌سازند، و نیز از یک تاخوردگی هموردای دارای نقص $\lambda-1$. ساده‌ترین تاخوردگی‌ای که این کار را انجام می‌دهد تاخوردگیِ زیر با نقص 1 است:
\[
 \pair{sT^\lambda}{\sigma p_\lambda}=\eps\pair{S^{2\lambda-1}}{R_{\lambda-1}p_\lambda}=\eps\pair{R^\lambda}{p_\lambda},
\]
که در آن قرار می‌دهیم:
\[
 R^{\lambda+1}=sT^\lambda,
 \quad s=S^1,
 \quad \sigma=S_1,
 \quad R^\lambda\sim\sigma R_{\lambda-1}.
\]
این تاخوردگی، بنا بر (65.) و(64.)، متشکل از $2\lambda-1$ تاخوردگی پایهٔ زیر است:
\srcnumdisplay[4.2em]{(67.)}{%
 T^\lambda S^\lambda,S^{\lambda-1},\ldots,S^1;
 \qquad R^{\lambda+1}S^{\lambda+1},S^{\lambda+2},\ldots,S^{2\lambda-1}.}
یک تاخوردگیِ دارای نقص $\lambda-1$، با در نظر گرفتن (21.) و(62.)، حاصل زیر را می‌دهد:
\[
 \pair{uS^\lambda}{\sigma R_{\lambda-1}x}=\eps\pair{R^{\lambda+1}}{S_\lambda x}
 =\eps\pair{sT^\lambda}{S^\lambda x}=\eps\pair{S^\lambda}{T_\lambda}.
\]
اکنون می‌توان ساختِ حالت خاص را برای دو ردیف $S^\rho,T^\rho$ و$n$ دلخواه، مستقیماً به حالت کلی تعمیم داد. به جای (67.)، $2\lambda-1$ تاخوردگی پایهٔ زیر قرار می‌گیرند:
\srcnumdisplay[4.2em]{(68.)}{%
 T^\rho S^\rho,S^{\rho-1},\ldots,S^{\rho-\lambda+1};
 \qquad R^{\rho+1}S^{\rho+1},S^{\rho+2},\ldots,S^{\rho+\lambda-1}.}
نیمهٔ نخستِ تاخوردگی‌ها فرم زیر را به دست می‌دهد:
\[
 \pair{q_{n-\rho-1}T^\rho}{S_{n-\rho+\lambda-1}p_{\rho-\lambda}},
\]
که از آن، با جایگذاری‌های
\[
 S_{n-\rho+\lambda-1}p_{\rho-\lambda}\sim w,
 \qquad T^\rho w=R^{\rho+1}
\]
و با اعمال نیمهٔ دومِ تاخوردگی‌ها، فرم زیر حاصل می‌شود:
\[
 \pair{q_{n-\rho-\lambda}S^{\rho+\lambda-1}}{R_{n-\rho-1}p_\rho}.
\]
اگر جایگذاری کنیم:
\[
 q_{n-\rho-\lambda}S^{\rho+\lambda-1}\sim y,
\]
آنگاه یک تاخوردگی هموردای دارای نقص $\lambda-1$ حاصل زیر را می‌دهد:
\srcnumdisplay[4.2em]{(69.)}{%
 \pair{q_{n-\rho-\lambda+1}S^\rho}{yR_{n-\rho-1}p_{\rho-\lambda+1}}.}
نشان می‌دهیم که این همان تاخوردگیِ مطلوبِ دارای نقص $\lambda$ است. اگر جایگذاری کنیم
\[
 R_{n-\rho-1}p_{\rho-\lambda+1}=R_{n-\lambda}
\]
و توجه کنیم که با بازکردنِ $y$
\[
 \pair{q_{n-\rho-\lambda+1}R^\lambda}{S_{n-\rho}y}
 =\eps\pair{q_{n-\rho-\lambda}S^{\rho+\lambda-1}}{S_{n-\rho}p'_{\rho-1}}=0
\]
به دست می‌آید، آنگاه بنا بر (20.) مقدار (69.) چنین است:
\[
 \eps\pair{S^\rho R^\lambda}{p_{\rho+\lambda-1}y}
 =\eps\pair{q_{n-\rho-\lambda}S^{\rho+\lambda-1}}{[S^\rho R^\lambda]_{n-\rho-\lambda}p_{\rho+\lambda-1}}.
\]
اما این تاخوردگی از نوع (43.) است و بنابراین هم‌ارز است با
\[
 \pair{q_{n-\rho-\lambda}S^\rho}{R_{n-\rho-1}p_{\rho-\lambda+1}}
 =\eps\pair{wT^\rho}{R_\lambda p_{\rho-\lambda+1}},
 \qquad \text{که در آن }R_\lambda\sim q_{n-\rho-\lambda}S^\rho.
\]
در اینجا $R_\lambda$ از پیش دارای فرم پایانیِ مطلوب است؛ این عبارت در حالت خاص با فرم $\pair{sT^\lambda}{S_\lambda x}$ متناظر است. اکنون اگر توجه کنیم که بازکردنِ $w$ و$R_\lambda$ حاصل زیر را می‌دهد:
\[
 \pair{wR^{n-\lambda}}{T_{n-\rho}p_{\rho-\lambda+1}}
 =\eps\pair{q'_{\lambda-1}R^{n-\lambda}}{S_{n-\rho+\lambda-1}p_{\rho-\lambda}}
 =\eps\pair{q''_{n-\rho-1}S^\rho}{S_{n-\rho+\lambda-1}p_{\rho-\lambda}}=0,
\]
آنگاه بنا بر (21.) مقدار زیر حاصل می‌شود:
\[
 \pair{wq_{n-\rho+\lambda-1}}{T_{n-\rho}R_\lambda}
 =\eps\pair{q_{n-\rho+\lambda-1}[T_{n-\rho}R_\lambda]^{n-\lambda}}{S_{n-\rho+\lambda-1}p_{\rho-\lambda}}
 \quad\text{هم‌ارز}\quad
 \pair{q_{n-\rho-\lambda}S^\rho}{T_{n-\rho}p_{\rho-\lambda}}.
\]
در سراسر این فرایند فقط خاصیت فرم نرمالِ $S$ به کار رفت، نه خاصیت $T$؛ بنابراین می‌توان تاخوردگی هموردای دارای نقص $\lambda-1$ را که برای تولید (69.) به کار رفت، با $2\lambda-3$ تاخوردگی پایه و یک تاخوردگی هموردای دارای نقص $\lambda-2$ جایگزین کرد؛ تاخوردگی اخیر نیز به نوبهٔ خود تجزیه‌ای به تاخوردگی‌های پایه می‌پذیرد، و به همین ترتیب ادامه می‌یابد. به‌طور مشابه، با استفاده از خاصیت فرم نرمالِ $T$ و اعمال فرایند دوگان، تولید به کمک تاخوردگی‌های پایهٔ زیر انجام می‌شود:
\[
 S^\rho T^\rho,\quad T^{\rho+1},\ldots,T^{\rho+\lambda-1};
 \qquad R^{\rho-1}T^{\rho-1},\quad T^{\rho-2},\ldots,T^{\rho-\lambda+1},
\]
و یک تاخوردگیِ دوگانِ (69.). اکنون تاخوردگیِ هم‌ارز با تاخوردگیِ پیشین پدید می‌آید:
\[
 \pair{q_{n-\rho-\lambda}T^\rho}{S_{n-\rho}p_{\rho-\lambda}}.
\]

اگر تجزیهٔ به‌دست‌آمده در 2) را نیز بیفزاییم، در واقع \emph{تولید کلی‌ترین تاخوردگی‌ها از تاخوردگی‌های پایه} را به دست آورده‌ایم. هنوز باید ثابت کنیم که این تولید، صرف‌نظر از جابه‌جاییِ ترتیب، \emph{یکتا} است؛ یعنی برای هر تاخوردگی دقیقاً یک عدد $\alpha_\rho$ وجود دارد که شمار تاخوردگی‌های پایهٔ نوع $\rho$ را، برای $\rho=1,2,\ldots,n-1$، بیان می‌کند.

فرض کنیم $m_\rho$ درجه‌های ردیف‌های متغیر $p_\rho$ در فرم آغازین $f$ و $l_\rho$ درجه‌های فرم حاصل از تاخوردگی باشند. هر تاخوردگیِ پایهٔ نوع $\rho$ درجه‌های $m_{\rho-1},m_\rho,m_{\rho+1}$ را به‌ترتیب به $m_{\rho-1}+1,m_\rho-2,m_{\rho+1}+1$ تبدیل می‌کند. پس $\alpha$ها دستگاه معادلات زیر را ارضا می‌کنند:
\srcnumdisplay[4.2em]{(70.)}{%
 m_\rho-l_\rho=-\alpha_{\rho-1}+2\alpha_\rho-\alpha_{\rho+1},
 \qquad (\rho=1,2,\ldots,n-1)}
که در آن باید $\alpha_0$ را صفر و نیز $\alpha_n=0$ قرار داد؛ دترمینان این دستگاه مقدار $n$ را دارد\srcfn{*)}{اگر $\Delta_n$ دترمینانِ مرتبهٔ $n$ را نشان دهد، رابطهٔ بازگشتی چنین است: $\Delta_n=2\Delta_{n-1}-\Delta_{n-2}$؛ یعنی $\Delta_n-\Delta_{n-1}=\Delta_{n-1}-\Delta_{n-2}=\cdots=\Delta_2-\Delta_1=1$، زیرا $\Delta_2=3$ و$\Delta_1=2$.}. بنابراین، $\alpha$ها به‌طور یکتا تعیین می‌شوند و بنا بر قضیهٔ هشتم اعداد صحیح مثبت‌اند. از این، شرط‌هایی برای تفاضل‌های $m_\rho-l_\rho$ به دست می‌آید.

نتایج این بند \emph{به‌موجب روابط (62.)} برقرارند؛ اما اگر روابط (62.) برقرار نباشند، عبارت‌های پایانی با عبارت‌های آغازین هم‌ارز نیستند. زیرا تعریف هم‌ارزیِ داده‌شده در \S\ 6 صورت‌بندی زیر را نیز می‌پذیرد\srcfn{**)}{دراین‌باره، نک. استدلالِ پایان بند بعد.}: \glqq{}دو فرم اگر و فقط اگر اختلافشان از فرم‌هایی با تاخوردگیِ بالاتر تشکیل شده باشد، یعنی از فرم‌هایی که کاهش وزن بیشتری دارند، هم‌ارزند.\grqq{} این کاهش وزنِ بیشتر همیشه هنگامی رخ می‌دهد که دو ردیفِ $q_{n-\sigma-\lambda},p_{\tau-\lambda}$ که وارد تاخوردگی می‌شوند واقعاً ردیف‌های متغیر باشند؛ چنان‌که در (41.) و(42.)، در تاخوردگی‌هایی که در این بند هم‌ارز شمرده شدند، و نیز در فرم‌های هم‌ارزِ قضیهٔ هفتم چنین بود. اما اگر این ردیف‌ها با جایگذاری از ردیف‌های نماد مشتق شده باشند، لازم نیست این کاهش رخ دهد؛ چنان‌که دربارهٔ فرم‌های این بند که به‌موجب (62.) برحسب یکدیگر بیان‌پذیر بودند چنین بود. به‌آسانی می‌توان دید که حتی برخی از فرم‌های صفرشونده وزن کل بیشتری دارند.

\begin{center}
\emph{\S\ 9. رشته‌های فرم. قضیه‌های کاهش}\srcfn{*)}{مقایسه کنید با بررسی‌های متناظر در حوزهٔ سه‌تایی. رساله، \S\S\ 1--3.}.
\end{center}
اکنون \emph{رشتهٔ فرم} را، که در پایان \S\ 6 معرفی شد، اندکی دقیق‌تر چنین تعریف می‌کنیم: \glqq{}مجموعهٔ ساخته‌های ناوردای خطی در ضرایب یک فرم آغازینِ معین؛ یعنی مجموعهٔ فرم‌هایی که با تاخوردگیِ فرم آغازین در خودش پدید آمده‌اند، \emph{در ترتیبی برحسب فرم‌های بالاتر}\grqq{}.

برای توضیح فرم‌های بالاتر، فرض کنیم بنا بر قضیهٔ هفتم، فرم آغازین به‌صورت حاصل‌ضرب دو فرم نرمالِ $S\cdot T$ داده شده باشد؛ آنگاه فرم‌های دارای تاخوردگیِ بیشتر در خود را \emph{فرم‌های بالاتر} می‌نامیم؛ و در میان فرم‌هایی با تاخوردگیِ یکسان در خود، فرم‌هایی را که به‌طور ویژه نرمال‌سازی شده‌اند نیز بالاتر می‌نامیم. دقیق‌تر بگوییم: فرض کنیم فرم $A$ با شمارِ $\alpha_\rho$ و فرم $B$ با شمارِ $\beta_\rho$ از تاخوردگی‌های پایهٔ نوع $\rho$ پدید آمده باشند. $B$ بالاتر از $A$ است:
\begin{enumerate}
\renewcommand{\labelenumi}{\arabic{enumi})}
\item اگر در دستگاه نامساوی‌های $\beta_\rho\geqq\alpha_\rho$، $\rho=1,2,\ldots,n-1$، برای دست‌کم یک مقدارِ $\rho$ علامت نامساوی برقرار باشد\srcfn{**)}{اگر در دستگاه نامساوی‌ها در برخی جاها علامت $>$ و در برخی جاها علامت $<$ برقرار باشد، فرم‌ها با یکدیگر مقایسه‌پذیر نیستند؛ زیرا فرمی که با تاخوردگی از فرم دیگری پدید آمده است همیشه با یک دستگاه مثبت از مقادیر تاخوردگی‌های پایه همراه است، یعنی در حالت ما با مقادیر مثبت یا صفرِ $\beta_\rho-\alpha_\rho$.}؛
\item اگر برای همهٔ مقادیرِ $\rho$ علامت تساوی برقرار باشد، اما ردیف‌های نمادِ کمتری در یک حاصل‌ضرب ماتریسی گرد آمده باشند. این نرمال‌سازی به‌سبب کاهش‌های \S\ 8 ضروری است. بر این اساس، برای نمونه، $\pair{S^{n-1}}{T_{n-1}}$ از $\pair{S^{n-1}}{[S^1T^\rho]_{n-\rho-1}p_\rho}$ بالاتر می‌شود.
\item اگر برای همهٔ مقادیرِ $\rho$ علامت تساوی برقرار و شمار ردیف‌های نمادِ گردآمده در یک حاصل‌ضرب ماتریسی نیز یکسان باشد، $B$ با نرمال‌سازی‌ای که فی‌نفسه دلخواه است اما یک‌بار برای همیشه تثبیت می‌شود، از $A$ بالاتر شمرده می‌شود. این نرمال‌سازی به‌سبب متناهی‌بودنِ شمار فرم‌های متعلق به یک دستگاه مقادیرِ $\alpha$ همیشه ممکن است؛ برای نمونه، می‌توان آن را با ممتازکردنِ فرم $S$ به دست آورد (شمار بیشتر ردیف‌های نمادِ $S$، مجموع بالاتر اندیس‌های وزنیِ بالاییِ ردیف‌های $S$، و مانند این‌ها).
\end{enumerate}
بهتر است رشتهٔ فرم‌ها را چنین تصور کنیم که به‌صورت آرایه‌ای $(n-1)$-بعدی از نقاط شبکه سامان یافته است و جهت‌های $n-1$گانه با تاخوردگی‌های پایهٔ $n-1$گانه متناظرند. در این صورت، به هر فرم دقیقاً یک دستگاه مقادیرِ $\alpha$، یعنی یک نقطهٔ شبکه با مختصات صحیح مثبت، نسبت داده می‌شود؛ در حالی که نرمال‌سازیِ بالا ترتیب فرم‌های متفاوتِ متعلق به یک دستگاه مقادیرِ $\alpha$ را به‌طور یکتا تعیین می‌کند. هر فرم دلخواهِ رشتهٔ فرم‌ها، فرم آغازینِ رشتهٔ تازه‌ای از فرم‌ها را به دست می‌دهد که به‌عنوان بخشی از کل رشتهٔ فرم‌ها در آن جای دارد؛ دقیق‌تر، در آن فرم‌هایی جای دارد که با حرکت از فرم آغازینِ تازه در جهت‌های مثبتِ $n-1$گانه به دست می‌آیند.

اکنون بنا بر قضیهٔ هشتم و نظریهٔ عمومیِ بسط‌های سری، دقیقاً همان قضیه‌های کاهشِ حوزه‌های دوتایی و سه‌تایی برای رشته‌های فرم برقرارند (مقایسه کنید با رساله، \S\ 3)؛ قضیهٔ اصلیِ آن‌ها را نیز به‌اختصار استخراج می‌کنیم. نخست تعریف را می‌آوریم:

\glqq{}در یک دستگاه معین از فرم‌ها، شمار متناهیِ معینی از فرم‌ها در یک پیمانه گرد آورده شده باشند؛ فرمی را \emph{کاهش‌پذیر} می‌نامیم اگر بتوان آن را با فرم‌هایی بیان کرد که ناورداها را به‌عنوان عامل دارند، یا با \glqq{}فرم‌های بالاتر\grqq{}؛ یعنی با فرم‌های بالاترِ کل رشتهٔ فرم‌هایی که آن فرم به آن تعلق دارد، یا با فرم‌هایی که نمادهای پیمانه را در مرتبه‌ای بالاتر دربر دارند. رشتهٔ فرمی را که کاهش‌پذیر باشد \emph{کاهشگر} می‌نامیم.\grqq{} آنگاه داریم

\emph{قضیهٔ نهم: اگر فرم آغازینِ یک رشتهٔ فرم‌ها از آن رو کاهش‌پذیر باشد که با تاخوردگی با یک کاهشگر پدید آمده است، آنگاه کل رشتهٔ فرم‌هایی که از این فرم آغازین پدید می‌آید کاهش‌پذیر است.}

برای اثبات، توجه می‌کنیم که تاخوردگیِ فرم $f$ با فرم دومِ $g$، فرم $h$ را پدید می‌آورد که بنا بر (45.) می‌تواند فرمی دانسته شود که چند ردیف متغیرِ هموردای $p_\rho,p'_\rho$ را دربر دارد، یعنی همان فرم $f$. اما طبق نظریهٔ عمومیِ بسط سری، چنین فرم‌هایی را می‌توان به‌صورت قطبی‌های همورداهای بنیادیِ $f$ نمایش داد؛ به این ترتیب که بسط سری را پیاپی بر هر دو ردیف متغیرِ هموردای $p_\rho,p'_\rho$ اعمال کنیم. بنا بر (38.)، همورداهای بنیادی‌ای که در این میان ظاهر می‌شوند، فرم‌های تولیدشده با تاخوردگیِ هموردا هستند. از آنجا که ترتیبِ به‌کارگیریِ بسط سری هنوز دلخواه است، می‌توانیم به‌ویژه ترتیبی را برگزینیم که با تولید کلی‌ترین تاخوردگی‌ها از تاخوردگی‌های پایه متناظر باشد. دستگاه حاصلِ همورداهای بنیادی آنگاه با رشتهٔ فرم‌های $f$ یکسان می‌شود؛ و فرم‌های پدیدآمده با تاخوردگیِ هموردای دارای نقص بالاتر، در میان فرم‌های پدیدآمده با تاخوردگی‌های پایه جای می‌گیرند. با هر ترتیب دلخواهِ بسط سری نیز به قطبی‌های رشتهٔ فرم‌های $f$ می‌رسیم؛ فرض کنیم دستگاه دومی از همورداهای بنیادی با آن ترتیب متناظر باشد. چون رشتهٔ فرم‌ها همهٔ ساخته‌های ناوردا را فرامی‌گیرد، هر فرمِ دستگاه دوم را می‌توان به‌صورت خطی برحسب قطبی‌های رشتهٔ فرم‌ها بیان کرد، و دقیق‌تر برحسب قطبی‌های فرم‌هایی که کاهش وزن یکسان یا بیشتری نشان می‌دهند. این گزارهٔ اخیر از آنجا نتیجه می‌شود که قطبی‌ها را می‌توان (نک. \S\ 7) ناورداهای هم‌زمانِ یک رشتهٔ فرم‌های $H$ در (52.)، با درجه‌های متناظر با فرم $h$، و رشتهٔ فرم‌های $f$ دانست. هر تاخوردگیِ $H$ در خودش با کاهشی در وزن متناظر است که پس از انجام جایگذاریِ (11.) به ردیف‌های نماد و بدین‌وسیله به فرم‌های $f$ منتقل می‌شود\srcfn{*)}{این استدلال مبنای صورت‌بندیِ دومِ مفهوم هم‌ارزی (پایان \S\ 8) نیز هست. شرح‌های بیشتر دربارهٔ دستگاه‌های مختلفِ همورداهای بنیادی در حوزهٔ سه‌تایی نزد شتودی، همان اثر، II. \S\ 7، قضیه‌های 7) و8)، آمده است. بنا بر قضیهٔ هفتمِ ما، همان مطالب برای $n$ متغیر نیز برقرارند.}.

هر فرم دلخواهِ رشتهٔ فرم‌های $h$ با تاخوردگیِ $h$ در خودش پدید آمده است؛ یعنی از یک سو با تاخوردگیِ بالاترِ $g$ با $f$، و از سوی دیگر با تاخوردگیِ $f$ یا $g$ در خودش. در هر دو حالت، بسط سری، همان‌گونه که در بالا برای فرم آغازینِ $h$ نشان داده شد، به قطبی‌های رشتهٔ فرم‌های $f$ می‌انجامد. اگر به‌ویژه $f$ کاهشگر باشد، بسطی برحسب قطبی‌های کاهشگر به دست می‌آید؛ پس رشتهٔ فرم‌های $h$ کاهش‌پذیر است.

کاربردهای این قضیه در کاهش دستگاه‌های کاملِ فرم‌ها نیز همانند حوزه‌های دوتایی و سه‌تایی به دست می‌آیند.


\clearpage
\setcounter{footnote}{0}
\providecommand{\Kfield}{\mathfrak{K}}
\providecommand{\Ssys}{\mathfrak{S}}
\providecommand{\Mfield}{\mathfrak{M}}
\providecommand{\tq}{\tau}
\providecommand{\ds}{\displaystyle}

\begin{center}
{\Large\bfseries \srcLTR{5.}~\srcRTL{میدان‌های توابع گویا.}\srcfn{1)}{سخنرانی‌ای که در گردهمایی طبیعت‌پژوهان وین در سال \srcLTR{1913} ایراد شد.}}\par
\vspace{1.2em}
\srcRTL{نوشتهٔ \srcspaced{امی نوتر} در ارلانگن.}\par
\medskip
\srcLTR{J. Ber. d. DMV 22 (1913), S. 316--319}
\end{center}

\vspace{2em}

پرسش‌های زیر در اصل به گفت‌وگوهایی با ا. فیشر بازمی‌گردند. افزون بر این، ا. اشتاینیتس پیش‌تر برخی از این پرسش‌ها را در حالت ویژهٔ یک متغیر نامعین ــ و آن هم با فرض‌هایی کلی‌تر دربارهٔ حوزهٔ ضرایب ــ مطرح و حل کرده است.\srcfn{2)}{\srcLTR{E. Steinitz: Algebraische Theorie der Körper} (به‌ویژه \S\ 24). \srcLTR{Crelles Journal 137. 1910.}}

از \glqq{}میدان توابع گویا\grqq{} میدانی را مراد می‌کنم که اعضایش توابع گویای $n$ متغیر نامعین‌اند و ضرایبشان از میدانی از اعدادِ ازپیش‌داده‌شده گرفته می‌شود؛ این میدان می‌تواند، به‌ویژه، همهٔ اعداد مختلط را نیز در بر گیرد. میدان توابع متقارنِ $n$ کمیت، یا به‌طور کلی‌تر کلِ توابع گویای $n$ متغیر نامعین که جایگشت‌های گروه معینی را می‌پذیرند (میدان‌های گونهٔ لاگرانژی)، و نیز میدان ناورداها، نمونه‌هایی از چنین میدان‌هایی‌اند.

میان دو میدان نخست و میدان اخیر تفاوتی شاخص وجود دارد: دو میدان نخست \emph{$n$ تابعِ مستقلِ جبری}، یعنی توابع متقارنِ مقدماتی، در بر دارند؛ حال‌آنکه شمار ناورداهای مستقلِ جبری همیشه از شمار متغیرهای نامعین کمتر است. ازاین‌رو مهم است که بتوان \emph{به‌طور کلی میان چنین میدان‌هایی از نوع دوم و میدان‌های نوع اول تناظری یک‌به‌یک برقرار کرد}، بدین‌گونه که برخی از متغیرهای نامعین را با اعداد جایگزین کنیم. پس در ادامه خود را \emph{به میدان‌های نوع اول، یعنی میدان‌های دارای $n$ تابع مستقلِ جبری} محدود می‌کنم؛ در پرتو این تناظر، نتایج سپس به‌طور کلی برقرارند.

پرسش‌ها پیرامون سه مفهوم بنیادیِ پایه سامان می‌یابند: پایهٔ گویا، پایهٔ کمینه و پایهٔ تمامیت.

1. از \glqq{}پایهٔ گویا\grqq{} شمار متناهی‌ای از توابع میدان را مراد می‌کنم، چنان‌که هر تابعِ میدان را بتوان با ضرایبی از حوزهٔ ضرایبِ ازپیش‌داده‌شده به‌صورت ترکیبی گویا از آن شمار متناهی نمایش داد. با استدلال‌هایی ساده ــ که در حالت یک متغیر نامعین نزد ا. اشتاینیتس نیز در همان موضع آمده‌اند ــ وجود پایهٔ گویا برای هر میدان توابع گویا ثابت می‌شود. برای نمونه، بنا بر قضیهٔ لاگرانژ، توابع متقارنِ مقدماتی همراه با تابعی متعلق به گروه، پایه‌ای گویا برای میدان‌های گونهٔ لاگرانژیِ یادشده تشکیل می‌دهند.

2. هر پایهٔ گویا (برای میدان‌های نوع اول) باید دست‌کم $n$ تابع داشته باشد؛ اگر دقیقاً $n$ تابع داشته باشد، که در آن صورت ناگزیر مستقلِ جبری‌اند، آن را \glqq{}پایهٔ کمینه\grqq{} می‌نامم. با قضیه‌های لوروت، کاستلنووو و انریکس می‌توان به پرسش وجود پایهٔ کمینه پاسخی جزئی داد.\srcfn{1)}{\srcLTR{J. Lüroth: Beweis eines Satzes über rationale Kurven. Math. Ann. 9. 1875. --- G. Castelnuovo: Sulla razionalità delle involuzioni piane. Math. Ann. 44. 1893. --- F. Enriques: Sopra una involuzione non razionale dello spazio. Rend. Acc. Linc. vol. 21. 21. Jan. 1912.}} بنا بر این قضیه‌ها، برای میدانِ یک متغیر نامعین همیشه پایه‌ای کمینه وجود دارد: تابع لوروتِ میدان که تا یک تبدیل خطی-کسری به‌طور یکتا معین است (نک. اشتاینیتس، \S\ 24). برای میدان‌های دو متغیر نامعین، هرگاه توسیع جبریِ حوزهٔ ضرایب را مجاز بدانیم همیشه پایه‌ای کمینه وجود دارد و عموماً فقط در همین صورت وجود دارد؛ حال‌آنکه برای سه یا بیش از سه متغیر نامعین، پایهٔ کمینه عموماً اصلاً وجود ندارد. هنوز تلاشی برای مشخص‌کردن میدان‌های ویژهٔ دارای پایهٔ کمینه صورت نگرفته است.

اکنون پرسش پایهٔ کمینه را در میدان‌های گونهٔ لاگرانژی بیشتر پی گرفته‌ام. این پرسش در اینجا معنای زیر را می‌یابد: \glqq{}اگر میدان گونهٔ لاگرانژیِ متعلق به گروه $G$ پایه‌ای کمینه داشته باشد، می‌توان کلِ معادله‌های دارای گروه مقررِ $G$ را به‌طور گویا و با نمایش پارامتری ساخت؛ آن هم بر روی هر میدان از اعدادی که حوزهٔ ضرایبِ پایهٔ کمینه را در بر گیرد ــ پس به‌ویژه بر روی هر میدان دلخواهی از اعداد، اگر این حوزهٔ ضرایب میدان اعداد گویا باشد.\grqq{}

گروه متقارن ساده‌ترین نمونه را به دست می‌دهد؛ در اینجا توابع متقارنِ مقدماتی پایه‌ای کمینه با ضرایب گویا تشکیل می‌دهند. از این امر، معادله‌ای با ضرایب نامعین به‌سادگی به‌عنوان نمایش پارامتری حاصل می‌شود. قضیهٔ تحویل‌ناپذیریِ هیلبرت نشان می‌دهد که چگونه می‌توان از این معادله، بر روی هر میدان از اعداد، به شمار دلخواهی از معادله‌های دارای گروه گالوای متقارن گذر کرد.

اکنون از قضیه‌های لوروت و کاستلنووو بی‌درنگ وجود پایهٔ کمینه برای همهٔ گروه‌هایی که در معادله‌های درجهٔ سوم و چهارم پدید می‌آیند نتیجه می‌شود؛ همچنین روشن می‌شود که این پایهٔ کمینه ضرایبی از اعداد گویا دارد. بنابراین بر روی هر میدان دلخواهی از اعداد می‌توان کلِ معادله‌های درجهٔ سوم و چهارمِ دارای گروه مقرر را به‌طور گویا ساخت. برای گروه‌های دوری و دووجهی پایه‌ای کمینه وجود دارد که حوزهٔ ضرایبش میدان ریشه‌های $n$اُمِ واحد است؛ همین حکم برای برخی گروه‌های فرادوریِ دیگر نیز برقرار است. این پرسش که آیا وجود پایهٔ کمینه اصلاً ویژگیِ مشخص‌کنندهٔ برخی گروه‌هاست یا نه، ناگزیر بی‌پاسخ می‌ماند.

3. برای رسیدن به آخرین مفهوم پایه، کلِ چندجمله‌ای‌های موجود در میدان را در نظر می‌گیرم. از \glqq{}پایهٔ تمامیت\grqq{} شمار متناهی‌ای از این چندجمله‌ای‌ها را مراد می‌کنم، چنان‌که هر چندجمله‌ایِ میدان را بتوان با ضرایبی از حوزهٔ ضرایبِ ازپیش‌داده‌شده به‌صورت ترکیبی چندجمله‌ای از آن شمار متناهی نمایش داد. پرسش وجود پایهٔ تمامیت ــ که برای نمونه متناهی‌بودن دستگاه ناورداها را به‌عنوان حالتی ویژه در بر می‌گیرد ــ با صورت‌بندی‌ای اندکی متفاوت، از سوی هیلبرت در \glqq{}مسائل ریاضی\grqq{} او\srcfn{1)}{\srcLTR{D. Hilbert: Mathematische Probleme. Vortrag Paris 1900. Problem 14. Göttinger Nachr. 1900.}} به‌عنوان مسئلهٔ توابعِ تمامِ نسبی مطرح شده است.

فعلاً فقط می‌توانم یک رده از میدان‌هایی را ذکر کنم که برایشان پایهٔ تمامیت وجود دارد. اگر میدان دستگاهی از $n$ چندجمله‌ای در بر داشته باشد، به‌گونه‌ای که برآیندِ همگنِ جزءهای همگنِ دارای بالاترین درجهٔ این چندجمله‌ای‌ها ناصفر باشد، آنگاه پایهٔ تمامیت وجود دارد؛ این پایه را ضرایبِ آن معادلهٔ تحویل‌ناپذیرِ میدان برای فرم خطی زیر، نسبت به همهٔ متغیرهای $u$، به دست می‌دهند:
\[
  u_0=u_1x_1+u_2x_2+\cdots+u_nx_n.\srcfnmark{2)}
\]
\srcfntext{2)}{این معادلهٔ تحویل‌ناپذیر، در حالت یک متغیر نامعین، نزد ا. اشتاینیتس در برهان قضیهٔ لوروت آمده است.}
اگر، به‌ویژه، مقدار برآیند برابر با یکه‌ای از حوزهٔ ضرایب باشد، تمام‌بودنِ نمایش نیز تضمین می‌شود. میدان‌های گونهٔ لاگرانژی نمونه‌ای از این ردهٔ میدان‌ها هستند، زیرا برآیندِ توابع متقارنِ مقدماتی برابر با $\pm 1$ است. نمونه‌ای دیگر (بدون تمام‌بودن) همهٔ میدان‌هایی‌اند که فقط یک چندجمله‌ایِ مستقلِ جبری در بر دارند؛ پایهٔ تمامیت در اینجا با تابع لوروتِ کوچک‌ترین زیرمیدانی که همهٔ چندجمله‌ای‌ها را در بر دارد یکسان است. بدین‌سان، چنان‌که می‌دانیم، همهٔ ناورداهای پروژکتیوِ گویای چندجمله‌ایِ یک فرم درجه‌دو با $n$ متغیر، تابع‌هایی چندجمله‌ای از ممیزند.

شرط یادشده فقط برای وجود پایهٔ تمامیت کافی است و به‌هیچ‌وجه ضروری نیست. به‌آسانی می‌توان میدان‌هایی ساخت که در آن‌ها برآیندِ هر $n$ چندجمله‌ای صفر است، اما بااین‌همه پایهٔ تمامیتی دارند که ضرایبِ همان معادلهٔ تحویل‌ناپذیر آن را به دست می‌دهند. این پرسش پیش می‌آید که آیا در عام‌ترین حالت نیز ضرایبِ همان معادله پایهٔ تمامیت را به دست می‌دهند یا نه.

\clearpage

\begin{center}
{\Large\bfseries \srcLTR{6.}~\srcRTL{میدان‌ها و دستگاه‌های توابع گویا}}\par
\vspace{0.6em}
\srcLTR{Emmy Noether (Erlangen).}\par
\vspace{1.2em}
\srcLTR{Math. Ann. 76 (1915), S. 161--196}
\end{center}

\vspace{2em}

این مقاله به پرسش‌های پایه در دستگاه‌های دلخواهِ توابع گویا و توابعِ گویایِ تمام (چندجمله‌ای‌ها) می‌پردازد؛ روش‌های به‌کاررفتهٔ نظریهٔ میدان نشان می‌دهند که امر اساسی، حل این پرسش‌ها برای میدان‌هایی از توابع گویا ــ یعنی میدان‌های توابع گویا\srcfn{*)}{نک. گزارش مقدماتی‌ای به مناسبت گردهمایی طبیعت‌پژوهان وین: «میدان‌های توابع گویا» ــ \srcLTR{Jahresber. d. D. Math.-Ver. 22 (1913)} ــ، که در آن مروری بر پرسش‌های مطرح‌شده و نتایج مربوط به میدان‌های توابع آمده است.} ــ است، حال‌آنکه تعمیم نتایج به دستگاه‌های دلخواه به‌صورت نتیجه‌ای حاصل می‌شود.

از پرسش‌های پایه برای دستگاه‌های عمومی، تاکنون فقط وجود پایهٔ مدولی (\srcLTR{Modulbasis}) شناخته شده است که قضیهٔ هیلبرت (\srcLTR{Math. Ann. 36}) آن را تضمین می‌کند. در ادامه، پرسش نمایش‌پذیریِ گویا با اثبات وجود پایهٔ گویا برای هر دستگاه دلخواه به‌طور کامل پاسخ داده می‌شود (\S\ 7)؛ پایهٔ گویای میدان‌ها پیش‌تر در \S\ 4 به دست می‌آید. وجود پایهٔ گویا اجازه می‌دهد که در همه‌جا از میدان یا دستگاهی که به‌طور مجرد تعریف شده است آغاز کنیم و بدین‌سان از دشواری‌هایی بپرهیزیم که نه از خود دستگاه، بلکه صرفاً از انتخاب ویژهٔ پایهٔ گویا ناشی می‌شوند؛ مانند پدیدآمدن مخرج‌های ویژه یا نقاط پایهٔ توابع پایه.

برای میدان‌ها، پرسش پایهٔ کمینه، یعنی پایه‌ای گویا از توابع مستقلِ جبری، نیز بررسی می‌شود (\S\ 6). روش‌های نظریهٔ میدان همچنین به پایه‌ای گویای ممتاز، یعنی پایهٔ اینولوشن (\srcLTR{Involutionsbasis})\srcfn{**)}{این نام از آن رو برگزیده شد که میدان توابع گویا، در تعبیر هندسی، عمومی‌ترین اینولوشن را در یک فضای خطی نمایش می‌دهد.} می‌انجامند (\S\ 5)؛ این پایه به‌ویژه برای دامنه‌های صحیحِ ساخته‌شده از چندجمله‌ای‌ها اساسی است (\S\ 8). در اینجا نمایشی با مخرج ثابت به دست می‌دهد (توانی از تابعی که دامنهٔ صحیح تعیین می‌کند)، همان‌گونه که برای نمونه در حالت ویژهٔ نمایش تیپیکِ ناورداها شناخته شده است.

اکنون از نمایش‌پذیریِ گویا، تا حد امکان نتایجی فراگیر برای نمایش‌پذیریِ گویایِ تمام، یعنی برای پرسش متناهی‌بودن به معنای خاص‌تر، استخراج می‌شود.

همهٔ قضیه‌های متناهی‌بودن که تاکنون شناخته شده‌اند بر این واقعیت استوارند که قضیهٔ هیلبرت دربارهٔ پایهٔ مدولی، متناهی‌بودن را برای همهٔ دستگاه‌هایی تضمین می‌کند که در آن‌ها نمایشی
\[
  F=A_1f_1+A_2f_2+\cdots+A_kf_k
\]
نمایش دومی
\[
  F=B_1f_1+B_2f_2+\cdots+B_kf_k
\]
را در پی دارد، که در آن $B_i$ها به دستگاه تعلق دارند و درجه‌شان از درجهٔ $F$ کمتر است.

در مقابل، قضیهٔ پایهٔ گویا و روش‌های نظریهٔ میدان اجازه می‌دهند که تحت فرض‌هایی از نوعی کاملاً متفاوت به متناهی‌بودن برسیم.

بدین‌سان، در پاسخی جزئی به یکی از مسائل هیلبرت (\srcLTR{Math. Probleme 14})، رده‌ای از توابع تمامِ نسبی (دامنه‌های تمامِ نسبی از نوع اول) به دست می‌آید که دامنه‌های صحیحِ متناهی‌مولدند و پایهٔ تمامیتشان را پایهٔ اینولوشنِ یادشده به دست می‌دهد (\S\ 10). همچنین، به قیاس با برهان هیلبرت برای قضیهٔ کرونکر دربارهٔ دستگاه بنیادیِ کمیت‌های تمامِ جبری (\srcLTR{Volle Invariantensysteme, \S\ 2; Math. Ann. 42})، می‌توان نشان داد که همهٔ دستگاه‌های منظمِ چندجمله‌ای‌ها ــ یعنی همهٔ دستگاه‌هایی که می‌توان آن‌ها را به دستگاه‌های بی‌نقطهٔ پایه نگاشت ــ پایهٔ تمامیت دارند (\S\ 12)، حال‌آنکه به‌آسانی می‌توان دستگاه‌های نامنظمی بدون پایهٔ تمامیت به دست داد. به‌عنوان نمونه‌ای ساده از این واقعیت، قضیهٔ زیر را یاد می‌کنیم: \glqq{}در هر دستگاه دلخواه از بی‌نهایت چندجمله‌ایِ یک متغیر نامعین، شمار متناهی‌ای از همین چندجمله‌ای‌ها وجود دارد، چنان‌که هر چندجمله‌ایِ دستگاه را بتوان به‌صورت ترکیبی چندجمله‌ای از آن شمار متناهی نمایش داد.\grqq{} نمونه‌های بیشتر در \S\ 13 آمده‌اند.

سرانجام، بندهای پایانی (\S\S\ 14 و15) قضیه‌های پایه را از حیث تمامیت تقویت می‌کنند.

اصل انتقالِ \S\ 3 ابزاری اساسی برای بررسی فراهم می‌آورد؛ بنا بر آن، کافی است همهٔ پرسش‌های پایهٔ مطرح‌شده در اینجا را برای دستگاه‌هایی بررسی کنیم که در آن‌ها شمار متغیرهای نامعین با شمار توابع مستقلِ جبریِ دستگاه برابر است.

گفت‌وگوهایی با آقای ا. فیشر، به‌ویژه پرسش او دربارهٔ پایهٔ کمینهٔ میدان‌های گونهٔ لاگرانژی، انگیزهٔ این مقاله را فراهم آورد. بررسی‌های مربوط، که به ساخت معادله‌هایی با گروه مقرر انجامیدند (نک. گزارش یادشده دربارهٔ \glqq{}میدان‌های توابع گویا\grqq{}، بخش 2)، برای انتشار دیگری محفوظ مانده‌اند.

گسترش اساسیِ مقالهٔ حاضر نسبت به گزارش نخستین در این است که قضیه‌های پایه، که در آنجا فقط برای میدان‌ها یا دامنه‌های صحیحِ ویژه بیان شده بودند، به دستگاه‌های دلخواه انتقال یافته‌اند. این انتقال با مفهوم کوچک‌ترین میدان دربرگیرندهٔ یک دستگاه دلخواه، به‌عنوان تعمیم میدان کسرهای یک دامنهٔ صحیح، به روشی ساده انجام می‌شود (\S\ 7). آنچه مرا به ساختن این مفهوم واداشت این بود که آقای ک. هنتزلت، پس از آگاهی از پایهٔ گویای میدان‌ها، توانست وجود پایهٔ گویا برای دستگاه‌های دلخواه را از راه غیرمستقیمِ پایهٔ مدولیِ هیلبرت ثابت کند. همچنین ملاحظهٔ ک. هنتزلت، مبنی بر اینکه استدلال‌هایی که من بر دامنه‌های صحیح اعمال کرده بودم برای دستگاه‌های دلخواهِ تابع همان فرض‌ها نیز معتبرند، مرا به قضیهٔ پایهٔ تمامیت برای دستگاه‌های منظم رهنمون شد؛ افزون بر این، چند ملاحظهٔ کوچک و پراکندهٔ دیگر را نیز وامدار آقای هنتزلت‌ام.

سرانجام باید یاد کرد که در حالت یک متغیر نامعین، پایهٔ گویا و پایهٔ کمینهٔ میدان‌ها نزد ا. اشتاینیتس در «نظریهٔ جبری میدان‌ها» (\srcLTR{Algebraische Theorie der Körper})، \S\ 24 (\srcLTR{Crelles Journ. 137}) آمده‌اند؛ در آنجا حوزهٔ ضرایب میدانی به عمومی‌ترین معنا فرض شده است. اینکه قضیه‌های ارائه‌شده در اینجا تا چه اندازه تحت این فرض‌های کلی‌تر پابرجا می‌مانند، در ادامه بررسی نشده است؛ به‌هرحال برهان‌ها نیازمند تعدیل‌اند، زیرا برای نمونه ابزارهای نظریهٔ دترمینان‌های تابعی در میدان‌هایی با مشخصهٔ $p$ از کار می‌افتند.

\section*{\S\ 1. میدان‌های $\Kfield_{n\rho}$ و دستگاه‌های $\Ssys_{n\rho}$.}

در ادامه، دستگاه‌های توابع گویای $n$ متغیر نامعین، و به‌ویژه میدان‌های توابع گویا، بررسی می‌شوند.

بنابراین از $f(x_1\cdots x_n)$، $g(x_1\cdots x_n)$، \ldots\ یا به‌اختصار $f(x)$، $g(x)$، \ldots\ همواره توابع گویای $x_1\cdots x_n$ مراد است؛ این توابع در نمایش کاهش‌یافته ــ یعنی با صورت و مخرج نسبت به هم اول ــ فرض می‌شوند. توابع گویای تمام (چندجمله‌ای‌ها) با حروف بزرگ، $F(x)$، $G(x)$، \ldots\ نشان داده خواهند شد. حوزهٔ ضرایب $\Omega$ میدانی از اعداد، هرچه باشد، فرض می‌شود و ازاین‌رو می‌تواند، به‌ویژه، همهٔ اعداد مختلط را نیز در بر گیرد؛ همچنین $\Omega$ می‌تواند شمار متناهی‌ای پارامتر داشته باشد.

کمیت‌های متعلق به $\Omega$، یعنی همهٔ توابع درجهٔ صفر، همواره متعلق به دستگاه شمرده می‌شوند.

پس از این قراردادها می‌توان میدان ساخته‌شده از توابع گویا ــ میدان توابع گویا ــ را به‌طور مجرد تعریف کرد.

\emph{تعریف \srcLTR{I}: دستگاهی از توابع گویا میدان نامیده می‌شود اگر شرط‌های زیر را برآورده کند:}
\begin{enumerate}
\item \emph{اگر $f(x)$ را در بر داشته باشد، برای هر کمیت $c$ از $\Omega$، $c\cdot f(x)$ را نیز در بر دارد.}
\item \emph{اگر $f(x)$ و $g(x)$ را در بر داشته باشد، همواره $f(x)+g(x)$ و $f(x)\cdot g(x)$، و ــ اگر $g(x)\ne0$ ــ خارج‌قسمت $f(x):g(x)$ را نیز در بر دارد.}\srcfn{*)}{$f(x)$ و $g(x)$ در اینجا می‌توانند درجهٔ صفر، یعنی کمیت‌هایی از $\Omega$، نیز باشند.}
\end{enumerate}
از میدان‌های ویژه، میدان‌هایی‌اند که با الحاق شمار متناهی‌ای از توابع گویا
\[
  f_1(x),\ f_2(x),\ldots, f_k(x)
\]
به $\Omega$ پدید می‌آیند؛ آن‌ها را با
\[
  \Omega\bigl(f_1(x)\cdots f_k(x)\bigr)
  \quad\text{یا}\quad
  \Omega(f_1\cdots f_k)
\]
نشان می‌دهیم.\srcfn{**)}{در \S\ 4 روشن می‌شود که همین \glqq{}میدان‌های ویژه\grqq{} همهٔ میدان‌ها را در بر می‌گیرند.} از میان این میدان‌های ویژه، میدان حاصل از الحاق $x_1,\ldots,x_n$ به $\Omega$، یعنی
\[
  \Omega(x_1\cdots x_n)
\]
را نیز یاد می‌کنیم؛ این میدان عینِ مجموعهٔ همهٔ توابع گویای $x_1\cdots x_n$ با ضرایبی از $\Omega$ است.

همچنین (به پیروی از ا. اشتاینیتس) مفهوم میدان واسط را وارد می‌کنیم:
\begin{quote}
\glqq{}میدان $\Mfield$ میان $\Omega_1$ و $\Omega_2$ قرار دارد اگر $\Omega_1$ را در بر داشته و در $\Omega_2$ مندرج باشد\grqq{}؛
\end{quote}
آنگاه تعریف \srcLTR{I} را می‌توان به صورت زیر نیز بیان کرد:

میدان توابع گویای $n$ متغیر نامعین، میدانی واسط میان $\Omega$ و $\Omega(x_1\cdots x_n)$ است که در دست‌کم یک تابع، هر $n$ متغیر نامعین در آن واقعاً حضور دارند.

در کنار شمار متغیرهای نامعین، عدد مشخصهٔ دیگری برای دستگاه‌های توابع گویا به دست می‌آید: شمار توابع مستقلِ جبری، یا رتبهٔ جبری، که تعریف زیر آن را تعیین می‌کند:

\emph{تعریف \srcLTR{II}: دستگاهی از توابع گویا رتبهٔ جبری $\rho$ دارد اگر بتوان $\rho$ تابع از دستگاه یافت که مستقلِ جبری باشند، اما هر مجموعهٔ $(\rho+1)$ تابع از دستگاه وابستهٔ جبری باشد.}\srcfn{***)}{$\rho$ تابع $f_1(x)\cdots f_\rho(x)$ مستقلِ جبری نامیده می‌شوند اگر هر رابطهٔ
\[
F\bigl(f_1(x)\cdots f_\rho(x)\bigr)=0\quad\text{همانی نسبت به }x_1\cdots x_n
\]
نتیجه دهد که
\[
F(\lambda_1\cdots \lambda_\rho)=0\quad\text{همانی نسبت به }\lambda_1\cdots \lambda_\rho.
\]
در حالت مقابل، آن‌ها وابستهٔ جبری نامیده می‌شوند.}

دستگاه‌های متشکل از شمار متناهی یا نامتناهیِ توابع گویای $n$ متغیر نامعین و دارای رتبهٔ جبری $\rho$ را با زیرنویس دوگانهٔ
\[
  \Ssys_{n\rho}
\]
نشان می‌دهیم. به‌ویژه، از
\[
  \Kfield_{n\rho}
\]
میدانی با $n$ متغیر نامعین و رتبهٔ جبری $\rho$ مراد است. نامساوی زیر برقرار است:
\[
  1\leq \rho\leq n.
\]

چند نمونه از میدان‌های $\Kfield_{nn}$ و $\Kfield_{n\rho}$ نیز می‌آوریم:

1. میدان‌های گونهٔ لاگرانژی نمونه‌هایی از میدان‌های $\Kfield_{nn}$ هستند؛ می‌توان آن‌ها را چنین تعریف کرد:
\begin{quote}
\glqq{}مجموعهٔ همهٔ توابع گویا (و توابع گویای تمام) در $x_1\cdots x_n$ که جایگشت‌های یک گروه جایگشتی $G$ روی $x_1\cdots x_n$ را می‌پذیرند.\grqq{}
\end{quote}
این میدان‌ها همواره میدان توابع متقارن، و در نتیجه $n$ تابع متقارنِ مقدماتیِ مستقلِ جبری را نیز در بر دارند.

2. میدان‌های ناورداهای پروژکتیو نیز نمونه‌هایی از میدان‌های $\Kfield_{n\rho}$ هستند؛ این میدان‌ها از مجموعهٔ همهٔ ناورداهای گویا (و گویای تمام) یک یا چند فرم بنیادی ساخته می‌شوند.\srcfn{*)}{مناسب است که تنها تبدیل‌های با دترمینان $1$ را در نظر بگیریم؛ در این صورت هر ترکیب گویای هر شمار از ناورداها نیز تبدیل را می‌پذیرد و به‌راستی میدانی به دست می‌آید. ناورداهای همگن و هم‌وزن، اگر به‌تنهایی در نظر گرفته شوند، میدان تشکیل نمی‌دهند، اما دستگاهی $\Ssys_{n\rho}$ تشکیل می‌دهند.} داریم $\rho<n$، زیرا شمار ناورداهای مستقلِ جبری همواره از شمار ضرایب کمتر است.

حکم متناظر برای میدان‌های ناوردا نسبت به زیرگروه‌های گروه پروژکتیو نیز برقرار است.

\section*{\S\ 2. لمی دربارهٔ توابع وابستهٔ جبری.}

به‌کمک لم این بند، در \S\ 3 معلوم خواهد شد که رتبهٔ جبری یک دستگاه، عددی است که به‌راستی مشخصهٔ آن است. زیرا این لم نشان می‌دهد که با تعیین پیاپی مقادیر عددی برای برخی متغیرهای نامعین، به‌نحوی که $\rho$ تابع مستقل جبری همچنان مستقل جبری بمانند، هیچ تابعی که به این $\rho$ تابع وابستهٔ جبری است نه می‌تواند به‌طور همانی صفر شود و نه می‌تواند نامعین گردد.

پس دستگاه زیر را در نظر بگیریم:
\[
\begin{array}{@{}r@{\quad}l@{}}
(1)& g_1(x),\ g_2(x),\ldots,g_\rho(x)
\end{array}
\]
که دارای رتبهٔ جبری $\rho$ --- با $\rho<n$ --- است، و $h(x)$ تابعی گویا باشد که به توابع (1) وابستهٔ جبری است. بنا بر قضایای شناخته‌شده، رتبهٔ ماتریس ژاکوبیِ دستگاه (1) برابر $\rho$ است. بنابراین، برای مثال، متغیرهای نامعین را چنان نام‌گذاری کنیم که
\[
\begin{array}{@{}r@{\quad}l@{}}
(2)& \displaystyle
\frac{\partial(g_1\ g_2\cdots g_\rho)}{\partial(x_1\ x_2\cdots x_\rho)}
=\frac{\Gamma(x)}{G(x)^{\rho+1}}\ne0,
\end{array}
\]
که در آن، چنان‌که می‌دانیم، $G(x)$ مخرج مشترک توابع (1) است.

\emph{اکنون عدد $a$ را چنان برگزینیم که حاصل‌ضرب $G(x)\cdot\Gamma(x)$ بر $(x_i-a)$ بخش‌پذیر نباشد}، که در آن $i$ یکی از اعداد $\rho+1,\rho+2,\ldots,n$ است. پس در $x_i=a$ مخرج مشترک توابع (1) نمی‌تواند صفر شود، و $\rho$ تابع
\[
\begin{array}{@{}r@{\quad}l@{}}
(3)& [g_1(x)]_{x_i=a},\ldots,[g_\rho(x)]_{x_i=a}
\end{array}
\]
نیز مستقل جبری‌اند. به تعبیر ما، رتبهٔ جبری دستگاه (1) با جایگذاری $x_i=a$ حفظ می‌شود.

معادلهٔ تحویل‌ناپذیری که $h(x)$ نسبت به توابع (1) ارضا می‌کند، به صورت زیر باشد:
\[
\begin{array}{@{}r@{\quad}l@{}}
(4)& A_0(g_1\cdots g_\rho)\cdot h(x)^\tau
 +A_1(g_1\cdots g_\rho)\cdot h(x)^{\tau-1}
 +\cdots+A_\tau(g_1\cdots g_\rho)=0,
\end{array}
\]
که در آن $A_0(g_1\cdots g_\rho)$ و $A_\tau(g_1\cdots g_\rho)$ به‌طور همانی صفر نمی‌شوند. برای آنکه از این رابطه به یک همانی میان چندجمله‌ای‌ها برسیم، می‌گذاریم:
\[
\begin{array}{@{}r@{\quad}l@{}}
(5)& g_i(x)=G_i(x):G(x);\qquad h(x)=H_1(x):H_2(x),
\end{array}
\]
و یادآور می‌شویم که --- به‌سبب فرض کاهش‌یافته بودن نمایش $h(x)$ --- $H_1(x)$ و $H_2(x)$ نسبت به هم اول‌اند. به‌موجب (5)، رابطهٔ (4) به صورت زیر درمی‌آید:
\[
\begin{array}{@{}r@{\quad}l@{}}
(6)& \begin{aligned}[t]
&B_0(G\,G_1\cdots G_\rho)\cdot G(x)^{\sigma_0}\cdot H_1(x)^\tau
+B_1(G\,G_1\cdots G_\rho)\cdot G(x)^{\sigma_1}\cdot H_1(x)^{\tau-1}H_2(x)\\
&\quad+\cdots+B_\tau(G\,G_1\cdots G_\rho)\cdot G(x)^{\sigma_\tau}\cdot H_2(x)^\tau=0,
\end{aligned}
\end{array}
\]
که در آن $B_k$ها فرم‌های همگن در $G,G_i$ هستند و دست‌کم یکی از توان‌های نامنفی $\sigma_k$ باید صفر باشد.

اکنون فرض کنیم $H_1(x)$ بر $(x_i-a)$ بخش‌پذیر باشد؛ آنگاه بنا بر (6)، عبارت
\[
  B_\tau(G\,G_1\cdots G_\rho)\cdot G(x)^{\sigma_\tau}\cdot H_2(x)^\tau
\]
نیز بر $(x_i-a)$ بخش‌پذیر است. این امر طبق فرض برای $G(x)$ منتفی است؛ بنابراین بخش‌پذیری $B_\tau$، بخش‌پذیریِ
\[
  A_\tau=B_\tau:G^\lambda
\]
را نیز در پی می‌داشت و میان توابع (3) رابطهٔ $A_\tau=0$ برقرار می‌شد، برخلاف استقلال جبری مفروض.
سرانجام $H_2(x)$ نیز چون با $H_1(x)$ نسبت به هم اول است، نمی‌تواند بر $(x_i-a)$ بخش‌پذیر باشد\srcfn{*)}{این استنتاج اخیر که بر کاهش‌یافته بودن نمایش تکیه دارد، در صورت انجام یک جایگذاری هم‌زمانِ $x_i=a$ و $x_k=b$ با حفظ سایر فرض‌ها دیگر ممکن نبود. در آن صورت ممکن بود $H_1(x)$ و $H_2(x)$ هم‌زمان صفر شوند و خارج‌قسمت بر اثر جایگذاری نامعین $=\frac00$ گردد؛ چنان‌که مثلاً
\[
 h(x)=\frac{x_3(\alpha x_1+\beta x_2)+x_4(\gamma x_1+\delta x_2)}
 {x_3(\alpha' x_1+\beta' x_2)+x_4(\gamma' x_1+\delta' x_2)}
\]
با جایگذاری‌های $x_3=0$ و $x_4=0$.}؛ پس فرض بخش‌پذیری $H_1(x)$ بر $(x_i-a)$ به تناقض می‌انجامد. به همین نحو نشان می‌دهیم که $H_2(x)$ نیز نمی‌تواند بر $(x_i-a)$ بخش‌پذیر باشد. بنابراین تابع $[h(x)]_{x_i=a}=k(x)$ نه می‌تواند به‌طور همانی صفر شود و نه می‌تواند نامعین گردد.

استدلال بالا را می‌توان دوباره بر دستگاه (3) و تابع $k(x)$ به کار برد، به شرط آنکه این تابع باز در نمایش کاهش‌یافته نوشته شده باشد؛ و ادامهٔ این روش سرانجام به لم زیر می‌انجامد:

\emph{لم: دو دستگاه}
\[
\begin{array}{c}
g_1(x),\ g_2(x),\ldots,\ g_\rho(x),\\
g_1(x),\ g_2(x),\ldots,\ g_\rho(x),\ h(x)
\end{array}
\]
\emph{هر یک دارای رتبهٔ جبری $\rho$ --- با $\rho<n$ --- باشند، و مثلاً}
\[
 \frac{\partial(g_1\cdots g_\rho)}
 {\partial(x_1\cdots x_\rho)}
 =\frac{\Gamma(x)}{G(x)^{\rho+1}}\ne0.
\]
\emph{اعداد}
\[
 a_n,\ a_{n-1},\ldots,\ a_{\rho+1}
\]
\emph{چنان انتخاب شده باشند که در اثر جایگذاری}
\[
 x_n=a_n,\quad x_{n-1}=a_{n-1},\ldots,\quad x_{\rho+1}=a_{\rho+1}
\]
\emph{حاصل‌ضرب $G(x)\cdot\Gamma(x)$ صفر نشود --- یعنی رتبهٔ جبری دستگاه (1) حفظ شود. در جایگذاری‌های پیاپی}
\[
\begin{array}{c}
[h(x)]_{x_n=a_n}=h^{(n-1)}(x);\quad
[h^{(n-1)}(x)]_{x_{n-1}=a_{n-1}}=h^{(n-2)}(x);\quad\cdots;\\[0.25em]
[h^{(\rho+1)}(x)]_{x_{\rho+1}=a_{\rho+1}}=h^*(x_1,x_2,\ldots,x_\rho)
\end{array}
\]
\emph{فرض کنیم هر یک از توابع $h(x),h^{(n-1)}(x),\ldots,h^{(\rho+1)}(x)$ در نمایش کاهش‌یافته نوشته شده باشد. تحت این فرض‌ها در تابع $h^*(x_1\cdots x_\rho)$ نه صورت و نه مخرج می‌تواند به‌طور همانی صفر شود، و تابع نیز نمی‌تواند $\frac00$ گردد.}\srcfn{**)}{جایگذاری پیاپی، برای مثال، در تابع $h(x)$ از یادداشت پیشین نتیجهٔ زیر را می‌دهد:
\[
 [h(x)]_{x_4=0}=h^{(3)}(x)=\frac{\alpha x_1+\beta x_2}{\alpha' x_1+\beta' x_2},\qquad
 h^*(x_1,x_2)=\frac{\alpha x_1+\beta x_2}{\alpha' x_1+\beta' x_2}.
\]}

\section*{\S\ 3. تناظر یک‌به‌یکِ دستگاه‌های $\Ssys_{n\rho}$ با دستگاه‌های $\Ssys^*_{\rho\rho}$.}

از لم بی‌درنگ تناظری میان دستگاه‌های $\Ssys_{n\rho}$ و دستگاه‌های $\Ssys^*_{\rho\rho}$ به دست خواهیم آورد، بدین‌گونه که برخی از متغیرهای نامعین با اعداد جایگزین شوند؛ به‌این‌ترتیب رتبهٔ جبری $\rho$ به‌منزلهٔ عددی که به‌واقع مشخصه است پدیدار می‌شود.

چون دستگاه $\Ssys_{n\rho}$ دارای رتبهٔ جبری $\rho$ است، می‌توان $\rho$ تابع از توابع $\Ssys_{n\rho}$ یافت، یعنی
\[
\begin{array}{@{}l@{\quad}c@{}}
(1)& g_1(x)\cdots g_\rho(x),
\end{array}
\]
که مستقل جبری‌اند. افزون بر این، اگر $f(x)$ تابعی دلخواه از $\Ssys_{n\rho}$ باشد، آنگاه دستگاه
\[
\begin{array}{@{}l@{\quad}c@{}}
(2)& g_1(x)\cdots g_\rho(x),\ f(x)
\end{array}
\]
نیز رتبهٔ جبری $\rho$ دارد؛ و دستگاه‌های (1) و(2) شرایط لم را برآورده می‌کنند.\srcfn{*)}{در صورت نیاز، متغیرهای نامعین دوباره شماره‌گذاری می‌شوند.}

پس اعداد
\[
 a_n,\ a_{n-1},\ldots,\ a_{\rho+1}
\]
را مطابق لم چنان تعیین می‌کنیم که رتبهٔ جبری $\rho$ دستگاه (1) بر اثر جایگذاری حفظ شود ــ و این کار به هر شمار دلخواه از شیوه‌ها ممکن است ــ؛ آنگاه در تابعی که با
\[
\begin{array}{c}
[f(x)]_{x_n=a_n}=f^{(n-1)}(x);\quad
[f^{(n-1)}(x)]_{x_{n-1}=a_{n-1}}=f^{(n-2)}(x);\quad\cdots;\\[0.25em]
[f^{(\rho+1)}(x)]_{x_{\rho+1}=a_{\rho+1}}=f^*(x_1\cdots x_\rho)
\end{array}
\]
تعریف می‌شود، یعنی در تابع
\[
\begin{array}{@{}l@{\quad}c@{}}
(3)& f^*(x_1\cdots x_\rho)
\end{array}
\]
نه صورت و نه مخرج می‌تواند به‌طور همانی صفر شود.

اکنون $f(x)$ همهٔ توابع $\Ssys_{n\rho}$ را، چه تعدادشان متناهی باشد و چه نامتناهی، طی کند. دستگاهی را در نظر می‌گیریم که از مجموعهٔ همهٔ توابع (3) تشکیل شده و خواص زیر را دارد:

1. رتبهٔ جبری آن $\rho$ است؛ زیرا توابع حاصل از (1) را در بر دارد:
\[
g_1^*(x_1\cdots x_\rho),\ldots,g_\rho^*(x_1\cdots x_\rho),
\]
که به سبب انتخاب $a_n,a_{n-1},\ldots,a_{\rho+1}$ مستقل جبری‌اند. ازاین‌رو این دستگاه را
\[
\Ssys^*_{\rho\rho}
\]
می‌نامیم.\srcfn{**)}{از $f^*(x)$ و$g^*(x)$ و\ldots{} از این پس همواره توابع تصویر، یعنی توابعی از $x_1\cdots x_\rho$، مراد است.}

2. تناظر میان توابع $f(x)$ و$f^*(x)$ بی‌هیچ استثنا یک‌به‌یک است.

زیرا اگر $f_1(x)$ و$f_2(x)$ به $\Ssys_{n\rho}$ تعلق داشته باشند، تفاضل
\[
d(x)=f_1(x)-f_2(x)
\]
نیز به دستگاه (1) وابستهٔ جبری است؛ همچنین داریم:
\[
d^*(x)=f_1^*(x)-f_2^*(x).
\]
چون $d(x)$ همراه با دستگاه (1) شرایط لم را برآورده می‌کند، از
\[
d(x)\ne0
\]
ضرورتاً نتیجه می‌شود:
\[
d^*(x)\ne0;
\]
یعنی توابع متفاوت از $\Ssys_{n\rho}$ با توابع متفاوت از $\Ssys^*_{\rho\rho}$ متناظرند.

3. از هر رابطه میان توابع $\Ssys^*_{\rho\rho}$:
\[
F\bigl(f_1^*(x),f_2^*(x),\ldots,f_k^*(x)\bigr)=0
\quad\text{همانی نسبت به }x_1\cdots x_\rho
\]
برای توابعی از $\Ssys_{n\rho}$ که به‌طور یکتا با آن‌ها متناظرند، رابطهٔ زیر نتیجه می‌شود:
\[
F\bigl(f_1(x),f_2(x),\ldots,f_k(x)\bigr)=0
\quad\text{همانی نسبت به }x_1\cdots x_n.
\]
زیرا همراه با $f_1(x),f_2(x),\ldots,f_k(x)$، هر ترکیب گویای این توابع نیز به دستگاه (1) وابستهٔ جبری است.

همچنین نتیجه می‌شود که رابطهٔ نخست زیر، رابطهٔ دوم را در پی دارد:
\[
\begin{array}{@{}l@{\quad}l@{}}
(4)& h(x)=F\bigl(f_1(x)\cdots f_k(x)\bigr):\\
& h^*(x)=F\bigl(f_1^*(x)\cdots f_k^*(x)\bigr).
\end{array}
\]
بنابر لم، از برقرار بودن شرط ناصفر نخست در نمایش زیر نتیجه می‌شود که شرط دوم نیز برقرار است:
\[
h(x)\ne0\quad\text{همچنین}\quad h^*(x)\ne0;
\]
و این همان چیزی بود که باید ثابت می‌شد.

در جمع‌بندی، قضیهٔ زیر را به دست می‌آوریم:

\emph{قضیهٔ I: به هر دستگاه $\Ssys_{n\rho}$ متشکل از تعداد متناهی یا نامتناهی تابع گویا می‌توان ــ به هر شمار دلخواه از شیوه‌ها ــ دستگاهی $\Ssys^*_{\rho\rho}$ را به‌طور یک‌به‌یک متناظر کرد، چنان‌که همهٔ روابط برقرار میان توابع $\Ssys^*_{\rho\rho}$ در $\Ssys_{n\rho}$ نیز حفظ شوند و برعکس. به‌ویژه، مطابق (4)، به هر میدان $\Kfield_{n\rho}$ نیز میدانی $\Kfield^*_{\rho\rho}$ متناظر می‌شود.}

از قضیهٔ I نتیجهٔ زیر را نیز، که همهٔ برهان‌های بعدی را ساده می‌کند، به دست می‌آوریم: \emph{خواص دستگاه‌های $\Ssys_{n\rho}$ که با تعداد متناهی یا نامتناهی رابطه میان توابع دستگاه مشخص می‌شوند، برای دستگاه‌های $\Ssys_{n\rho}$ برقرارند، به‌محض آنکه برای یک دستگاه تصویر $\Ssys^*_{\rho\rho}$ برقرار باشند.}

این نتیجه را به‌اختصار \glqq{}اصل انتقال\grqq{} می‌نامیم.

ساده‌ترین مثال این تناظر را نظریهٔ ناورداهای جبری به دست می‌دهد. در واقع همواره می‌توان با یک تبدیل خطی، برخی ضرایب فرم بنیادی را چنان عددی معین کرد که همهٔ روابط برقرار برای فرم خاص، در حالت کلی نیز حفظ شوند.

\section*{\S\ 4. پایهٔ گویای میدان‌های $\Kfield_{n\rho}$.}

در بررسی‌های بعدی که پیرامون مفهوم‌های پایه سامان یافته‌اند، در همه‌جا از \glqq{}اصل انتقال\grqq{} بند \S\ 3 استفاده می‌کنیم. نخست وجود پایه را برای میدان‌ها ثابت می‌کنیم، اما تعریف‌ها را به‌طور کلی می‌آوریم:

\emph{تعریف سوم: شمار متناهی‌ای از توابع $\Ssys_{n\rho}$ پایهٔ گویا نامیده می‌شود هرگاه هر تابع از $\Ssys_{n\rho}$ را بتوان با ضرایبی از $\Omega$ به‌صورت ترکیبی گویا از این شمار متناهی تابع نمایش داد.}

پایهٔ گویا باید $\rho$ تابع مستقل جبری در بر داشته باشد، زیرا رتبهٔ جبریِ دستگاهی که از یک پایهٔ گویا به دست می‌آید با رتبهٔ جبریِ خود پایهٔ گویا برابر است.

نخست نشان می‌دهیم که برهان وجود پایهٔ گویا اصل انتقال را می‌پذیرد. زیرا تعریف سوم را می‌توان چنین صورت‌بندی کرد:

توابع $\Ssys_{n\rho}$
\[
g_1(x),\ g_2(x),\ldots,g_k(x)
\]
درست در صورتی یک پایهٔ گویا تشکیل می‌دهند که رابطه‌های نامتناهی زیر برقرار باشند:
\[
\begin{array}{@{}l@{\quad}c@{}}
(1)& f(x)=\gamma\bigl(g_1(x)\cdots g_k(x)\bigr),
\end{array}
\]
که در آن $f(x)$ همهٔ توابع $\Ssys_{n\rho}$ را طی می‌کند و $\gamma$ تابعی گویاست که با $f$ تغییر می‌کند و $k$ آرگومان دارد.

پس اگر در دستگاه تصویر $\Ssys^*_{\rho\rho}$ پایه‌ای گویا با توابع زیر داده شده باشد،
\[
g_1^*(x),\ g_2^*(x),\ldots,g_k^*(x),
\]
که رابطه‌های زیر را نتیجه می‌دهد،
\[
f^*(x)=\gamma\bigl(g_1^*(x)\cdots g_k^*(x)\bigr)
\]
آنگاه بنا بر قضیهٔ نخست، رابطه‌های (1) برای $\Ssys_{n\rho}$ برقرارند؛ یعنی \emph{اصل انتقال برای پایهٔ گویا} چنین است:

\emph{از وجود پایهٔ گویا برای یک دستگاه تصویر $\Ssys^*_{\rho\rho}$، وجود پایهٔ گویا برای دستگاه اصلی $\Ssys_{n\rho}$ نتیجه می‌شود.}\srcfn{*)}{اصل انتقال برای هر پرسش مربوط به پایه که در اینجا بررسی می‌شود، به شیوهٔ متناظر صورت‌بندی می‌گردد.}

پس برای اثبات وجود پایهٔ گویا برای همهٔ میدان‌های $\Kfield_{n\rho}$ می‌توانیم خود را به میدان‌های $\Kfield^*_{\rho\rho}$ محدود کنیم؛ در اینجا تعمیم مستقیم شیوهٔ استدلالی که ا. اشتاینیتس به کار برده است ما را به مقصود می‌رساند.\srcfn{*)}{ا. اشتاینیتس، \srcLTR{Algebraische Theorie der Körper}، \S\ 24، \srcLTR{Crelles Journ. 137 (1909).}}

فرض کنیم
\[
\begin{array}{@{}l@{\quad}c@{}}
(2)& g_1^*(x_1\cdots x_\rho),\ldots,g_\rho^*(x_1\cdots x_\rho)
\end{array}
\]
دستگاهی از $\rho$ تابع مستقل جبری در $\Kfield^*_{\rho\rho}$ باشد. با حذف $x_1\cdots x_\rho$، $\rho$ رابطه به دست می‌آید:
\[
\begin{array}{@{}l@{\quad}c@{}}
(3)& F_1\bigl(g_1^*(x)\cdots g_\rho^*(x),x_1\bigr)=0,\ldots,
F_\rho\bigl(g_1^*(x)\cdots g_\rho^*(x),x_\rho\bigr)=0,
\end{array}
\]
به‌گونه‌ای که در هر رابطهٔ $F_i=0$، متغیر نامعین $x_i$ واقعاً ظاهر می‌شود؛ زیرا در غیر این صورت، برخلاف فرض، میان توابع (2) وابستگی جبری برقرار می‌بود. رابطه‌های (3) بیان می‌کنند که
\[
x_1,x_2,\ldots,x_\rho
\]
عناصری جبری نسبت به
\[
\Omega\bigl(g_1^*(x)\cdots g_\rho^*(x)\bigr)
\]
هستند، و میدانی که با الحاق $x_1\cdots x_\rho$ پدید می‌آید، یعنی
\[
\Omega\bigl(g_1^*\cdots g_\rho^*,x_1\cdots x_\rho\bigr)=\Omega(x_1\cdots x_\rho)
\]
یک توسیع جبریِ متناهی از $\Omega(g_1^*\cdots g_\rho^*)$ است. بنا بر قضیه‌ای شناخته‌شده،\srcfn{**)}{برای نمونه نک. \srcLTR{Weber, Lehrbuch d. Algebra 1}، \S\ 144 (چاپ نخست)، یا \S\ 55 (چاپ مختصر).} خود $\Kfield^*_{\rho\rho}$، چون میدان واسط میان $\Omega(g_1^*\cdots g_\rho^*)$ و$\Omega(x_1\cdots x_\rho)$ است، یک توسیع جبری از $\Omega(g_1^*\cdots g_\rho^*)$ است؛ درحالی‌که $\Omega(x_1\cdots x_\rho)$ توسیعی جبری از $\Kfield^*_{\rho\rho}$ است. پس $\Kfield^*_{\rho\rho}$ از $\Omega(g_1^*\cdots g_\rho^*)$ با الحاق شمار متناهی‌ای از عناصر $\Kfield^*_{\rho\rho}$ پدید می‌آید، یا نیز با الحاق \emph{یک} تابع که به‌گونه‌ای مناسب انتخاب شده باشد؛ یعنی
\[
\Kfield^*_{\rho\rho}
=\Omega(g_1^*\cdots g_\rho^*,g_{\rho+1}^*\cdots g_k^*)
=\Omega(g_1^*\cdots g_\rho^*g_0^*).
\]

پس به کمک اصل انتقال داریم:

\emph{قضیهٔ دوم: هر میدان $\Kfield_{n\rho}$ پایه‌ای گویا با رتبهٔ جبری $\rho$ دارد؛ به‌ویژه می‌توان پایهٔ گویا را چنان برگزید که حداکثر $\rho+1$ تابع داشته باشد.}

فرض کنیم دستگاهی دلخواه از $\rho$ تابع مستقل جبری در $\Kfield_{n\rho}$ با $h_1(x)\cdots h_\rho(x)$ داده شده باشد؛ به روش بالا می‌توان آن را به پایهٔ گویای زیر گسترش داد:
\[
h_1(x)\cdots h_\rho(x);\ h_{\rho+1}(x)\cdots h_\sigma(x)
\]
بنا بر ویژگی تعریفی، رابطه‌های زیر وجود دارند:
\[
\begin{array}{ll}
h_i(x)=\chi_i\bigl(g_1(x)\cdots g_k(x)\bigr)&(i=1,2,\ldots,\sigma),\\
g_j(x)=\psi_j\bigl(h_1(x)\cdots h_\sigma(x)\bigr)&(j=1,2,\ldots,k),
\end{array}
\]
که در آن $\chi_i$ و$\psi_j$ توابعی گویا از آرگومان‌های خودند؛ یعنی:

از هر پایهٔ گویای دلخواه، هر پایهٔ گویای دیگر با یک تبدیل دوگویا ــ یعنی تبدیلی گویا با وارونِ گویا ــ به دست می‌آید؛ و نیز:

اگر
\[
h_1(x)\cdots h_\rho(x)
\]
دستگاهی دلخواه از $\rho$ تابع مستقل جبری در $\Kfield_{n\rho}$ باشد، آنگاه $\Kfield_{n\rho}$ یک توسیع جبریِ متناهی از میدان زیر است:
\[
\Omega\bigl(h_1(x)\cdots h_\rho(x)\bigr).
\]

\section*{\S\ 5. فرم اینولوشن و پایهٔ اینولوشنِ میدان‌های $\Kfield_{n\rho}$.}

هر پایهٔ گویایی که در \S\ 4 ساخته شد این عیب را داشت که به توابع آغازینِ برگزیده‌شده به‌دلخواه وابسته بود؛ اکنون وجود \emph{پایهٔ گویایی را نشان می‌دهیم که خود میدان آن را به‌طور یکتا تعیین می‌کند}.

بنا بر \S\ 4، $\Omega(x_1\cdots x_\rho)$ میدانی جبری ــ مثلاً با درجهٔ $i$ ــ بر روی $\Kfield^*_{\rho\rho}$ است و، به‌سبب
\[
\Kfield^*_{\rho\rho}(x_1\cdots x_\rho)=\Omega(x_1\cdots x_\rho)
\]
با الحاق عناصری که نسبت به $\Kfield^*_{\rho\rho}$ جبری‌اند، یعنی $x_1\cdots x_\rho$، پدید می‌آید. پس چندجمله‌ای تحویل‌ناپذیر نسبت به $\Kfield^*_{\rho\rho}$، یعنی $\Phi^*(z,u)$، را در نظر می‌گیریم که ریشهٔ زیر را دارد:
\[
z=u_1x_1+u_2x_2+\cdots+u_\rho x_\rho=u(x),
\]
و درجه‌اش $i$ در $z$ است. افزون بر این، $\Phi^*(z,u)$، چون حاصل‌ضرب کمیت‌های مزدوج با $[z-u(x)]$ است، فرمی همگن از درجهٔ $i$ در $z,u_1\cdots u_\rho$ است؛ ازاین‌رو ضریب $z^i$ در آن تابعی مستقل از $u$ و متعلق به $\Kfield^*_{\rho\rho}$ است. اگر با تقسیم، این ضریبِ بالاترین درجه را برابر $1$ کنیم، چندجمله‌ای تحویل‌ناپذیر $\Phi^*(z,u)$ به‌طور یکتا تعیین می‌شود. این فرم همگن که بدین‌سان تعریف شده است،
\[
\begin{array}{@{}l@{\quad}l@{}}
(1)& \displaystyle \Phi^*(z,u)=z^i+\sum{}'g^*_{\alpha\alpha_1\ldots\alpha_\rho}(x_1\cdots x_\rho)\cdot
z^\alpha u_1^{\alpha_1}\cdots u_\rho^{\alpha_\rho}\\
& \displaystyle\qquad(\alpha+\alpha_1+\cdots+\alpha_\rho=i)
\end{array}
\]
را \emph{فرم اینولوشن}\srcfn{*)}{دلیل این نام‌گذاری در پایان این بند روشن می‌شود؛ $\sum'$ یعنی جمع باید بر همهٔ جمله‌ها به‌جز $z^i$ گسترده شود.} میدان $\Kfield^*_{\rho\rho}$ می‌نامیم؛ از (1) نیز
\[
\begin{array}{@{}l@{\quad}l@{}}
(2)& \displaystyle \Phi(z,u)=z^i+\sum{}'g_{\alpha\alpha_1\ldots\alpha_\rho}(x)\cdot
z^\alpha u_1^{\alpha_1}\cdots u_\rho^{\alpha_\rho}\\
& \displaystyle\qquad(\alpha+\alpha_1+\cdots+\alpha_\rho=i)
\end{array}
\]
به‌عنوان یک \emph{فرم اینولوشن} میدان $\Kfield_{n\rho}$ به دست می‌آید.



نشان خواهیم داد که ضرایب فرم اینولوشن $\Phi^*(z,u)$ یک پایهٔ گویای $\Kfield^*_{\rho\rho}$ را تشکیل می‌دهند؛ یعنی رابطهٔ زیر برقرار است:
\[
\begin{array}{@{}l@{\quad}c@{}}
(3)& \Omega\bigl(g^*_{\alpha\alpha_1\ldots\alpha_\rho}(x_1\cdots x_\rho)\bigr)=\Kfield^*_{\rho\rho}\qquad(\alpha+\alpha_1+\cdots+\alpha_\rho=i).
\end{array}
\]
برای اثبات، توجه می‌کنیم که $\Phi^*(z,u)$ ضرایبی از $\Omega(g^*_{\alpha\alpha_1\ldots\alpha_\rho}(x))$ دارد و، چون نسبت به $\Kfield^*_{\rho\rho}$ تحویل‌ناپذیر است، به طریق اولی نسبت به این زیرمیدان نیز تحویل‌ناپذیر است. پس $\Phi^*(z,u)$ معادلهٔ تحویل‌ناپذیر $u(x)$ نسبت به $\Omega(g^*_{\alpha\alpha_1\ldots\alpha_\rho}(x))$ را نمایش می‌دهد. در نتیجه $\Omega(x_1\cdots x_\rho)$ میدانی جبری ــ با درجهٔ $i$ ــ بر روی $\Omega(g^*_{\alpha\alpha_1\ldots\alpha_\rho}(x))$ است. چون $\Kfield^*_{\rho\rho}$ میان $\Omega(g^*_{\alpha\alpha_1\ldots\alpha_\rho}(x))$ و$\Omega(x_1\cdots x_\rho)$ قرار دارد، چنان‌که می‌دانیم، برابری درجه‌ها برابری میدان‌ها را نتیجه می‌دهد؛ پس رابطهٔ (3) برقرار است. بنا بر اصل انتقال، ضرایب $\Phi(z,u)$ نیز به‌طور متناظر پایه‌ای گویا برای $\Kfield_{n\rho}$ به دست می‌دهند؛ یعنی قضیهٔ زیر را داریم:

\emph{قضیهٔ سوم: ضرایب فرم اینولوشن پایه‌ای گویا تشکیل می‌دهند که میدان آن را تعیین می‌کند؛ این پایه، پایهٔ اینولوشن است. برای میدان‌های $\Kfield^*_{\rho\rho}$ این پایه به‌طور یکتا تعیین می‌شود.}\srcfn{*)}{برای میدان‌های $\Kfield_{n\rho}$ ممکن است فرم اینولوشن به نگاشت وابسته باشد؛ یکتایی را می‌توان با نگاشتی دارای ضرایب نامعین به دست آورد.}

اکنون تعبیری ناگویا از این واقعیت به دست می‌دهیم که پیوند آن را با نظریهٔ هندسی اینولوشن برقرار می‌کند.

فرض کنیم تجزیهٔ $\Phi^*(z,u)$ به عامل‌های خطی در یک میدان توسیع به صورت زیر داده شده باشد:
\[
\begin{array}{@{}l@{\quad}l@{}}
(4)& \Phi^*(z,u)=(z-u_1x_1-\cdots-u_\rho x_\rho)(z-u_1x_1^{(1)}-\cdots-u_\rho x_\rho^{(1)})\cdots\\[0.25em]
& \qquad\qquad\quad (z-u_1x_1^{(i-1)}-\cdots-u_\rho x_\rho^{(i-1)}).
\end{array}
\]
مقایسه با (1) نشان می‌دهد که توابع
\[
g^*_{\alpha\alpha_1\ldots\alpha_\rho}(x_1\cdots x_\rho)\qquad(\alpha+\alpha_1+\cdots+\alpha_\rho=i)
\]
همان توابع متقارنِ مقدماتیِ ردیف‌های کمیت زیرند:
\[
\begin{array}{@{}l@{\quad}cccc@{}}
(5)&x_1&x_2&\cdots&x_\rho\\
&x_1^{(1)}&x_2^{(1)}&\cdots&x_\rho^{(1)}\\
&\vdots&\vdots&&\vdots\\
&x_1^{(i-1)}&x_2^{(i-1)}&\cdots&x_\rho^{(i-1)}.
\end{array}
\]
پس هر تابع از $\Kfield^*_{\rho\rho}$ ترکیبی گویا از این توابع مقدماتی و در ردیف‌های کمیت (5) متقارن است؛ و برعکس، هر تابع متقارنِ ردیف‌های کمیت (5) به $\Kfield^*_{\rho\rho}$ تعلق دارد.

بنابراین اگر
\[
f^*(x)=\frac{F^*(x)}{H^*(x)}
\]
همهٔ توابع $\Kfield^*_{\rho\rho}$ را طی کند، دستگاهی از رابطه‌های نامتناهیِ زیر برقرار است:
\[
\begin{array}{@{}l@{\quad}c@{}}
(6)&\displaystyle \frac{F^*(x)}{H^*(x)}=\frac{F^*(x^{(1)})}{H^*(x^{(1)})}=\cdots=\frac{F^*(x^{(i-1)})}{H^*(x^{(i-1)})};
\end{array}
\]
یا به‌صورت فشرده:
\[
F^*(x^{(k)})H^*(x^{(l)})-F^*(x^{(l)})H^*(x^{(k)})=0\qquad(k,l=0,1,\ldots,(i-1));
\]
که در آن نه $F^*(x^{(k)})$ و نه $H^*(x^{(k)})$ به‌طور همانی صفر نمی‌شوند، زیرا به‌ترتیب با $F^*(x)$ و$H^*(x)$ مزدوج و به‌طور صوری متقارن‌اند.

برعکس، برای هر ردیف کمیت $\xi$ که رابطه‌های نامتناهیِ
\[
\begin{array}{@{}l@{\quad}c@{}}
(7)&F^*(x)H^*(\xi)-H^*(x)F^*(\xi)=0,\qquad H^*(\xi)\ne0
\end{array}
\]
را برآورده کند، تعلق به ردیف‌های کمیت (5) نتیجه می‌شود.

زیرا اگر $G^*(x)$ مخرج مشترک ضرایب $\Phi^*(z,u)$ باشد، بنا بر فرض (7)، تابع چندجمله‌ایِ زیر که در $\xi$ نیز چندجمله‌ای است، به‌طور همانی صفر نمی‌شود:
\[
G^*(\xi)\cdot\Phi^*(z,u;\xi)
\]
ــ یعنی ضرایب $\Phi^*(z,u;\xi)$ نامعین نمی‌شوند ــ و ریشهٔ زیر را دارد:
\[
z=u_1\xi_1+u_2\xi_2+\cdots+u_\rho\xi_\rho.
\]
پس به‌سبب (7)، $\Phi^*(z,u;x)$ نیز ریشهٔ $z=u(\xi)$ را دارد؛ یعنی $\xi$ به ردیف‌های کمیت (5) تعلق دارد. بنا بر (6)، حکم متناظر هنگامی نیز برقرار است که $\xi$ دستگاه معادلاتی را نسبت به یک ردیف کمیت مزدوج برآورده کند:
\[
F^*(x^{(k)})H^*(\xi)-H^*(x^{(k)})F^*(\xi)=0;\qquad H^*(\xi)\ne0,
\]
یعنی ردیف‌های کمیت (5) به‌صورت جواب‌های $\xi$ دستگاه معادلات زیر تعریف می‌شوند:
\[
F^*(x)H^*(\xi)-H^*(x)F^*(\xi)=0;\qquad H^*(\xi)\ne0,
\]
که در آن $\frac{F^*(x)}{H^*(x)}$ همهٔ توابع $\Kfield^*_{\rho\rho}$ را طی می‌کند؛ در اینجا می‌توان ردیف $x$ را نیز با هر ردیف مزدوجی جایگزین کرد.

از دید هندسی، اگر هر ردیف $x$ را یک نقطه تعبیر کنیم، این مطلب چنین است: جواب‌های (7) در یک فضای خطیِ $\rho$بعدی $\infty^\rho$ گروه‌هایی هریک شامل $i$ نقطه را نمایش می‌دهند، به‌گونه‌ای که هر گروه معین با هر نقطهٔ آن گروه به‌طور یکتا تعیین می‌شود. چنین دستگاهی از گروه‌های نقطه‌ای را \glqq{}\emph{اینولوشن در یک فضای خطی}\grqq{} می‌نامند.\srcfn{*)}{جواب‌های کنارگذاشته‌شدهٔ (7) که در آن‌ها $F^*(\xi)=0$ و$H^*(\xi)=0$ می‌شود ــ و همچنین جواب‌های کنارگذاشته‌شدهٔ دستگاه معادلات متناظر با پایهٔ اینولوشن ــ نقاط بنیادیِ اینولوشن نامیده می‌شوند؛ نقاطی که با همگن‌سازی به دست می‌آیند و \glqq{}در بی‌نهایت قرار دارند\grqq{} نیز باید در شمار آیند.}

به‌طور متناظر، یک میدان $\Kfield_{n\rho}$ اینولوشنی را در یک فضای خطی مشخص می‌کند که از منحنی‌ها، رویه‌ها و مانند این‌ها تشکیل شده است.

چون بنا بر مطالب بالا درجهٔ $i$ از $\Omega(x_1\cdots x_\rho)$ نسبت به $\Kfield^*_{\rho\rho}$ برابر شمار دستگاه‌های جواب (7) با شرط جانبی $H^*(\xi)\ne0$ است ــ که می‌توانیم آن‌ها را دستگاه‌های جواب $\Kfield^*_{\rho\rho}$ بنامیم ــ استدلالی که برای اثبات وجود پایهٔ اینولوشن به کار رفت صورت‌بندی ناگویای زیر را نیز می‌پذیرد:

\glqq{}اگر $\Lsys^*$ زیرمیدانی از $\Kfield^*_{\rho\rho}$ باشد و هر دستگاه جوابِ میدان $\Lsys^*$ دستگاه جوابِ $\Kfield^*_{\rho\rho}$ نیز باشد، آنگاه $\Lsys^*$ و$\Kfield^*_{\rho\rho}$ یکسان‌اند.\grqq{}

حکم متناظر، بنا بر اصل انتقال، برای زیرمیدان‌های $\Lsys$ از $\Kfield_{n\rho}$ برقرار است.

\section*{\S\ 6. پایهٔ کمینهٔ میدان‌های $\Kfield_{n\rho}$.}

این بند مروری بر قضیه‌های شناخته‌شده به دست می‌دهد، اما در صورت‌بندیِ بسط‌یافته‌ای که اصل انتقال و وجود پایهٔ گویا آن را ممکن ساخته‌اند.

\emph{تعریف چهارم: پایهٔ گویای $\Kfield_{n\rho}$ را پایهٔ کمینه می‌نامیم اگر فقط کمترین شمار ممکن، یعنی $\rho$ تابع را دربر داشته باشد؛ این توابع در آن صورت باید مستقلِ جبری باشند.}

از قضیه‌های لوروت، کاستلنووو و انریکس\srcfn{*)}{\srcLTR{J. Lüroth: Beweis eines Satzes über rationale Kurven, Math. Ann. 9 (1875). --- G. Castelnuovo: Sulla razionalità delle involuzioni piane, Math. Ann. 44 (1893). --- F. Enriques: Sopra una involuzione non razionale dello spazio, Rend. Acc. Linc, Vol. 21, 21. Jan. 1912.}} حکم‌های زیر نتیجه می‌شوند:

I. هر میدان $\Kfield_{n1}$ پایه‌ای کمینه دارد: تابع لوروتِ میدان که تا یک تبدیل خطی-کسری به‌طور یکتا معین است.

II. برای میدان‌های $\Kfield_{n2}$، اگر توسیع جبریِ حوزهٔ ضرایب $\Omega$ را مجاز بدانیم، همیشه پایه‌ای کمینه وجود دارد و عموماً فقط در همین صورت وجود دارد.

III. برای میدان‌های $\Kfield_{n\rho}$ با $\rho\ge3$ عموماً پایهٔ کمینه وجود ندارد. هنوز تلاشی برای مشخص‌کردن میدان‌های ویژهٔ $\Kfield_{n\rho}$ که پایهٔ کمینه دارند صورت نگرفته است.

اشاره‌ای کوتاه می‌کنیم به برهانی که ا. اشتاینیتس\srcfn{**)}{همان مأخذ، \S\ 24.} برای قضیهٔ لوروت در مورد میدان‌های $\Kfield^*_{11}$ داده است:

فرض کنیم $\Phi^*(z,u)$ فرم اینولوشنِ $\Kfield^*_{11}$ باشد و $\psi^*(x)$ ضریبی دلخواه باشد که متغیر $x$ واقعاً در آن ظاهر می‌شود؛ یعنی ضریبی از $\Phi^*(z,u)$. معلوم می‌شود درجهٔ $\Omega(x)$ نسبت به $\Omega(\psi^*(x))$ با درجهٔ $\Omega(x)$ نسبت به $\Kfield^*_{11}$ برابر است و از این برابری یکسانیِ دو میدان نتیجه می‌شود. افزون بر این، اگر $\Omega(\chi^*(x))$ برابر $\Omega(\psi^*(x))$ باشد، از وابستگیِ گویای دوسویهٔ $\chi^*(x)$ و$\psi^*(x)$، وابستگیِ خطی-کسریِ دو تابع نتیجه می‌شود.

پس بنابر اصل انتقال، قضیهٔ لوروت صورت‌بندیِ زیر را نیز می‌پذیرد:

\emph{پایهٔ کمینهٔ هر میدان $\Kfield_{n1}$ از یک تابع دلخواهِ پایهٔ اینولوشن تشکیل می‌شود؛ و هر دو تابع از پایهٔ اینولوشنِ $\Kfield_{n1}$ به‌طور خطی-کسری به یکدیگر وابسته‌اند.}

گ. کاستلنووو حکم II را با نظریهٔ هندسیِ اینولوشن برای میدان‌های $\Kfield^*_{22}$ ثابت می‌کند که از یک پایهٔ گویای سه‌تابعی به دست آمده‌اند. چون اصل انتقال در توسیع جبریِ متناهیِ حوزهٔ ضرایب $\Omega$ نیز برقرار می‌ماند، وجود پایهٔ گویا این برهان را برای همهٔ میدان‌های $\Kfield_{n2}$ کامل می‌کند.

در مورد III باید یادآور شد که انریکس اینولوشنی ناگویا در فضای سه‌بعدی ساخته است؛ یعنی میدانی $\Kfield_{33}$ که پایهٔ کمینه ندارد.

\section*{\S\ 7. دستگاه دلخواه $\Ssys_{n\rho}$، خانوادهٔ خطی $\Lsys_{n\rho}$ و دامنهٔ صحیح $\Jdom_{n\rho}$ از توابع گویا. وجود پایهٔ گویا.}

اکنون از وجود پایهٔ گویای میدان‌های $\Kfield_{n\rho}$، وجود پایهٔ گویا برای دستگاه‌های دلخواه از توابع گویا و گویای تمام (یعنی چندجمله‌ای‌ها) به دست خواهد آمد. این دستگاه‌ها را به ترتیبی سامان می‌دهیم که با عمل‌های ابتداییِ جمع، ضرب و تقسیم، یکی پس از دیگری از هم پدید می‌آیند.

I. مانند همیشه، $\Ssys_{n\rho}$ دستگاهی دلخواه از توابع گویا (یا گویای تمام) در $n$ متغیر نامعین، با رتبهٔ جبری $\rho$، باشد. کمیت‌های متعلق به $\Omega$ نیز از اعضای دستگاه شمرده می‌شوند.

II. دستگاه $\Lsys_{n\rho}$ را \emph{خانوادهٔ خطی} می‌نامیم هرگاه شرط‌های زیر را برآورده کند:

1. $\Lsys_{n\rho}$ همراه با $f(x)$، همواره $c\cdot f(x)$ را نیز دربر دارد، که در آن $c$ کمیتی دلخواه از $\Omega$ است؛

2. $\Lsys_{n\rho}$ همراه با $f(x)$ و $g(x)$، همواره $f(x)+g(x)$ را نیز دربر دارد.

III. خانوادهٔ خطی $\Jdom_{n\rho}$ را، با شرط افزودهٔ زیر، \emph{دامنهٔ صحیح} می‌نامیم:

3. $\Jdom_{n\rho}$ همراه با $f(x)$ و $g(x)$، همواره $f(x)\cdot g(x)$ را نیز دربر دارد.

IV. دامنهٔ صحیح $\Kfield_{n\rho}$ را، با شرط افزودهٔ زیر، \emph{میدان} می‌نامیم:

4. $\Kfield_{n\rho}$ همراه با $f(x)$ و $g(x)$ ــ که $g(x)\ne0$ ــ همواره خارج‌قسمت $f(x):g(x)$ را نیز دربر دارد.

شرط‌های 1 تا 4 همان شرط‌هایی‌اند که در \S\ 1 برای میدان بیان شدند.\srcfn{*)}{$f(x)$ و $g(x)$ می‌توانند از درجهٔ صفر نیز باشند، یعنی کمیت‌هایی از $\Omega$.}

اکنون این ساختارها را یکی پس از دیگری از هم به دست می‌آوریم.

a) یک دستگاه دلخواه $\Ssys_{n\rho}$ با شرط زیر به خانوادهٔ خطی $\Lsys_{n\rho}$ گسترش می‌یابد:

$\Lsys_{n\rho}$ افزون بر $\Ssys_{n\rho}$، همهٔ توابع $h(x)$ و فقط آن‌هایی را دربر دارد که بتوان آن‌ها را با ضرایبی از $\Omega$، به صورت ترکیبی خطی از شمار متناهی‌ای از توابع $\Ssys_{n\rho}$ نمایش داد:
\[
h(x)=c_1f_1(x)+c_2f_2(x)+\cdots+c_kf_k(x),
\]
که در آن $f_1(x)\cdots f_k(x)$ در میان همهٔ توابع $\Ssys_{n\rho}$ و $c_1\cdots c_k$ در میان همهٔ کمیت‌های $\Omega$ تغییر می‌کنند.

زیرا افزودن ترکیب‌های گویایی از توابع دستگاه، رتبهٔ جبریِ دستگاه را تغییر نمی‌دهد. $\Lsys_{n\rho}$ شرط‌های 1 و 2 را برآورده می‌کند و بنابراین خانواده‌ای خطی است. افزون بر این، چون هر خانوادهٔ خطیِ دربرگیرندهٔ $\Ssys_{n\rho}$، بنابر 1 و 2، $\Lsys_{n\rho}$ را نیز دربر می‌گیرد، $\Lsys_{n\rho}$ کوچک‌ترین خانوادهٔ خطیِ دربرگیرندهٔ $\Ssys_{n\rho}$ است.

b) به همین ترتیب، از یک خانوادهٔ خطی $\Lsys_{n\rho}$، کوچک‌ترین دامنهٔ صحیحِ دربرگیرندهٔ آن، یعنی $\Jdom_{n\rho}$، با شرط زیر پدید می‌آید:

$\Jdom_{n\rho}$ افزون بر $\Lsys_{n\rho}$، همهٔ توابع $h(x)$ و فقط آن‌هایی را دربر دارد که بتوان آن‌ها را به صورت \emph{ترکیبی گویای تمام} ــ یعنی ترکیبی چندجمله‌ای ــ با ضرایبی از $\Omega$، از شمار متناهی‌ای از توابع $\Lsys_{n\rho}$، یعنی $f_1(x)\cdots f_k(x)$، نمایش داد؛ در اینجا $f_1(x)\cdots f_k(x)$ در میان همهٔ توابع $\Lsys_{n\rho}$ تغییر می‌کنند.

c) سرانجام، از یک دامنهٔ صحیح $\Jdom_{n\rho}$، کوچک‌ترین میدانِ دربرگیرندهٔ آن، یعنی $\Kfield_{n\rho}$، با شرط زیر پدید می‌آید: $\Kfield_{n\rho}$ افزون بر $\Jdom_{n\rho}$، همهٔ توابع $h(x)$ و فقط آن‌هایی را دربر دارد که بتوان آن‌ها را به صورت خارج‌قسمت دو تابع از $\Jdom_{n\rho}$ نمایش داد.

اگر این سه گسترش را در یک گام گرد آوریم، نتیجه چنین است:

\emph{هر دستگاه $\Ssys_{n\rho}$ را می‌توان با شرط زیر به کوچک‌ترین میدانِ دربرگیرندهٔ آن، یعنی $\Kfield_{n\rho}$، گسترش داد:}

\emph{$\Kfield_{n\rho}$ افزون بر $\Ssys_{n\rho}$، همهٔ توابع $h(x)$ و فقط آن‌هایی را دربر دارد که بتوان آن‌ها را به صورت ترکیبی گویا ــ با ضرایبی از $\Omega$ ــ از شمار متناهی‌ای از توابع $\Ssys_{n\rho}$، یعنی $f_1(x)\cdots f_k(x)$، نمایش داد؛ در اینجا $f_1(x)\cdots f_k(x)$ در میان همهٔ توابع $\Ssys_{n\rho}$ تغییر می‌کنند.}

از حقیقت اخیر، وجود پایهٔ گویا برای ساختارهای دلخواه بی‌درنگ نتیجه می‌شود:

$\Kfield_{n\rho}$ کوچک‌ترین میدانِ دربرگیرندهٔ $\Ssys_{n\rho}$ باشد و یک پایهٔ گویای $\Kfield_{n\rho}$ با توابع زیر داده شده باشد:
\[
\begin{array}{@{}l@{\quad}c@{}}
(1)&h_1(x),\ h_2(x),\ldots,h_\sigma(x).
\end{array}
\]
بنابر تعریف $\Kfield_{n\rho}$، رابطه‌های زیر وجود دارند:
\[
\begin{array}{@{}l@{\quad}c@{}}
(2)&h_1(x)=r_1(g_1(x)\cdots g_\lambda(x));\quad\ldots;\quad h_\sigma(x)=r_\sigma(g_\mu(x)\cdots g_\nu(x)),
\end{array}
\]
که در آن $r_1\cdots r_\sigma$ توابعی گویا در آرگومان‌های خود هستند، و
\[
\begin{array}{@{}l@{\quad}c@{}}
(3)&g_1(x)\cdots g_\lambda(x),\ldots,g_\mu(x)\cdots g_\nu(x)
\end{array}
\]
همگی به $\Ssys_{n\rho}$ تعلق دارند. به سبب رابطه‌های (2)، (3) نیز، همانند (1)، یک پایهٔ گویای $\Kfield_{n\rho}$ است؛ و چون توابع (3) به $\Ssys_{n\rho}$ تعلق دارند، هم‌زمان پایهٔ گویای $\Ssys_{n\rho}$ نیز هست. پس داریم:

\emph{قضیهٔ چهارم: هر دستگاه دلخواه $\Ssys_{n\rho}$ از توابع گویا (یا گویای تمام) دارای پایه‌ای گویاست که از شمار متناهی‌ای از توابع همان دستگاه تشکیل شده و رتبهٔ جبریِ آن $\rho$ است.}

اکنون، به‌طور ویژه، $\Lsys_{n\rho}$ خانواده‌ای خطی و
\[
\begin{array}{@{}l@{\quad}c@{}}
(4)&g_1(x)\cdots g_\rho(x);\ g_{\rho+1}(x)\cdots g_\nu(x)
\end{array}
\]
یک پایهٔ گویای $\Lsys_{n\rho}$ باشد. فرض کنیم، برای نمونه،
\[
g_1(x)\cdots g_\rho(x)
\]
از نظر جبری مستقل باشند. آنگاه اگر بگذاریم
\[
g(x)=c_1g_{\rho+1}(x)+c_2g_{\rho+2}(x)+\cdots+c_{\nu-\rho}g_\nu(x),
\]
می‌توان $c_i$ها را چنان به‌صورت اعدادی از $\Omega$ برگزید که
\[
\begin{array}{@{}l@{\quad}c@{}}
(5)&g_1(x)\cdots g_\rho(x),\ g(x)
\end{array}
\]
یک پایهٔ گویای کوچک‌ترین میدانِ دربرگیرنده، یعنی $\Kfield_{n\rho}$، باشد. اما چون $\Lsys_{n\rho}$ خانواده‌ای خطی است، $g(x)$ به $\Lsys_{n\rho}$ تعلق دارد و (5) یک پایهٔ گویای $\Lsys_{n\rho}$ است؛ یعنی:

\emph{قضیهٔ پنجم: پایهٔ گویای هر خانوادهٔ خطی $\Lsys_{n\rho}$ از توابع گویا (یا گویای تمام)، و در نتیجه به‌ویژه پایهٔ گویای هر دامنهٔ صحیح $\Jdom_{n\rho}$ و هر میدان $\Kfield_{n\rho}$، را می‌توان چنان برگزید که حداکثر $\rho+1$ تابع (و حداقل $\rho$ تابع) داشته باشد.}

پس کران شمار توابعِ پایهٔ گویا که در \S\ 4 برای میدان به دست آمد، بر این ویژگیِ میدان تکیه دارد که خانواده‌ای خطی است؛ در حالی که محدودشدن به $\rho$ تابع در صورت وجود پایهٔ کمینه، همهٔ فرض‌های میدان را به کار می‌گیرد.

همچنین یادآور می‌شویم که بنابر \S\ 3 (4)، همهٔ ویژگی‌های شاخصِ دستگاه‌ها و کوچک‌ترین ساختارهای دربرگیرندهٔ آن‌ها تحت نگاشت حفظ می‌شوند.


\section*{\S\ 8. پایهٔ اینولوشنِ دامنه‌های صحیح $\Jdom_{n\rho}$ ساخته از چندجمله‌ای‌ها.}

در بندهای بعدی با دستگاه‌های $\Ssys_{n\rho}$ سروکار داریم که اعضایشان توابع گویای تمام (چندجمله‌ای‌ها) هستند؛ نخست دامنه‌های صحیح $\Jdom_{n\rho}$ ساخته از چندجمله‌ای‌ها را بررسی می‌کنیم. ابتدا باید مفاهیم بخش‌پذیری را برای این دامنه‌های $\Jdom_{n\rho}$ تثبیت کنیم.

$F(x),G(x)$ دو چندجمله‌ای از $\Jdom_{n\rho}$ و $L(x)$ مقسوم‌علیه مشترکی برای آن‌ها باشند، به‌طوری‌که:
\[
F(x)=L(x)\cdot F_1(x);\qquad G(x)=L(x)\cdot G_1(x).
\]
$L(x)$ آنگاه و تنها آنگاه \emph{مقسوم‌علیه مشترک در $\Jdom_{n\rho}$} نامیده می‌شود که هر سه چندجمله‌ای $L(x),F_1(x),G_1(x)$ به $\Jdom_{n\rho}$ تعلق داشته باشند.

دو چندجمله‌ای از $\Jdom_{n\rho}$ را \emph{هم‌اول در $\Jdom_{n\rho}$} می‌نامیم هرگاه در $\Jdom_{n\rho}$ هیچ مقسوم‌علیه مشترکِ ناواحدی نداشته باشند.

$\Kfield_{n\rho}$ کوچک‌ترین میدانِ دربرگیرندهٔ $\Jdom_{n\rho}$ باشد. آنگاه از تعریف، برای هر تابع از $\Kfield_{n\rho}$، نمایشی به صورت زیر به دست می‌آید:
\[
f(x)=F_1(x):F_2(x),
\]
که در آن $F_1(x),F_2(x)$ به $\Jdom_{n\rho}$ تعلق دارند و در $\Jdom_{n\rho}$ هم‌اول‌اند. به همین ترتیب، چند تابع از $\Kfield_{n\rho}$ را می‌توان در $\Jdom_{n\rho}$ به مخرج مشترک برد:
\[
f_1(x)=\frac{F_1(x)}{F(x)}\cdots f_k(x)=\frac{F_k(x)}{F(x)}.
\]
$F(x)$ آنگاه و تنها آنگاه \emph{یک مخرج مشترکِ کمینه در $\Jdom_{n\rho}$} نامیده می‌شود که چندجمله‌ای‌های متعلق به $\Jdom_{n\rho}$، یعنی $F(x),F_1(x),\ldots,F_k(x)$، در $\Jdom_{n\rho}$ هم‌اول باشند.

چنین نمایشی ممکن است به شیوه‌های گوناگون امکان‌پذیر باشد، زیرا قوانین تجزیهٔ یکتا به توابع تحویل‌ناپذیر عموماً در $\Jdom_{n\rho}$ برقرار نیستند.

اکنون معنای پایهٔ اینولوشنِ $\Kfield_{n\rho}$ را برای $\Jdom_{n\rho}$ بررسی می‌کنیم؛ برای این کار، مفاهیم فرم اینولوشن و پایهٔ اینولوشن را به دامنه‌های صحیح انتقال می‌دهیم. فرض کنیم
\[
\begin{gathered}
\Phi(z,u)=z^i+\sum{}' g_{\alpha\alpha_1\ldots\alpha_\rho}(x)z^\alpha u_1^{\alpha_1}\cdots u_\rho^{\alpha_\rho}\\
(\alpha+\alpha_1+\cdots+\alpha_\rho=i)
\end{gathered}
\]
یک فرم اینولوشنِ $\Kfield_{n\rho}$ و $G(x)$ یک مخرج مشترکِ کمینه در $\Jdom_{n\rho}$ برای ضرایب $\Phi(z,u)$ باشد. با ضرب در $G(x)$، $\Phi(z,u)$ به فرم $\Psi(z,u)$ تبدیل می‌شود که ضرایبش چندجمله‌ای‌هایی از $\Jdom_{n\rho}$ و در $\Jdom_{n\rho}$ هم‌اول‌اند:
\[
\begin{gathered}
G(x)\cdot\Phi(z,u)=\Psi(z,u)=G(x)z^i+\sum{}'G_{\alpha\alpha_1\ldots\alpha_\rho}(x)z^\alpha u_1^{\alpha_1}\cdots u_\rho^{\alpha_\rho}\\
(\alpha+\alpha_1+\cdots+\alpha_\rho=i).
\end{gathered}
\]

\emph{هر فرم $\Psi(z,u)$ که از یک فرم اینولوشنِ $\Phi(z,u)$ متعلق به $\Kfield_{n\rho}$ با ضرب در چندجمله‌ای‌ای از $\Jdom_{n\rho}$ به‌گونه‌ای پدید آید که ضرایب $\Psi(z,u)$ چندجمله‌ای‌هایی از $\Jdom_{n\rho}$ و در $\Jdom_{n\rho}$ هم‌اول باشند، فرم اینولوشنِ $\Jdom_{n\rho}$ نامیده می‌شود؛ و ضرایب $\Psi(z,u)$ یک پایهٔ اینولوشنِ متناظر نامیده می‌شوند.}

نخست، از معنای پایهٔ اینولوشنِ $\Kfield_{n\rho}$ روشن است که هر پایهٔ اینولوشنِ $\Jdom_{n\rho}$ هم‌زمان یک پایهٔ گویاست. برای بررسی بیشتر، دامنهٔ تصویر $\Jdom^*_{\rho\rho}$ را در نظر می‌گیریم که برای کوچک‌ترین میدانِ دربرگیرندهٔ آن، $\Kfield^*_{\rho\rho}$، معادلهٔ تحویل‌ناپذیرِ دارای ریشهٔ
\[
z=u_1x_1+\cdots+u_\rho x_\rho
\]
با $\Phi^*(z,u)=0$ داده شده است.

بنا بر \S\ 5، ضرایب
\[
\begin{gathered}
\Phi^*(z,u)=z^i+\sum{}'g^*_{\alpha\alpha_1\ldots\alpha_\rho}(x)z^\alpha u_1^{\alpha_1}\cdots u_\rho^{\alpha_\rho}\\
(\alpha+\alpha_1+\cdots+\alpha_\rho=i)
\end{gathered}
\]
توابع متقارنِ مقدماتیِ ردیف‌های کمیت زیرند که نسبت به $\Kfield^*_{\rho\rho}$ مزدوج‌اند:
\[
x,\ x^{(1)}\cdots x^{(i-1)}
\]
$F^*(x)$ چندجمله‌ای دلخواهی از $\Kfield^*_{\rho\rho}$ باشد؛ آنگاه از
\[
F^*(x)=\frac1i\left(F^*(x)+F^*(x^{(1)})+\cdots+F^*(x^{(i-1)})\right),
\]
نتیجه می‌شود که $F^*(x)$ تابعی متقارن و تمام، یعنی چندجمله‌ای، در این ردیف‌های کمیت است؛ و بنابراین می‌توان آن را به صورت ترکیبی گویای تمام ــ یعنی ترکیبی چندجمله‌ای ــ از توابع متقارن مقدماتی، یعنی توابع
\[
g^*_{\alpha\alpha_1\ldots\alpha_\rho}(x)
\]
بیان کرد. اکنون اگر
\[
\begin{gathered}
G^*(x)\cdot\Phi^*(z,u)=\Psi^*(z,u)=G^*(x)z^i+
\sum{}'G^*_{\alpha\alpha_1\ldots\alpha_\rho}(x)z^\alpha u_1^{\alpha_1}\cdots u_\rho^{\alpha_\rho}\\
(\alpha+\alpha_1+\cdots+\alpha_\rho=i)
\end{gathered}
\]
یک فرم اینولوشنِ $\Jdom^*_{\rho\rho}$ باشد، آنگاه هر چندجمله‌ای از $\Kfield^*_{\rho\rho}$، و به‌ویژه هر چندجمله‌ای از $\Jdom^*_{\rho\rho}$، ترکیبی گویای تمام از خارج‌قسمت‌های زیر خواهد بود:
\[
G^*_{\alpha\alpha_1\ldots\alpha_\rho}(x):G^*(x).
\]

با در نظر گرفتن اصل انتقال، نتیجهٔ زیر به دست می‌آید:

\emph{قضیهٔ ششم: برای هر دامنهٔ صحیحِ ساخته از چندجمله‌ای‌ها، $\Jdom_{n\rho}$، دست‌کم یک پایهٔ گویای ممتاز، یعنی پایهٔ اینولوشن، وجود دارد؛ به‌گونه‌ای که در نمایش گویای همهٔ چندجمله‌ای‌های $\Jdom_{n\rho}$ با توابع پایهٔ اینولوشن، تنها توانی از یک تابعِ معین (ضریبِ بالاترین درجهٔ فرم اینولوشنِ متناظر) در مخرج ظاهر می‌شود.}

\section*{\S\ 9. دامنه‌های تمامِ نسبیِ $\Gdom_{n\rho}$ ساخته از چندجمله‌ای‌ها.}

در ادامه، دامنه‌های صحیحِ ویژه‌ای ساخته از چندجمله‌ای‌ها را بررسی می‌کنیم که مرحله‌ای میانی میان دامنهٔ صحیح و میدان‌اند.

یک دامنهٔ صحیح ساخته از چندجمله‌ای‌ها، $\Gdom_{n\rho}$، به شرط زیر «دامنهٔ تمامِ نسبی» نامیده می‌شود:

\emph{4الف) $\Gdom_{n\rho}$ افزون بر $F(x)$ و$G(x)$ --- که در آن $G(x)\ne0$ ---، خارج‌قسمت $F(x):G(x)$ را همواره اگر و تنها اگر دربر دارد که این خارج‌قسمت در متغیرهای نامعین تمام، یعنی یک چندجمله‌ای، باشد.}

نخست باید همانیِ دامنه‌های تمامِ نسبی با مجموعهٔ همهٔ چندجمله‌ای‌های یک میدان $\Kfield_{n\sigma}$ ($\sigma\ge \rho$) از توابع گویا اثبات شود.

فرض کنید $\Kfield_{n\rho}$ کوچک‌ترین میدان دربرگیرندهٔ $\Gdom_{n\rho}$ باشد. هر تابع $f(x)$ از $\Kfield_{n\rho}$ نمایشی به صورت خارج‌قسمت دو چندجمله‌ای از $\Gdom_{n\rho}$ دارد:
\[
f(x)=F_1(x):F_2(x),
\]
که در آن چندجمله‌ای‌های درجهٔ صفر نیز مجازند. هرگاه $f(x)$ یک چندجمله‌ای باشد، بنا بر شرط 4الف به $\Gdom_{n\rho}$ تعلق دارد؛ و چون $\Gdom_{n\rho}$ نمی‌تواند هیچ چندجمله‌ای دیگری را نیز دربر داشته باشد، \emph{هر دامنهٔ تمامِ نسبیِ $\Gdom_{n\rho}$ با مجموعهٔ همهٔ چندجمله‌ای‌های کوچک‌ترین میدان دربرگیرندهٔ خود یکسان است.}

برعکس، فرض کنید $\Kfield_{n\sigma}$ یک میدان دلخواه از توابع گویا باشد (پس نه کوچک‌ترین میدان دربرگیرندهٔ دستگاهی داده‌شده) و $\mathfrak J$ مجموعهٔ همهٔ چندجمله‌ای‌های موجود در $\Kfield_{n\sigma}$ باشد. $\mathfrak J$ شرط‌های 1، 2 و 3 دامنهٔ صحیح و شرط 4الف را برآورده می‌کند و بنابراین یک دامنهٔ تمامِ نسبی است؛ یعنی \emph{مجموعهٔ همهٔ چندجمله‌ای‌های موجود در یک میدان $\Kfield_{n\sigma}$ یک دامنهٔ تمامِ نسبیِ $\Gdom_{n\rho}$ ($\rho\le\sigma$) تشکیل می‌دهد.}\srcfn{*)}{ممکن است $\rho=0$ نیز باشد؛ یعنی مجموعهٔ چندجمله‌ای‌های $\Kfield_{n\sigma}$ فقط از حوزهٔ ضرایب $\Omega$ تشکیل شود.}

همچنین همانیِ دامنه‌های تمامِ نسبی را با «دامنه‌های صحیحِ ساخته از توابع تمامِ نسبی» که هیلبرت معرفی کرده است نشان می‌دهیم.\srcfn{**)}{د. هیلبرت: مسائل ریاضی. سخنرانی پاریس 1900. مسئلهٔ 14 (اخبار گوتینگن 1900).}

فرض کنید
\[
\begin{array}{@{}l@{\quad}c@{}}
(1)& G_1(x)\cdots G_k(x)
\end{array}
\]
یک پایهٔ گویای $\Gdom_{n\rho}$ باشد؛ بنا بر شرط‌های 1، 2، 3 و 4الف، هر ترکیب گویای چندجمله‌ای‌های (1) که در متغیرهای نامعین تمام، یعنی یک چندجمله‌ای، باشد به $\Gdom_{n\rho}$ تعلق دارد. برعکس، چون طبق مفهوم پایهٔ گویا هر چندجمله‌ای از $\Gdom_{n\rho}$ را می‌توان به‌طور گویا با چندجمله‌ای‌های (1) بیان کرد، $\Gdom_{n\rho}$ \emph{با مجموعهٔ آن ترکیب‌های گویای چندجمله‌ای‌های (1) که در متغیرهای نامعین تمام، یعنی چندجمله‌ای، هستند یکسان است.} این همان تعریف هیلبرت است که از پایهٔ گویای ویژه آغاز می‌کند، در حالی که تعریف ما تنها خواص مجرد دامنه را به کار می‌برد.

برای دامنه‌های تمامِ نسبی، تعریف مقسوم‌علیه مشترک که در \S\ 8 برای دامنه‌های صحیحِ عمومی داده شد ساده‌تر می‌شود:

فرض کنید $F(x),G(x)$ دو چندجمله‌ای از $\Gdom_{n\rho}$ و $L(x)$ یک مقسوم‌علیه مشترک این دو چندجمله‌ای باشد. آنگاه $L(x)$ اگر و تنها اگر یک \emph{مقسوم‌علیه مشترک در $\Gdom_{n\rho}$} نامیده می‌شود که $L(x)$ به $\Gdom_{n\rho}$ تعلق داشته باشد.

زیرا در این صورت تعلق خارج‌قسمت‌های $F(x):L(x)$ و$G(x):L(x)$ از شرط 4الف نتیجه می‌شود.

یکی دیگر از ویژگی‌های اساسی دامنه‌های تمامِ نسبی این است که مفهوم «تمامِ جبری نسبت به $\Gdom_{n\rho}$» را می‌توان دقیقاً همان‌گونه با یک معادلهٔ دلخواه تعریف کرد که برای میدان‌های اعداد جبری یا میدان $\Omega(x_1\cdots x_n)$ شناخته شده است.

\emph{تعریف V: کمیتی مانند $\xi$ تمامِ جبری نسبت به $\Gdom_{n\rho}$ --- یا نیز تمامِ جبریِ نسبی --- نامیده می‌شود اگر در معادله‌ای صدق کند که ضریب بالاترین درجهٔ آن واحد و دیگر ضرایبش چندجمله‌ای‌هایی از $\Gdom_{n\rho}$ باشند.}

فرض کنید کمیت تمامِ جبریِ نسبیِ $\xi$، برای نمونه، در معادلهٔ
\[
L(z)=z^\lambda+G_1(x)z^{\lambda-1}+\cdots+G_\lambda(x)=0
\]
صدق کند. $L(z)$ بر تابع تمامِ تحویل‌ناپذیر نسبت به کوچک‌ترین میدان دربرگیرندهٔ $\Kfield_{n\rho}$، یعنی
\[
M(z)=z^\mu+H_1(x)z^{\mu-1}+\cdots+H_\mu(x)
\]
بخش‌پذیر است. اما از این بخش‌پذیری، بنا بر تعمیم شناخته‌شدهٔ قضیهٔ گاوس\srcfn{*)}{برای نمونه بنگرید به ی. کونیگ: درآمدی بر نظریهٔ عمومی کمیت‌های جبری (لایپزیگ، ب. گ. تویبنر، 1903)، فصل نهم، \S\ 2، یا وبر (ویرایش کوچک)، \S\ 20.}، نتیجه می‌شود که توابع گویای $H_1(x)\cdots H_\mu(x)$ چندجمله‌ای‌هایی از $\Gdom_{n\rho}$ هستند و بنابراین به $\Gdom_{n\rho}$ تعلق دارند. بدین‌ترتیب تعریف کمیت تمامِ جبریِ نسبیِ $\xi$ از معادلهٔ تحویل‌ناپذیر به دست می‌آید. به‌ویژه، هر چندجمله‌ای $F(x)$ از $\Gdom_{n\rho}$ نیز نسبت به $\Gdom_{n\rho}$ تمامِ جبری است، زیرا در معادلهٔ $z-F(x)=0$ صدق می‌کند.

همچنین یادآور می‌شویم که برای یک دامنهٔ صحیحِ چندجمله‌ایِ $\Jdom_{n\rho}$، همانند کوچک‌ترین میدان دربرگیرنده، کوچک‌ترین دامنهٔ تمامِ نسبیِ دربرگیرندهٔ $\Gdom_{n\rho}$ نیز با شرط زیر تعریف می‌شود:

$\Gdom_{n\rho}$ افزون بر $\Jdom_{n\rho}$، همهٔ چندجمله‌ای‌های $H(x)$ و فقط آن‌هایی را دربر دارد که بتوان آن‌ها را به صورت خارج‌قسمت دو چندجمله‌ای از $\Jdom_{n\rho}$ نمایش داد.

از این تعریف همچنین نتیجه می‌شود که از هر فرم اینولوشن و پایهٔ اینولوشنِ $\Jdom_{n\rho}$، فرم و پایهٔ متناظری برای $\Gdom_{n\rho}$ به دست می‌آید؛ کافی است حداکثر یک مقسوم‌علیه مشترک در $\Gdom_{n\rho}$ از همهٔ چندجمله‌ای‌های پایه جدا شود.

زیرا $\Jdom_{n\rho}$ و$\Gdom_{n\rho}$ کوچک‌ترین میدان دربرگیرندهٔ یکسانِ $\Kfield_{n\rho}$ را دارند؛ و هر فرم اینولوشنِ $\Jdom_{n\rho}$ از ضرب یک فرم اینولوشنِ $\Kfield_{n\rho}$ در چندجمله‌ای‌ای از $\Jdom_{n\rho}$ به دست می‌آید؛ بنابراین ضرایبی از $\Gdom_{n\rho}$ دارد که بااین‌حال ممکن است در $\Gdom_{n\rho}$ مقسوم‌علیه مشترکی داشته باشند.\srcfn{*)}{برای نمونه، دامنهٔ صحیحی که از همهٔ ترکیب‌های گویای تمامِ $x_1^2$ و$x_1x_2$ تشکیل شده است، یعنی $\Jdom_{22}$، فرم اینولوشن زیر را دارد:
\[
\Psi(z,u)=x_1^2z^2-x_1^4u_1^2-2x_1^3x_2u_1u_2-x_1^2x_2^2u_2^2.
\]
چون $x_1^2$ به $\Jdom_{22}$، و پس به طریق اولی به کوچک‌ترین دامنهٔ تمامِ نسبیِ دربرگیرندهٔ $\Gdom_{22}$ تعلق دارد و در نتیجه در $\Gdom_{22}$ مقسوم‌علیه مشترک است، فرم اینولوشنِ $\Gdom_{22}$ چنین می‌شود:
\[
\Psi(z,u)=z^2-x_1^2u_1^2-2x_1x_2u_1u_2-x_2^2u_2^2.
\]}

سرانجام یادآور می‌شویم که \emph{دامنهٔ تصویر} یک دامنهٔ تمامِ نسبیِ $\Gdom_{n\rho}$، که بنا بر \S\ 7 خود دامنهٔ صحیحِ چندجمله‌ایِ $\Jdom^*_{\rho\rho}$ است، لزوماً دامنه‌ای تمامِ نسبی نیست. زیرا اگر $F(x)$ و$G(x)$ دو چندجمله‌ای از $\Gdom_{n\rho}$ باشند که خارج‌قسمتشان چندجمله‌ای نیست و در نتیجه به $\Gdom_{n\rho}$ تعلق ندارد، خارج‌قسمت $F^*(x):G^*(x)$ نیز به $\Jdom^*_{\rho\rho}$ تعلق ندارد. اما این خارج‌قسمت که از تخصصی‌کردن برخی متغیرهای نامعین پدید می‌آید می‌تواند به‌خوبی یک چندجمله‌ای باشد؛ در آن صورت شرط 4الف برای $\Jdom^*_{\rho\rho}$ برقرار نیست.\srcfn{**)}{برای نمونه، فرض کنید $\Gdom_{32}$ از همهٔ چندجمله‌ای‌هایی تشکیل شده باشد که ترکیب‌های گویای $F=x_1^2x_2$ و$G=x_1^2(1+x_3)$ هستند. تخصص مجاز $x_3=0$ رابطهٔ $F^*:G^*=x_2$ را می‌دهد، در حالی که $F:G$ یک چندجمله‌ای نیست. پس $\Jdom^*_{22}$ دامنه‌ای تمامِ نسبی نیست.}

\section*{\S\ 10. دامنه‌های تمامِ نسبیِ نوع اول و پایهٔ تمامیت آن‌ها.}

مفهوم «تمامِ جبریِ نسبی» که در \S\ 9 (تعریف V) معرفی شد، امکان می‌دهد رده‌ای از دامنه‌های تمامِ نسبی مشخص کنیم که دامنه‌های صحیحِ متناهی‌مولدند؛ و بدین‌سان، دست‌کم برای این رده، به پرسشی که د. هیلبرت (مسائل ریاضی، مسئلهٔ 14) مطرح کرده بود پاسخی مثبت بدهیم. اکنون

\emph{تعریف VI: شمار متناهی‌ای از چندجمله‌ای‌های $\Gdom_{n\rho}$ «پایهٔ تمامیت» نامیده می‌شود اگر بتوان هر چندجمله‌ایِ $\Gdom_{n\rho}$ را، با ضرایبی از $\Omega$، به صورت ترکیبی گویای تمام از اعضای آن شمار متناهی نمایش داد. دامنهٔ صحیحی از چندجمله‌ای‌ها که دارای پایهٔ تمامیت باشد «دامنهٔ صحیحِ متناهی‌مولد» نامیده می‌شود.}

ردهٔ دامنه‌های تمامِ نسبی‌ای را که باید بررسی شوند و پایهٔ اینولوشنشان پایهٔ تمامیت از کار درخواهد آمد، بنابر تعریف زیر «دامنه‌های تمامِ نسبیِ نوع اول» می‌نامیم:

\emph{تعریف VII: دامنهٔ تمامِ نسبیِ $\Gdom_{n\rho}$ «نوع اول» نامیده می‌شود اگر به‌عنوان دامنهٔ تصویر، دامنهٔ تمامِ نسبیِ $\Gdom^*_{\rho\rho}$ را چنان داشته باشد که هر چندجمله‌ایِ دلخواه در $x_1\cdots x_\rho$ (با ضرایبی از $\Omega$) به‌طور تمامِ جبری به $\Gdom^*_{\rho\rho}$ وابسته باشد.}

زیرا بنابر این تعریف، متغیرهای نامعین $x_1\cdots x_\rho$ و در نتیجه فرم خطی $u(x)=u_1x_1+\cdots+u_\rho x_\rho$ نیز به‌طور تمامِ جبری به $\Gdom^*_{\rho\rho}$ وابسته‌اند. پس تابع تمامِ تحویل‌ناپذیر با ریشهٔ $z=u(x)$ --- یعنی فرم اینولوشنِ کوچک‌ترین میدان دربرگیرندهٔ $\Kfield^*_{\rho\rho}$ --- ضریب بالاترین درجه‌اش واحد و دیگر ضرایبش چندجمله‌ای‌هایی از $\Gdom^*_{\rho\rho}$ هستند؛ بنابراین، هم‌زمان فرم اینولوشنِ $\Gdom^*_{\rho\rho}$ است و در $\Gdom^*_{\rho\rho}$ هیچ فرم اینولوشن دیگری نمی‌تواند وجود داشته باشد. پس بنابر اصل انتقال، $\Gdom_{n\rho}$ نیز فرم اینولوشنی دارد که ضریب بالاترین درجه‌اش واحد است؛ و چون بنابر قضیهٔ VI، در نمایش گویای همهٔ چندجمله‌ای‌های $\Gdom_{n\rho}$ به‌وسیلهٔ پایهٔ اینولوشن متناظر، فقط همین ضریب بالاترین درجه وارد مخرج می‌شود، پایهٔ اینولوشن به پایهٔ تمامیت تبدیل می‌شود؛ یعنی:

\emph{قضیهٔ VII: هر دامنهٔ تمامِ نسبیِ نوع اول، دامنهٔ صحیحِ متناهی‌مولد است؛ پایهٔ تمامیت آن را پایهٔ اینولوشنی به دست می‌دهد که از دامنهٔ تصویرِ مشخص‌کنندهٔ آن حاصل شده است. به‌ویژه، پایهٔ تمامیتِ دامنه‌های تمامِ نسبیِ نوع اول $\Gdom_{\rho\rho}$ را پایهٔ اینولوشنِ به‌طور یکتا تعریف‌شدهٔ $\Gdom_{\rho\rho}$ به دست می‌دهد.}

معیاری نیز برای دامنه‌های تمامِ نسبیِ نوع اول می‌دهیم. فرض کنید $\Gdom_{n\rho}$ از نوع اول باشد و پایهٔ تمامیتی برای $\Gdom_{n\rho}$ به صورت زیر داده شده باشد:
\[
\begin{array}{@{}l@{\quad}c@{}}
(1)& F_1(x),F_2(x),\ldots,F_k(x).
\end{array}
\]
آنگاه بنابر لمی از هیلبرت\srcfn{*)}{د. هیلبرت: \srcLTR{Über die vollen Invariantensysteme}، \srcLTR{Math. Ann. 42 (1893)}، \S\ 1. (لمی که در آنجا فقط برای چندجمله‌ای‌های همگن بیان شده، برای چندجمله‌ای‌های ناهمگن نیز به همان‌سان برقرار است.)}، $\rho$ چندجمله‌ایِ مستقلِ جبری از $\Gdom_{n\rho}$ وجود دارند:
\[
\begin{array}{@{}l@{\quad}c@{}}
(2)& G_1(x),G_2(x),\ldots,G_\rho(x)
\end{array}
\]
به‌گونه‌ای که چندجمله‌ای‌های (1)، و در نتیجه هر چندجمله‌ایِ $\Gdom_{n\rho}$، به‌طور تمامِ جبری به چندجمله‌ای‌های (2) وابسته باشند. همین امر برای هر چندجمله‌ایِ دامنهٔ تصویر $\Gdom^*_{\rho\rho}$ نسبت به چندجمله‌ای‌های متناظر زیر برقرار است:
\[
\begin{array}{@{}l@{\quad}c@{}}
(3)& G_1^*(x),G_2^*(x),\ldots,G_\rho^*(x).
\end{array}
\]


پس بنابر تعریف VII، هر چندجمله‌ایِ دلخواه در $x_1\rcdots x_\rho$ نیز به‌طور تمامِ جبری به چندجمله‌ای‌های (3) وابسته است.

اکنون برعکس، فرض کنیم دامنهٔ تصویر $\Jdom^*_{\rho\rho}$ یک دامنهٔ تمامِ نسبیِ دلخواهِ $\Gdom_{n\rho}$، $\rho$ چندجمله‌ایِ (3) را در بر داشته باشد که هر چندجمله‌ای در $x_1\rcdots x_\rho$ به‌طور تمامِ جبری به آن‌ها وابسته باشد. آنگاه نشان می‌دهیم که از این فرض، $\Jdom^*_{\rho\rho}$ به‌صورت دامنهٔ تمامِ نسبیِ $\Gdom^*_{\rho\rho}$ به دست می‌آید، و نیز $\Gdom^*_{\rho\rho}$ و در نتیجه $\Gdom_{n\rho}$ از نوع اول‌اند.

در واقع، فرض کنید $F^*(x)$ چندجمله‌ای‌ای از کوچک‌ترین میدانِ دربرگیرندهٔ $\Jdom^*_{\rho\rho}$، یعنی $\Kfield^*_{\rho\rho}$، باشد؛ پس بنابر فرض، به‌طور تمامِ جبری به چندجمله‌ای‌های (3) وابسته است. چون تابع متناظر در $\Kfield_{n\rho}$ آنگاه به‌طور تمامِ جبری به $\rho$ چندجمله‌ای از $\Gdom_{n\rho}$ وابسته است، یک چندجمله‌ایِ $F(x)$ است و در نتیجه به $\Gdom_{n\rho}$ تعلق دارد؛ پس $F^*(x)$ به $\Jdom^*_{\rho\rho}$ تعلق دارد. بنابراین دامنهٔ تصویر $\Jdom^*_{\rho\rho}$ همهٔ چندجمله‌ای‌های $\Kfield^*_{\rho\rho}$ را در بر دارد و دامنهٔ تمامِ نسبیِ $\Gdom^*_{\rho\rho}$ است. اما برای $\Gdom^*_{\rho\rho}$ و ازاین‌رو برای $\Gdom_{n\rho}$ نیز، بنابر تعریف‌های V و VII، بی‌درنگ از فرض نتیجه می‌شود که از نوع اول‌اند؛ یعنی:

\emph{دامنه‌های تمامِ نسبیِ نوع اول $\Gdom_{n\rho}$ با این ویژگی مشخص می‌شوند که در یک دامنهٔ تصویر $\Jdom^*_{\rho\rho}$، $\rho$ چندجمله‌ای وجود دارند که هر چندجمله‌ایِ دلخواه در $x_1\rcdots x_\rho$ به‌طور تمامِ جبری به آن‌ها وابسته است. دامنهٔ تصویر از این فرض به‌صورت دامنه‌ای تمامِ نسبی به دست می‌آید.}\srcfn{*)}{برای نمونه، در دامنهٔ $\Gdom_{32}$ که در یادداشت پایانیِ \S\ 9 ذکر شد، تخصص مجاز $x_1=1$ دو چندجمله‌ایِ $F^*=x_2$ و$G^*=1+x_3$ را به دست می‌دهد که هر چندجمله‌ایِ دلخواه در $x_2,x_3$ به‌طور تمامِ جبری به آن‌ها وابسته است. پس $\Gdom_{32}$ از نوع اول است و $F(x)$ و$G(x)$ پایهٔ تمامیت آن‌اند.}

\section*{\S\ 11. لمی دربارهٔ وابستگیِ تمامِ جبری.}

اکنون معیاری به دست می‌دهیم تا هر چندجمله‌ای دلخواه در $x_1\rcdots x_\rho$ به $\rho$ چندجمله‌ای زیر در $x_1\rcdots x_\rho$
\[
\begin{array}{@{}l@{\quad}c@{}}
(1)& G_1^*(x)\rcdots G_\rho^*(x)
\end{array}
\]
به‌طور تمامِ جبری وابسته باشد.

اگر چندجمله‌ای‌های (1)، به‌طور خاص، فرم‌های همگن باشند، آنگاه برای هر دستگاه مقادیری که (1) را صفر کند، همهٔ فرم‌هایی که به‌طور تمامِ جبری به آن‌ها وابسته‌اند نیز هم‌زمان صفر می‌شوند؛ پس به‌ویژه $x_1\rcdots x_\rho$ نیز صفر می‌شوند.

بنابراین فرم‌های (1) در هیچ دستگاه مقادیری جز
\[
x_1=0,\ldots,x_\rho=0
\]
صفر نمی‌شوند؛ برآیند آن‌ها باید ناصفر باشد.

برعکس، این شرط کافی نیز هست و طبق لم زیر می‌توان آن را به چندجمله‌ای‌های ناهمگن تعمیم داد:

\emph{لم: شرطی کافی برای اینکه هر چندجمله‌ای دلخواه در $x_1\rcdots x_\rho$ به‌طور تمامِ جبری به $\rho$ چندجمله‌ای در $x_1\rcdots x_\rho$
\[
\begin{array}{@{}l@{\quad}c@{}}
(1)& G_1^*(x)\rcdots G_\rho^*(x)
\end{array}
\]
وابسته باشد، ناصفربودنِ برآیندِ ــ نسبت به $x_1\rcdots x_\rho$ ــ جزءهای همگنِ دارای بالاترین درجه}
\[
H_1^*(x)\rcdots H_\rho^*(x)
\]
\emph{از چندجمله‌ای‌های (1) است:}
\[
\begin{array}{@{}l@{\quad}c@{}}
(2)& \displaystyle R_0=\left[\begin{array}{c} H_1^*\rcdots H_\rho^*\\ x_1\rcdots x_\rho\end{array}\right]\ne0.
\end{array}
\]
\emph{برای فرم‌های همگن (1)، شرط (2) ضروری است.}\srcfn{*)}{دربارهٔ نظریهٔ برآیند بنگرید به ف. مرتنس، \srcLTR{Zur Theorie der Elimination (I. u. II. Teil)}، \srcLTR{Sitzungsber. d. k. Ak. d. Wiss. Wien, Math.-naturw. Kl. Abt. IIa, Bd. 108 (1899)}. (بخشی از آن نزد ی. کونیگ، \srcLTR{Theorie d. algebr. Größen}، فصل ششم، بازآورده شده است).}

برای اثبات، چندجمله‌ای‌های زیر را می‌سازیم:
\[
\begin{array}{@{}l@{\quad}c@{}}
(3)& \begin{gathered}
\bar G_1^*(x)=G_1^*(x)-\lambda_1,\ldots,\bar G_\rho^*(x)=G_\rho^*(x)-\lambda_\rho,\\
\bar\Phi^*(x)=\Phi^*(x)-\mu,
\end{gathered}
\end{array}
\]
که در آن $\lambda_1\rcdots\lambda_\rho,\mu$ متغیرهای نامعین و $\Phi^*(x)$ یک چندجمله‌ای دلخواه در $x_1\rcdots x_\rho$ هستند. برآیندِ چندجمله‌ای‌های (3) نسبت به $x_1\rcdots x_\rho,1$\srcfn{**)}{یعنی برآیند همگن نسبت به $x_1\rcdots x_\rho x_0$، هنگامی که چندجمله‌ای‌های (3) را با یک متغیر نامعین اضافی $x_0$ همگن می‌کنیم، به‌گونه‌ای که درجه نسبت به $x$ حفظ شود.} را به صورت زیر در نظر می‌گیریم:
\[
\begin{array}{@{}l@{\quad}c@{}}
(4)& \displaystyle R(\lambda,\mu)=\left[\begin{array}{cc} \bar G_1^*\rcdots \bar G_\rho^*\bar\Phi^*\\ x_1\rcdots x_\rho\quad 1\end{array}\right];
\end{array}
\]
این یک تابع گویای تمام از $\lambda_1\rcdots\lambda_\rho,\mu$ است. اکنون بنا بر ف. مرتنس، در بسط (4) برحسب توان‌های $\mu$ می‌توان ضریبِ بالاترین توان را به‌صراحت تعیین کرد\srcfn{***)}{همان‌جا، \S\ 3 (برهان قضیهٔ اول)، و به‌صراحت \S\ 6 (تعیین مجموعهٔ همهٔ جمله‌های برآیند $R$ که حاصل‌توانِ معیّن $a_{\lambda\lambda}^{p_\lambda}a_{\mu\mu}^{p_\mu}\cdots a_{\sigma\sigma}^{p_\sigma}$ را به‌عنوان عامل در بر دارند). در حالت ما $p_\lambda=0$، و $p_\mu=0,\ldots,p_\sigma=\tau$، و $a_{\sigma\sigma}=(-\mu)$ گذاشته می‌شود. نزد کونیگ، فصل ششم، \S\ 9، بازآورده شده است.}؛ نتیجه چنین است:
\[
\begin{array}{@{}l@{\quad}c@{}}
(5)& (-1)^\tau R(\lambda,\mu)=R_0^\sigma\cdot\mu^\tau+A_1(\lambda)\mu^{\tau-1}+\cdots+A_\tau(\lambda),
\end{array}
\]
که در آن $A_1(\lambda)\rcdots A_\tau(\lambda)$ تابع‌های چندجمله‌ای با ضرایب صحیح در $\lambda_1\rcdots\lambda_\rho$ و در ضرایب چندجمله‌ای‌های (3) هستند. این بسط نخست نشان می‌دهد که با فرض (2)، یعنی $R_0\ne0$، $R(\lambda,\mu)$ نیز نمی‌تواند به‌طور همانی صفر باشد. هم‌نهشتیِ زیر در $x_1\rcdots x_\rho,\lambda_1\rcdots\lambda_\rho\mu$:
\[
R(\lambda,\mu)\equiv0\pmod{\bar G_1^*(x),\ldots,\bar G_\rho^*(x),\bar\Phi^*(x)}
\]
با جایگذاری زیر
\[
\lambda_i=G_i^*(x),\qquad \mu=\Phi^*(x)
\]
به رابطهٔ زیر تبدیل می‌شود:
\[
\begin{array}{@{}l@{\quad}l@{}}
(6)& \displaystyle R(G_i^*(x),\Phi^*(x))=R_0^\sigma\cdot\Phi^{*\tau}(x)+A_1(G_i^*(x))\Phi^{*\tau-1}(x)+\cdots\\
& \displaystyle\qquad +A_\tau(G_i(x))=0,
\end{array}
\]
و بدین‌سان لم اثبات می‌شود، زیرا $R_0$ بنا بر فرض عددی از $\Omega$ است.

\section*{\S\ 12. پایهٔ تمامیتِ دستگاه‌های منظم $\Ssys_{n\rho}$ از چندجمله‌ای‌ها.}

لمِ \S\ 11، همراه با نتیجهٔ پایانیِ \S\ 10، بی‌درنگ واقعیت زیر را به دست می‌دهد:

\emph{اگر در یک دامنهٔ تصویر $\Jdom^*_{\rho\rho}$ از $\Gdom_{n\rho}$، $\rho$ چندجمله‌ای وجود داشته باشد، چنان‌که برآیند همگنِ جزءهای دارای بالاترین درجه به‌طور همانی صفر نباشد --- $R_0\ne0$ ---، آنگاه $\Gdom_{n\rho}$ از نوع اول و در نتیجه دامنهٔ صحیحی متناهی‌مولد است، و پایهٔ اینولوشنِ آن یک پایهٔ تمامیت است.}

اکنون شرط $R_0\ne0$ به‌کمک لم، برهان دومی برای وجود پایهٔ تمامیت ممکن می‌سازد که اساساً از د. هیلبرت است\srcfn{*)}{\srcLTR{Über die vollen Invariantensysteme}، \S\ 2. نزد هیلبرت، با فرض اینکه وجود پایهٔ تمامیتِ میدان ناورداها از راهی دیگر به دست آمده است، سخن فقط بر سر تفسیر این پایه به‌وسیلهٔ دستگاه بنیادیِ کرونکر برای کمیت‌های تمامِ جبریِ یک میدان است.}؛ این برهان اگرچه پایهٔ اینولوشن را از آن حیث بازنمی‌شناساند، در عوض به‌طور یکنواخت برای همهٔ دستگاه‌های $\Ssys_{n\rho}$ با ویژگیِ یادشده برقرار است.

به‌طور مشخص، فرض کنیم آن $\rho$ چندجمله‌ای از $\Jdom^*_{\rho\rho}$ که برایشان $R_0\ne0$ است و بنابراین مستقلِ جبری‌اند، عبارت باشند از
\[
\begin{array}{@{}l@{\quad}c@{}}
(1)& G_1^*(x)\rcdots G_\rho^*(x).
\end{array}
\]
بنا بر لم، هر چندجمله‌ای دلخواه در $x_1\rcdots x_\rho$ به‌طور تمامِ جبری به چندجمله‌ای‌های (1) وابسته است؛ پس این حکم به‌ویژه برای هر چندجمله‌ایِ متعلق به کوچک‌ترین میدانِ دربرگیرندهٔ $\Jdom^*_{\rho\rho}$، یعنی $\Kfield^*_{\rho\rho}$، برقرار است. اگر چندجمله‌ای‌های متناظر با آن‌ها در $\Ssys_{n\rho}$ چنین باشند:
\[
\begin{array}{@{}l@{\quad}c@{}}
(2)& G_1(x)\rcdots G_\rho(x),
\end{array}
\]
آنگاه، بنا بر اصل انتقال، هر چندجمله‌ایِ کوچک‌ترین میدانِ دربرگیرندهٔ $\Ssys_{n\rho}$، یعنی $\Kfield_{n\rho}$، و در نتیجه به‌ویژه هر چندجمله‌ای از $\Ssys_{n\rho}$، به‌طور تمامِ جبری به چندجمله‌ای‌های (2) وابسته است. برعکس، هر تابع از $\Kfield_{n\rho}$ که به‌طور تمامِ جبری به چندجمله‌ای‌های (2) وابسته باشد نیز یک چندجمله‌ای است. بنابراین اگر، مطابق \S\ 4، $\Kfield_{n\rho}$ را میدانی جبری بر $\Omega(G_1\rcdots G_\rho)$ در نظر بگیریم، مجموعهٔ همهٔ چندجمله‌ای‌های $\Kfield_{n\rho}$ --- یعنی دامنهٔ تمامِ نسبی $\Gdom_{n\rho}$ --- با مجموعهٔ کمیت‌های تمامِ جبریِ این میدان یکسان است. فرض کنیم $G(x)$ چندجمله‌ای‌ای از $\Kfield_{n\rho}$ باشد که (2) را به یک پایهٔ گویا تکمیل می‌کند، و $D(x)$ ممیّزِ معادلهٔ تحویل‌ناپذیرِ درجهٔ $k$ باشد که $G(x)$ نسبت به $\Omega(G_1\rcdots G_\rho)$ در آن صدق می‌کند. آنگاه هر چندجمله‌ای از $\Kfield_{n\rho}$، و به‌ویژه هر چندجمله‌ای $F(x)$ از $\Ssys_{n\rho}$، به‌عنوان کمیتی تمامِ جبری، نمایش زیر را می‌پذیرد:
\[
\begin{array}{@{}l@{\quad}c@{}}
(3)& \displaystyle F(x)=\frac{\Gamma_1(G_i)\cdot G(x)^{k-1}+\Gamma_2(G_i)\cdot G(x)^{k-2}+\cdots+\Gamma_k(G_i)}{D(x)}.
\end{array}
\]

اکنون بگذاریم $F(x)$ همهٔ چندجمله‌ای‌های $\Ssys_{n\rho}$ را طی کند. در این صورت، کاربرد قضیهٔ I هیلبرت دربارهٔ پایهٔ مدولی (\srcLTR{Math. Ann. 36}) بر دستگاهی که باید از تابع‌های متناظر $v_1\Gamma_1(G_i)+v_2\Gamma_2(G_i)+\cdots+v_k\Gamma_k(G_i)$ ساخته شود، وجود شمار متناهی‌ای از چندجمله‌ای‌های $\Ssys_{n\rho}$ را نشان می‌دهد:
\[
\begin{array}{@{}l@{\quad}c@{}}
(4)& F_1(x),F_2(x),\ldots,F_\lambda(x)
\end{array}
\]
به‌گونه‌ای که هر چندجمله‌ایِ $\Ssys_{n\rho}$ نمایش زیر را می‌پذیرد:
\[
\begin{array}{@{}l@{\quad}c@{}}
(5)& F(x)=A_1(G_i)F_1(x)+A_2(G_i)\cdot F_2(x)+\cdots+A_\lambda(G_i)F_\lambda(x),
\end{array}
\]
که در آن $A_1\rcdots A_\lambda$ چندجمله‌ای‌هایی در $G_i(x)$ هستند. بنابراین چندجمله‌ای‌های (2) و (4) یک پایهٔ تمامیت برای $\Ssys_{n\rho}$ می‌سازند؛ یعنی:

\emph{قضیهٔ VIII: اگر در یک دستگاه تصویر $\Jdom^*_{\rho\rho}$ از دستگاه چندجمله‌ای‌ها $\Ssys_{n\rho}$، $\rho$ چندجمله‌ای وجود داشته باشد، چنان‌که برآیندِ جزءهای دارای بالاترین درجه به‌طور همانی صفر نباشد --- $R_0\ne0$ ---، آنگاه $\Ssys_{n\rho}$ دارای پایهٔ تمامیت است.}

این قضیه تعمیم مختصری نیز می‌پذیرد. کافی است آن $\rho$ چندجمله‌ایِ (1) که برایشان $R_0\ne0$ است، به کوچک‌ترین دامنهٔ صحیحِ دربرگیرندهٔ $\Ssys^*_{\rho\rho}$، یعنی $\Jdom^*_{\rho\rho}$، تعلق داشته باشند. در واقع، استدلال‌های منتهی به قضیهٔ VIII نشان می‌دهند که در این حالت نیز نمایش (5) همچنان برقرار است. اما بنا بر تعریف، هر چندجمله‌ای $L(x)$ از کوچک‌ترین دامنهٔ صحیحِ دربرگیرنده $\Jdom_{n\rho}$ را می‌توان به‌صورت ترکیبی چندجمله‌ای از شمار متناهی‌ای --- که با $L(x)$ تغییر می‌کنند --- از چندجمله‌ای‌های $\Ssys_{n\rho}$ نمایش داد؛ بنابراین $\sigma$ چندجمله‌ای از $\Ssys_{n\rho}$ وجود دارند:
\[
\begin{array}{@{}l@{\quad}c@{}}
(6)& K_1(x),K_2(x),\ldots,K_\sigma(x)
\end{array}
\]
چنان‌که چندجمله‌ای‌های $G_i(x)$ ترکیب‌های چندجمله‌ایِ (6) می‌شوند. پس (6) و (4) یک پایهٔ تمامیت به دست می‌دهند. با معرفی تعریف زیر:

\emph{تعریف VIII: دستگاه $\Ssys_{n\rho}$ از چندجمله‌ای‌ها منظم نامیده می‌شود هرگاه در یک دامنهٔ تصویر $\Jdom^*_{\rho\rho}$ از کوچک‌ترین دامنهٔ صحیحِ دربرگیرنده، $\rho$ چندجمله‌ای وجود داشته باشد، چنان‌که برآیندِ جزءهای دارای بالاترین درجه به‌طور همانی صفر نباشد --- $R_0\ne0$ ---}

در نتیجه داریم:

\emph{قضیهٔ IX: هر دستگاه منظم $\Ssys_{n\rho}$ از چندجمله‌ای‌ها دارای پایهٔ تمامیت است.}\srcfn{*)}{هیلبرت با مثالی نشان داده است که دستگاه‌های نامنظمِ بدون پایهٔ تمامیت وجود دارند (\srcLTR{Gött. Nachr. 1891, S. 233}). مثال زیر اندکی ساده‌تر است:
\[
\Ssys_{22}=x_1,
x_1x_2,
x_1x_2^2,
\ldots,
x_1x_2^\nu,
\ldots\quad\text{in inf.}
\]
وجود پایهٔ تمامیت در اینجا هم‌ارز با وجود عدد صحیح مثبت $\nu$ای است که همانی‌های زیر برقرار باشند:
\[
x_1x_2^{\nu+\sigma}=\sum c\cdot x_1^{i_0}(x_1x_2)^{i_1}(x_1x_2^2)^{i_2}\cdots(x_1x_2^\nu)^{i_\nu},
\qquad(\sigma=1,2,\ldots).
\]
اما مقایسهٔ درجه‌ها در $x_1,x_2$، حتی برای $\sigma=1$، به دو معادلهٔ شرطیِ زیر برای توان‌ها می‌انجامد:
\[
1=i_0+i_1+i_2+\cdots+i_\nu,
\qquad
\nu+1=i_1+2i_2+\cdots+\nu i_\nu,
\]
که آشکارا در اعداد صحیح نامنفی جوابی ندارند. پس $\Ssys_{22}$ پایهٔ تمامیت ندارد.}

همچنین، همانند کاری که در \S\ 5 برای اینولوشن انجام دادیم، تفسیری هندسی از دستگاه‌های منظم به دست می‌دهیم. برای این منظور، مانند (7) در آنجا، دستگاه معادلات زیر را در نظر می‌گیریم:
\[
\begin{array}{@{}l@{\quad}c@{}}
(7)& L^*(x)-L^*(\xi)=0,
\end{array}
\]
که در آن $L^*(x)$ همهٔ چندجمله‌ای‌های $\Jdom^*_{\rho\rho}$ را طی می‌کند. اگر با استفاده از $x_0,\xi_0$ همگن‌سازی کنیم، فرض می‌کنیم (7) به دستگاه معادلات زیر میان فرم‌های همگن تبدیل شود:
\[
\begin{array}{@{}l@{\quad}c@{}}
(8)& \xi_0^m\cdot M^*(x)-x_0^m\cdot M^*(\xi)=0,
\end{array}
\]
و اگر $H^*(x)$ جزء همگنِ دارای بالاترین درجهٔ $L^*(x)$ را نشان دهد، رابطهٔ زیر برقرار است:
\[
\begin{array}{@{}l@{\quad}c@{}}
(9)& H^*(x)\equiv M^*(x)\pmod{x_0}.
\end{array}
\]

\emph{نقاط بنیادیِ} $\Jdom^*_{\rho\rho}$ را (مقایسه کنید با \S\ 5، یادداشت) جواب‌های (8) می‌نامیم که از $x$ مستقل‌اند --- یعنی با ردیف‌های مزدوج با $x$ تفاوت دارند. پس نقاط بنیادی همان دستگاه‌های مقادیری‌اند که با دستگاه $\xi_1=0,\ldots,\xi_\rho=0$ تفاوت دارند و در آن‌ها هم‌زمان برقرار است:
\[
\xi_0^m=0,
\qquad M^*(\xi)=0
\]
یا، بنا بر (9):
\[
\xi_0=0,
\qquad H^*(\xi)=0.
\]

اکنون اگر $\rho$ فرم $H^*(\xi)$ وجود داشته باشند که برایشان $R_0\ne0$ است، و بنابراین دستگاه جواب مشترکی ندارند، آنگاه $\Jdom^*_{\rho\rho}$ نمی‌تواند نقطهٔ بنیادی داشته باشد. به‌طور خاص، اگر $\Ssys^*_{\rho\rho}$ از فرم‌های همگن تشکیل شده باشد، با $R_0\ne0$ خودِ $\Ssys^*_{\rho\rho}$ نیز نمی‌تواند نقطهٔ بنیادی داشته باشد، زیرا چنین نقطه‌ای هم‌زمان نقطهٔ بنیادیِ $\Jdom^*_{\rho\rho}$ نیز می‌بود. اما عکس این حکم نیز برقرار است. برای این منظور دستگاهی را که از مجموعهٔ همهٔ فرم‌های $H^*(x)$ ساخته شده است، یعنی $\Jdom(H^*)$، در نظر می‌گیریم؛ این دستگاه نیز یک دامنهٔ صحیح می‌سازد.

ازاین‌رو، از قضیهٔ I هیلبرت دربارهٔ پایهٔ مدولی و لم هیلبرت (مقایسه کنید با \S\ 10، ارجاع) وجود $\rho$ فرم از $\Jdom(H^*)$ نتیجه می‌شود که صفرشدنشان، صفرشدن همهٔ فرم‌های $\Jdom(H^*)$ را در پی دارد. اگر $\Jdom^*_{\rho\rho}$ نقطهٔ بنیادی نداشته باشد، این $\rho$ فرم نمی‌توانند هم‌زمان صفر شوند؛ پس برآیندشان $R_0\ne0$ است.

بنابراین:

\emph{دستگاه‌های منظم $\Ssys_{n\rho}$ از چندجمله‌ای‌ها با این ویژگی مشخص می‌شوند که یک دامنهٔ تصویر $\Jdom^*_{\rho\rho}$ از کوچک‌ترین دامنهٔ صحیحِ دربرگیرنده، نقطهٔ بنیادی ندارد. در دستگاه‌های منظم $\Ssys^*_{\rho\rho}$ از فرم‌های همگن، خودِ دستگاه نیز نقطهٔ بنیادی ندارد.}\srcfn{*)}{دستگاه نامنظم $\Ssys_{22}$ که در یادداشت پیشین ذکر شد، نقطهٔ بنیادیِ زیر را دارد:
\[
\xi_1:\xi_2:\xi_0=0:1:0.
\]}

\section*{\S\ 13. مثال‌هایی از دامنه‌های تمامِ نسبیِ نوع اول و دستگاه‌های منظم.}

1. به‌عنوان نخستین نمونه از یک دامنهٔ تمامِ نسبیِ نوع اول $\Gdom_{nn}$، \emph{کل چندجمله‌ای‌های $x_1\rcdots x_n$ را که یک گروه جایگشتیِ معین $G$ را می‌پذیرند} در نظر می‌گیریم --- یعنی کل چندجمله‌ای‌های یک میدانِ گونهٔ لاگرانژی. از آنجا که $x_1\rcdots x_n$ به‌موجب همانی
\[
\begin{array}{@{}l@{\quad}c@{}}
(1)& x_k^n-\sigma_1(x)x_k^{n-1}+\sigma_2(x)x_k^{n-2}+\cdots\pm\sigma_n(x)=0\quad(k=1,2,\ldots,n)
\end{array}
\]
به توابع متقارن مقدماتیِ متعلق به $\Gdom_{nn}$ وابستگیِ تمامِ جبری دارند، و بنابراین بنا بر تعریف V به $\Gdom_{nn}$ وابستگیِ تمامِ جبری دارند، $\Gdom_{nn}$ از نوع اول است. همچنین $\Gdom_{nn}$، به‌عنوان دامنه‌ای از نوع اول، چون اعضایش فرم‌های همگن‌اند و $n=\rho$ است، (بنا بر پایان \S\ 10) دستگاهی منظم است. در واقع، به‌سبب (1)، برآیندِ توابع متقارن مقدماتی $R_0$ نمی‌تواند به‌طور همانی صفر شود؛ $R_0$ بر پایهٔ قواعد محاسبهٔ نظریهٔ برآیند برابر با $\pm1$ به دست می‌آید. \emph{فرم اینولوشنِ} $\Gdom_{nn}$ نیز --- چون $\Gdom_{nn}$ از نوع اول و $n=\rho$ است --- به‌طور یکتا همان فرم اینولوشنِ میدانِ گونهٔ لاگرانژیِ متناظر تعیین می‌شود. از آنجا که کمیت‌های مزدوج با $x_1\rcdots x_n$ نسبت به این میدان از جایگشت‌های گروه $G$ حاصل می‌شوند، این فرم اینولوشن به‌صورت زیر است:
\[
\begin{array}{@{}l@{\quad}c@{}}
(2)&
\begin{aligned}
  \Psi(z,u)&=\prod_G
  (z-u_1x_{i_1}-u_2x_{i_2}-\cdots-u_nx_{i_n})\\
  &=z^i+\sum{}'G_{\alpha\alpha_1\ldots\alpha_n}(x)
  z^\alpha u_1^{\alpha_1}\cdots u_n^{\alpha_n},
\end{aligned}
\end{array}
\]
که در آن حاصل‌ضرب باید بر همهٔ جایگشت‌های گروه $G$ گسترده شود. (2) \glqq{}فرم گالوای گروه $G$\grqq{} را نمایش می‌دهد؛ پس حکم زیر برقرار است:

\emph{همهٔ چندجمله‌ای‌های $x_1\rcdots x_n$ که گروه جایگشتی $G$ را می‌پذیرند، ترکیب‌های چندجمله‌ایِ ضرایب $G_{\alpha\alpha_1\ldots\alpha_n}(x)$ فرم گالوای گروه‌اند.}

2. به‌عنوان نمونه‌های \emph{دستگاه‌های منظم}، دستگاه‌های $\Ssys_{n1}$ متشکل از چندجمله‌ای‌ها را در نظر می‌گیریم؛ یعنی دستگاه‌هایی که تنها یک چندجمله‌ایِ مستقلِ جبری در بر دارند.

در واقع، فرض کنیم $\Ssys^*_{11}$ یک دستگاه تصویرِ $\Ssys_{n1}$ باشد و
\[
  F^*(x)=a_0x^k+a_1x^{k-1}+\cdots+a_k\qquad(a_0\ne0),
\]
چندجمله‌ای‌ای از $\Ssys^*_{11}$ باشد؛ آنگاه
\[
\begin{array}{@{}l@{\quad}c@{}}
(3)& R_0=a_0\ne0\text{\srcfnmark{*)}}
\end{array}
\]
\srcfntext{*)}{مرتنس، همان مأخذ، \S\ 5.}
پس $\Ssys_{n1}$ بنا بر تعریف VIII منظم است و بنابراین:

\emph{هر دستگاه $\Ssys_{n1}$ از چندجمله‌ای‌ها، به‌ویژه هر دستگاهی از بی‌نهایت چندجمله‌ای در یک متغیر نامعین، پایهٔ تمامیت دارد.}\srcfn{**)}{پس ساده‌ترین دستگاه‌های نامنظم ممکنِ فاقد پایهٔ تمامیت، دستگاه‌های $\Ssys_{22}$ هستند. مثال هیلبرت و مثال آمده در یادداشت \S\ 12 نشان می‌دهند که چنین دستگاه‌هایی واقعاً وجود دارند.}

3. اکنون دستگاه‌های $\Ssys_{n1}$ را به حالت دامنه‌های تمامِ نسبی $\Gdom_{n1}$ تخصیص می‌دهیم؛ این دامنه‌ها چون دستگاه‌هایی منظم‌اند، بنا بر \S\ 12 دامنه‌های تمامِ نسبیِ نوع اول‌اند. بنابراین پایهٔ تمامیت آن‌ها یک پایهٔ اینولوشن است که هم‌زمان پایهٔ اینولوشنِ کوچک‌ترین میدان دربرگیرنده $\Kfield_{n1}$ نیز هست. اکنون بنا بر \S\ 6، هر تابع دلخواهی از پایهٔ اینولوشنِ $\Kfield_{n1}$ در عین حال پایه‌ای کمینه است --- آن را تابع لوروتِ میدان نامیدیم --- و هر تابع از پایهٔ اینولوشن به‌طور خطی-کسری به این تابع لوروت وابسته است. در حالت ما، تابع لوروت چون به $\Gdom_{n1}$ تعلق دارد چندجمله‌ای است؛ بنابراین هر چندجمله‌ایِ دیگرِ پایهٔ اینولوشن به‌طور خطی-کسری، و بنا بر (3) به‌صورت تمامِ جبری، به این تابع لوروت $L(x)$ وابسته است و در نتیجه وابستگی‌اش به آن چندجمله‌ای و خطی است؛ پس همین تابع پایهٔ تمامیت می‌شود. فرم اینولوشنِ $\Gdom^*_{11}$ نیز، با قرار دادن $(u_1=1)$، چنین به دست می‌آید:
\[
  \Psi^*(z)=L^*(z)-L^*(x),
\]
که از آن دوباره روشن می‌شود $L(x)$ پایهٔ تمامیت است. بنابراین:

\emph{دامنه‌های تمامِ نسبی $\Gdom_{n1}$ از نوع اول‌اند؛ پایهٔ تمامیتشان از یک چندجمله‌ای، یعنی تابع لوروتِ کوچک‌ترین میدان دربرگیرنده، تشکیل می‌شود.}

4. به‌عنوان نمونهٔ خاص دیگری از یک دامنهٔ تمامِ نسبی $\Gdom_{n1}$، ناورداهای پروژکتیوِ گویای تمام --- یعنی چندجمله‌ای --- یک فرم درجه‌دو با $s$ متغیر را یاد می‌کنیم. در سازگاری با مطالب پیشین، چنان‌که می‌دانیم، این ناورداها را می‌توان به‌صورت ترکیبی چندجمله‌ای از دترمینانِ $|a_{ik}|$ فرم نمایش داد؛ و این دترمینان در عین حال تابع لوروتِ میدان ناورداهای متناظر است.

\section*{\S\ 14. دستگاه‌هایی از چندجمله‌ای‌های صحیح.}

اکنون می‌خواهیم قضیه‌های پایهٔ به‌دست‌آمده را از نظر صحیح‌بودنِ ضرایب تقویت کنیم.

ازاین‌رو، در ادامه حوزهٔ ضرایب $\Omega$ را یک میدان عددی جبریِ (متناهی) فرض می‌کنیم و $[\Omega]$ را مجموعهٔ اعداد صحیح جبریِ $\Omega$ می‌گیریم که به‌اختصار آن‌ها را اعداد صحیح خواهیم نامید. منظور از چندجمله‌ای صحیح، چندجمله‌ای‌ای با ضرایب متعلق به $[\Omega]$ است؛ از این پس $F(x),G(x),\ldots$ فقط چندجمله‌ای‌های صحیح را نشان می‌دهند.

به قیاس با \S\ 7، دستگاه‌های گوناگونِ ساخته از چندجمله‌ای‌های صحیح را در نظر می‌گیریم:

هر دستگاه دلخواه از چندجمله‌ای‌های صحیح را \emph{دستگاه صحیح‌ضرایب} $\Ssys_{n\rho}$ می‌نامیم.

خانوادهٔ خطیِ صحیح‌ضرایب $\Lsys_{n\rho}$ دستگاهی صحیح‌ضرایب است که همراه با $F_1(x)$ و $F_2(x)$، همواره $c_1F_1(x)+c_2F_2(x)$ را نیز دربر دارد؛ در اینجا $c_1,c_2$ اعداد صحیحِ دلخواهی از $[\Omega]$ هستند.

دامنهٔ صحیحِ صحیح‌ضرایب $\Jdom_{n\rho}$ خانواده‌ای خطی و صحیح‌ضرایب است که همراه با $F(x)$ و $G(x)$، همواره $F(x)\cdot G(x)$ را نیز دربر دارد.

دامنهٔ تمامِ نسبیِ صحیح‌ضرایب $\Gdom_{n\rho}$ دامنه‌ای صحیح است که همراه با $F(x)$ و $G(x)$ ــ که $G(x)\ne0$ ــ خارج‌قسمت $F(x):G(x)$ را همواره درست در صورتی دربر دارد که این خارج‌قسمت چندجمله‌ای‌ای صحیح باشد.

همانند \S\ 7، کوچک‌ترین دامنهٔ صحیحِ صحیح‌ضرایب $\Jdom_{n\rho}$ که دستگاه صحیح‌ضرایب $\Ssys_{n\rho}$ را دربر دارد، از همهٔ چندجمله‌ای‌های $H(x)$ تشکیل می‌شود که ترکیب‌های چندجمله‌ای با ضرایب صحیح از شمار متناهی‌ای از چندجمله‌ای‌های $\Ssys_{n\rho}$ هستند؛

و کوچک‌ترین دامنهٔ تمامِ نسبیِ صحیح‌ضرایب $\Gdom_{n\rho}$ که آن را دربر دارد، از همهٔ چندجمله‌ای‌های صحیح $H(x)$ تشکیل می‌شود که ترکیب‌های گویایی از شمار متناهی‌ای از چندجمله‌ای‌های $\Ssys_{n\rho}$ هستند.

اکنون به انتقال قضیه‌های پایه می‌پردازیم. دربارهٔ پایهٔ گویا باید توجه کرد که مفهوم‌های میدان، کوچک‌ترین میدانِ دربرگیرنده و پایهٔ گویا، که فقط بر عملیات «گویا» تکیه دارند، هیچ تغییری نمی‌کنند. افزون بر این، دستگاه تصویر $\Ssys^*_{\rho\rho}$ یک دستگاه صحیح‌ضرایب $\Ssys_{n\rho}$ را می‌توان به هر تعداد شیوه، باز به صورت دستگاهی صحیح‌ضرایب برگزید. چون استدلال‌های منتهی به کران شمار اعضا برای خانوادهٔ خطیِ صحیح‌ضرایب نیز برقرار می‌مانند، نتیجه می‌شود:

\emph{قضیهٔ پنجم$'$: هر دستگاه صحیح‌ضرایب $\Ssys_{n\rho}$ پایه‌ای گویا دارد؛ به‌ویژه می‌توان پایهٔ گویای یک خانوادهٔ خطیِ صحیح‌ضرایب $\Lsys_{n\rho}$ را چنان برگزید که حداکثر $\rho+1$ چندجمله‌ای داشته باشد.}

اکنون نظیرِ پایهٔ اینولوشن برای دامنه‌های صحیحِ ساخته از چندجمله‌ای‌ها (\S\ 8) را بررسی می‌کنیم. برای این منظور، قضیهٔ مربوط به \emph{توابع متقارنِ ردیف‌های کمیت} را از نظر \emph{صحیح‌بودنِ ضرایب} با لم زیر گسترش می‌دهیم:

\emph{لم. توابع متقارنِ چندجمله‌ای و صحیح‌ضرایبِ $n$ ردیف کمیت $(yz\cdots t)$:}
\[
\begin{array}{@{}l@{\quad}c@{}}
(1)&
\begin{array}{cccc}
  y_1&z_1&\cdots&t_1\\
  y_2&z_2&\cdots&t_2\\
  \vdots&\vdots&&\vdots\\
  y_n&z_n&\cdots&t_n
\end{array}
\end{array}
\]
\emph{توابعی چندجمله‌ای و صحیح‌ضرایب از شمار متناهی‌ای از میان خود هستند؛ یعنی از توابع متقارن مقدماتی}
\[
\begin{array}{@{}l@{\quad}c@{}}
(2)&
\begin{gathered}
  \sigma_1(y),\sigma_2(y),\ldots,\sigma_n(y);\quad
  \sigma_1(z),\sigma_2(z),\ldots,\sigma_n(z);\quad\ldots;\\
  \sigma_1(t),\ldots,\sigma_n(t);
\end{gathered}
\end{array}
\]
\emph{و توابع}
\[
\begin{array}{@{}l@{\quad}c@{}}
(3)&
  \chi_1(y,z,\ldots,t),\ldots,\chi_k(y,z,\ldots,t),
\end{array}
\]
\emph{که وابستگی‌شان به توابع (2) هم تمامِ جبری و هم صحیح‌ضرایب است.}

در واقع، وابستگیِ $n$ کمیتِ $y_1,
\ldots,y_n$ به $\sigma_1(y),
\ldots,\sigma_n(y)$ هم تمامِ جبری و هم صحیح‌ضرایب است؛ حکم متناظر برای $z_1,
\ldots,z_n$ نسبت به $\sigma_1(z),
\ldots,\sigma_n(z)$ و جز آن نیز برقرار است. پس میدان توابع متقارنِ ردیف‌های کمیت (1) میدانی جبری بر میدانِ تعیین‌شده به‌وسیلهٔ توابع (2) است، به‌گونه‌ای که کمیت‌های تمامِ جبری و صحیح‌ضرایبِ این میدان عیناً همان توابع متقارنِ چندجمله‌ای و صحیح‌ضرایبِ ردیف‌های کمیت (1) هستند. بنابراین، اگر توابع (3) را به‌عنوان دستگاه بنیادیِ کرونکر تعیین کنیم، لم ثابت می‌شود.

اکنون $\Jdom_{n\rho}$ دامنه‌ای صحیح و صحیح‌ضرایب، و $\Jdom^*_{\rho\rho}$ دامنهٔ تصویری صحیح‌ضرایب باشد، و
\[
  \Psi^*(z,u)=G^*(x)\cdot z^i+
  \sum{}'G^*_{\alpha\alpha_1\ldots\alpha_\rho}(x)\cdot
  z^\alpha u_1^{\alpha_1}\cdots u_\rho^{\alpha_\rho}
\]
یک فرم اینولوشنِ $\Jdom^*_{\rho\rho}$ باشد؛ در اینجا $G^*(x)$ می‌تواند از بُعد صفر نیز باشد، یعنی عددی از $[\Omega]$. آنگاه توابع متقارن مقدماتیِ متناظر با توابع (2)، با بعضی از خارج‌قسمت‌های زیر داده می‌شوند:
\[
\begin{array}{@{}l@{\quad}c@{}}
(4)&
  G^*_{\alpha\alpha_1\ldots\alpha_\rho}(x):G^*(x).
\end{array}
\]
چون بنا بر لم، وابستگیِ توابع (3) به (2) هم تمامِ جبری و هم صحیح‌ضرایب است، توابعِ متناظر با (3) در $\Kfield^*_{\rho\rho}$ ــ بنا بر تعمیم قضیهٔ گاوس به چندجمله‌ای‌های دارای ضرایب صحیح جبری ــ با عبارت‌های زیر داده می‌شوند:
\[
\begin{array}{@{}l@{\quad}c@{}}
(5)&
  K_\lambda^*(x):(G^*(x))^\sigma,
\end{array}
\]
که در آن $K_\lambda^*(x)$ چندجمله‌ای‌های صحیح‌ضرایب از $\Kfield^*_{\rho\rho}$ هستند. پس اگر با دامنه‌های تمامِ نسبی سروکار داشته باشیم، $K_\lambda^*(x)$ به دامنه تعلق دارند؛ و در حالت کلی خارج‌قسمتِ چندجمله‌ای‌هایی از $\Jdom^*_{\rho\rho}$ هستند.

اکنون اگر $F^*(x)$ چندجمله‌ای صحیحِ دلخواهی از $\Jdom^*_{\rho\rho}$ باشد، داریم
\[
  F^*(x)=\frac1i\Bigl(F^*(x)+F^*(x^{(1)})+\cdots+F^*(x^{(i-1)})\Bigr)
\]
و بنابراین $i\cdot F^*(x)$ تابعی متقارن، چندجمله‌ای و صحیح‌ضرایب در ردیف‌های کمیت $x,x^{(1)},\ldots,x^{(i-1)}$ است.

پس بنا بر لم و اصل انتقال، نتیجه می‌شود:

\emph{قضیهٔ ششم$'$: برای دامنه‌های صحیحِ صحیح‌ضرایب $\Jdom_{n\rho}$، هر فرم اینولوشن پایهٔ گویای ممتازی را تعیین می‌کند، به‌گونه‌ای که هر چندجمله‌ای $F(x)$ از $\Jdom_{n\rho}$ نمایش زیر را می‌پذیرد:}
\[
  F(x)=\frac{\Gamma(H(x),H_1(x),\ldots,H_\sigma(x))}{i\cdot H(x)^k},
\]
\emph{که در آن $\Gamma$ تابعی چندجمله‌ای و صحیح‌ضرایب است. برای دامنه‌های تمامِ نسبیِ صحیح‌ضرایب $\Gdom_{n\rho}$ که دامنهٔ تصویرشان $\Gdom^*_{\rho\rho}$ نیز تمامِ نسبی و صحیح‌ضرایب است، مخرج $H(x)$ با ضریبِ بالاترین درجهٔ فرم اینولوشنِ متناظر داده می‌شود.}

\section*{\S\ 15. دامنه‌های تمامِ نسبیِ صحیح‌ضرایبِ نوع اول و دستگاه‌های منظمِ صحیح‌ضرایب.}

در دستگاه‌های صحیح‌ضرایب، پایهٔ تمامیت جای خود را به پایهٔ تمامیتِ صحیح‌ضرایب می‌دهد؛ یعنی شمار متناهی‌ای از چندجمله‌ای‌های صحیح‌ضرایبِ دستگاه، به‌گونه‌ای که بتوان هر چندجمله‌ایِ صحیح‌ضرایبِ دستگاه را به‌صورت ترکیبی چندجمله‌ای و صحیح‌ضرایب از اعضای آن شمار متناهی نمایش داد.

قضیه‌های مربوط به پایهٔ تمامیتِ صحیح‌ضرایب در اصل هم‌راستا با قضیه‌های پیشین‌اند و فقط برهان‌ها به اصلاحات اندکی نیاز دارند.

نخست می‌توان تعریف V (\S\ 9) را مستقیماً منتقل کرد.

\emph{تعریف V$'$: کمیت $\xi$ نسبت به دامنهٔ تمامِ نسبیِ صحیح‌ضرایبِ $\Gdom_{n\rho}$ تمامِ جبری و صحیح‌ضرایب نامیده می‌شود اگر در معادله‌ای صدق کند که ضریب بالاترین درجهٔ آن واحدی از $[\Omega]$ و دیگر ضرایبش چندجمله‌ای‌هایی از $\Gdom_{n\rho}$ باشند.}

زیرا تعمیم قضیهٔ گاوس به چندجمله‌ای‌هایی با ضرایب صحیح جبری، همانند \S\ 9، نشان می‌دهد که از این تعریف می‌توان تعریف مبتنی بر معادلهٔ تحویل‌ناپذیر را به دست آورد.

از تعریف V$'$ انتقال تعریف VII نیز به دست می‌آید.

\emph{تعریف VII$'$: دامنهٔ تمامِ نسبیِ صحیح‌ضرایبِ $\Gdom_{n\rho}$ نوع اول نامیده می‌شود اگر به‌عنوان دامنهٔ تصویر، دامنهٔ تمامِ نسبیِ $\Gdom^*_{\rho\rho}$ را چنان داشته باشد که هر چندجمله‌ایِ صحیح‌ضرایبِ دلخواه در $x_1\rcdots x_\rho$ به‌طور تمامِ جبری و صحیح‌ضرایب به $\Gdom^*_{\rho\rho}$ وابسته باشد.}

برای اثبات متناهی‌مولدبودن این دامنه‌های نوع اول، نخست توجه می‌کنیم که همانند \S\ 10 از تعریف نتیجه می‌شود که ضریب بالاترین درجهٔ فرم اینولوشن واحدی از $[\Omega]$ است.

پس بنا بر قضیهٔ VI$'$ (\S\ 14)، هر چندجمله‌ایِ $F(x)$ از $\Gdom_{n\rho}$ نمایش زیر را می‌پذیرد:
\[
  F(x)=\frac{\Gamma(H_1(x)\rcdots H_\sigma(x))}{i}.
\]
اگر در اینجا قضیهٔ دوم هیلبرت دربارهٔ پایهٔ مدولیِ صحیح‌ضرایب (Math. Ann. 36) --- که برای ضرایب از اعداد صحیح معمولی بیان شده است --- را، پس از تعمیم آن به چندجمله‌ای‌هایی با ضرایب صحیح جبری، بر دستگاه ساخته‌شده از چندجمله‌ای‌های صحیح‌ضرایبِ $\Gamma(H_1\rcdots H_\sigma)$ اعمال کنیم،\srcfn{*)}{اگر $w_1\rcdots w_k$ دستگاهی بنیادی از اعداد صحیح جبریِ $\Omega$ باشد و بگذاریم
\[
  \Gamma(H_1\rcdots H_\sigma)=\Gamma_1(H)w_1+\cdots+\Gamma_k(H)w_k,
\]
آنگاه $\Gamma_1(H)\rcdots \Gamma_k(H)$ ضرایبی از اعداد صحیح معمولی دارند. اعمال قضیهٔ دوم بر دستگاه ساخته‌شده از فرم‌های $v_1\Gamma_1(H)+\cdots+v_k\Gamma_k(H)$ بی‌درنگ اعتبار این تعمیم را نشان می‌دهد.} وجود پایهٔ تمامیتِ صحیح‌ضرایب از آن آشکار می‌شود؛ یعنی:

\emph{قضیهٔ VII$'$: هر دامنهٔ تمامِ نسبیِ صحیح‌ضرایبِ نوع اول دارای پایهٔ تمامیتِ صحیح‌ضرایب است.}

اکنون تعریف دستگاه‌های منظم را نیز منتقل می‌کنیم:

\emph{تعریف VIII$'$: دستگاه صحیح‌ضرایبِ $\Ssys_{n\rho}$ منظمِ صحیح‌ضرایب نامیده می‌شود هرگاه در یک دامنهٔ تصویر $\Jdom^*_{\rho\rho}$ از کوچک‌ترین دامنهٔ صحیحِ صحیح‌ضرایبِ دربرگیرنده، $\rho$ چندجمله‌ای وجود داشته باشد، چنان‌که برآیندِ جزءهای دارای بالاترین درجهٔ آن‌ها واحدی از $[\Omega]$ باشد --- $R_0=\eps$.}

برای انتقال برهان وجود پایهٔ تمامیت، توجه می‌کنیم که لمِ \S\ 11 صورت قوی‌تر زیر را می‌پذیرد --- صورتی که از معادله‌های (5) و (6) در \S\ 11 و از این واقعیت به دست می‌آید که $A_1(\lambda)\rcdots A_\tau(\lambda)$ توابع چندجمله‌ایِ صحیح‌ضرایب در $\lambda_1\rcdots\lambda_\rho$ هستند:

\emph{شرط کافی برای آنکه هر چندجمله‌ایِ صحیح‌ضرایبِ دلخواه در $x_1\rcdots x_\rho$ به‌طور تمامِ جبری و صحیح‌ضرایب به $\rho$ چندجمله‌ایِ صحیح‌ضرایب وابسته باشد، این است که برآیندِ جزءهای دارای بالاترین درجهٔ این چندجمله‌ای‌ها واحدی از $[\Omega]$ باشد --- $R_0=\eps$.}

از این لم، درست همانند \S\ 12، برای هر چندجمله‌ایِ $F(x)$ از یک دستگاه منظمِ صحیح‌ضرایب، نمایش (3) در \S\ 12 نتیجه می‌شود؛ با این تفاوت که اکنون $\Gamma_1(G_i)\rcdots\Gamma_k(G_i)$ چندجمله‌ای‌های صحیح‌ضرایب در $G_i$ هستند. سپس، همانند \S\ 12، اعمال قضیهٔ دوم هیلبرت (Math. Ann. 36) در تعمیمش به چندجمله‌ای‌هایی با ضرایب صحیح جبری، همراه با تعریف کوچک‌ترین دامنهٔ صحیحِ صحیح‌ضرایبِ دربرگیرنده، وجود پایهٔ تمامیتِ صحیح‌ضرایب را نتیجه می‌دهد؛ یعنی:

\emph{قضیهٔ IX$'$: هر دستگاه منظمِ صحیح‌ضرایب دارای پایهٔ تمامیتِ صحیح‌ضرایب است.}

به‌عنوان نمونه‌ای از یک دامنهٔ تمامِ نسبیِ صحیح‌ضرایبِ نوع اول --- به قیاس با مثال 1 در \S\ 13 --- به مجموعهٔ $\Gdom_{nn}$ از چندجمله‌ای‌های صحیح‌ضرایبِ یک میدانِ گونهٔ لاگرانژی اشاره می‌کنیم. نوع اول و صحیح‌ضرایب‌بودنِ $\Gdom_{nn}$ از همانی زیر نتیجه می‌شود:
\[
  x_k^n-\sigma_1(x)x_k^{n-1}+\cdots\pm\sigma_n(x)=0
  \qquad(k=1,2,
\ldots,n);
\]
و چون برآیندِ توابع متقارن مقدماتی مقدار $\pm1$ دارد، $\Gdom_{nn}$ نیز یک دستگاه منظمِ صحیح‌ضرایب است.

نمونه‌ای دیگر، بنا بر لمِ \S\ 14، از مجموعهٔ همهٔ توابع متقارنِ چندجمله‌ای و صحیح‌ضرایب در $n$ ردیفِ کمیت به دست می‌آید.

ارلانگن، مه ۱۹۱۴.

\clearpage

\providecommand{\Hgrp}{\mathfrak{H}}
\providecommand{\GaloisRes}{\Phi}
\providecommand{\pol}[2]{\left(#1\frac{\partial}{\partial #2}\right)}
\providecommand{\Deltaop}{\Delta\!\left(\frac{\partial}{\partial x},\frac{\partial}{\partial y},\frac{\partial}{\partial z}\right)}
\providecommand{\Omegaop}{\Omega}
\providecommand{\NoetherSrcNote}[2]{\textsuperscript{#1}\begingroup\renewcommand{\thefootnote}{}\footnotetext{\textsuperscript{#1}~#2}\endgroup}


\clearpage
\setcounter{footnote}{0}
\section*{۷. قضیهٔ متناهی‌بودنِ ناورداهای گروه‌های متناهی}
\begin{center}
Math. Ann. ۷۷ (۱۹۱۶)، صص. ۸۹--۹۲
\end{center}

\begin{center}
{\Large قضیهٔ متناهی‌بودنِ ناورداهای گروه‌های متناهی.}\par
\vspace{1.2em}
از\par
\vspace{0.8em}
\srcspaced{امی نوتر} در ارلانگن.
\end{center}

در ادامه برهانی کاملاً مقدماتی برای متناهی‌بودنِ ناورداهای گروه‌های متناهی عرضه می‌کنیم --- برهانی که فقط بر نظریهٔ توابع متقارن تکیه دارد --- و هم‌زمان دستگاه کامل را واقعاً به دست می‌دهد؛ حال آنکه برهان معمول، که بر قضیهٔ هیلبرت دربارهٔ پایهٔ مدولی (Ann. 36) استوار است، تنها برهان وجود است.\NoetherSrcNote{*)}{برای نمونه، نک. Weber, \emph{Lehrbuch der Algebra}، ویرایش دوم، جلد ۲، \S\ 57.}

فرض کنیم گروه متناهی $\Hgrp$ از $h$ تبدیل خطی $A_1,\ldots,A_h$ با دترمینان ناصفر تشکیل شده باشد و $A_k$ نمایانگر تبدیل خطیِ
\[
  x_1^{(k)}=\sum_{\nu=1}^{n}a_{1\nu}^{(k)}x_\nu,\ldots,
  x_n^{(k)}=\sum_{\nu=1}^{n}a_{n\nu}^{(k)}x_\nu
\]
یا، به اختصار، $(x^{(k)})=A_k(x)$ باشد. بنابراین گروه $\Hgrp$ ردیف $(x)$ با عناصر $x_1,\ldots,x_n$ را به ردیف‌های $(x^{(k)})$ با عناصر $x_1^{(k)},\ldots,x_n^{(k)}$ می‌برد. چون تبدیل همانی باید در میان $A_1,\ldots,A_h$ باشد، ردیف‌های $(x^{(k)})$ خود ردیف $(x)$ را نیز در بر دارند. منظور از ناوردای مطلقِ تمامِ گویا برای گروه، تابعی تمامِ گویا در $x_1,\ldots,x_n$، یعنی تابعی چندجمله‌ای، است که با اعمال $A_1,\ldots,A_h$ عیناً بی‌تغییر بماند؛ پس برای چنین ناوردایی $f(x)$ داریم:
\begin{equation}
  f(x)=f(x^{(1)})=\cdots=f(x^{(h)})=
  \frac1h\sum_{k=1}^{h}f(x^{(k)}).
\tag{1}
\end{equation}

\subsection*{۱.}
فرمول (1) بیان می‌کند که $f(x)$ تابعی متقارن و تمامِ گویا، یعنی چندجمله‌ای، در ردیف‌های کمیت $(x^{(1)}),\ldots,(x^{(h)})$ است؛ و چون هر جملهٔ $f(x^{(k)})$ فقط همان یک ردیف $(x^{(k)})$ را در بر دارد، با ساده‌ترین حالت روبه‌رو هستیم که معمولاً \emph{تک‌ردیفی} نامیده می‌شود. بنابراین بنا بر قضیهٔ شناخته‌شده دربارهٔ توابع متقارنِ ردیف‌های کمیت،\NoetherSrcNote{*)}{نک. یادداشتِ بند ۲.} می‌توان $f(x)$ را به‌صورت چندجمله‌ای برحسب توابع متقارن مقدماتیِ این ردیف‌ها نمایش داد؛ یعنی برحسب ضرایب $G_{\alpha\alpha_1\ldots\alpha_n}(x)$ در «رزولونت گالوا»:
\[
\begin{aligned}
  \Phi(z,u)
  &=\prod_{k=1}^{h}\bigl(z+u_1x_1^{(k)}+\cdots+u_nx_n^{(k)}\bigr) \\
  &=z^h+
    \sum G_{\alpha\alpha_1\ldots\alpha_n}(x)
    z^{\alpha}u_1^{\alpha_1}\cdots u_n^{\alpha_n},
\end{aligned}
\qquad
\left(\begin{array}{c}
\alpha+\alpha_1+\cdots+\alpha_n=h\\
\alpha\ne h
\end{array}\right).
\]
در اینجا $G_{\alpha\alpha_1\ldots\alpha_n}(x)$ ناورداهایی از درجهٔ $\alpha_1+\cdots+\alpha_n$ در $x$ هستند. پس ثابت شد:

\emph{ضرایب $G_{\alpha\alpha_1\ldots\alpha_n}(x)$ رزولونت گالوا دستگاهی کامل از ناورداهای گروه تشکیل می‌دهند، به‌گونه‌ای که هر ناوردا را می‌توان به‌صورت چندجمله‌ای برحسب این شمار متناهی ناوردا نمایش داد.}

\subsection*{۲.}
همچنین می‌توان با آغاز از (1)، استدلال مقدماتی‌تر زیر را پیش برد،\NoetherSrcNote{**)}{این در واقع به برهانی برای قضیهٔ توابع متقارنِ ردیف‌های کمیت در حالت تک‌ردیفیِ یادشده در بالا می‌انجامد.} که هم‌زمان به دستگاه کامل دومی می‌رسد. قرار می‌دهیم
\[
  f(x)=a+b x_1^{\mu_1}\cdots x_n^{\mu_n}
       +\cdots+c x_1^{\nu_1}\cdots x_n^{\nu_n},
\]
که در آن $a,b,\ldots,c$ ثابت‌اند؛ آنگاه بنا بر (1) داریم:
\[
\begin{aligned}
 h\cdot f(x)
 &=h\cdot a+b\sum_{k=1}^{h}(x_1^{(k)})^{\mu_1}\cdots(x_n^{(k)})^{\mu_n}
     +\cdots
     +c\sum_{k=1}^{h}(x_1^{(k)})^{\nu_1}\cdots(x_n^{(k)})^{\nu_n} \\
 &=h\cdot a+b\cdot J_{\mu_1\ldots\mu_n}+\cdots+c\cdot J_{\nu_1\ldots\nu_n}.
\end{aligned}
\]
پس هر ناوردا را می‌توان به‌صورت خطی و بدون مخرج در این مولدها از ناورداهای ویژهٔ
\[
  J_{\mu_1\ldots\mu_n}
  =\sum_{k=1}^{h}(x_1^{(k)})^{\mu_1}\cdots(x_n^{(k)})^{\mu_n}
\]
ساخت، و کافی است ادعا را برای همین ناورداهای ویژه ثابت کنیم. اما $J_{\mu_1\ldots\mu_n}$، تا یک عامل عددی، ضریب $u_1^{\mu_1}\cdots u_n^{\mu_n}$ در عبارتِ
\[
  S_\mu=\sum_{k=1}^{h}
  \bigl(u_1x_1^{(k)}+\cdots+u_nx_n^{(k)}\bigr)^{\mu},
\]
است، که در آن $\mu=\mu_1+\cdots+\mu_n$ گذاشته‌ایم؛ این عبارت مجموع توانیِ مرتبهٔ $\mu$ برای $h$ فرم خطیِ
\[
  \xi_1=u_1x_1^{(1)}+\cdots+u_nx_n^{(1)},\ldots,
  \xi_h=u_1x_1^{(h)}+\cdots+u_nx_n^{(h)}
\]
است. چنان‌که می‌دانیم، بی‌نهایت مجموع توانی $S_\mu$ ترکیب‌های چندجمله‌ایِ
\[
  S_1,S_2,\ldots,S_h,
\]
هستند که ضرایبشان با ناورداهای
\[
  J_{\mu_1\ldots\mu_n}\qquad(\mu_1+\cdots+\mu_n\le h)
\]
داده می‌شوند؛ پس دستگاه کامل دومی به دست آمده است:

\emph{یک دستگاه کامل از ناورداهای گروه را همهٔ ناورداهای $J_{\mu_1\ldots\mu_n}$ با شرط $\mu_1+\cdots+\mu_n\le h$ تشکیل می‌دهند، که در آن $h$ مرتبهٔ گروه است.}

پیوند با بند ۱ از این ملاحظه به دست می‌آید که می‌توان مجموع‌های توانی $S_1,\ldots,S_h$ را با توابع متقارن مقدماتیِ $\xi_1,\ldots,\xi_h$ جایگزین کرد؛ این کار به همان دستگاه کاملِ $G_{\alpha\alpha_1\ldots\alpha_n}(x)$ می‌انجامد که آنجا داده شد. هر دو نتیجه نشان می‌دهند که همهٔ ناورداها را می‌توان به‌صورت چندجمله‌ای برحسب ناورداهایی بیان کرد که درجه‌شان در $x$ از مرتبهٔ $h$ گروه بیشتر نیست.

\subsection*{۳.}
از آنچه همین اکنون ثابت شد می‌توان نتیجه‌هایی برای نمایش گویا گرفت. هر ناوردای مطلقِ گویا را می‌توان --- چنان‌که با بسط صورت و مخرج از راه ضرب آن‌ها در مزدوج‌های مخرج نسبت به گروه بی‌درنگ دیده می‌شود --- به‌صورت خارج‌قسمت دو ناوردای مطلقِ چندجمله‌ای نمایش داد که لزوماً نسبت به هم اول نیستند. از این رو هر ناوردای مطلقِ گویا را می‌توان به‌صورت تابعی گویا از ضرایب $G_{\alpha\alpha_1\ldots\alpha_n}(x)$ رزولونت گالوای $\Phi(z,u)$ بیان کرد. این قضیه را به چند شیوه و بدون گذر از قضیهٔ متناهی‌بودن نیز می‌توان به‌آسانی ثابت کرد؛ صورت آن پیش‌تر در Weber II, \S\ 58 آمده است.\NoetherSrcNote{*)}{بااین‌حال، برهان ارائه‌شده در آنجا معتبر نیست؛ زیرا فقط نشان داده شده است که تابع $\Psi(t)$ در فرمول (7) ناورداهایی به‌عنوان ضریب دارد، نه اینکه این ناورداها را بتوان به‌صورت گویا برحسب ضرایب $\Phi(t)$ بیان کرد. اگر برای نمایش $x_i^{(k)}$ برحسب $\xi_k$، به‌جای فرمول درون‌یابی لاگرانژ، روشی شناخته‌شده مبتنی بر مشتق‌گیری به کار رود، این شکاف برطرف می‌شود. رابطه‌ای که با مشتق‌گیری همانی $\Phi(-\xi_k,u)=0$ نسبت به همهٔ $u$ها به دست می‌آید چنین است:
\[
 \left[\frac{\partial\Phi}{\partial u_i}
      -x_i^{(k)}\frac{\partial\Phi}{\partial z}\right]_{z=-\xi_k}=0.
\]
پس به‌جای فرمول‌های (7) و (8) وبر، برای هر تابع گویای $\omega(x)$ باید فرمول زیر قرار گیرد:
\[
 \omega(x_1^{(k)},\ldots,x_n^{(k)})=
 \left.
 \omega\!\left(
   \frac{\partial\Phi/\partial u_1}{\partial\Phi/\partial z},\ldots,
   \frac{\partial\Phi/\partial u_n}{\partial\Phi/\partial z}
 \right)\right|_{z=-\xi_k},
\]
که دیگر فقط ضرایب $\Phi(z,u)$ را در بر دارد؛ جمع آن بر همهٔ $k$ها، هنگامی که $\omega(x)$ ناوردا باشد، نمایش مطلوب را به دست می‌دهد.}

\subsection*{۴.}
سرانجام یادآور می‌شویم که نتیجه‌های به‌دست‌آمده در اینجا به‌طور ضمنی در مقالهٔ من «Körper und Systeme rationaler Funktionen» آمده‌اند؛\NoetherSrcNote{*)}{Math. Ann. 76, S. 161 (1915).} دستگاه کاملِ بند ۱ در قضیه‌های VI و VII، و برهانی برای نمایش گویا که از این قضیهٔ متناهی‌بودن مستقل است در قضیهٔ III آمده است.\NoetherSrcNote{**)}{برای ناورداهای نسبی، قضیه‌های VIII و IX برهانی برای متناهی‌بودن در بر دارند که آن نیز بر مبنایی متفاوت از مبنای معمول است، اما فقط برهان وجود است؛ علت این امر آن است که ناورداهای نسبی یک میدان تشکیل نمی‌دهند.} بااین‌همه، مسیر برهان و نتیجه در حالت ویژهٔ حاضر نسبت به نظریهٔ عمومی بسیار ساده‌تر می‌شوند. ایدهٔ به‌کاربستن پژوهش‌های عمومی دربارهٔ ناورداهای گروه‌های متناهی در گفت‌وگو با آقای ا. فیشر به ذهنم رسید؛ همچنین مضمونِ برهانِ داده‌شده در بند ۲ --- هرچند دستگاه کاملِ پدیدارشده در آنجا در حالت عمومی به‌طور گویا تعریف‌پذیر نیست --- و یادداشت بند ۳ از اوست.

ارلانگن، مه ۱۹۱۵.

\clearpage
\section*{۸. دربارهٔ نمایش چندجمله‌ایِ ناورداهای دستگاهی متشکل از تعداد دلخواهی فرم بنیادی}
\begin{center}
Math. Ann. ۷۷ (۱۹۱۶)، صص. ۹۳--۱۰۲
\end{center}

\begin{center}
{\Large دربارهٔ نمایش چندجمله‌ایِ ناورداهای دستگاهی\\
متشکل از تعداد دلخواهی فرم بنیادی.}\par
\vspace{1.2em}
از\par
\vspace{0.8em}
\srcspaced{امی نوتر} در ارلانگن.
\end{center}

د. هیلبرت در پایان مقاله‌اش در جشن‌نامهٔ شوارتس، «\srcLTR{Über die Invarianten eines Systems von beliebig vielen Grundformen}»، این حدس را مطرح می‌کند که \emph{این ناورداها را می‌توان به‌طور تمام و گویا، یعنی به‌صورت چندجمله‌ای، برحسب ناورداهای متناهیِ دستگاه $(J,PJ)$ نمایش داد}؛ در اینجا $J$ دستگاه کاملِ ناورداهای آن دسته از فرم‌های بنیادی است که مستقل خطی‌اند، و ناورداهای $PJ$ با فرایندهای قطبی از $J$ پدید می‌آیند.

در ادامه، در I نشان می‌دهم که این حدس واقعاً درست است و به‌سادگی از قضیهٔ کاهش نتیجه می‌شود؛ قضیه‌ای که می‌گوید \emph{هر فرم در تعداد دلخواهی ردیف، هر یک با $n$ متغیر، را می‌توان به‌صورت مجموع قطبی‌های فرم‌هایی نمایش داد که فقط $n$ ردیف ثابت دارند و خود با فرایندهای قطبی از فرم آغازین مشتق شده‌اند}. این قضیهٔ کاهش حالت ویژه‌ای است که نخست ف. مرتنس صورت‌بندی کرد، از تعمیمی که آ. کاپلی، ف. مرتنس و ژ. درویتس برای بسط سلسله‌ای کلبش--گوردان به فرم‌های دارای تعداد دلخواهی ردیف، هر یک با $n$ متغیر، داده‌اند؛ تعمیمی که با تبدیل صوریِ فرایندهای مشتق‌گیری اثبات شده است.\NoetherSrcNote{*)}{F. Mertens: «\srcLTR{Über eine Formel der Determinantentheorie}» (فرمول ۵)، \srcLTR{Sitzb. d. Ak. d. Wiss. Wien}، جلد ۹۱، بخش II (۱۸۸۵)؛ و با تفصیل بیشتر، همراه با کاربردها، در: «\srcLTR{Über ganze Funktionen von $m$ Systemen von je $n$ Unbestimmten}»، \srcLTR{Monatsh. f. Math. u. Phys.}، سال ۴ (۱۸۹۳). برای برخی از برهان‌های کاپلی، از ۱۸۸۲ به بعد، برای نمونه نک. \srcLTR{Math. Ann.} ۳۷، یا درس‌گفتارهای تکثیرشدهٔ «\srcLTR{Lezioni sulla teoria delle forme algebriche}» (۱۹۰۲). J. Deruyts: \emph{\srcLTR{Essai d'une theorie générale des formes algébriques}}، \srcLTR{Mém. d. l. Soc. Roy. des Sciences de Liège} (2) 17 (1892).}

در II نیز برهانی مفهومی‌تر برای قضیهٔ کاهش می‌دهم؛ بدین‌صورت که مسئله را مسئلهٔ هم‌ارزیِ خانواده‌های خطیِ فرم‌ها در نظر می‌گیرم. این دیدگاه و نیز بخش اساسی‌ای از استدلال‌های به‌کاررفته را از مقالهٔ «\srcLTR{Über algebraische Modulsysteme}»\NoetherSrcNote{*)}{E. Fischer: \srcLTR{Über algebraische Modulsysteme und lineare homogene partielle Differentialgleichungen mit konstanten Koeffizienten}، \S\ 1--4، \srcLTR{J. f. M.} 140 (1911).} و قضیه‌های منتشرنشدهٔ پس از آن از ا. فیشر گرفته‌ام؛ او با نهایت لطف اجازهٔ استفاده از آن‌ها را به من داد.

سرانجام در III نشان می‌دهم که استدلال‌های متناظر چگونه به عمومی‌ترین بسط‌های سلسله‌ای نیز می‌انجامند.

\subsection*{I.}
فرض کنیم $\theta$ فرمی باشد که در هر یک از $N>n$ ردیفِ
\[
  A_1,A_2,\ldots,A_N
\]
همگن است؛ به‌طور کلی، ردیف $A_k$ می‌تواند از $n$ عنصرِ
\[
  A_k^{(1)},A_k^{(2)},\ldots,A_k^{(n)}
\]
تشکیل شده باشد. ضرایب $\theta$ می‌توانند نامعین‌ها یا کمیت‌هایی از یک میدانِ معین باشند. همچنین $P$ را تابعی تمامِ گویا، یعنی چندجمله‌ای --- \emph{با ضرایب عددیِ گویا} --- در فرایندهای قطبی می‌گیریم:
\[
  P_{hk}=\left(A_h\frac{\partial}{\partial A_k}\right)
  =A_h^{(1)}\frac{\partial}{\partial A_k^{(1)}}+
   A_h^{(2)}\frac{\partial}{\partial A_k^{(2)}}+
   \cdots+
   A_h^{(n)}\frac{\partial}{\partial A_k^{(n)}},
\]
که در آن، از حاصل‌ضرب $P_{hk}P_{ij}\theta$ اعمال $P_{hk}$ بر $P_{ij}\theta$ مراد است.

آنگاه \emph{قضیهٔ کاهشِ} یادشده با همانیِ زیر داده می‌شود:
\begin{equation}
  \theta=\sum PZ,
\tag{1}
\end{equation}
که در آن فرم‌های $Z$ دیگر فقط ردیف‌های
\[
  A_1,A_2,\ldots,A_n
\]
را دربر دارند و از $\theta$ با فرایندهای قطبیِ $P$ مشتق شده‌اند.

اکنون دستگاه فرم‌های بنیادیِ $(F')$ را در نظر می‌گیریم که از $N$ فرم هم‌درجه با شمار یکسانی از متغیرها، یعنی
\[
  F_1,F_2,\ldots,F_N,
\]
تشکیل شده است و ضرایبشان، به‌ترتیب، همان $N$ ردیفِ
\[
  A_1,A_2,\ldots,A_N
\]
گرفته می‌شوند. فرض کنیم $(J)$ دستگاه کاملِ $F_1,\ldots,F_n$ --- متشکل از شمار متناهی ناوردا --- باشد، به‌گونه‌ای که هر ناوردای $F_1,\ldots,F_n$ را بتوان به‌طور تمام و گویا، یعنی به‌صورت چندجمله‌ای، برحسب ناورداهای $(J)$ بیان کرد. \emph{آنگاه حدس هیلبرت که باید اثبات شود این است که برای ناورداهای هم‌زمانِ $S$ از $F_1,\ldots,F_N$، ناورداهای متناهیِ $(J,PJ)$ دستگاهی کامل به دست می‌دهند؛ در اینجا $PJ$ به همهٔ ناورداهایی گفته می‌شود که می‌توان آن‌ها را با فرایندهای قطبیِ $P$ از $J$ مشتق کرد.}

در واقع، هر ناوردای هم‌زمانِ $S$ از $(F')$ که ضرایب عددیِ گویا دارد، فرمی است که در هر یک از $N$ ردیفِ $A_1,\ldots,A_N$ همگن است --- و ضرایب عددیِ گویا دارد ---؛ پس همانی (1) را می‌توان بر آن اعمال کرد. چون اعمال فرایندهای قطبیِ $P$ بر $S$، چنان‌که می‌دانیم، باز ناوردا پدید می‌آورد، فرم‌های $Z$ نیز ناوردا هستند؛ دقیق‌تر، ناورداهای هم‌زمانِ $F_1,
\ldots,F_n$ هستند، زیرا دیگر فقط ردیف‌های $A_1,
\ldots,A_n$ را دربر دارند. بنابراین بنا بر فرض، آن‌ها توابعی تمامِ گویا، یعنی چندجمله‌ای، از ناورداهای $(J)$ هستند. در نتیجه همانی (1) به صورت زیر درمی‌آید:
\[
  S=\sum PY,
\]
که در آن $Y$ها حاصل‌ضرب‌های توانیِ ناورداهای $(J)$ هستند. اما اجرای واقعیِ فرایندهای $P$ بر $Y$، بنا بر تعریف $P$، عبارت است از اجرای پیاپیِ عملیات سادهٔ قطبیِ $P_{hk}$ که تنها شمار متناهی‌ای تکرار می‌شوند؛ این عملیات، طبق قاعدهٔ مشتق‌گیریِ حاصل‌ضرب، به مجموع‌هایی از حاصل‌ضرب‌های ناورداهای $(J)$ و $(PJ)$ می‌انجامند. همین حکم برای اعمال $P_{hk}$ بر حاصل‌ضربی از $(J)$ و $(PJ)$ نیز برقرار است؛ پس $PY$ نیز فقط ترکیب‌های تمامِ گویا، یعنی چندجمله‌ای، از $(J,PJ)$ به دست می‌دهد. \emph{بدین‌ترتیب، نمایش‌پذیریِ تمام و گویای همهٔ ناورداهای $S$ برحسب $(J,PJ)$ ثابت شد.}

\subsection*{II.}
پیش از برهان قضیهٔ کاهش، مطالبی را دربارهٔ خانواده‌های خطیِ فرم‌ها بیان می‌کنیم. در اینجا از \emph{خانوادهٔ خطیِ فرم‌ها در $K$} دستگاهی از فرم‌های هم‌درجه را مراد می‌کنیم که به معنای زیر بسته است: هرگاه $\theta_1$ و $\theta_2$ به آن تعلق داشته باشند، فرم $c_1\theta_1+c_2\theta_2$ نیز به آن تعلق دارد؛ در اینجا $c_1,c_2$ کمیت‌هایی از میدانِ معین $K$ هستند و $K$ میدان ضرایب فرم‌های $\theta$ را دربر می‌گیرد. در میان این فرم‌ها همواره تعدادی متناهی --- مثلاً $\rho$ فرم --- روی $K$ مستقل خطی‌اند، در حالی که هر $\rho+1$ فرم رابطه‌ای خطی با ضرایبی از $K$ را برآورده می‌کنند؛ این خانواده دارای \emph{رتبهٔ $\rho$} نسبت به $K$ است.\NoetherSrcNote{*)}{خانواده‌های خطیِ عام‌تری از فرم‌ها --- که لازم نیست رتبه‌شان متناهی باشد --- با کنارگذاشتن قید هم‌درجه‌بودن، یا نیز با کنارگذاشتن قیدِ اینکه $K$ میدان ضرایب فرم‌های $\theta$ را دربر گیرد، به دست می‌آیند.} \emph{واقعیت} روشن زیر برقرار است: اگر $\mathcal L$ و $\mathcal T$ دو خانوادهٔ خطی در $K$ با رتبهٔ یکسان $\rho$ باشند و $\mathcal T$ زیرخانواده‌ای از $\mathcal L$ باشد، آنگاه $\mathcal L$ و $\mathcal T$ یکسان‌اند (هم‌ارزند).

با این مقدمات، به قضیهٔ کاهش می‌پردازیم. اکنون همانیِ (در I)
\[
  \theta=\sum PZ
\]
به همانیِ مشابهی برای $P\theta$ تبدیل می‌شود، اگر همان فرایند قطبیِ $P$ را بر دو طرف اعمال کنیم؛ پس قضیهٔ کاهش هم‌زمان نمایشی از \emph{همهٔ} قطبی‌های $\theta$ به صورت مجموع‌هایی از قطبی‌های ویژهٔ $PZ$ به دست می‌دهد. از سوی دیگر، چون این قطبی‌های ویژه بی‌گمان در مجموعهٔ همهٔ قطبی‌ها قرار دارند، \emph{قضیهٔ کاهش هم‌ارزی این دو خانوادهٔ خطیِ فرم‌ها را بیان می‌کند}؛ به‌ویژه، هم‌ارزی آن زیرخانواده‌هایی را نیز بیان می‌کند که فرم‌هایشان در هر ردیف \emph{به‌تنهایی} همان درجهٔ $\theta$ را دارند. برای اثبات این هم‌ارزی، یک خانوادهٔ میانی را نیز وارد می‌کنیم که فرم‌های $Z$ آن را تعیین می‌کنند. برای سادگی، \emph{نخست برهان را در حالت سه ردیف دوتاییِ $(N=3,n=2)$} می‌دهیم و سپس برهان عمومیِ موازی را به‌اختصار نشان می‌دهیم.

پس فرض کنیم $\theta(x,y,z)$ به‌ترتیب دارای درجه‌های $\alpha,\beta,p$ در ردیف‌های
\[
  x_1x_2;\quad y_1y_2;\quad z_1z_2;
\]
باشد؛ فعلاً ضرایب $\theta$ نامعین‌های $a$ (یا نیز عددهای گویا) باشند. آنگاه سه خانوادهٔ خطیِ فرم‌ها را در نظر می‌گیریم:

۱. خانوادهٔ $\mathcal L$، متشکل از همهٔ فرم‌های $f(x,y,z)$ که از $\theta$ با فرایندهای قطبیِ $P$ حاصل شده‌اند و در هر یک از ردیف‌های $x,y,z$ همان درجهٔ $\theta$ را دارند؛ $\mathcal L$ به‌ویژه خودِ $\theta$ را نیز دربر می‌گیرد. اگر $f(x,y,z)$ را فرم‌هایی در $x,y,z$ و نامعین‌های $a$ در نظر بگیریم، میدان ضرایب برابر با میدان $R$ اعداد گویاست؛ برای $K$ نیز باید $R$ را برگزید، زیرا تنها هنگامی که $c_1,c_2$ گویا باشند، $c_1P_1+c_2P_2$ دوباره فرایندی از نوع $P$ می‌شود. فرض کنیم $\mathcal L$ دارای رتبهٔ $\rho$ نسبت به $R$ باشد.

۲. خانوادهٔ $\mathcal S$، متشکل از همهٔ فرم‌های $g(\lambda;x,y)$ که با رابطهٔ زیر تعریف می‌شوند:
\[
  g(\lambda;x,y)=f(x,y,\lambda_1x+\lambda_2y),
\]
یا، به زبان فرایندهای قطبی:
\[
\begin{aligned}
  g(\lambda;x,y)
  &=\frac1{p!}\left[(\lambda_1x+\lambda_2y)\frac{\partial}{\partial z}\right]^p f(x,y,z)\srcfnmark{*)} \\
  &=g_0(xy)\lambda_1^p+g_1(x,y)\lambda_1^{p-1}\lambda_2+\cdots+g_p(xy)\lambda_2^p,
\end{aligned}
\]
که در آن:
\[
  i!(p-i)!\,g_i(xy)=
  \left(y\frac{\partial}{\partial z}\right)^i
  \left(x\frac{\partial}{\partial z}\right)^{p-i}f(xyz).
\]
\srcfntext{*)}{منظور، همانند I، این است:
\[
 \left(t\frac{\partial}{\partial x}\right)=t_1\frac{\partial}{\partial x_1}+t_2\frac{\partial}{\partial x_2},
\]
و به‌طور خاص:
\[
 \left[(\lambda_1x+\lambda_2y)\frac{\partial}{\partial z}\right]
 = (\lambda_1x_1+\lambda_2y_1)\frac{\partial}{\partial z_1}
 + (\lambda_1x_2+\lambda_2y_2)\frac{\partial}{\partial z_2},
\]
\[
 \left[z\left(\frac{\partial}{\partial\lambda_1}\frac{\partial}{\partial x}
 +\frac{\partial}{\partial\lambda_2}\frac{\partial}{\partial y}\right)\right]
 =z_1\left(\frac{\partial}{\partial\lambda_1}\frac{\partial}{\partial x_1}
 +\frac{\partial}{\partial\lambda_2}\frac{\partial}{\partial y_1}\right)
 +z_2\left(\frac{\partial}{\partial\lambda_1}\frac{\partial}{\partial x_2}
 +\frac{\partial}{\partial\lambda_2}\frac{\partial}{\partial y_2}\right).
\]}
بدین‌ترتیب $g_i(xy)$ فرم‌های $Z$ در همانیِ (۱) را به دست می‌دهند. فرم‌های $g$ در $x,y,\lambda$ و نامعین‌های $a$ دارای ضرایب عددیِ گویا هستند؛ فرض کنیم $\mathcal S$ دارای رتبهٔ $\sigma$ نسبت به $R$ باشد.

۳. خانوادهٔ $\mathcal T$ که از $\mathcal S$ به‌وسیلهٔ فرایند قطبیِ معکوس به دست می‌آید، همان‌گونه که $\mathcal S$ از $\mathcal L$ به دست آمد؛ فرم‌های $h$ در $\mathcal T$ با رابطهٔ زیر تعریف می‌شوند:
\[
\begin{aligned}
 h(x,y,z)
 &=\frac1{p!}\left[z\left(\frac{\partial}{\partial\lambda_1}\frac{\partial}{\partial x}
 +\frac{\partial}{\partial\lambda_2}\frac{\partial}{\partial y}\right)\right]^p g(\lambda;xy)\\
 &=h_0(xyz)+h_1(xyz)+\cdots+h_p(xyz),
\end{aligned}
\]
که بنا بر ۲ داریم:
\[
  h_i(xyz)=
  \left(z\frac{\partial}{\partial x}\right)^{p-i}
  \left(z\frac{\partial}{\partial y}\right)^i g_i(xy).
\]
پس هر $h_i$ یک فرم منفرد است و در نتیجه خودِ $h$ مجموعی از این فرم‌هاست؛ اولی فرمی از نوع $PZ$ و دومی مجموع فرم‌هایی از این نوع است؛ فرض کنیم $\mathcal T$ دارای رتبهٔ $\tau$ نسبت به $R$ باشد.

چنان‌که ساخت فرم‌های $h(x,y,z)$ در $\mathcal T$ نشان می‌دهد، این فرم‌ها با فرایندهای قطبیِ $P$ از $\theta$ حاصل شده‌اند و در هر یک از ردیف‌های $x,y,z$، به‌تنهایی، همان درجهٔ $\theta$ را دارند؛ پس \emph{خانوادهٔ $\mathcal T$ زیرخانواده‌ای از $\mathcal L$ را تشکیل می‌دهد}. بنا بر \emph{واقعیت} یادشده در بالا، برای اثبات هم‌ارزی فقط باید برابری رتبه‌ها را ثابت کرد. در این اثبات دیگر از ساختار ویژهٔ $f(xyz)$ --- یعنی قطبی‌بودن آنها از یک فرم واحدِ $\theta$ --- استفاده نمی‌شود.

\emph{الف) برای مقایسهٔ رتبهٔ $\mathcal L$ و$\mathcal S$ از لمِ کمکیِ (الف) استفاده می‌کنیم: از $f(x,y,\lambda_1x+\lambda_2y)=0$ ضرورتاً $f(x,y,z)=0$ نتیجه می‌شود.} این حکم بی‌درنگ از رابطهٔ همانیِ میان سه ردیف دوتایی به دست می‌آید:
\[
  (xy)z+(yz)x+(zx)y=0,
  \qquad \text{که در آن }(xy)=(x_1y_2-x_2y_1),\ldots,
\]
و این رابطه نتیجه می‌دهد:
\[
  f[x,y,(zy)x+(xz)y]=(xy)^p f(x,y,z).
\]
پس از $f(x,y,\lambda_1x+\lambda_2y)=0$ (به‌طور همانی در $\lambda$)، چنان‌که جایگذاریِ $\lambda_1=(zy)$ و $\lambda_2=(xz)$ نشان می‌دهد، و چون $(xy)\ne0$، ضرورتاً $f(x,y,z)=0$ به دست می‌آید.

بنا بر این لم، \emph{میان فرم‌های $\mathcal S$ و فرم‌های $\mathcal L$ تناظری یک به یک برقرار است}. زیرا از
\[
  f_1(x,y,\lambda_1x+\lambda_2y)=g(\lambda;xy),\qquad
  f_2(x,y,\lambda_1x+\lambda_2y)=g(\lambda;xy)
\]
ضرورتاً نتیجه می‌شود:
\[
  f_1(xyz)=f_2(xyz).
\]
تحت این تناظر، \emph{هر رابطه‌ای میان فرم‌های $\mathcal S$ برای فرم‌های یک‌به‌یک متناظرِ آنها در $\mathcal L$ نیز برقرار می‌ماند (و برعکس)}. زیرا از
\[
  c_1g^{(1)}(\lambda;xy)+c_2g^{(2)}(\lambda;xy)+\cdots+c_rg^{(r)}(\lambda;xy)=0
\]
بنا بر لم کمکی ضرورتاً نتیجه می‌شود:
\[
  c_1f^{(1)}(x,y,z)+c_2f^{(2)}(x,y,z)+\cdots+c_rf^{(r)}(x,y,z)=0.
\]
\emph{بدین‌ترتیب برابری رتبهٔ $\mathcal L$ و$\mathcal S$ ثابت شد؛ $\rho=\sigma$.}

\emph{ب) برای مقایسهٔ رتبهٔ $\mathcal S$ و$\mathcal T$ نیز به‌طور متناظر از لمِ کمکیِ (ب) استفاده می‌کنیم: از $h(x,y,z)=0$ ضرورتاً $g(\lambda;xy)=0$ نتیجه می‌شود.} برای اثبات، توجه می‌کنیم که بنا بر ۳، $h(x,y,z)=0$ هم‌ارزِ صفرشدن حاصل‌ضرب‌های توانیِ جداگانهٔ زیر است:
\[
 \left(\frac{\partial}{\partial\lambda_1}\frac{\partial}{\partial x_1}
 +\frac{\partial}{\partial\lambda_2}\frac{\partial}{\partial y_1}\right)^{p_1}
 \left(\frac{\partial}{\partial\lambda_1}\frac{\partial}{\partial x_2}
 +\frac{\partial}{\partial\lambda_2}\frac{\partial}{\partial y_2}\right)^{p_2}g(\lambda;x,y),
 \qquad p_1+p_2=p,
\]
از این امر، با مشتق‌گیریِ بیشتر، صفرشدن حاصل‌ضرب‌های توانیِ زیر نیز نتیجه می‌شود:
\[
\begin{aligned}
&\left(\frac{\partial}{\partial x_1}\right)^{\alpha_1}
 \left(\frac{\partial}{\partial x_2}\right)^{\alpha_2}
 \left(\frac{\partial}{\partial y_1}\right)^{\beta_1}
 \left(\frac{\partial}{\partial y_2}\right)^{\beta_2} \\
&\quad\cdot
 \left(\frac{\partial}{\partial\lambda_1}\frac{\partial}{\partial x_1}
 +\frac{\partial}{\partial\lambda_2}\frac{\partial}{\partial y_1}\right)^{p_1}
 \left(\frac{\partial}{\partial\lambda_1}\frac{\partial}{\partial x_2}
 +\frac{\partial}{\partial\lambda_2}\frac{\partial}{\partial y_2}\right)^{p_2}
 g(\lambda;xy)
\end{aligned}
\]
پس هر ترکیب خطی از این حاصل‌ضرب‌های توانی نیز صفر می‌شود و بنابراین هر فرمِ زیر صفر است:
\[
  f\left(\frac{\partial}{\partial x},\frac{\partial}{\partial y},
  \frac{\partial}{\partial\lambda_1}\frac{\partial}{\partial x}
  +\frac{\partial}{\partial\lambda_2}\frac{\partial}{\partial y}\right)g(\lambda,xy).
\]
اگر $f$ را به‌طور خاص چنان برگزینیم که
\[
  f(x,y,\lambda_1x+\lambda_2y)=g(\lambda;xy),
\]
آنگاه این رابطه نیز به دست می‌آید:
\[
  g\left(\frac{\partial}{\partial\lambda};
    \frac{\partial}{\partial x},
    \frac{\partial}{\partial y}\right)g(\lambda,x,y)=0;
\]
یا، اگر بنویسیم
\[
  g(\lambda,xy)=\sum G_i\lambda_1^{\varkappa_1}\lambda_2^{\varkappa_2}
  x_1^{\varkappa_3}x_2^{\varkappa_4}y_1^{\varkappa_5}y_2^{\varkappa_6},
\]
که در آن $G_i$ فرم‌های خطیِ نامعین‌های $a$ با ضرایب عددیِ گویا هستند، آنگاه
\[
  \sum \varkappa_1!\varkappa_2!\cdots\varkappa_6!\,G_i^2=0
\]
و در نتیجه $g(\lambda;xy)=0$.

از لم کمکیِ (ب)، درست مانند (الف)، \emph{برابری رتبهٔ $\mathcal S$ و$\mathcal T$ نتیجه می‌شود؛ $\sigma=\tau$.}

پس بنا بر (الف) و (ب) داریم: $\rho=\tau$؛ و بدین‌ترتیب، چون $\mathcal T$ زیرخانواده‌ای از $\mathcal L$ است، هم‌ارزی این دو خانوادهٔ خطی ثابت می‌شود. بنابراین هر فرم از $\mathcal L$، و به‌ویژه خودِ $\theta$، را می‌توان به صورت ترکیبی خطی با ضرایب عددیِ گویا از $\rho$ فرم مستقل خطیِ $h(x,y,z)$ نمایش داد:
\[
  \theta=\sum c_i h^{(i)}(x,y,z)=\sum PZ.
\]
اما این همانی که برای ضرایب نامعینِ $a$ در $\theta$ به دست آمد، با جایگزین‌کردن نامعین‌های $a$ با کمیت‌هایی از هر میدان دلخواه همچنان برقرار می‌ماند؛ \emph{پس قضیهٔ کاهش در حالت سه ردیف دوتایی به‌طور کلی ثابت شد.}

اکنون برهان عمومی برای $N$ ردیف، که هر یک $n$ متغیر دارند، کاملاً به‌طور موازی پیش می‌رود. فرض کنیم $\theta(A)$ دارای درجهٔ $\alpha_i$ در ردیف $A_i$ باشد. سپس دوباره سه خانوادهٔ خطیِ فرم‌ها را در نظر می‌گیریم:

۱. خانوادهٔ $\mathcal L$، متشکل از همهٔ فرم‌های $f(A)$ که از $\theta(A)$ با فرایندهای قطبیِ $P$ به دست آمده‌اند و در هر ردیفِ منفردِ $A_i$ همان درجهٔ $\theta(A)$ را دارند؛

۲. خانوادهٔ $\mathcal S$، متشکل از همهٔ فرم‌های $g(\lambda;A_1\cdots A_n)$ که از $f(A)$ با جایگذاری‌های زیر به دست می‌آیند:
\[
\begin{aligned}
 A_{n+1}&=\lambda_{n+1}^{(1)}A_1+\lambda_{n+1}^{(2)}A_2+
        \cdots+\lambda_{n+1}^{(n)}A_n,\\
 &\vdotswithin{=}\\
 A_N&=\lambda_N^{(1)}A_1+\lambda_N^{(2)}A_2+
        \cdots+\lambda_N^{(n)}A_n
\end{aligned}
\]
در اینجا ضرایب حاصل‌ضرب‌های توانیِ جداگانهٔ $\lambda$ در $g(\lambda;A_1\cdots A_n)$، فرم‌های $Z$ در همانی (1) را به دست می‌دهند؛

۳. خانوادهٔ $\mathcal T$، متشکل از همهٔ فرم‌های $h(A)$ که به صورت زیر تعریف می‌شوند:
\[
\begin{aligned}
 h(A)=&\left[A_{n+1}\left(
 \frac{\partial}{\partial\lambda_{n+1}^{(1)}}\frac{\partial}{\partial A_1}
 +\cdots+
 \frac{\partial}{\partial\lambda_{n+1}^{(n)}}\frac{\partial}{\partial A_n}\right)\right]^{\alpha_{n+1}} \\
&\cdots
\left[A_N\left(
 \frac{\partial}{\partial\lambda_N^{(1)}}\frac{\partial}{\partial A_1}
 +\cdots+
 \frac{\partial}{\partial\lambda_N^{(n)}}\frac{\partial}{\partial A_n}\right)\right]^{\alpha_N}
 g(\lambda;A_1\cdots A_n);
\end{aligned}
\]
و $h(A)$ مجموعی از فرم‌های $PZ$ است.

اکنون به سبب رابطهٔ همانیِ موجود میان هر $n+1$ ردیف، لمِ کمکیِ (الف) همچنان برقرار است؛ لمِ کمکیِ (ب) نیز برقرار می‌ماند، زیرا شمار متغیرها اصلاً در برهان آن دخالت نداشت. پس $\mathcal L$ و$\mathcal T$ هم‌رتبه‌اند و — چون $\mathcal T$ زیرخانواده‌ای از $\mathcal L$ است — یکسان‌اند. بنابراین همانیِ زیر برقرار است:
\[
  \theta=\sum PZ
\]
برای ضرایب نامعین، و در نتیجه برای مقادیری از هر میدان دلخواه به‌عنوان ضرایب نیز؛ \emph{بدین‌ترتیب قضیهٔ کاهش در حالت عمومی ثابت شده است.}

\subsection*{III.}
اکنون به‌اختصار نشان می‌دهیم که با استدلال‌های مشابه، \emph{بسط عمومی برحسب ردیف‌ها} (برای $N\le n$) نیز چگونه به دست می‌آید. فرض کنیم $Z$ فرمی باشد که در هر یک از $n$ ردیف $A_1,
\ldots,A_n$، هرکدام شامل $n$ کمیت، جداگانه همگن است؛ و $\Delta$ دترمینان این ردیف‌ها باشد. نخست همانیِ متناظر با (1) را ثابت می‌کنیم:
\begin{equation}
  Z\equiv\sum PH \pmod{\Delta},
\tag{2}
\end{equation}
که در آن فرم‌های $H$ دیگر فقط ردیف‌های $A_1,\ldots,A_{n-1}$ را دربر دارند و از $Z$ با فرایندهای $P$ به دست آمده‌اند.

برهان بر تبدیل رابطه‌های به‌دست‌آمده در II به هم‌نهشتی‌هایی به پیمانهٔ $\Delta$ استوار است. برای سادگی، \emph{سه ردیف سه‌تایی} $x,y,z$ را مبنا می‌گیریم و ضرایب را نامعین فرض می‌کنیم.

سه خانوادهٔ خطیِ $\mathcal L,\mathcal S,\mathcal T$ را مانند حالت ردیف‌های دوتایی در II تعریف می‌کنیم؛ رتبه‌هایشان را $\rho,\sigma,\tau$ می‌نامیم، در حالی که $\rho^*$ و$\tau^*$ رتبهٔ $\mathcal L$ و$\mathcal T$ به پیمانهٔ $\Delta$ را نشان می‌دهند (یعنی $\rho^*$ فرم، ولی نه $\rho^*+1$ فرم، از $\mathcal L$ به پیمانهٔ $\Delta$ مستقل خطی‌اند). آنگاه مانند II لم کمکیِ زیر برقرار است:

\emph{لم کمکی (الف): از $f(x,y,\lambda_1x+\lambda_2y)=0$ ضرورتاً $f(xyz)\equiv0\pmod{\Delta}$ نتیجه می‌شود (و برعکس).} زیرا به‌موجب رابطهٔ موجود میان چهار ردیف سه‌تایی،
\[
  \Delta_1x-\Delta_2y+\Delta_3z=\Delta t,
  \qquad \text{که در آن }\Delta_1=(yzt),\ldots,
\]
همواره میان سه ردیف سه‌تایی رابطه‌ای خطی به پیمانهٔ $\Delta$ برقرار است:
\[
  \Delta_1x-\Delta_2y+\Delta_3z\equiv0\pmod{\Delta},
\]
و در نتیجه، مانند II، داریم:
\[
  f(x,y,-\Delta_1x+\Delta_2y)\equiv\Delta_3^p f(x,y,z)\pmod{\Delta}.
\]
از $f(x,y,\lambda_1x+\lambda_2y)=0$ (به‌طور همانی نسبت به $\lambda$)، چون $\Delta$ تحویل‌ناپذیر است و $\Delta_3$ بر $\Delta$ بخش‌پذیر نیست، ضرورتاً $f(x,y,z)\equiv0\pmod{\Delta}$ نتیجه می‌شود. عکس حکم روشن است، زیرا $\Delta(x,y,\lambda_1x+\lambda_2y)=0$. از اینجا مانند II نتیجه می‌گیریم: $\rho^*=\sigma$.

لم کمکی (ب) همراه با برهانش همچنان برقرار است؛ پس $\sigma=\tau$.

برای اثبات $\tau=\tau^*$ از لم کمکیِ زیر استفاده می‌کنیم:

\emph{لم کمکی (ج): از $h(x,y,z)\equiv0\pmod{\Delta}$ ضرورتاً $h(x,y,z)=0$ نتیجه می‌شود.}
زیرا $h(x,y,z)$، چون یک فرم قطبی است، با اعمال فرایند شناخته‌شدهٔ $\Omega$ صفر می‌شود؛ این امر با معادلهٔ دیفرانسیل زیر بیان می‌شود:
\[
  \Delta\!\left(\frac{\partial}{\partial x},\frac{\partial}{\partial y},\frac{\partial}{\partial z}\right)h(xyz)
  -\left(\frac{\partial}{\partial x}\frac{\partial}{\partial y}\frac{\partial}{\partial z}\right)h(x,y,z)=0.\srcfnmark{*)}
\]
\srcfntext{*)}{طبق 3. داریم:
\[
 c\cdot\Omega(h)=-\left(
\frac{\partial}{\partial \xi}\frac{\partial}{\partial \eta}
\left(\frac{\partial}{\partial\lambda_1}\frac{\partial}{\partial \xi}
+\frac{\partial}{\partial\lambda_2}\frac{\partial}{\partial \eta}\right)
\right)
\left(x\frac{\partial}{\partial \xi}\right)^{\alpha-1}
\left(y\frac{\partial}{\partial \eta}\right)^{\beta-1}
\left[
z\left(\frac{\partial}{\partial\lambda_1}\frac{\partial}{\partial \xi}
+\frac{\partial}{\partial\lambda_2}\frac{\partial}{\partial \eta}\right)
\right]^{p-1}\cdot g(\lambda;\xi,\eta)
\]
و بنابراین $\Omega(h)$ به سبب صفرشدن عامل دترمینانی صفر می‌شود.}
اگر اکنون $h(x,y,z)=\Delta(x,y,z)\cdot k(x,y,z)$ باشد، با مشتق‌گیری بیشتر، مقدار زیر نیز صفر می‌شود:
\[
  k\!\left(\frac{\partial}{\partial x},\frac{\partial}{\partial y},\frac{\partial}{\partial z}\right)
  \Delta\!\left(\frac{\partial}{\partial x},\frac{\partial}{\partial y},\frac{\partial}{\partial z}\right)
  =
  h\!\left(\frac{\partial}{\partial x},\frac{\partial}{\partial y},\frac{\partial}{\partial z}\right)h(x,y,z),
\]
و در نتیجه $h(x,y,z)$ نیز صفر می‌شود. بنابراین $\rho^*=\sigma=\tau=\tau^*$ و چون $\mathcal T$ زیرخانواده‌ای از $\mathcal L$ است، \emph{همانیِ (2) ثابت شده است}؛ استدلالی متناظر برهان را در حالت $n$ متغیر به دست می‌دهد.

با اعمال همانی (2) بر عاملی که در $\Delta$ ضرب شده است، بسط زیر به دست می‌آید:
\[
  Z=\varphi_0+\Delta\varphi_1+\Delta^2\varphi_2+\cdots+\Delta^\mu\varphi_\mu,
\]
که هر یک از $\varphi_i$ها با اعمال فرایند $\Omega$ صفر می‌شود. بنابراین تکرار $i$بارهٔ فرایند --- زیرا به‌طور کلی داریم: $\Omega(\Delta\cdot\psi)=c_1\psi+c_2\Delta\cdot\Omega(\psi)$ --- نتیجه می‌دهد:
\[
  c\cdot \Omega^i(Z)\equiv\varphi_i\pmod{\Delta};
\]
از سوی دیگر، بنابر (2) داریم:
\[
  c\cdot\Omega^i(Z)\equiv\sum PH_i\pmod{\Delta},
\]
که در آن $H_i$ همان‌گونه از فرم $\Omega^i(Z)$ به دست می‌آید که $H$ از $Z$. پس بنابر لم کمکی (ج) (در تعمیم آن به $n$ متغیر) داریم: $\varphi_i=\sum PH_i$، و در نتیجه:
\begin{equation}
  Z=\sum PH+\Delta\cdot\sum PH_1+\Delta^2\cdot\sum PH_2+
  \cdots+\Delta^\mu\cdot\sum PH_\mu
\tag{3}
\end{equation}
\emph{به‌عنوان عمومی‌ترین بسط برحسب ردیف‌ها برای فرمی با $n$ ردیف، هر یک دارای $n$ متغیر، برحسب قطبی‌های فرم‌هایی که فقط $n-1$ ردیف دارند.}

چنان‌که می‌دانیم، بسط برحسب ردیف‌ها برای فرم‌های دارای $N$ ردیف از $n$ متغیر $(N<n)$ با یک جایگذاری ساده از این نتیجه به دست می‌آید. می‌گذاریم (برای $N=3$ و$n>3$):
\[
  u=\xi_1x+\xi_2y+\xi_3z,
  \quad v=\eta_1x+\eta_2y+\eta_3z,
  \quad w=\zeta_1x+\zeta_2y+\zeta_3z
\]
و $H(u,v,w)=Z(\xi,\eta,\zeta)$ را بر پایهٔ (3) بسط می‌دهیم. آنگاه، چون $\xi,\eta,\zeta$ فقط در ترکیب‌های $u,v,w$ ظاهر می‌شوند، داریم:
\[
  \left(\eta\frac{\partial}{\partial \xi}\right)=
  \left(v\frac{\partial}{\partial u}\right)\cdots,
  \qquad
  \Omega=\sum_{ikl}\left(\frac{\partial}{\partial u}
  \frac{\partial}{\partial v}
  \frac{\partial}{\partial w}\right)_{ikl}(xyz)_{ikl}
  =\nabla_{uvw}(xyz).
\]
به‌موجب تخصص $\xi_1=\eta_2=\zeta_3=1$؛ $\xi_2=\xi_3=\cdots=\zeta_2=0$، $H(u,v,w)$ به $H(x,y,z)$ و $\nabla_{uvw}(xyz)$، تا یک عامل عددی، به $\nabla_{xyz}(xyz)=\nabla_{xyz}$ تبدیل می‌شود، زیرا اعمال $(y\partial/\partial x)$ و مانند آن بر دترمینان‌های $(xyz)_{ikl}$ سبب صفرشدن آن‌ها می‌شود. چون حکم متناظر برای $N$ ردیف نیز برقرار است، از (3) بسط زیر نتیجه می‌شود:
\begin{equation}
  H=\sum P\Phi+\sum P\Phi_1+\cdots+\sum P\Phi_\mu,
\tag{4}
\end{equation}
که در آن $\Phi_i$ها از $\nabla^i_{A_1\ldots A_N}H$ به‌وسیلهٔ فرایندهای $P$ به دست آمده‌اند و فقط ردیف‌های $A_1,
\ldots,A_{N-1}$ را دربر دارند، جدا از ترکیب‌های دترمینانیِ پدیدارشونده در $\nabla^i$ --- که چنان‌که گفته شد هنگام قطبی‌سازی به حساب نمی‌آیند.

بنابراین، اگر این دترمینان‌ها در $\nabla^i$ با متغیرهای نامعین $p$ جایگزین شوند، می‌توان فرم‌های حاصل را دوباره بر پایهٔ (4) بسط داد؛ و سرانجام بسط زیر را به دست آورد:
\begin{equation}
  H=\left[\sum P\psi(p)\right]_{p=\text{دترمینان‌های }A_1\cdots A_N},
\tag{5}
\end{equation}
که در آن $\psi$ها از فرم آغازین به‌وسیلهٔ فرایندهای $P$ و $\nabla_{A_1\ldots A_N}(p)$، $\nabla_{A_1\ldots A_{N-1}}(p)$ و مانند آن به دست می‌آیند. با این بسط‌ها و ترکیب‌هایشان --- مثلاً با اعمال (3) بر فرم‌های $Z$ ظاهرشده در (1)، و سپس اعمال (4) و (5) --- \emph{همهٔ بسط‌های شناخته‌شده برای فرم‌های دارای هر تعداد دلخواه ردیف، هر یک با $n$ متغیر هموردا، استنفاد می‌شوند}.

ارلانگن، ۵ ژانویهٔ ۱۹۱۵.



\clearpage
\setcounter{footnote}{0}

\section*{9. عمومی‌ترین دامنه‌های ساخته‌شده از اعداد متعالیِ صحیح}
\begin{center}
{\Large\bfseries 9. عمومی‌ترین دامنه‌های ساخته‌شده از اعداد متعالیِ صحیح}\par
\vspace{0.75em}
نوشتهٔ امی نوتر در ارلانگن.\par
\vspace{1.2em}
\srcLTR{Math. Ann. 77 (1916)}، ص. ۱۰۳--۱۲۸
\end{center}

ا. تسرملو به‌موجب قضیهٔ خوش‌ترتیبی دامنه‌ای از
\glqq{}اعداد متعالیِ صحیح\grqq{} ساخته است؛ یعنی دامنه‌ای از اعداد
$\mathfrak G$ که دارای ویژگی‌های \emph{انتزاعیِ} زیر است:
\begin{enumerate}
\item[I.] مجموع، تفاضل و حاصل‌ضربِ هر دو عدد از $\mathfrak G$ باز هم عددی از
$\mathfrak G$ است.
\item[II.] هر عدد حقیقی یا مختلط خارج‌قسمتِ دو عدد از $\mathfrak G$ است.
\item[III.] هر عدد صحیحِ گویا (یا عدد صحیحِ جبری، برحسب مورد) عضوی از $\mathfrak G$ است.
\item[IV.] هیچ عدد گویای ناصحیح (یا عدد جبریِ ناصحیح، برحسب مورد) به
$\mathfrak G$ تعلق ندارد.\srcfn{*)}{\srcLTR{E. Zermelo, Über ganze
transzendente Zahlen, Math. Ann. 75}، ص. ۴۳۴ (۱۹۱۴).}
\end{enumerate}

این ساخت بر آن استوار است که قضیهٔ خوش‌ترتیبی وجود یک \emph{پایهٔ جبری}
برای همهٔ اعداد را آشکار می‌سازد؛ یعنی \emph{دستگاهی} $H$ \emph{از اعداد}
$\eta$ که \emph{هیچ رابطهٔ جبری‌ای میانشان برقرار نیست، اما همهٔ اعداد دیگر
را می‌توان به‌طور جبری برحسب آن‌ها بیان کرد}. به‌سبب استقلال جبریِ اعداد پایهٔ
$\eta$، می‌توان آن‌ها را \emph{متغیرهای نامعین} در نظر گرفت؛ ازاین‌رو دامنهٔ
مورد نظر به \emph{دامنه‌ای صحیح از توابع گویا و جبریِ متغیرهای نامعین}
$\eta$ تبدیل می‌شود،\srcfn{**)}{از همین رو بررسی‌های مقالهٔ حاضر تا اندازه‌ای
به موازات بررسی‌های مقالهٔ پیشینِ نویسنده، «میدان‌ها و دستگاه‌های توابع گویا»
(\srcLTR{Körper und Systeme rationaler Funktionen, Math. Ann. 76, S. 161
(1915)}) پیش می‌روند.} که ضرایبشان به میدان $K$ همهٔ اعداد جبری تعلق دارند.
دامنهٔ ویژهٔ $\mathfrak G_\eta$ که تسرملو ساخته است، آنگاه با
\emph{ویژگی‌های بنیادیِ} زیر مشخص می‌شود:

$\mathfrak G_\eta$ موارد زیر را در بر دارد:
\begin{enumerate}
\item[1)] همهٔ اعداد پایهٔ $\eta$،
\item[2)] همهٔ چندجمله‌ای‌های $\eta$ با ضرایب صحیح\srcfn{***)}{در سراسر
مقاله، مراد از \glqq{}صحیح\grqq{} این است که ضرایب به دامنهٔ صحیح $[K]$ همهٔ
اعداد صحیح جبری تعلق دارند. برای تعریف $\mathfrak G_\eta$ در بندهای 1) تا 5)
کافی بود فقط ضرایبی از اعداد صحیحِ گویا را در نظر بگیریم؛ اما در آن صورت
ویژگی‌های بنیادیِ دامنه‌های عمومی‌تر ساده‌بودنِ کمتری می‌داشتند.}،
\item[3)] همهٔ توابع تمامِ جبریِ $\eta$ با ضرایب صحیح.
\end{enumerate}
$\mathfrak G_\eta$ موارد زیر را در بر نمی‌گیرد:
\begin{enumerate}
\item[4)] همهٔ توابعِ گویای کسریِ $\eta$ با ضرایب صحیح،
\item[5)] همهٔ توابعِ جبریِ کسریِ $\eta$ با ضرایب صحیح.\srcfn{*)}{تسرملو،
همان‌جا، § ۳.}
\end{enumerate}

اکنون به‌آسانی آشکار می‌شود (نک. مثال‌های § ۲ و۳) که در برابر دامنهٔ تسرملو
$\mathfrak G_\eta$، دامنه‌هایی $\mathfrak G$ وجود دارند که با آن
\emph{تفاوتی اساسی} دارند --- یعنی دامنه‌هایی $\mathfrak G$ که در
\emph{هیچ خوش‌ترتیبی‌ای} همهٔ ویژگی‌های بنیادیِ 1) تا 5) را ندارند و در
نتیجه با هیچ خوش‌ترتیبی‌ای نمی‌توان آن‌ها را به‌وسیلهٔ ساخت تسرملو پدید آورد؛
به بیان دیگر: دامنه‌هایی $\mathfrak G$ وجود دارند که نمی‌توان آن‌ها را
به‌طور یک‌به‌یک و یکریخت بر $\mathfrak G_\eta$ نگاشت. ازاین‌رو پرسش دربارهٔ
\emph{عمومی‌ترین دامنه‌های ساخته‌شده از اعداد متعالیِ صحیح} پیش می‌آید؛
پرسشی که در ادامه به‌طور کامل پاسخ داده خواهد شد.

برای این منظور نخست باید معلوم کرد کدام‌یک از ویژگی‌های بنیادی پیامد
ویژگی‌های تعریف انتزاعی‌اند و کدام‌یک از ساخت ویژه ناشی می‌شوند. آشکار می‌شود
(§ ۱) که برای هر دامنهٔ داده‌شدهٔ $\mathfrak G$ خوش‌ترتیبی‌ای وجود دارد که
ویژگی‌های بنیادیِ 1) و2) را برآورده می‌کند؛ بنابراین \emph{هر دامنهٔ
$\mathfrak G$ را می‌توان به‌وسیلهٔ یک پایهٔ جبریِ همهٔ اعداد ساخت}. از
ویژگی‌های بنیادیِ دیگر هیچ‌یک لازم نیست برقرار باشد (§ ۲ و۳). به‌ویژه از این
امر نتیجه می‌شود که ویژگی انتزاعیِ دیگری از دامنهٔ تسرملو در همهٔ دامنه‌های
$\mathfrak G$ برقرار نیست: \emph{بسته‌بودن نسبت به وابستگیِ تمامِ جبری}.
این ویژگی را می‌توان چنین صورت‌بندی کرد:
\begin{enumerate}
\item[V.] هر عددی که به شمار متناهی‌ای از اعداد $\mathfrak G$ به‌طور
تمامِ جبری و صحیح‌ضرایب وابسته باشد، به $\mathfrak G$ تعلق دارد.
\end{enumerate}

دامنه‌های ویژه‌ای را که افزون بر I--IV دارای V نیز هستند، دامنه‌های
$\mathfrak H$ ساخته‌شده از \emph{اعداد متعالیِ تمامِ جبری} می‌نامیم. برای
عمومی‌ترین دامنه‌های $\mathfrak H$ نتیجه‌های سادهٔ زیر به دست می‌آیند (§ ۴):
\begin{enumerate}
\item[a)] \emph{اشتراک همهٔ دامنه‌های قابل‌ساخت به‌وسیلهٔ پایهٔ یکسان
$H$ در ردهٔ دامنه‌های $\mathfrak H$، با مجموعهٔ همهٔ توابع تمامِ جبریِ پایه داده
می‌شود --- یعنی با دامنهٔ تسرملو $\mathfrak G_\eta$ که خود یک دامنهٔ
$\mathfrak H$ است.} (بدین‌ترتیب، هم‌زمان توصیفی انتزاعی از دامنهٔ تسرملو
به دست می‌آید.)
\item[b)] \emph{هر دامنهٔ $\mathfrak H$ از راه الحاقِ تمامِ
جبری\srcfn{**)}{توضیح این عبارت در § ۴ آمده است.} یک دستگاه دلخواه
$\mathfrak S(\eta)$ به $\mathfrak G_\eta$ پدید می‌آید؛ دستگاه
$\mathfrak S(\eta)$ فقط باید این یک شرط را برآورده کند که بر اثر الحاق، هیچ
عدد جبریِ کسری‌ای وارد دامنه نشود.}
\end{enumerate}

نتیجه‌ها برای عمومی‌ترین دامنه‌های $\mathfrak G$، هنگامی که یک پایهٔ جبری
مبنای کار قرار گیرد، کمتر ساده‌اند. در اینجا اشتراک همهٔ دامنه‌های
$\mathfrak G$ خود یک دامنهٔ $\mathfrak G$ نیست؛ ازاین‌رو احاطه بر ساخت
عمومی‌ترین دامنه‌ها نیز دشوارتر است (§ ۵).

برای به‌دست‌آوردن نتیجه‌هایی همانند نتیجه‌های دامنه‌های $\mathfrak H$، پایهٔ
جبری با یک \emph{پایهٔ گویا}ی $\Theta$ جایگزین می‌شود --- یعنی با
\emph{دستگاهی} $\Theta$ \emph{از اعداد} $\vartheta$ که \emph{اجازه می‌دهد
همهٔ اعداد را به‌طور گویا بیان کنیم، در حالی که در دست‌کم یک خوش‌ترتیبی، هیچ
عدد پایه‌ای را نمی‌توان به‌طور گویا برحسب شمار متناهی‌ای از اعداد پیش از آن
نمایش داد}. بر این پایهٔ گویای $\Theta$ باید شرط جانبیِ دیگری نیز تحمیل شود:
دامنهٔ صحیحی که از همهٔ چندجمله‌ای‌های صحیح‌ضرایبِ $\vartheta$ها تشکیل
می‌شود، یعنی $[\Theta]$، هیچ عدد جبریِ کسری‌ای را دربر نداشته باشد. (ساختن چنین پایه‌ای
با شرط جانبی در § ۷ نشان داده می‌شود.) آنگاه برای عمومی‌ترین دامنه‌های
$\mathfrak G$ --- که می‌توان آن‌ها را دامنه‌های ساخته‌شده از \emph{اعداد
متعالیِ تمامِ نسبی} نیز نامید --- باز نتیجه‌های سادهٔ زیر به دست می‌آیند (§ ۶):
\begin{enumerate}
\item[a)] \emph{اشتراک همهٔ دامنه‌های قابل‌ساخت به‌وسیلهٔ پایهٔ گویای یکسان
$\Theta$ در ردهٔ دامنه‌های $\mathfrak G$، با مجموعهٔ همهٔ توابع تمامِ نسبیِ
پایه داده می‌شود --- یعنی با دامنهٔ $[\Theta]$ که خود یک دامنهٔ
$\mathfrak G$ است.}
\item[b)] \emph{هر دامنهٔ $\mathfrak G$ از راه الحاقِ تمامِ نسبیِ یک
دستگاه دلخواه $\mathfrak S(\vartheta)$ به $[\Theta]$ پدید می‌آید؛ دستگاه
$\mathfrak S(\vartheta)$ فقط باید این یک شرط را برآورده کند که بر اثر الحاق،
هیچ عدد جبریِ کسری‌ای وارد دامنه نشود.}
\end{enumerate}

برای آنکه قیاس با دامنه‌های $\mathfrak H$ کامل باشد، می‌بایست اعداد جبری را
در ویژگی‌های تعریف III وIV برای دامنه‌های $\mathfrak G$ وارد نمی‌کردیم. با این
تغییر در III وIV، ساخت عناصر تمام به‌وسیلهٔ یک پایهٔ گویا به \emph{هر} میدان
دلخواهی انتقال‌پذیر است (§ ۱۰)، حال‌آنکه ساخت تسرملو بسته‌بودن جبریِ میدان را
فرض می‌کند.

\emph{اگر $H$ همهٔ پایه‌های جبری و، به‌ترتیب، $\Theta$ همهٔ پایه‌های گویای
برآورندهٔ شرط جانبی را بپیمایند، نتیجه‌های a) وb) کل دامنه‌های $\mathfrak H$
و$\mathfrak G$ را به دست می‌دهند.} اگر همهٔ دامنه‌های یکریخت را در یک کلاس
گرد آوریم، می‌توان از این کل، زیرمجموعه‌ای برگزید که نمایندهٔ کل کلاس‌ها باشد
--- کلاس‌هایی که بنا بر § ۹ شمارشان دست‌کم بی‌نهایتِ شماراست --- (§ ۸).


\end{document}
